% !TeX program = xelatex
% Compile twice: xelatex -interaction=nonstopmode main.tex
\documentclass[UTF8,10.5pt,a4paper,oneside,fontset=fandol,scheme=plain]{ctexbook}
% ==================== Art-Book Editorial Preamble for AI Series Volume 6 ====================
\usepackage[top=24mm,bottom=24mm,left=23mm,right=23mm,headheight=22pt,headsep=7mm]{geometry}
\usepackage{amsmath,amssymb,mathtools,amsthm}
\usepackage{booktabs,tabularx,array,multirow,longtable}
\usepackage{graphicx}
\usepackage{xcolor}
\usepackage{tikz}
\usepackage{pgfplots}
\pgfplotsset{compat=1.18}
\usetikzlibrary{arrows.meta,positioning,shapes.geometric,calc,decorations.pathreplacing,backgrounds}
\usepackage[most]{tcolorbox}
\usepackage{fancyhdr}
\usepackage{listings}
\usepackage{enumitem}
\usepackage{hyperref}

% ==================== Curated Gallery Color Palette ====================
\definecolor{ink}{HTML}{1B2430}        % Deep charcoal-slate ink
\definecolor{teal}{HTML}{0B5555}       % Refined verdigris teal
\definecolor{tealDark}{HTML}{083B3B}   % Deep gallery spine teal
\definecolor{tealLight}{HTML}{F2F8F8}  % Delicate ice-teal parchment
\definecolor{tealLine}{HTML}{A3C6C5}   % Soft architectural rule
\definecolor{tealMist}{HTML}{DCEBEB}   % Watermark chapter numeral tint
\definecolor{navy}{HTML}{1D3B5E}       % Muted Prussian indigo
\definecolor{navyLight}{HTML}{F2F5FA}  % Soft indigo tint
\definecolor{ochre}{HTML}{9E4A22}      % Burnished terracotta-ochre
\definecolor{gold}{HTML}{B97A24}       % Fine burnished amber-gold accent
\definecolor{ochreLight}{HTML}{FCF5F0} % Warm parchment tint
\definecolor{plum}{HTML}{5B3A6B}       % Muted editorial plum
\definecolor{plumLight}{HTML}{F8F4FB}  % Soft plum tint
\definecolor{slate}{HTML}{526071}      % Muted architectural gray
\definecolor{codebg}{HTML}{F7F9FB}     % Clean light code canvas
\definecolor{codeborder}{HTML}{D4DCE5}

\hypersetup{
  unicode=true,
  colorlinks=true,
  urlcolor=teal,
  linkcolor=teal,
  citecolor=teal,
  pdftitle={天书中的智能 · 第6卷：强化学习、大模型后训练与AI系统},
  pdfauthor={天书中的智能：人工智能的七卷手记}
}

% ==================== Editorial Chapter Openers & Typography ====================
\newcommand{\chapWatermarkNum}{\ifnum\value{chapter}<10 0\fi\thechapter}
\ctexset{
  contentsname={目\quad 录},
  chapter={
    name={第,章},
    number=\arabic{chapter},
    format=\LARGE\sffamily\bfseries\color{tealDark},
    nameformat=\small\sffamily\bfseries\color{gold},
    aftername=\hfill{\fontsize{38}{38}\selectfont\sffamily\bfseries\color{tealMist}\chapWatermarkNum}\par\vspace{-4mm},
    aftertitle=\par\vspace{2.8mm}{\noindent\color{gold}\rule{26mm}{1.5pt}}{\color{tealLine}\rule{\dimexpr\linewidth-26mm\relax}{0.55pt}}\par,
    beforeskip=6pt,
    afterskip=10pt
  },
  section={
    format=\Large\sffamily\bfseries\color{teal},
    beforeskip=14pt plus 3pt minus 2pt,
    afterskip=6pt plus 1pt
  },
  subsection={
    format=\large\sffamily\bfseries\color{ink},
    beforeskip=10pt plus 2pt minus 1pt,
    afterskip=4pt plus 1pt
  },
  subsubsection={
    format=\normalsize\sffamily\bfseries\color{navy},
    beforeskip=7pt,
    afterskip=2pt
  }
}
\renewcommand{\lstlistingname}{代码}
\renewcommand{\figurename}{图}
\renewcommand{\tablename}{表}

% Artisanal Chapter Epigraph macro (2 args: quote + right-aligned attribution)
\newcommand{\chapEpigraph}[2]{%
  \vspace{-2mm}%
  \begin{tcolorbox}[
    enhanced,
    colback=ochreLight!45!white,
    colframe=ochreLight!45!white,
    borderline west={1.8pt}{0pt}{gold},
    boxrule=0pt,
    arc=0.8mm,
    left=4mm,right=3.5mm,top=2mm,bottom=2mm,
    before skip=0pt,after skip=8pt
  ]
    {\kaishu\normalsize\color{ink!85}“#1”}\par\vspace{1mm}%
    \hfill{\footnotesize\sffamily\color{slate}——\,#2}%
  \end{tcolorbox}%
}

% Exercise tier header macro
\newcommand{\exTier}[2]{%
  \par\vspace{3.8mm}\noindent
  \colorbox{tealLight}{\strut\;\sffamily\bfseries\footnotesize\color{tealDark}#1\;}\;%
  {\small\sffamily\bfseries\color{slate}#2}\par\nopagebreak\vspace{1.2mm}%
}
\newcommand{\exHint}[1]{\par\smallskip\noindent{\footnotesize\sffamily\color{navy}\textbf{【提示】}\,#1}}

\setlength{\parindent}{0pt}
\setlength{\parskip}{5pt plus 1.5pt minus 0.5pt}
\linespread{1.26}
\setlist[itemize]{leftmargin=1.8em,topsep=3pt,itemsep=2.5pt,parsep=1pt}
\setlist[enumerate]{leftmargin=2.0em,topsep=3pt,itemsep=2.5pt,parsep=1pt}
\setlength{\tabcolsep}{5.5pt}
\renewcommand{\arraystretch}{1.22}

% ==================== Running Headers & Footers ====================
\pagestyle{fancy}
\fancyhf{}
\fancyhead[L]{\small\sffamily\color{slate}天书中的智能 \;\textcolor{gold}{$\cdot$}\; 第 6 卷}
\fancyhead[R]{\small\sffamily\color{slate}\nouppercase{\leftmark}}
\fancyfoot[C]{\small\sffamily\color{slate}\textcolor{gold}{$\diamond$}\quad\textcolor{tealDark}{\thepage}\quad\textcolor{gold}{$\diamond$}}
\renewcommand{\headrulewidth}{0.4pt}
\renewcommand{\headrule}{\hbox to\headwidth{\color{tealLine}\leaders\hrule height \headrulewidth\hfill}}
\renewcommand{\chaptermark}[1]{\markboth{第 \thechapter\ 章\quad #1}{}}

\fancypagestyle{plain}{
  \fancyhf{}
  \fancyfoot[C]{\small\sffamily\color{slate}\textcolor{gold}{$\diamond$}\quad\textcolor{tealDark}{\thepage}\quad\textcolor{gold}{$\diamond$}}
  \renewcommand{\headrulewidth}{0pt}
}

% ==================== Theorem & Pedagogical Box Environments ====================
\theoremstyle{definition}
\newtheorem{theorem}{定理}[chapter]
\newtheorem{lemma}[theorem]{引理}
\newtheorem{proposition}[theorem]{命题}
\newtheorem{corollary}[theorem]{推论}
\newtheorem{definition}[theorem]{定义}
\newtheorem{example}[theorem]{例}
\newtheorem{remark}[theorem]{注记}
\renewcommand{\proofname}{\textbf{\sffamily\color{teal}证明}}

\newtcolorbox{thmBox}[1][]{
  enhanced,breakable,
  colback=tealLight,
  colframe=tealLine,
  borderline west={2.4pt}{0pt}{teal},
  boxrule=0.45pt,
  arc=1.5mm,
  left=4mm,right=3.8mm,top=2.8mm,bottom=2.8mm,
  before skip=6.5pt,after skip=6.5pt,
  #1
}

\newtcolorbox{defBox}[1][]{
  enhanced,breakable,
  colback=navyLight,
  colframe=navy!30,
  borderline west={2.4pt}{0pt}{navy},
  boxrule=0.45pt,
  arc=1.5mm,
  left=4mm,right=3.8mm,top=2.8mm,bottom=2.8mm,
  before skip=6.5pt,after skip=6.5pt,
  #1
}

\newtcolorbox{insightBox}[1]{
  enhanced,breakable,
  colback=white,
  colframe=navy!35,
  borderline west={2.4pt}{0pt}{gold},
  coltitle=navy,
  fonttitle=\sffamily\bfseries\small,
  attach title to upper={\quad},
  title={【直觉与下游坐标】#1},
  boxrule=0.45pt,
  arc=1.2mm,
  left=4mm,right=3.8mm,top=2.6mm,bottom=2.6mm,
  before skip=6.5pt,after skip=6.5pt
}

\newtcolorbox{recallBox}[1]{
  enhanced,breakable,
  colback=codebg,
  colframe=codeborder,
  borderline west={2.2pt}{0pt}{slate},
  coltitle=slate,
  fonttitle=\sffamily\bfseries\small,
  attach title to upper={\quad},
  title={【前置速忆】#1},
  boxrule=0.4pt,
  arc=1.2mm,
  left=3.8mm,right=3.5mm,top=2.4mm,bottom=2.4mm,
  before skip=5.5pt,after skip=5.5pt
}

\newtcolorbox{failBox}[1]{
  enhanced,breakable,
  colback=ochreLight,
  colframe=ochre!35,
  borderline west={2.4pt}{0pt}{ochre},
  coltitle=ochre,
  fonttitle=\sffamily\bfseries\small,
  attach title to upper={\quad},
  title={【反例、数值陷阱与排错】#1},
  boxrule=0.45pt,
  arc=1.2mm,
  left=4mm,right=3.8mm,top=2.6mm,bottom=2.6mm,
  before skip=6.5pt,after skip=6.5pt
}

\newtcolorbox{extBox}[1]{
  enhanced,breakable,
  colback=plumLight,
  colframe=plum!30,
  borderline west={2.4pt}{0pt}{plum},
  coltitle=plum,
  fonttitle=\sffamily\bfseries\small,
  attach title to upper={\quad},
  title={【选读拓展 · 初读可跳过】#1},
  boxrule=0.45pt,
  arc=1.2mm,
  left=4mm,right=3.8mm,top=2.6mm,bottom=2.6mm,
  before skip=6.5pt,after skip=6.5pt
}

\newtcolorbox{chapIntro}[1]{
  enhanced,
  colback=tealLight,
  colframe=tealLine,
  borderline west={2.8pt}{0pt}{tealDark},
  boxrule=0.45pt,
  arc=1.5mm,
  left=4.2mm,right=4mm,top=3mm,bottom=3mm,
  before skip=2pt,after skip=10pt,
  title={\sffamily\bfseries\small\color{tealDark}#1},
  colbacktitle=tealLight,
  attach boxed title to top left={yshift=-1mm,xshift=2mm},
  boxed title style={empty,size=minimal}
}

% ==================== Code Listing Style ====================
\lstdefinestyle{pycode}{
  language=Python,
  basicstyle=\footnotesize\ttfamily\color{ink},
  keywordstyle=\bfseries\color{teal},
  commentstyle=\itshape\color{slate},
  stringstyle=\color{ochre},
  backgroundcolor=\color{codebg},
  frame=single,
  rulecolor=\color{codeborder},
  framesep=5pt,
  xleftmargin=4pt,
  xrightmargin=4pt,
  numbers=left,
  numberstyle=\tiny\ttfamily\color{slate!80},
  numbersep=6pt,
  columns=fullflexible,
  keepspaces=true,
  showstringspaces=false,
  breaklines=true,
  aboveskip=6.5pt,
  belowskip=6.5pt
}
\lstset{style=pycode}

% ==================== Math Shortcuts ====================
\newcommand{\R}{\mathbb{R}}
\newcommand{\N}{\mathbb{N}}
\newcommand{\E}{\mathbb{E}}
\newcommand{\Prob}{\mathbb{P}}
\newcommand{\Var}{\operatorname{Var}}
\newcommand{\DKL}{D_{\mathrm{KL}}}
\newcommand{\clip}{\operatorname{clip}}
\newcommand{\softmax}{\operatorname{softmax}}
\newcommand{\diag}{\operatorname{diag}}
\newcommand{\argmin}{\operatorname*{arg\,min}}
\newcommand{\argmax}{\operatorname*{arg\,max}}
\newcommand{\norm}[1]{\left\|#1\right\|}
\newcommand{\abs}[1]{\left|#1\right|}
\newcommand{\inner}[2]{\left\langle #1,\, #2 \right\rangle}

\makeatletter
\setlength{\@fpsep}{16pt}
\makeatother


\begin{document}
\color{ink}

% ==================== Cover Page (Gallery Art-Book Edition) ====================
\frontmatter
\thispagestyle{empty}
\begin{tikzpicture}[remember picture, overlay]
  % Architectural deep teal spine + burnished gold hairline + vertical monograph label
  \fill[tealDark] (current page.north west) rectangle ([xshift=11mm]current page.south west);
  \fill[gold] ([xshift=11mm]current page.north west) rectangle ([xshift=12.2mm]current page.south west);
  \node[rotate=90, anchor=center, font=\sffamily\bfseries\scriptsize, color=tealMist, opacity=0.85]
    at ([xshift=5.5mm]current page.west) {INTELLIGENCE FROM THE BOOK \quad $\cdot$ \quad VOLUME VI};

  % Bespoke fine-art geometric Policy Trust-Region & Roofline Architecture motif (top-right)
  \begin{scope}[shift={([xshift=-95mm, yshift=-20mm]current page.north east)}, scale=0.63, transform shape]
    % Subtle coordinate axes & Roofline knee curve (compute ceiling + memory slope)
    \draw[teal, opacity=0.14, line width=0.55pt] (-2.25,-1.15) -- (2.25,-1.15);
    \draw[navy, opacity=0.22, line width=0.95pt] (-2.15,-1.15) -- (-0.55,0.85) -- (2.15,0.85);
    \draw[teal, opacity=0.16, line width=0.6pt, dashed] (-0.55,-1.15) -- (-0.55,0.85);
    \fill[gold, opacity=0.55] (-0.55,0.85) circle (2.2pt);
    % Nested KL Trust-Region policy rings
    \draw[teal, opacity=0.18, line width=0.75pt] (0.55, -0.10) circle (1.42cm);
    \draw[gold, opacity=0.45, line width=1.0pt] (0.55, -0.10) circle (0.92cm);
    \draw[teal, opacity=0.14, line width=0.6pt, dashed] (0.55, -0.10) circle (0.46cm);
    % Smooth policy optimization trajectory ascending into the trust region
    \draw[ochre, opacity=0.55, line width=1.25pt, -{Stealth[length=2.1mm]}]
      (-1.55,-0.88) to[out=32, in=-150] (-0.35,-0.48)
      to[out=28, in=-135] (0.45,-0.12)
      to[out=42, in=-120] (1.12,0.42);
    \foreach \x/\y in {-1.55/-0.88, -0.35/-0.48, 0.45/-0.12, 1.12/0.42} {
      \fill[tealDark, opacity=0.50] (\x,\y) circle (2.0pt);
    }
  \end{scope}
\end{tikzpicture}
\vspace*{1mm}
\noindent{\small\sffamily\bfseries\color{teal}天书中的智能 \;\textcolor{gold}{$\cdot$}\; 人工智能的七卷手记 \hfill 第 6 卷 \quad 2026 年 11 月版}
\vspace{1.5mm}
\noindent{\color{gold}\rule{30mm}{1.8pt}}{\color{tealLine}\rule{\dimexpr\linewidth-30mm\relax}{0.7pt}}
\vspace{6.5mm}

{\fontsize{25}{33}\selectfont\sffamily\bfseries\color{ink}强化学习、大模型后训练与 AI 系统}\par
\vspace{2.5mm}
{\large\sffamily\color{teal}从时序差分、PPO 与 DPO 对齐，走向 Roofline 硬件、FlashAttention 与分布式推理}\par
\vspace{4.5mm}
{\small\color{slate}对标 UC Berkeley CS285 $\cdot$ Stanford CS234 / CS336 $\cdot$ CMU 10-414 / 10-703 $\cdot$ 清华智班与南大《强化学习》}\par

\vspace{5.5mm}

\begin{tcolorbox}[
  enhanced,
  colback=tealLight,
  colframe=tealLine,
  borderline west={2.6pt}{0pt}{tealDark},
  arc=2mm,
  boxrule=0.5pt,
  left=4.5mm,right=4.5mm,top=3mm,bottom=3mm
]
\textbf{\sffamily\color{tealDark}本卷定位：贯通“交互决策—人类对齐—底层系统”的现代 AI 工程脊梁}\par\vspace{1.2mm}
\small
预训练赋予模型世界知识，而真正让智能体学会复杂决策、对齐人类意图并在千万卡集群上高速运转的，是本卷的两大支柱——\textbf{现代强化学习（RL）}与\textbf{AI 计算系统（AI Systems）}：
\begin{itemize}
  \item \textbf{深度强化学习主线}：从无模型 MC/TD(0)、SARSA 与 Q-Learning 出发，攻克 DQN 致命三元组、策略梯度定理、基线降方差与 GAE，再到信任域 TRPO、PPO-Clip 与去 Critic 的 GRPO。
  \item \textbf{大模型后训练与对齐}：严格推导 KL 正则 RLHF 的闭式最优策略，一步消去配分函数导出 DPO 损失；深入剖析奖励黑客（Reward Hacking）与可验证奖励推理（RLVR）。
  \item \textbf{硬件与算子优化}：直接衔接计算机组成与操作系统，用 Roofline 模型定量区分 Prefill 计算受限与 Decode 访存受限，完整证明 FlashAttention 的在线 Softmax 分块递推。
  \item \textbf{大规模分布式与推理压缩}：建立 ZeRO/FSDP、张量并行与流水线并行的精确显存与通信账本，结合 LoRA 低秩微调、INT8/INT4 量化与 PagedAttention 实测验证。
\end{itemize}
\end{tcolorbox}

\vspace{4.5mm}

\noindent\textbf{\sffamily\color{tealDark}第 6 卷章节导览}
\vspace{1.5mm}

{\footnotesize\renewcommand{\arraystretch}{1.15}
\begin{tabularx}{\linewidth}{@{}p{15mm}p{47mm}X@{}}
\toprule
\textbf{章节} & \textbf{标题} & \textbf{核心理论、完整证明与实测实验} \\
\midrule
序\quad 章 & 本卷坐标与八校考点覆盖对照 & 从预训练迈向决策与系统、八校对标表与本卷知识依赖 DAG \\
第 1 章 & 无模型强化学习：TD、SARSA 与 Q-Learning & MC vs TD(0)、Robbins--Monro、悬崖行走安全/最优分化与 Double Q \\
第 2 章 & 深度价值、策略梯度与 Actor-Critic & 致命三元组与 DQN、策略梯度定理、基线零偏降方差证明与 GAE($\lambda$) \\
第 3 章 & 信任域、PPO 裁剪与 GRPO 演进 & 重要性采样、TRPO KL 约束、PPO-Clip 梯度分情形剖析与 GRPO 群体优势 \\
第 4 章 & 大模型后训练：SFT、RLHF、DPO 与 RLVR & Bradley--Terry、RLHF 闭式最优解、DPO 精确推导、奖励黑客与推理缩放 \\
第 5 章 & AI 硬件、Roofline 模型与 FlashAttention & GPU 访存墙、Prefill/Decode 算术强度、显存账本与在线 Softmax 分块实测 \\
第 6 章 & 分布式训练、LoRA、量化与高效推理 & DDP/ZeRO/TP/PP 通信显存公式、LoRA、INT8/INT4 量化与 PagedAttention \\
附\quad 录 & 全卷 48 道分层习题完整精解 & 第 1--6 章 A/B/C/D 四层习题分步推导、反例验算与代码自检基准 \\
\bottomrule
\end{tabularx}}

\vfill
\noindent{\color{tealLine}\rule{\linewidth}{0.5pt}}
{\footnotesize\color{slate}配套实验脚本 \texttt{code/vol6\_experiments.py} 与实测数据随本卷提供。上一卷：第 5 卷；下一卷：第 7 卷《AI 科研实践、严谨评测、安全伦理与顶会复现》。}

% ==================== Table of Contents ====================
\clearpage
\tableofcontents

% ==================== Main Matter ====================
\mainmatter
\chapter*{序章：从表征学习走向交互决策与系统工程}
\addcontentsline{toc}{chapter}{序章：从表征学习走向交互决策与系统工程}
\markboth{序章\quad 本卷坐标与八校考点对照}{}

\chapEpigraph{你不必相信上帝，但你应当相信“天书”的存在。}{Paul Erdős（保罗·埃尔德什，谈论记录最美证明的 \emph{THE BOOK}）}

\begin{chapIntro}{本章定位与系统闭环}
完成第 5 卷后，你已经能用自动微分与 Transformer 训练出强大的生成与表征模型。但在真实世界与前沿人工智能实验室中，还横亘着两座同样决定成败的大山：
\begin{enumerate}
  \item \textbf{怎么让模型在动态反馈中学会做决策、与人类价值观及客观正确性对齐？}——这属于\textbf{强化学习（RL）与大模型后训练}（第 1--4 章）。
  \item \textbf{当模型参数突破十亿、千亿，显存和访存带宽成了唯一瓶颈时，如何让算法与硬件同频共振？}——这属于\textbf{AI 计算系统与推理工程}（第 5--6 章）。
\end{enumerate}
本卷将这两条常被传统教材割裂的主线融为一炉：前四章把第 3 卷已知模型的 MDP 推进到无模型 TD/DQN、策略梯度、PPO、GRPO 与 DPO；后两章则直接打通计算机组成、操作系统分页与 GPU 硬件瓶颈，让你亲手算清每一字节显存与每一次 HBM 访存。
\end{chapIntro}

\section*{0.1\quad 第 6 卷八校核心考点覆盖对照表}
\addcontentsline{toc}{section}{0.1\quad 第 6 卷八校核心考点覆盖对照表}

\begin{center}
\footnotesize
\begin{tabularx}{\linewidth}{@{}p{35mm}p{35mm}p{15mm}X@{}}
\toprule
\textbf{八校核心考点} & \textbf{主要对标课程} & \textbf{本卷位置} & \textbf{下游直接用途} \\
\midrule
MC、TD(0)、SARSA、Q-Learning 与 Double Q & Berkeley CS285 / Stanford CS234 & 第 1 章 & 无模型时序差分控制与安全/最优策略权衡 \\
致命三元组、DQN、策略梯度定理与 GAE($\lambda$) & CS285 / CMU 10-703 & 第 2 章 & 深度强化学习与方差缩减理论基石 \\
重要性采样、TRPO、PPO-Clip 与 GRPO & CS285 / Stanford CS336 & 第 3 章 & 连续控制与现代推理模型强化学习的核心算法 \\
SFT、Bradley--Terry、RLHF、DPO 与 RLVR & Stanford CS336 / CS224n & 第 4 章 & 大模型指令微调、偏好对齐与推理时计算缩放 \\
GPU 内存墙、Roofline 模型与 FlashAttention & CMU 10-414 / Stanford CS336 & 第 5 章 & 识别算力/带宽瓶颈、显存核算与 IO 感知算子设计 \\
DDP/ZeRO/TP/PP、LoRA、INT4 量化与 PagedAttention & Stanford CS336 / MIT 6.5940 & 第 6 章 & 千亿参数分布式训练与高吞吐低延迟线上推理 \\
\bottomrule
\end{tabularx}
\end{center}

\section*{0.2\quad 第 6 卷知识阶梯（零先修倒挂）}
\addcontentsline{toc}{section}{0.2\quad 第 6 卷知识阶梯（零先修倒挂）}

\begin{center}
\resizebox{0.98\linewidth}{!}{%
\begin{tikzpicture}[
  node distance=6.5mm and 7.5mm,
  ch/.style={
    rectangle, rounded corners=1.8mm, draw=teal, thick, fill=tealLight,
    text width=34mm, minimum height=14.5mm, align=center, font=\sffamily\footnotesize, text=ink, inner sep=4pt
  },
  goal/.style={
    rectangle, rounded corners=1.8mm, draw=ochre, very thick, fill=ochreLight,
    text width=37mm, minimum height=14.5mm, align=center, font=\sffamily\footnotesize, text=ink, inner sep=4pt
  },
  arr/.style={-{Stealth[length=2.2mm]}, thick, teal}
]
  \node[ch] (c1) {\textbf{第 1 章}\\无模型 RL：MC/TD\\SARSA 与 Q-Learning};
  \node[ch, right=of c1] (c2) {\textbf{第 2 章}\\DQN、策略梯度定理\\与 Actor-Critic/GAE};
  \node[ch, right=of c2] (c3) {\textbf{第 3 章}\\信任域 TRPO、\\PPO-Clip 与 GRPO};

  \node[ch, below=9mm of c1] (c4) {\textbf{第 4 章}\\大模型后训练：\\RLHF、DPO 与 RLVR};
  \node[ch, below=9mm of c2] (c5) {\textbf{第 5 章}\\GPU 内存墙、Roofline\\与 FlashAttention};
  \node[goal, below=9mm of c3] (c6) {\textbf{\color{ochre}第 6 章（系统闭环）}\\分布式并行、LoRA、\\量化与 PagedAttention};

  \draw[arr] (c1) -- (c2);
  \draw[arr] (c2) -- (c3);
  \draw[arr] ([xshift=-9mm]c3.south) -- ([xshift=9mm]c4.north);
  \draw[arr] (c4) -- (c5);
  \draw[arr] (c5) -- (c6);
\end{tikzpicture}}
\end{center}

\chapter[无模型强化学习：TD、SARSA 与 Q-Learning]{无模型强化学习基础：蒙特卡洛、TD(0)、SARSA 与 Q-Learning}
\label{ch:model-free-rl}

\chapEpigraph{与其试图编写一个模拟成人心智的程序，何不先编写一个模拟孩童心智的程序？如果再施以恰当的教育，我们就能得到成人的心智。}{Alan M. Turing（艾伦·图灵，1950）}

\begin{chapIntro}{本章导读：从“手握地图的规划”跃入“黑暗中的试错学习”}
在第 3 卷第 6 章中，我们已经完整建立了马尔可夫决策过程（MDP）的数学骨架，并借助价值迭代与策略迭代，在\textbf{完全已知环境转移概率 $P(s'\mid s,a)$ 与奖励函数 $R(s,a)$} 的前提下，证明了贝尔曼最优算子的压缩不动点存在且唯一。

然而，真实世界几乎从不把“物理引擎源码”交到智能体手上：机器人不知晓每个关节的精确摩擦系数，围棋程序面对的是对手难以穷举的落子分布，大语言模型更不可能拥有刻画“人类满意度”的闭式转移矩阵。本章正式开启\textbf{无模型强化学习（Model-Free RL）}的大门，聚焦三个层层递进的核心问题：
\begin{enumerate}
  \item \textbf{无模型评估}：当只能通过与环境交互采集轨迹样本时，蒙特卡洛（MC）与时序差分（TD(0)）如何分别用“完整回报”与“单步自举”估计状态价值？两者在偏差、方差与收敛条件上有何本质差异？
  \item \textbf{无模型控制与同/异策略分化}：在缺乏转移模型时为何必须转向动作价值 $Q(s,a)$？同策略的 \textbf{SARSA} 与异策略的 \textbf{Q-Learning} 仅在目标项相差一个 $\max$，为何会在悬崖行走中走出截然相反的轨迹？
  \item \textbf{最大化偏差与双估计解耦}：为何带噪声的 $\max$ 算子必然系统性高估价值？\textbf{Double Q-Learning} 如何用一次优雅的解耦消除高估偏差？
\end{enumerate}
\end{chapIntro}

\begin{recallBox}{第 3 卷第 6 章 MDP 与贝尔曼方程核心记号}
记有限马尔可夫决策过程为五元组 $\mathcal{M} = (\mathcal{S}, \mathcal{A}, P, R, \gamma)$，其中 $\gamma \in [0, 1]$ 为折扣因子。在策略 $\pi(a \mid s)$ 下，从时刻 $t$ 起累积的折扣回报定义为
\[
  G_t = \sum_{k=0}^{T-t-1} \gamma^k r_{t+k} = r_t + \gamma G_{t+1} \qquad (\text{持续任务取 } T = \infty).
\]
状态价值函数与动作价值函数分别满足\textbf{贝尔曼期望方程}：
\[
  V^\pi(s) = \E_\pi[G_t \mid S_t = s] = \sum_{a \in \mathcal{A}} \pi(a \mid s)\, Q^\pi(s, a), \qquad
  Q^\pi(s, a) = \E\!\left[r_t + \gamma V^\pi(S_{t+1}) \mid S_t = s, A_t = a\right],
\]
而最优动作价值函数 $Q^*(s, a) = \max_\pi Q^\pi(s, a)$ 满足\textbf{贝尔曼最优方程}：
\[
  Q^*(s, a) = \E\!\left[r_t + \gamma \max_{a' \in \mathcal{A}} Q^*(S_{t+1}, a') \mid S_t = s, A_t = a\right].
\]
\end{recallBox}

\section{从已知模型到无模型交互：样本回溯的诞生}
\label{sec:rl-model-free-gap}

在已知模型的动态规划（DP）中，策略评估的每一次扫描都执行\textbf{全宽期望回溯（Full Expected Backup）}：
\[
  V_{k+1}(s) \leftarrow \sum_{a \in \mathcal{A}} \pi(a \mid s) \sum_{s' \in \mathcal{S}} P(s' \mid s, a)\!\left[R(s, a, s') + \gamma V_k(s')\right].
\]
一旦拿掉环境模型 $(P, R)$，上式立即陷入双重停摆：
\begin{itemize}
  \item \textbf{评估受阻}：内层对所有后继状态 $s'$ 的加权求和需要确切知道转移核 $P(s' \mid s, a)$，而智能体只能把动作 $a$ 送入黑盒环境，观测到\textbf{唯一一个}随机落点 $(r_t, s_{t+1})$。
  \item \textbf{决策受阻}：即便天降神谕把精确的 $V^\pi(s)$ 全部写进表格，智能体仍无法执行一步前瞻的贪心策略改进 $\argmax_a \sum_{s'} P(s'\mid s,a)[R + \gamma V^\pi(s')]$，因为缺少 $P$ 就无法预知动作 $a$ 会把我们带到哪里。
\end{itemize}

无模型强化学习用一条朴素却深远的统计原则化解了第一个难题：\textbf{用经验样本回溯（Sample Backup）替代全宽期望回溯}。既然环境本身就是依照真实分布 $P(\cdot \mid s, a)$ 吐出 $(r_t, s_{t+1})$ 的完美采样器，我们便无需显式重构 $P$，只需沿着实际经历的轨迹反复向单次观测目标做小步长增量更新。图~\ref{fig:backup-diagrams} 直观对比了动态规划、蒙特卡洛与单步时序差分三种回溯结构在“采样宽度”与“回溯深度”两个维度上的根本分歧。

\begin{figure}[htbp]
\centering
\resizebox{0.96\linewidth}{!}{%
\begin{tikzpicture}[
  state/.style={circle, draw=tealDark, thick, fill=tealLight, minimum size=6.2mm, inner sep=0pt, font=\scriptsize\sffamily},
  act/.style={circle, draw=navy, fill=navy, minimum size=2.4mm, inner sep=0pt},
  term/.style={rectangle, draw=ochre, thick, fill=ochreLight, minimum size=5.2mm, inner sep=0pt, font=\scriptsize\sffamily},
  edge/.style={-{Stealth[length=1.7mm]}, thick, slate},
  hl/.style={-{Stealth[length=1.9mm]}, very thick, teal},
  title/.style={font=\small\sffamily\bfseries, color=tealDark}
]
  % Panel 1: DP Full Backup
  \begin{scope}[xshift=0cm]
    \node[title] at (0, 1.35) {(a) 动态规划（全宽单步）};
    \node[state] (s0) at (0, 0.55) {$s$};
    \node[act] (a1) at (-1.15, -0.35) {};
    \node[act] (a2) at (1.15, -0.35) {};
    \draw[edge] (s0) -- (a1) node[midway, left, font=\tiny, color=slate] {$a_1$};
    \draw[edge] (s0) -- (a2) node[midway, right, font=\tiny, color=slate] {$a_2$};
    \foreach \x/\lbl in {-1.75/$s'_1$, -0.55/$s'_2$} {
      \node[state] (n) at (\x, -1.45) {\lbl};
      \draw[edge] (a1) -- (n);
    }
    \foreach \x/\lbl in {0.55/$s'_3$, 1.75/$s'_4$} {
      \node[state] (n) at (\x, -1.45) {\lbl};
      \draw[edge] (a2) -- (n);
    }
    \node[font=\scriptsize, color=slate, align=center] at (0, -2.25) {需遍历全部 $a, s'$（依赖已知 $P$）\\自举一步 $V(s')$};
  \end{scope}

  % Panel 2: TD(0) Sample Backup
  \begin{scope}[xshift=5.6cm]
    \node[title] at (0, 1.35) {(b) 时序差分 TD(0)（单样本单步）};
    \node[state] (t0) at (0, 0.55) {$s_t$};
    \node[act] (ta) at (0, -0.35) {};
    \node[state] (t1) at (0, -1.45) {$s_{t+1}$};
    \draw[hl] (t0) -- (ta) node[midway, right, font=\scriptsize, color=teal] {$a_t \sim \pi$};
    \draw[hl] (ta) -- (t1) node[midway, right, font=\scriptsize, color=ochre] {$r_t$};
    \node[font=\scriptsize, color=slate, align=center] at (0, -2.25) {仅采一条单步转移（无模型）\\用当前估计 $V(s_{t+1})$ 自举};
  \end{scope}

  % Panel 3: MC Sample Backup
  \begin{scope}[xshift=11.4cm]
    \node[title] at (0, 1.35) {(c) 蒙特卡洛 MC（单样本全轨迹）};
    \node[state] (m0) at (0, 0.55) {$s_t$};
    \node[act] (ma0) at (0, -0.05) {};
    \node[state] (m1) at (0, -0.75) {$s_{t+1}$};
    \node[font=\scriptsize, color=slate] (dots) at (0, -1.56) {$\vdots$};
    \node[term] (mt) at (0, -2.36) {$T$};
    \draw[hl] (m0) -- (ma0);
    \draw[hl] (ma0) -- (m1) node[midway, right, font=\scriptsize, color=ochre] {$r_t$};
    \draw[edge] (m1.south) -- (dots.north);
    \draw[edge] (dots.south) -- (mt.north);
    \node[font=\scriptsize, color=slate, align=center] at (0, -3.10) {采完整回合至终止（无模型）\\无自举，直接用真实回报 $G_t$};
  \end{scope}
\end{tikzpicture}}
\caption{三种价值更新机制的回溯图对比。空心圆代表状态，实心小圆代表动作，方形代表回合终止。强化学习的方法光谱本质上是在“是否采样（宽 vs 窄）”与“是否自举（浅 vs 深）”所张成的二维平面上权衡偏差、方差与计算代价。}
\label{fig:backup-diagrams}
\end{figure}

\section{蒙特卡洛策略评估：无偏但高方差的基准}
\label{sec:rl-mc}

\subsection{首次访问与每次访问蒙特卡洛}

面对一个未知模型但可反复重启的\textbf{回合制（Episodic）}环境，估算固定策略 $\pi$ 的状态价值 $V^\pi(s) = \E_\pi[G_t \mid S_t = s]$ 最直接的办法便是依照大数定律：反复运行策略 $\pi$ 收集大量完整回合 $\tau^{(1)}, \dots, \tau^{(M)}$，统计每次经过状态 $s$ 之后实际收获的折扣回报 $G_t$，然后取经验均值。这一思路称为\textbf{蒙特卡洛（Monte Carlo, MC）策略评估}。

如果同一个状态 $s$ 在单条轨迹中被多次访问，存在两种自然变体：
\begin{defBox}
\begin{definition}[首次访问与每次访问 MC 评估]
给定由策略 $\pi$ 独立生成的 $M$ 条完整回合，对任意状态 $s \in \mathcal{S}$：
\begin{itemize}
  \item \textbf{首次访问 MC（First-Visit MC）}：在每条回合中，仅取状态 $s$ \emph{第一次}出现时的时刻 $t$ 对应的回报 $G_t$ 计入样本集 $\mathcal{G}_{\mathrm{first}}(s)$，估计值 $\hat{V}_{\mathrm{FV}}(s) = \frac{1}{|\mathcal{G}_{\mathrm{first}}(s)|}\sum_{G \in \mathcal{G}_{\mathrm{first}}(s)} G$。
  \item \textbf{每次访问 MC（Every-Visit MC）}：将回合中状态 $s$ \emph{每一次}出现对应的回报 $G_t$ 全部计入样本集 $\mathcal{G}_{\mathrm{every}}(s)$，取其算术平均 $\hat{V}_{\mathrm{EV}}(s)$。
\end{itemize}
若按增量形式在线更新，设第 $k$ 次采样到状态 $s$ 的回报为 $G^{(k)}$，则算术平均等价于增量式
\begin{equation}
  V(s) \leftarrow V(s) + \alpha_k \!\left(G^{(k)} - V(s)\right),
  \label{eq:mc-inc}
\end{equation}
其中取 $\alpha_k = 1/k$ 给出精确经验均值；在非平稳环境或后续控制任务（策略 $\pi$ 随时间演化）中，常改用固定常数步长 $\alpha \in (0, 1)$，使旧回报权重按 $(1-\alpha)^m$ 指数遗忘。
\end{definition}
\end{defBox}

\begin{thmBox}
\begin{theorem}[首次访问 MC 的无偏性与均方一致性]
\label{thm:mc-unbiased}
设回合间相互独立，且在策略 $\pi$ 下状态 $s$ 的折扣回报方差有限：$\Var_\pi(G_t \mid S_t = s) = \sigma^2(s) < \infty$。记 $N_M(s)$ 为前 $M$ 条回合中至少访问过一次 $s$ 的回合数，在条件 $N_M(s) = n \ge 1$ 下，首次访问 MC 估计 $\hat{V}_n(s) = \frac{1}{n}\sum_{j=1}^n G^{(j)}$ 满足：
\[
  \E\!\left[\hat{V}_n(s)\right] = V^\pi(s), \qquad
  \Var\!\left(\hat{V}_n(s)\right) = \frac{\sigma^2(s)}{n} \xrightarrow{n \to \infty} 0.
\]
\end{theorem}
\begin{proof}
在每条访问过状态 $s$ 的回合 $j$ 中，设 $T_j = \inf\{t \ge 0 : S_t^{(j)} = s\}$ 为首次击中时刻。由于 $T_j$ 是关于轨迹自然滤子 $\mathcal{F}_t = \sigma(S_0, A_0, \dots, S_t)$ 的停时，由马尔可夫过程的强马尔可夫性，条件于 $\{T_j < \infty\}$ 时，从 $T_j$ 起截取的后缀轨迹与“直接从 $S_0 = s$ 出发”的轨迹同分布，因此
\[
  \E\!\left[G^{(j)} \mid T_j < \infty\right] = \E_\pi\!\left[G_0 \mid S_0 = s\right] = V^\pi(s), \qquad
  \Var\!\left(G^{(j)} \mid T_j < \infty\right) = \sigma^2(s).
\]
又因不同回合的随机轨迹彼此独立，故 $G^{(1)}, \dots, G^{(n)}$ 为独立同分布（i.i.d.）随机变量。由期望的线性性与独立和方差的可加性立得
\[
  \E\!\left[\hat{V}_n(s)\right] = \frac{1}{n}\sum_{j=1}^n V^\pi(s) = V^\pi(s), \qquad
  \Var\!\left(\hat{V}_n(s)\right) = \frac{1}{n^2}\sum_{j=1}^n \sigma^2(s) = \frac{\sigma^2(s)}{n}.
\]
当 $n \to \infty$ 时，均方误差 $\mathrm{MSE}(\hat{V}_n(s)) = \Var(\hat{V}_n(s)) \to 0$，且由强大数定律知 $\hat{V}_n(s) \xrightarrow{\text{a.s.}} V^\pi(s)$。
\end{proof}
\end{thmBox}

\begin{remark}[每次访问 MC 为何在有限样本下可能有偏？]
在同一回合中，若轨迹按 $s \to s \to \text{终止}$ 连续访问两次 $s$，第一次访问的回报 $G_0 = r_0 + \gamma G_1$ 与第二次访问的回报 $G_1$ 不仅高度相关，而且每个回合贡献的样本数 $K_j$ 本身就是与回报相关的随机变量！由于 $\E\big[\frac{\sum_{j=1}^M \sum_{i=1}^{K_j} G_{j,i}}{\sum_{j=1}^M K_j}\big] \neq \frac{\E[\sum_i G_{1,i}]}{\E[K_1]}$，每次访问 MC 在有限样本下通常存在轻微的比率偏差（Ratio Bias）。不过 Singh 与 Sutton（1996）证明了它仍是渐近一致的，且均方误差同样以 $O(1/M)$ 收敛；在实现上甚至省去了记录“本回合是否已访问过 $s$”的集合开销。
\end{remark}

\begin{example}[三状态带环小链的 MC 手算对比]
考虑含终止态 $\mathsf{T}$ 的微型 MDP，折扣因子 $\gamma = 1$。策略 $\pi$ 产生了两条回合：
\begin{itemize}
  \item 回合 1：$s_A \xrightarrow{r=0} s_A \xrightarrow{r=2} s_B \xrightarrow{r=4} \mathsf{T}$。对应各步回报：$G_0 = 6$（$s_A$ 首次），$G_1 = 6$（$s_A$ 二次），$G_2 = 4$（$s_B$）。
  \item 回合 2：$s_A \xrightarrow{r=1} s_B \xrightarrow{r=0} \mathsf{T}$。对应各步回报：$G_0 = 1$（$s_A$ 首次），$G_1 = 0$（$s_B$）。
\end{itemize}
按\textbf{首次访问 MC}：$\hat{V}_{\mathrm{FV}}(s_A) = \frac{6 + 1}{2} = 3.5$，$\hat{V}_{\mathrm{FV}}(s_B) = \frac{4 + 0}{2} = 2.0$。按\textbf{每次访问 MC}：$\hat{V}_{\mathrm{EV}}(s_A) = \frac{6 + 6 + 1}{3} \approx 4.33$，$\hat{V}_{\mathrm{EV}}(s_B) = 2.0$。注意在回合 1 中，$s_A$ 的两次回报完全耦合在同一条好运轨迹上，使每次访问估计明显偏向更高分。
\end{example}

尽管首次访问 MC 拥有无可挑剔的无偏性，它在工程实践中却暴露出两大致命短板：
\begin{enumerate}
  \item \textbf{必须苦等回合结束}：只有走到终止状态才能反推 $G_t$，因此完全无法直接用于永不终止的持续任务（如机器人平衡、高频交易），在长视野任务中学习反馈也极其迟缓。
  \item \textbf{方差沿时间步雪崩累积}：$G_t = r_t + \gamma r_{t+1} + \dots + \gamma^{T-t-1} r_{T-1}$ 叠加了未来所有随机动作与随机环境转移的噪声；即便 $s_t$ 之后的某一步只是碰巧遭遇了一次罕见暴击，也会连带污染前面几十个状态的梯度与价值估计。
\end{enumerate}

\section{时序差分学习 TD(0) 与随机逼近理论}
\label{sec:rl-td0}

\subsection{TD(0) 更新律与自举思想}

Sutton（1988）提出的\textbf{时序差分（Temporal-Difference, TD）学习}被誉为强化学习领域最具原创性的核心思想——它把动态规划的\textbf{自举（Bootstrapping）}与蒙特卡洛的\textbf{样本驱动}完美焊接在一起。

观察贝尔曼期望方程 $V^\pi(s_t) = \E_\pi\!\left[r_t + \gamma V^\pi(s_{t+1}) \mid s_t\right]$。为了在仅仅走完\textbf{一步} $(s_t, a_t, r_t, s_{t+1})$ 后就立即更新 $V(s_t)$，我们不必等到回合结束去获取真实回报 $G_t$，只需用当前手头的价值估计 $V(s_{t+1})$ 顶替未知的剩余期望回报 $\E[G_{t+1} \mid s_{t+1}]$：

\begin{defBox}
\begin{definition}[TD(0) 目标、TD 误差与更新规则]
在每观测到单步转移 $(s_t, a_t, r_t, s_{t+1})$ 后，\textbf{单步时序差分算法 TD(0)} 定义\textbf{TD 目标（TD Target）}与\textbf{TD 误差（TD Error）}如下：
\begin{equation}
  y_t^{\mathrm{TD}} = r_t + \gamma V(s_{t+1}), \qquad
  \delta_t = y_t^{\mathrm{TD}} - V(s_t) = r_t + \gamma V(s_{t+1}) - V(s_t),
  \label{eq:td-error}
\end{equation}
其中若 $s_{t+1}$ 为终止状态，严格规定 $V(s_{t+1}) \equiv 0$。随后立即执行增量更新：
\begin{equation}
  V(s_t) \leftarrow V(s_t) + \alpha_t \,\delta_t = V(s_t) + \alpha_t \!\left[r_t + \gamma V(s_{t+1}) - V(s_t)\right].
  \label{eq:td0-update}
\end{equation}
\end{definition}
\end{defBox}

对比式~\eqref{eq:mc-inc} 与式~\eqref{eq:td0-update}，两者形式如出一辙，唯一差别是用“一步真实奖励 $+$ 折扣后继估计” $r_t + \gamma V(s_{t+1})$ 替换了“完整轨迹真实回报” $G_t$。这种“用估计去更新估计”的机制正是\textbf{自举（Bootstrapping）}。有趣的是，TD 误差 $\delta_t$ 不仅支配算法收敛，神经科学研究（Schultz, Dayan \& Montague, 1997）还发现灵长类动物中脑多巴胺神经元的脉冲发放模式，竟与“意外奖励预测误差” $\delta_t$ 高度吻合。

\subsection{偏差--方差权衡的严格刻画}

“用尚未准确的 $V(s_{t+1})$ 充当教学目标，难道不会越学越歪吗？”——这是初学者面对 TD(0) 时最敏锐的质疑。下面的命题精确量化了 MC 与 TD(0) 在偏差与方差上的消长关系。

\begin{thmBox}
\begin{proposition}[MC 目标与 TD(0) 目标的偏差--方差比较]
\label{prop:mc-td-bv}
设当前价值估计表为确定性函数 $V: \mathcal{S} \to \R$，策略 $\pi$ 固定，从状态 $S_t = s$ 出发执行一步转移 $(A_t, R_t, S_{t+1})$ 并继续演化得到完整回报 $G_t = R_t + \gamma G_{t+1}$。则：
\begin{enumerate}
  \item \textbf{偏差对比}：MC 目标 $G_t$ 严格无偏，即 $\E_\pi[G_t \mid S_t = s] - V^\pi(s) = 0$；而 TD(0) 目标 $y_t^{\mathrm{TD}} = R_t + \gamma V(S_{t+1})$ 的条件偏差为
  \begin{equation}
    \operatorname{Bias}\!\left(y_t^{\mathrm{TD}} \mid S_t = s\right)
    = \E_\pi\!\left[R_t + \gamma V(S_{t+1}) \mid S_t = s\right] - V^\pi(s)
    = \gamma\, \E_\pi\!\left[V(S_{t+1}) - V^\pi(S_{t+1}) \mid S_t = s\right].
    \label{eq:td-bias}
  \end{equation}
  只有当后继状态的价值估计无误（$V = V^\pi$）时，TD 目标的偏差才归零。
  \item \textbf{方差分解}：当 $V$ 接近真实价值 $V^\pi$ 时，理想单步目标 $y_t^* = R_t + \gamma V^\pi(S_{t+1})$ 与完整回报 $G_t$ 的条件方差满足精确恒等式：
  \begin{equation}
    \Var_\pi\!\left(G_t \mid S_t = s\right)
    = \Var_\pi\!\left(R_t + \gamma V^\pi(S_{t+1}) \mid S_t = s\right)
    + \gamma^2\, \E_\pi\!\left[\Var_\pi(G_{t+1} \mid S_{t+1}) \mid S_t = s\right].
    \label{eq:mc-td-var}
  \end{equation}
  因第二项恒非负，只要后续轨迹存在随机性，单步 TD 目标的方差便严格小于 MC 回报的方差。
\end{enumerate}
\end{proposition}
\begin{proof}
第一部分直接由贝尔曼期望方程 $V^\pi(s) = \E_\pi[R_t + \gamma V^\pi(S_{t+1}) \mid S_t = s]$ 相减即得：
\[
  \E_\pi\!\left[R_t + \gamma V(S_{t+1}) \mid S_t = s\right] - V^\pi(s)
  = \E_\pi\!\left[\gamma\big(V(S_{t+1}) - V^\pi(S_{t+1})\big) \mid S_t = s\right].
\]
第二部分对随机变量 $G_t = R_t + \gamma G_{t+1}$ 应用全方差公式（条件于 $\mathcal{F}_{t,1} = (S_t, A_t, R_t, S_{t+1})$，在给定 $S_t = s$ 下）：
\begin{align*}
  \Var_\pi(G_t \mid S_t = s)
  &= \Var_\pi\!\Big(\E_\pi[G_t \mid S_t = s, A_t, R_t, S_{t+1}] \;\Big|\; S_t = s\Big) \\
  &\quad + \E_\pi\!\Big[\Var_\pi(G_t \mid S_t = s, A_t, R_t, S_{t+1}) \;\Big|\; S_t = s\Big].
\end{align*}
利用马尔可夫性，条件于 $(S_t, A_t, R_t, S_{t+1})$ 时 $R_t$ 为常量，且 $G_{t+1}$ 的条件分布仅依赖于 $S_{t+1}$，故
\[
  \E_\pi[R_t + \gamma G_{t+1} \mid S_t = s, A_t, R_t, S_{t+1}] = R_t + \gamma V^\pi(S_{t+1}),
\]
\[
  \Var_\pi(R_t + \gamma G_{t+1} \mid S_t = s, A_t, R_t, S_{t+1}) = \gamma^2 \Var_\pi(G_{t+1} \mid S_{t+1}).
\]
代回全方差公式，即证式~\eqref{eq:mc-td-var}。
\end{proof}
\end{thmBox}

式~\eqref{eq:mc-td-var} 清楚地揭示了 MC 与 TD(0) 的分工：MC 忍受了后续全部步长的方差总和 $\gamma^2 \E[\Var(G_{t+1}\mid S_{t+1})]$ 以换取绝对零偏；TD(0) 则在训练初期承受由初值不准带来的自举偏差，但因为每一步只引入\textbf{单次转移}的随机性，方差大幅锐减。随着迭代推进，$V$ 向 $V^\pi$ 逼近，由式~\eqref{eq:td-bias} 可见 TD 偏差也在以几何级数自我消融。

\begin{example}[Sutton 经典双状态批处理反例：批量 MC 与批量 TD(0) 的本质分歧]
\label{ex:batch-mc-td}
设一个未知的马尔可夫过程包含两个非终止状态 $A, B$，折扣因子 $\gamma = 1$。我们只观测到如下 8 条完整回合：
\begin{center}
\small
回合 1：$A \xrightarrow{0} B \xrightarrow{0} \mathsf{T}$；\qquad
回合 2--7（共 6 条）：$B \xrightarrow{1} \mathsf{T}$；\qquad
回合 8：$B \xrightarrow{0} \mathsf{T}$。
\end{center}
假设我们把这 8 条数据当作固定训练集，反复执行批量更新直至收敛，问 $V(A)$ 与 $V(B)$ 分别应等于多少？
\begin{itemize}
  \item \textbf{批量 MC 的答案}：状态 $B$ 共出现 8 次，其中 6 次回报为 $1$、2 次为 $0$，故 $V_{\mathrm{MC}}(B) = \frac{6}{8} = \mathbf{0.75}$。状态 $A$ 仅在回合 1 出现 1 次，其完整回报为 $0 + 0 = 0$，故 $V_{\mathrm{MC}}(A) = \mathbf{0}$。批量 MC 最小化的是\emph{训练集观测回报的均方误差}。
  \item \textbf{批量 TD(0) 的答案}：从数据拟合出最大似然经验 MDP——从 $A$ 出发以概率 $1$ 零奖励转移到 $B$（$\hat{P}(B\mid A) = 1$），从 $B$ 出发有 $6/8 = 0.75$ 的概率获奖励 $1$ 终止、$2/8$ 获奖励 $0$ 终止。解该经验 MDP 的贝尔曼方程，得 $V_{\mathrm{TD}}(B) = \mathbf{0.75}$，且
  \[
    V_{\mathrm{TD}}(A) = 0 + 1 \cdot V_{\mathrm{TD}}(B) = \mathbf{0.75}.
  \]
\end{itemize}
当真实环境满足马尔可夫性质时，$A$ 之所以在回合 1 拿了 0 分，纯粹是因为到了 $B$ 之后倒霉抽中了 $1/4$ 的坏结局，而并非 $A$ 本身毫无价值。TD(0) 利用马尔可夫结构将 $B$ 上积累的 8 次统计证据“无缝传导”给了仅被观测 1 次的 $A$，因此在马尔可夫环境中往往比 MC 具有高得多的数据效率！
\end{example}

\subsection{Robbins--Monro 随机逼近条件与 TD(0) 收敛性}

无论是式~\eqref{eq:mc-inc} 还是式~\eqref{eq:td0-update}，其本质都是\textbf{在带随机噪声的观测中求解不动点方程 $F(v) = v$}。步长序列 $\{\alpha_t\}_{t=0}^\infty$ 该如何设计，才能保证迭代既不会“半途停摆”，又不会被噪声“永远晃晕”？1951 年 Robbins 与 Monro 给出了支配整个随机优化与强化学习收敛理论的黄金法则。

\begin{thmBox}
\begin{theorem}[Robbins--Monro 随机逼近条件]
\label{thm:robbins-monro}
考虑标量随机迭代 $\theta_{t+1} = \theta_t + \alpha_t (\mu + \xi_t - \theta_t)$，其中 $\mu \in \R$ 为未知真值，噪声序列 $\{\xi_t\}$ 独立同分布，满足 $\E[\xi_t] = 0$ 且 $0 < \Var(\xi_t) = \sigma^2 < \infty$，步长 $\alpha_t \in (0, 1)$。则 $\E[(\theta_t - \mu)^2] \to 0$ 的充分必要条件是步长序列同时满足
\begin{equation}
  \sum_{t=0}^{\infty} \alpha_t = \infty \qquad \text{且} \qquad \sum_{t=0}^{\infty} \alpha_t^2 < \infty.
  \label{eq:rm-cond}
\end{equation}
典型例子为 $\alpha_t = \frac{c}{(t + t_0)^p}$，其中 $\frac{1}{2} < p \le 1$（例如 $\alpha_t = \frac{1}{t+1}$ 满足调和级数发散而平方级数收敛于 $\pi^2/6$）。
\end{theorem}
\begin{proof}
令误差 $e_t = \theta_t - \mu$，则递推化为 $e_{t+1} = (1 - \alpha_t) e_t + \alpha_t \xi_t$。由于 $\xi_t$ 独立于 $e_t$ 且零均值，两边取平方再取期望得
\begin{equation}
  \E[e_{t+1}^2] = (1 - \alpha_t)^2 \E[e_t^2] + \alpha_t^2 \sigma^2.
  \label{eq:rm-mse-rec}
\end{equation}
\textbf{必要性}：首先，若 $\sum_{t=0}^\infty \alpha_t < \infty$，则由不等式 $1 - x \ge e^{-2x}$（对充分小的 $x > 0$）知 $\prod_{t=t_0}^\infty (1 - \alpha_t)^2 > 0$，于是初值误差的贡献 $e_0^2 \prod_{k=0}^{t-1}(1-\alpha_k)^2$ 无法衰减至 $0$（当 $e_0 \neq 0$ 时 $\E[e_t^2]$ 有严格正下界）。其次，若 $\alpha_t \not\to 0$，由式~\eqref{eq:rm-mse-rec} 知 $\E[e_{t+1}^2] \ge \alpha_t^2 \sigma^2 \not\to 0$。（一般情形下，将式~\eqref{eq:rm-mse-rec} 展开后，只有当 $\sum \alpha_t^2 < \infty$ 才能使累积噪声项 $\sigma^2 \sum_{k=0}^{t-1} \alpha_k^2 \prod_{j=k+1}^{t-1}(1-\alpha_j)^2$ 趋于 $0$——更严格地，对标量线性递推只需 $\alpha_t \to 0$ 与 $\sum \alpha_t = \infty$ 即可均方收敛，而几乎必然收敛 $\theta_t \xrightarrow{\text{a.s.}} \mu$ 以及一般非线性随机逼近的均方/几乎必然收敛需要更强的平方可和条件 $\sum \alpha_t^2 < \infty$。）

\textbf{充分性}：设条件~\eqref{eq:rm-cond} 成立。由 $\sum \alpha_t^2 < \infty$ 知 $\alpha_t \to 0$，故对任意给定 $\varepsilon > 0$，存在 $N$ 使得对一切 $t \ge N$ 有 $\sum_{k=N}^{\infty} \alpha_k^2 < \frac{\varepsilon}{2\sigma^2}$。对 $t > N$ 反复展开式~\eqref{eq:rm-mse-rec}，并利用 $1 - \alpha \le e^{-\alpha}$ 与 $(1-\alpha_j)^2 \le 1$：
\[
  \E[e_t^2]
  = \E[e_N^2] \prod_{k=N}^{t-1} (1 - \alpha_k)^2
  + \sigma^2 \sum_{k=N}^{t-1} \alpha_k^2 \prod_{j=k+1}^{t-1}(1 - \alpha_j)^2
  \le \E[e_N^2] \exp\!\left(-2\sum_{k=N}^{t-1}\alpha_k\right)
  + \sigma^2 \sum_{k=N}^{\infty} \alpha_k^2.
\]
由于 $\sum_{k=N}^\infty \alpha_k = \infty$，第一项当 $t \to \infty$ 时趋于 $0$，故存在 $T > N$ 使得对一切 $t \ge T$ 第一项小于 $\varepsilon/2$，从而 $\E[e_t^2] < \varepsilon$。因此 $\lim_{t \to \infty} \E[(\theta_t - \mu)^2] = 0$。
\end{proof}
\end{thmBox}

直观记忆这两条条件堪称毫不费力：\textbf{$\sum \alpha_t = \infty$ 保证步子总长足够大，无论初始猜测离真值多远都能走得到；$\sum \alpha_t^2 < \infty$ 保证步子最终足够小，让每步观测噪声的方差累加起来依然有限，从而稳稳停在真值处}。将此框架推广到贝尔曼算子 $\mathcal{T}^\pi$（它在 $\ell_\infty$ 范数下是 $\gamma$-压缩映射），Jaakkola、Jordan 与 Singh（1994）以及 Tsitsiklis（1994）证明：在表格型表示下，只要每个状态被无限次访问且步长序列 $\alpha_t(s)$ 对每个 $s$ 均满足 Robbins--Monro 条件，TD(0) 便以概率 1 收敛至精确的状态价值函数 $V^\pi$。

\section{无模型控制：动作价值函数与 $\varepsilon$-贪心探索}
\label{sec:rl-control-q}

学会评估给定策略 $\pi$ 后，我们的终极目标是\textbf{控制（Control）}——在未知环境中找到最优策略 $\pi^*$。回顾广义策略迭代（GPI）的两大交替步骤：策略评估与策略改进。在无模型设定下，两道新障碍随之浮现：

\begin{enumerate}
  \item \textbf{为什么 $V(s)$ 不够用，必须学习 $Q(s, a)$？} 若只知道状态价值 $V^\pi(s)$，贪心改进需要计算
  \[
    \pi'(s) = \argmax_{a \in \mathcal{A}} \sum_{s' \in \mathcal{S}} P(s' \mid s, a)\!\left[R(s, a, s') + \gamma V^\pi(s')\right],
  \]
  这再次要求已知模型 $(P, R)$。但如果我们直接缓存并学习\textbf{动作价值函数 $Q(s, a)$}，策略改进就退化为一次纯局部、完全无模型的查表取最大值操作：
  \[
    \pi'(s) = \argmax_{a \in \mathcal{A}} Q(s, a).
  \]
  \item \textbf{为什么纯粹贪心会困死在局部次优解？} 如果智能体从第一步起就死板地执行 $\argmax_a Q(s,a)$，那么任何从未尝试过、初值为 $0$ 的动作（哪怕真实奖励高达 $+100$），只要当前最优动作的估计值大于 $0$，就将\textbf{永生不被探索}！这正是强化学习最核心的\textbf{探索与利用权衡（Exploration--Exploitation Trade-off）}。
\end{enumerate}

化解这一矛盾最简洁且久经考验的方案是 \textbf{$\varepsilon$-贪心（$\varepsilon$-Greedy）策略}：对给定动作价值表 $Q$ 和探索率 $\varepsilon \in (0, 1)$，记动作空间大小 $m = |\mathcal{A}|$，则
\begin{equation}
  \pi_\varepsilon(a \mid s) =
  \begin{cases}
    1 - \varepsilon + \dfrac{\varepsilon}{m}, & \text{若 } a = a^*(s) := \argmax_{a'} Q(s, a'), \\[6pt]
    \dfrac{\varepsilon}{m}, & \text{否则}.
  \end{cases}
  \label{eq:eps-greedy}
\end{equation}
凡是对所有 $(s, a)$ 均满足 $\pi(a \mid s) \ge \varepsilon / m$ 的策略，统称为\textbf{$\varepsilon$-软策略（$\varepsilon$-Soft Policy）}。在继续之前，我们必须严格验证一个关键问题：强制保留 $\varepsilon$ 比例的随机探索后，策略改进定理是否依然成立？

\begin{thmBox}
\begin{theorem}[$\varepsilon$-软策略改进定理]
\label{thm:eps-soft-imp}
设 $\pi$ 为任意 $\varepsilon$-软策略，$\pi'$ 是关于 $Q^\pi$ 构造的 $\varepsilon$-贪心策略（即式~\eqref{eq:eps-greedy} 中取 $a^*(s) = \argmax_a Q^\pi(s, a)$）。则对一切状态 $s \in \mathcal{S}$，均有
\[
  V^{\pi'}(s) \ge V^\pi(s).
\]
\end{theorem}
\begin{proof}
对任意状态 $s \in \mathcal{S}$，直接计算在状态 $s$ 改用 $\pi'$ 一步、此后仍遵循 $\pi$ 的期望价值：
\begin{align*}
  Q^\pi(s, \pi'(s))
  &= \sum_{a \in \mathcal{A}} \pi'(a \mid s)\, Q^\pi(s, a) \\
  &= \frac{\varepsilon}{m} \sum_{a \in \mathcal{A}} Q^\pi(s, a) + (1 - \varepsilon) \max_{a \in \mathcal{A}} Q^\pi(s, a).
\end{align*}
因为 $\pi$ 是 $\varepsilon$-软策略，故对每个动作 $a$ 均有 $w(a) := \frac{\pi(a \mid s) - \varepsilon/m}{1 - \varepsilon} \ge 0$，且 $\sum_{a \in \mathcal{A}} w(a) = \frac{1 - \varepsilon}{1 - \varepsilon} = 1$。任何凸组合都不可能超过最大值，于是
\[
  \max_{a \in \mathcal{A}} Q^\pi(s, a) \ge \sum_{a \in \mathcal{A}} w(a)\, Q^\pi(s, a)
  = \frac{1}{1 - \varepsilon}\!\left(\sum_{a \in \mathcal{A}} \pi(a \mid s) Q^\pi(s, a) - \frac{\varepsilon}{m}\sum_{a \in \mathcal{A}} Q^\pi(s, a)\right).
\]
将此不等式代回上式，恰好消去 $\frac{\varepsilon}{m}\sum_a Q^\pi(s, a)$，立即得到
\[
  Q^\pi(s, \pi'(s)) \ge \sum_{a \in \mathcal{A}} \pi(a \mid s)\, Q^\pi(s, a) = V^\pi(s).
\]
随后与经典策略改进定理完全相同地沿贝尔曼期望算子 $\mathcal{T}^{\pi'}$ 递推展开：
\[
  V^\pi(s) \le \E_{\pi'}\!\left[r_t + \gamma V^\pi(S_{t+1}) \mid S_t = s\right]
  \le \E_{\pi'}\!\left[r_t + \gamma r_{t+1} + \gamma^2 V^\pi(S_{t+2}) \mid S_t = s\right]
  \le \cdots \le V^{\pi'}(s). \qedhere
\]
\end{proof}
\end{thmBox}

然而，若 $\varepsilon$ 永远保持固定正数，策略改进只会收敛到“所有 $\varepsilon$-软策略中的最优者”，而非真正的全局最优确定性策略 $\pi^*$。要兼顾“充分探索”与“最终纯粹最优”，策略序列必须满足 \textbf{GLIE（Greedy in the Limit with Infinite Exploration，无限探索下的极限贪心）} 条件：
\begin{enumerate}
  \item 所有状态--动作对都被访问无限多次：$\lim_{k \to \infty} N_k(s, a) = \infty$；
  \item 策略在极限意义下收敛至贪心策略：$\lim_{k \to \infty} \pi_k(a \mid s) = \mathbf{1}\!\left[a = \argmax_{a'} Q_k(s, a')\right]$。
\end{enumerate}
最常见的 GLIE 退火方案是在第 $k$ 个回合令 $\varepsilon_k = 1/k$（因为 $\sum_{k=1}^\infty 1/k = \infty$ 保证每个动作以概率 1 被采样无限次，而 $\lim_{k\to\infty} \varepsilon_k = 0$ 保证最终归于纯粹贪心）。

\section{同策略 SARSA 与异策略 Q-Learning}
\label{sec:rl-sarsa-q}

现在我们将 TD(0) 的单步自举思想直接搬到动作价值函数 $Q(s, a)$ 上。一个极其微妙的分岔口立刻出现：在当前状态 $s_t$ 执行动作 $a_t$、拿到奖励 $r_t$ 并落入 $s_{t+1}$ 之后，\textbf{构造自举目标时，后继状态 $s_{t+1}$ 的动作究竟应该取“我们接下来实际要做的探索动作 $a_{t+1}$”，还是取“假设自己将来完全理智时的最优动作 $\max_{a'} Q(s_{t+1}, a')$”？}

这两个看似微小的选择，分别催生了强化学习史上最著名的双子星——\textbf{同策略（On-Policy）的 SARSA} 与 \textbf{异策略（Off-Policy）的 Q-Learning}。

\begin{defBox}
\begin{definition}[SARSA 与 Q-Learning 更新规则]
设智能体按行为策略 $\mu$（通常取当前 $Q$ 表的 $\varepsilon$-贪心策略 $\pi_\varepsilon$）与环境交互：
\begin{itemize}
  \item \textbf{SARSA（同策略 TD 控制）}：在到达 $s_{t+1}$ 后，先由同一策略 $\pi_\varepsilon(\cdot \mid s_{t+1})$ 采样出下一步实际要执行的动作 $a_{t+1}$，利用五元组 $(s_t, a_t, r_t, s_{t+1}, a_{t+1})$ 更新：
  \begin{equation}
    Q(s_t, a_t) \leftarrow Q(s_t, a_t) + \alpha_t \!\left[r_t + \gamma Q(s_{t+1}, a_{t+1}) - Q(s_t, a_t)\right].
    \label{eq:sarsa}
  \end{equation}
  算法名正是五元组首字母 \texttt{S-A-R-S-A} 的缩写。
  \item \textbf{Q-Learning（异策略 TD 控制，Watkins, 1989）}：无需预先采样 $a_{t+1}$，直接用后继状态的贪心最大值作为目标更新：
  \begin{equation}
    Q(s_t, a_t) \leftarrow Q(s_t, a_t) + \alpha_t \!\left[r_t + \gamma \max_{a' \in \mathcal{A}} Q(s_{t+1}, a') - Q(s_t, a_t)\right].
    \label{eq:qlearning}
  \end{equation}
  随后智能体仍按 $\pi_\varepsilon(\cdot \mid s_{t+1})$ 与环境继续交互。
\end{itemize}
对两种算法，若 $s_{t+1}$ 为终止状态，自举项（$\gamma Q$ 或 $\gamma \max Q$）均严格置零。
\end{definition}
\end{defBox}

\begin{insightBox}{同策略 vs 异策略的数学实质：逼近的究竟是哪个贝尔曼算子？}
\begin{itemize}
  \item \textbf{SARSA（同策略）}：产生数据的行为策略与被评估的目标策略是\textbf{同一个} $\pi_\varepsilon$。在固定 $\varepsilon > 0$ 下，SARSA 逼近的是 $\pi_\varepsilon$ 的\textbf{贝尔曼期望不动点 $Q^{\pi_\varepsilon}$}。这意味着 SARSA 拥有“自知之明”——它在评估每一步的好坏时，已经把“我下一步有 $\varepsilon$ 的概率手滑乱走”这一风险老老实实地计入了账本！
  \item \textbf{Q-Learning（异策略）}：行为策略是带探索的 $\pi_\varepsilon$，而被学习的目标策略是纯粹贪心的 $\pi_{\mathrm{greedy}}(s) = \argmax_a Q(s, a)$。无论行为策略多么随机（只要覆盖所有状态--动作对），式~\eqref{eq:qlearning} 的期望更新方向始终直指\textbf{贝尔曼最优算子} $(\mathcal{T}^* Q)(s, a) = \E[r + \gamma \max_{a'} Q(s', a')]$！因此 Watkins 与 Dayan（1992）证明：只要每个 $(s, a)$ 被访问无限次且步长满足 Robbins--Monro 条件，\textbf{哪怕行为策略的 $\varepsilon$ 永远不衰减，Q-Learning 依然以概率 1 收敛到精确的最优动作价值 $Q^*$}。
\end{itemize}
\end{insightBox}

为了把两者的差异从公式落到直观数字上，不妨亲手算一笔“悬崖边的账”。

\begin{example}[悬崖边单步更新的数值分化]
设智能体站在悬崖正上方的格子 $s_1$，向右走到 $s_2$（仍在悬崖正上方）获步长惩罚 $r = -1$；而在 $s_2$ 共有 4 个动作：向东（沿悬崖边前行）估值为 $-11$，向北（远离悬崖）估值为 $-15$，向西为 $-13$，而向南一步便坠入悬崖，当前估值为 $Q(s_2, \text{南}) = -100$。取 $\gamma = 1, \varepsilon = 0.1$。
\begin{itemize}
  \item \textbf{Q-Learning 看 $s_2$}：它只看最优动作 $\max_{a'} Q(s_2, a') = Q(s_2, \text{东}) = -11$，因此给出目标 $y^{\mathrm{Q}} = -1 + (-11) = \mathbf{-12}$。坠崖风险被 $\max$ 算子彻底屏蔽——因为纯粹贪心策略绝不会主动往南跳！
  \item \textbf{SARSA 看 $s_2$}：它的下一步动作 $a_{t+1}$ 要从 $\varepsilon$-贪心策略中抽取。有 $0.925$ 的概率选东（$-11$），各有 $0.025$ 的概率选北（$-15$）、西（$-13$）与南（$-100$）！因此 SARSA 目标的期望值为
  \[
    \E[y^{\mathrm{SARSA}}] = -1 + \big(0.925(-11) + 0.025(-15) + 0.025(-13) + 0.025(-100)\big) = \mathbf{-14.375}.
  \]
\end{itemize}
这凭空多出来的 $-2.375$ 罚分会在反向回溯中逐格叠加：沿着 10 格长的悬崖边走一趟，每一步都有 $2.5\%$ 概率跌落 $-100$ 的深渊，期望额外代价高达约 $-25$ 分，远不如干脆多绕 4 步路（仅耗 $-4$ 分）走远离悬崖的顶部安全通道！表~\ref{tab:sarsa-vs-q} 系统汇总了两者的结构对比。
\end{example}

\begin{table}[htbp]
\centering
\caption{SARSA 与 Q-Learning 的核心理论与行为对比}
\label{tab:sarsa-vs-q}
\small
\begin{tabularx}{\linewidth}{@{}p{28mm}XX@{}}
\toprule
\textbf{对比维度} & \textbf{SARSA（同策略 On-Policy）} & \textbf{Q-Learning（异策略 Off-Policy）} \\
\midrule
自举目标 $y_t$ & $r_t + \gamma Q(s_{t+1}, a_{t+1})$，$a_{t+1} \sim \pi_\varepsilon$ & $r_t + \gamma \max_{a'} Q(s_{t+1}, a')$ \\
逼近算子与不动点 & 贝尔曼期望算子 $\mathcal{T}^{\pi_\varepsilon} \to Q^{\pi_\varepsilon}$ & 贝尔曼最优算子 $\mathcal{T}^* \to Q^*$ \\
收敛到 $Q^*$ 条件 & 需 GLIE（$\varepsilon_t \to 0$）+ Robbins--Monro & 只需无限访问全部 $(s,a)$ + Robbins--Monro \\
固定 $\varepsilon$ 下行为 & 主动规避探索失误风险（在线回报更高） & 紧贴理论最优路径（训练期偶发惨痛跌落） \\
数据复用潜力 & 严格要求当前策略样本，难以直接复用旧数据 & 天然支持经验回放池（直接孕育第 2 章 DQN） \\
\bottomrule
\end{tabularx}
\end{table}

\section{最大化偏差与 Double Q-Learning}
\label{sec:rl-double-q}

Q-Learning 凭借 $\max_{a'} Q(s_{t+1}, a')$ 赢得了异策略收敛到 $Q^*$ 的桂冠，但它暗藏着一个在带噪环境或神经网络近似下极具破坏力的统计缺陷——\textbf{最大化偏差（Maximization Bias / Overestimation Bias）}。

\subsection{Jensen 不等式证明：为何 $\max$ 必然高估价值？}

直觉上，如果你有 10 台期望中奖额完全为 $0$ 的老虎机，但由于有限次试玩存在随机波动，各台机器的经验均值有正有负，此时你挑出“经验均值最高的那一台”并把它的样本均值当作真实收益汇报，这个数字几乎必然显著大于 $0$——你其实是在给\textbf{正向噪声}颁奖。下面的定理把这一直觉写成了无懈可击的数学必然。

\begin{thmBox}
\begin{theorem}[最大化偏差定理（Jensen 不等式）]
\label{thm:max-bias}
设在某状态 $s$ 下，对各动作 $a \in \{1, \dots, K\}$ 的动作价值随机估计为 $\hat{Q}(s, a)$，且满足无偏性 $\E[\hat{Q}(s, a)] = q(s, a)$。则取最大值的期望必然不低于真实最大价值：
\begin{equation}
  \E\!\left[\max_{1 \le a \le K} \hat{Q}(s, a)\right] \ge \max_{1 \le a \le K} \E\!\left[\hat{Q}(s, a)\right] = \max_{1 \le a \le K} q(s, a).
  \label{eq:max-bias}
\end{equation}
更进一步，当所有真实价值均为 $0$、且 $\hat{Q}(s, 1), \hat{Q}(s, 2) \stackrel{\text{i.i.d.}}{\sim} \mathcal{N}(0, \sigma^2)$ 时，精确高估偏差为
\[
  \E\!\left[\max\!\big(\hat{Q}(s,1), \hat{Q}(s,2)\big)\right] = \frac{\sigma}{\sqrt{\pi}} \approx 0.564\,\sigma > 0.
\]
\end{theorem}
\begin{proof}
\textbf{初等直接证明}：对任意固定的 $j \in \{1, \dots, K\}$，随机变量间逐样本成立 $\max_{1 \le a \le K} \hat{Q}(s, a) \ge \hat{Q}(s, j)$。对两边取期望，由期望单调性得
\[
  \E\!\left[\max_{1 \le a \le K} \hat{Q}(s, a)\right] \ge \E\!\left[\hat{Q}(s, j)\right] = q(s, j).
\]
由于左端与指标 $j$ 无关，对右端关于 $j \in \{1, \dots, K\}$ 取最大值，立即得到式~\eqref{eq:max-bias}。（从凸分析视角看，函数 $f(x_1, \dots, x_K) = \max_a x_a$ 是仿射函数的逐点上确界，因而是 $\R^K$ 上的凸函数，故式~\eqref{eq:max-bias} 正是多元 Jensen 不等式 $\E[f(X)] \ge f(\E[X])$ 的特例。）

\textbf{双动作高斯精确算例}：利用恒等式 $\max(X_1, X_2) = \frac{1}{2}(X_1 + X_2 + |X_1 - X_2|)$。令 $D = X_1 - X_2$，因 $X_1, X_2 \stackrel{\text{i.i.d.}}{\sim} \mathcal{N}(0, \sigma^2)$，故 $D \sim \mathcal{N}(0, 2\sigma^2)$。利用对称性直接积分：
\[
  \E[\max(X_1, X_2)]
  = \frac{1}{2}\E[|D|]
  = \frac{1}{2} \cdot 2 \int_0^{\infty} x \cdot \frac{1}{\sqrt{4\pi\sigma^2}} e^{-x^2/(4\sigma^2)}\,\mathrm{d}x
  = \frac{1}{2\sigma\sqrt{\pi}} \left[-2\sigma^2 e^{-x^2/(4\sigma^2)}\right]_0^{\infty}
  = \frac{\sigma}{\sqrt{\pi}}. \qedhere
\]
\end{proof}
\end{thmBox}

当动作数 $K$ 增大时，独立标准正态最大值期望增长速度为 $\Theta(\sqrt{2 \ln K})$！在经典的两状态反例（Sutton \& Barto, Fig. 6.5）中：起点 $A$ 有两个动作——“向右”立即以确定的 $0$ 奖励终止；“向左”零奖励进入状态 $B$，而状态 $B$ 有 $K = 10$ 个动作，每个动作都以 $\mathcal{N}(-0.1, 1^2)$ 的随机奖励终止。显然状态 $B$ 的真实最优价值是 $-0.1 < 0$，“向右”才是绝对最优策略。但 Q-Learning 在 $B$ 中对 10 个带噪估计取 $\max$，算出严重虚高的正估值 $\max_a \hat{Q}(B, a) > 0$，并通过自举将虚假繁荣传回 $Q(A, \text{左})$，导致算法在早期绝大多数时候执迷不悟地往陷阱里走！

\subsection{Double Q-Learning：解耦“动作选择”与“价值评估”}

最大化偏差的病根在于\textbf{用同一份带噪数据既做“裁判选冠军”又做“会计发奖金”}：同一个 $\hat{Q}$ 既决定 $a^* = \argmax_{a'} \hat{Q}(s', a')$，又提供估值 $\hat{Q}(s', a^*)$。Hasselt（2010）提出的 \textbf{Double Q-Learning} 用外科手术般的精准将这两项职能彻底拆开：

\begin{defBox}
\begin{definition}[Double Q-Learning 更新算法]
同时维护两张独立初始化、用互不重叠样本更新的动作价值表 $Q_A$ 与 $Q_B$。行为策略基于两表之和 $Q_A(s, \cdot) + Q_B(s, \cdot)$ 执行 $\varepsilon$-贪心。每观测到单步转移 $(s_t, a_t, r_t, s_{t+1})$，以 $0.5$ 概率随机抛硬币选择其中一张表更新：
\begin{itemize}
  \item \textbf{更新 $Q_A$}：用 $Q_A$ \emph{选择}后继最优动作，用 $Q_B$ \emph{评估}该动作的价值：
  \begin{equation}
    a^* = \argmax_{a' \in \mathcal{A}} Q_A(s_{t+1}, a'), \qquad
    Q_A(s_t, a_t) \leftarrow Q_A(s_t, a_t) + \alpha_t \!\left[r_t + \gamma Q_B(s_{t+1}, a^*) - Q_A(s_t, a_t)\right].
    \label{eq:double-q-a}
  \end{equation}
  \item \textbf{更新 $Q_B$}：完全对称地用 $Q_B$ 选择动作、用 $Q_A$ 评估价值：
  \begin{equation}
    b^* = \argmax_{a' \in \mathcal{A}} Q_B(s_{t+1}, a'), \qquad
    Q_B(s_t, a_t) \leftarrow Q_B(s_t, a_t) + \alpha_t \!\left[r_t + \gamma Q_A(s_{t+1}, b^*) - Q_B(s_t, a_t)\right].
    \label{eq:double-q-b}
  \end{equation}
\end{itemize}
\end{definition}
\end{defBox}

为何这样就能治愈正向高估？因为 $Q_A$ 与 $Q_B$ 使用独立样本更新，若 $a^*$ 取自 $\argmax_{a'} Q_A(s', a')$，则给定 $Q_A$ 时独立估计表 $Q_B$ 对 $a^*$ 完全无偏：
\[
  \E\!\left[Q_B(s', a^*) \mid Q_A\right] = q(s', a^*) \le \max_{a' \in \mathcal{A}} q(s', a').
\]
再对 $Q_A$ 取全期望，立刻有 $\E[Q_B(s', a^*)] = \E[q(s', a^*)] \le \max_{a'} q(s', a')$！正向高估偏差彻底消失（取而代之的是在动作价值相近时轻微、无害的保守低估，不会引发自举爆炸）。这一思想将在第~\ref{ch:dqn-pg-ac} 章无缝升级为深度强化学习标配 \textbf{Double DQN}。

\section{实测验证：悬崖行走中的 SARSA 与 Q-Learning 分化}
\label{sec:exp-cliff}

现在我们在强化学习最经典的检验场——\textbf{$4 \times 12$ 悬崖行走（CliffWalking）}网格世界中，亲自验证同策略与异策略的戏剧性分化。环境配置见脚本 \texttt{code/\allowbreak vol6\_experiments.py} 中的 \texttt{exp1\_cliff\_\allowbreak sarsa\_vs\_q()}：
\begin{itemize}
  \item 网格高 $H=4$、宽 $W=12$，起点在左下角 $S = (3, 0)$，终点在右下角 $G = (3, 11)$；底部中间 10 格 $(3, 1)$ 至 $(3, 10)$ 为万丈悬崖。
  \item 动作空间 $\mathcal{A} = \{\text{上}, \text{下}, \text{左}, \text{右}\}$。普通每走一步（含抵达终点）获奖励 $-1.0$；一旦踩入悬崖格，立刻承受 $-100.0$ 巨额罚分并被送回起点 $(3, 0)$（回合继续，\texttt{done=False}）；到达终点 $(3, 11)$ 则回合结束（\texttt{done=True}）。
  \item 超参数：$\gamma = 1.0, \alpha = 0.5, \varepsilon = 0.1$，每回合最多 200 步，共训练 500 回合，跨 20 组随机种子取平均。
\end{itemize}

\begin{lstlisting}[caption={SARSA 与 Q-Learning 在悬崖行走中的核心对比循环（摘自 \texttt{vol6\_experiments.py}）}]
for ep in range(episodes):
    s = 3 * W + 0
    a = pick(s)  # eps-greedy(eps=0.1)
    tot = 0.0
    for _ in range(200):
        ns, rew, done = step(s, a)
        tot += rew
        if mode == "sarsa":
            na = pick(ns)  # 同策略：先采样下一步真实动作 na
            target = rew + (0.0 if done else gamma * Q[ns, na])
            Q[s, a] += alpha * (target - Q[s, a])
            s, a = ns, na
        else:              # 异策略 Q-Learning：用 max 构造最优目标
            target = rew + (0.0 if done else gamma * np.max(Q[ns]))
            Q[s, a] += alpha * (target - Q[s, a])
            s = ns
            a = pick(s)
        if done:
            break
\end{lstlisting}

图~\ref{fig:cliff-grid} 画出了两种算法各自学会的典型路径，而图~\ref{fig:cliff-curve} 直接读取实测记录文件 \texttt{data/rl\_cliff.csv}（每 10 回合滑动平均、20 种子均值），完整呈现了在线训练曲线与离线贪心评估的惊人反差：

\begin{figure}[htbp]
\centering
\resizebox{0.96\linewidth}{!}{%
\begin{tikzpicture}[scale=0.92]
  % Grid cells
  \foreach \r in {0,1,2,3} {
    \foreach \c in {0,1,...,11} {
      \draw[tealLine, thin] (\c, 3-\r) rectangle (\c+1, 4-\r);
    }
  }
  % Cliff cells (row 3, col 1..10 -> y=0..1)
  \fill[ochreLight] (1, 0) rectangle (11, 1);
  \draw[ochre, thick] (1, 0) rectangle (11, 1);
  \node[font=\sffamily\bfseries\small, color=ochre] at (6, 0.5) {悬崖区域（跌落奖励 $-100$，立即重置回起点 $S$）};

  % Start & Goal
  \fill[tealLight] (0, 0) rectangle (1, 1);
  \draw[tealDark, thick] (0, 0) rectangle (1, 1);
  \node[font=\sffamily\bfseries\normalsize, color=tealDark] at (0.5, 0.5) {$S$};

  \fill[navyLight] (11, 0) rectangle (12, 1);
  \draw[navy, thick] (11, 0) rectangle (12, 1);
  \node[font=\sffamily\bfseries\normalsize, color=navy] at (11.5, 0.5) {$G$};

  % SARSA Safe Path (top detour)
  \draw[teal, line width=2.0pt, -{Stealth[length=2.6mm]}, rounded corners=4pt]
    (0.35, 0.75) -- (0.35, 3.55) -- (11.65, 3.55) -- (11.65, 0.85);
  \node[font=\sffamily\bfseries\scriptsize, color=tealDark, fill=white, inner sep=1.8pt]
    at (6.0, 3.55) {SARSA 安全路径（远离悬崖，步长 $\ge 17$，在线均分 $-26.22$）};

  % Q-Learning Optimal Path (cliff edge)
  \draw[ochre, line width=2.0pt, dashed, -{Stealth[length=2.6mm]}, rounded corners=3pt]
    (0.68, 0.75) -- (0.68, 1.48) -- (11.32, 1.48) -- (11.32, 0.85);
  \node[font=\sffamily\bfseries\scriptsize, color=ochre, fill=white, inner sep=1.8pt]
    at (6.0, 1.88) {Q-Learning 最优路径（紧贴崖边，精确 13 步，贪心均分 $-13.0$）};

  % Row indices
  \foreach \r in {0,1,2,3} {
    \node[font=\scriptsize, color=slate] at (-0.35, 3.5-\r) {行 \r};
  }
\end{tikzpicture}}
\caption{$4 \times 12$ 悬崖行走（CliffWalking）路径分化示意图。SARSA 因在更新中计入 $\varepsilon=0.1$ 探索跌崖的代价，学会沿着顶行绕远的安全路线；Q-Learning 直接逼近最优动作价值 $Q^*$，学会紧贴悬崖边缘的 13 步最短路径。}
\label{fig:cliff-grid}
\end{figure}

\begin{figure}[htbp]
\centering
\begin{tikzpicture}
  \begin{axis}[
    width=0.90\linewidth, height=6.4cm,
    xlabel={训练回合数（Episode）},
    ylabel={在线回合总回报（10 回合均值，20 种子）},
    xmin=0, xmax=500, ymin=-110, ymax=-10,
    grid=both, grid style={tealLine!45, dashed},
    axis line style={slate},
    tick label style={font=\scriptsize},
    label style={font=\footnotesize\sffamily, color=ink},
    legend style={at={(0.98,0.12)}, anchor=south east, font=\scriptsize, fill=white, draw=tealLine},
    clip=true
  ]
    \addplot[teal, very thick, mark=*, mark size=1.6pt]
      table[x=ep, y=sarsa, col sep=comma]{data/rl_cliff.csv};
    \addlegendentry{SARSA（后 100 轮在线均值 $-26.22$，贪心评估 $-23.9$）}

    \addplot[ochre, very thick, dashed, mark=square*, mark size=1.5pt]
      table[x=ep, y=qlearning, col sep=comma]{data/rl_cliff.csv};
    \addlegendentry{Q-Learning（后 100 轮在线均值 $-51.90$，贪心评估 $-13.0$）}

    \draw[dotted, navy, thick] (axis cs:0, -13.0) -- (axis cs:500, -13.0);
    \node[font=\scriptsize\sffamily, color=navy, anchor=north east] at (axis cs:495, -13.5) {理论最优回报 $-13.0$};
  \end{axis}
\end{tikzpicture}
\caption{悬崖行走实测在线学习曲线（读取自 \texttt{data/rl\_cliff.csv}；注：首点 \texttt{ep=10} 处两者分别从 $-310.08$ 与 $-300.46$ 快速爬升，纵轴截断于 $[-110, -10]$ 以清晰展示稳态分化）。在线交互时 SARSA 远胜 Q-Learning，而关掉探索做贪心评估时 Q-Learning 拿下精确理论最优 $-13.0$。}
\label{fig:cliff-curve}
\end{figure}

实测数字揭示了教科书级、却在工程选型中常被误读的深刻规律：
\begin{enumerate}
  \item \textbf{在线训练期（带 $\varepsilon = 0.1$ 探索）}：前 10 回合两者都在反复试错摔崖（\texttt{ep=10} 均分分别为 $-310.08$ 与 $-300.46$）。但到第 90 回合起分化彻底定型：最后 100 回合的平均在线回报，\textbf{SARSA 稳定在 $-26.22$}，而 \textbf{Q-Learning 跌落在 $-51.90$}！原因一目了然——Q-Learning 坚持贴着悬崖走，在 10 步紧邻深渊的行走中，每步都有 $2.5\%$ 的概率因 $\varepsilon$-探索失足跌落 $-100$；SARSA 则自觉退守远离悬崖的上方安全走廊，几乎不再坠崖。
  \item \textbf{离线部署期（关掉探索，纯贪心 $\argmax_a Q(s, a)$）}：战局瞬间逆转！\textbf{Q-Learning 的贪心评估回报在全部 20 组种子上均精准达到理论最优值 $-13.0$}（上 1 步 $+$ 右 11 步 $+$ 下 1 步），而 \textbf{SARSA 的贪心评估均值仅为 $-23.9$}（绕顶行安全路至少耗 17 步，且在恒定高步长 $\alpha=0.5$ 与不衰减 $\varepsilon=0.1$ 下部分种子在上方安全区还带有迂回噪声）。
\end{enumerate}

由此可以提炼出一条极具实战价值的选型原则：\textbf{如果训练过程本身必须在昂贵的真实物理世界（如真机机器人、在线推荐/医疗策略）中承担探索代价，同策略的保守安全特性（SARSA 类）至关重要；如果允许在无成本模拟器中尽情摔跤、最终只把学成的最优策略部署上线，则异策略的 Q-Learning 更具优势。}

\begin{failBox}{表格型 TD/Q-Learning 实现中最隐蔽的两大 Bug}
\begin{enumerate}
  \item \textbf{忘记屏蔽终止状态的自举项}：在上面的代码中，\texttt{target = rew + (0.0 if done else gamma * np.max(Q[ns]))} 里的条件掩码 \texttt{0.0 if done else ...} 绝非可有可无！许多实现会漏写 \texttt{if done}，把终止状态 \texttt{ns} 对应的 $Q[\text{ns}]$ 也加进目标。虽然初始全零时看似侥幸无事，但一旦使用函数逼近（神经网络对终止态输出非零值）或环境将回合结束时的状态重置为起点，漏掉终止掩码会直接制造虚假的永动机循环，导致 $Q$ 值迅速发散。
  \item \textbf{探索率 $\varepsilon$ 设得过小或永不退火}：若 $\varepsilon$ 从一开始就设得极小（如 $0.01$），智能体在尚未找到终点前容易把少数负奖励动作越压越低，或在多臂环境中过早锁死于次优路径；反之，若在需要高在线回报的同策略任务中一直保持大 $\varepsilon$ 不衰减，即便已收敛到最优安全路线，也会持续支付不必要的探索税。标准工程实践是采用线性或指数退火（如在训练前 80\% 步数内从 $\varepsilon_0 = 1.0$ 平滑降至 $\varepsilon_{\min} = 0.05$ 或 $0.01$）。
\end{enumerate}
\end{failBox}

\begin{extBox}{$n$-步 TD、资格迹 TD($\lambda$) 与期望 SARSA}
MC（完整回报，零偏高方差）与 TD(0)（单步自举，有偏低方差）并非非此即彼的二元对立，而是同一条连续谱系的两个端点！
\begin{itemize}
  \item \textbf{$n$-步时序差分（$n$-Step TD）}：向前真实走 $n$ 步再自举，定义 $n$-步回报
  \[
    G_t^{(n)} = r_t + \gamma r_{t+1} + \cdots + \gamma^{n-1} r_{t+n-1} + \gamma^n V(s_{t+n}).
  \]
  当 $n=1$ 时退化为 TD(0)；当 $n \to \infty$ 时退化为 MC。实证表明，在绝大多数任务中，中等步数（如 $n \in [4, 16]$）往往同时击败 $n=1$ 与 $n=\infty$。
  \item \textbf{TD($\lambda$) 与资格迹（Eligibility Traces）}：与其硬选某个固定的 $n$，不如对所有 $n$-步回报做几何加权平均，定义\textbf{$\lambda$-回报（前向视角）}：
  \[
    G_t^\lambda = (1 - \lambda) \sum_{n=1}^{\infty} \lambda^{n-1} G_t^{(n)}, \qquad \lambda \in [0, 1].
  \]
  在工程实现中，它可通过\textbf{资格迹（后向视角）}每步在线增量完成：维护迹向量 $e_t(s) = \gamma \lambda e_{t-1}(s) + \mathbf{1}[s = s_t]$，并用当前单步误差 $\delta_t$ 同时更新所有近期访问过的状态：$V(s) \leftarrow V(s) + \alpha \delta_t e_t(s)$。这正是第~\ref{ch:dqn-pg-ac} 章广义优势估计 GAE($\lambda$) 的直系血亲！
  \item \textbf{期望 SARSA（Expected SARSA）}：将 SARSA 目标中的随机采样项 $Q(s_{t+1}, a_{t+1})$ 替换为对当前策略 $\pi$ 的精确期望
  \[
    y_t^{\mathrm{ExpSARSA}} = r_t + \gamma \sum_{a' \in \mathcal{A}} \pi(a' \mid s_{t+1}) Q(s_{t+1}, a').
  \]
  它彻底消除了因 $a_{t+1}$ 随机抽样引入的额外方差，因而允许使用更大步长 $\alpha$ 且更稳定；更美妙的是，当目标策略 $\pi$ 取贪心策略时，期望 SARSA 直接退化为 Q-Learning——由此统一了同策略与异策略的表格控制算法。
\end{itemize}
\end{extBox}

\section*{本章小结与核心脉络}
\addcontentsline{toc}{section}{本章小结与核心脉络}

本章完成了从“已知模型动态规划”到“无模型交互学习”的决定性跨越：
\begin{itemize}
  \item \textbf{评估侧}：蒙特卡洛（MC）用完整回报给出无偏但高方差的估计；TD(0) 用单步自举 $r_t + \gamma V(s_{t+1})$ 大幅压低方差，并能利用马尔可夫结构实现在线持续学习；Robbins--Monro 条件（$\sum \alpha_t = \infty, \sum \alpha_t^2 < \infty$）为一切随机逼近算法的收敛托底。
  \item \textbf{控制侧}：无模型决策要求直接学习动作价值 $Q(s, a)$，并借助 $\varepsilon$-贪心与 GLIE 平衡探索与利用；SARSA 诚实评估包含探索风险的当前策略，Q-Learning 则异策略直指最优价值 $Q^*$，两者在悬崖行走中完美演绎了“在线安全（$-26.22$）”与“贪心最优（$-13.0$）”的权衡；Double Q-Learning 则通过解耦选择与评估，根除了 Jensen 不等式注定的正向高估偏差。
\end{itemize}

然而，本章的所有算法都建立在一张有限大小的查表 $Q[s, a]$ 之上。当状态变成 $84 \times 84 \times 4$ 的游戏像素画面、或是机械臂的连续关节角度与速度向量时，查表彻底失效。下一章，我们将把神经网络引入强化学习，直面价值近似中的“致命三元组”与 DQN 稳定化技术，并推导另一条通往现代大模型后训练的康庄大道——\textbf{策略梯度定理与 Actor-Critic}。

\section*{本章分层习题}
\addcontentsline{toc}{section}{本章分层习题}

\exTier{A 层：概念与手算热身}{巩固 MC/TD 更新公式、SARSA 与 Q-Learning 差异}

\begin{enumerate}
  \item[\textbf{6.1.1}] \textbf{单条轨迹的 MC 与 TD(0) 逐步手算}：考虑含状态 $\{s_1, s_2, \mathsf{T}\}$ 的回合制任务，初值 $V(s_1) = 0, V(s_2) = 0$，折扣因子 $\gamma = 0.9$，固定步长 $\alpha = 0.5$。观测到第一条回合：
  \[
    s_1 \xrightarrow{r=2} s_2 \xrightarrow{r=4} \mathsf{T}.
  \]
  \begin{enumerate}
    \item 分别计算该回合结束后，首次访问 MC（用恒定步长 $\alpha=0.5$）与 TD(0)（按时间步在线更新）给出的新价值表 $V(s_1), V(s_2)$。
    \item 若紧接着观测到完全相同的第二条回合，再次计算 TD(0) 的更新结果，观察奖励信息是如何沿时间逆向“逐回合传导”给 $s_1$ 的。
  \end{enumerate}
  \exHint{注意 TD(0) 在处理第 1 步 $s_1 \to s_2$ 时，$V(s_2)$ 尚未被第 2 步更新，仍为初值 $0$；到第 2 条回合时，$V(s_2)$ 的非零信息才第一次通过自举传给 $V(s_1)$。}

  \item[\textbf{6.1.2}] \textbf{SARSA、Q-Learning 与期望 SARSA 单步目标对比}：观测到转移 $(s_t, a_t, r_t, s_{t+1}) = (s_1, \text{右}, -1, s_2)$；其中状态 $s_2$ 含 3 个动作，当前表值为 $Q(s_2, a_1) = 10$、$Q(s_2, a_2) = 2$、$Q(s_2, a_3) = -20$，采用 $\varepsilon = 0.15$ 的 $\varepsilon$-贪心策略，$\gamma = 0.9$。
  \begin{enumerate}
    \item 分别写出 Q-Learning 目标、期望 SARSA 目标的精确数值；
    \item 写出 SARSA 目标的可能取值分布，并计算其期望与方差。
  \end{enumerate}
  \exHint{由式~\eqref{eq:eps-greedy}，$m=3$ 时最优动作 $a_1$ 概率为 $1 - 0.15 + 0.05 = 0.90$，另两个动作各为 $0.05$；验证 SARSA 目标的期望恰好等于期望 SARSA 目标，而方差严格为正。}
\end{enumerate}

\exTier{B 层：定理推导与统计分析}{深入掌握 Robbins--Monro、最大化偏差与无偏性证明}

\begin{enumerate}
  \item[\textbf{6.1.3}] \textbf{Robbins--Monro 步长家族的边界检验}：考虑递减步长序列 $\alpha_t = (t + 1)^{-p}$（$t \ge 0$）。
  \begin{enumerate}
    \item 证明当且仅当 $\frac{1}{2} < p \le 1$ 时，$\{\alpha_t\}$ 同时满足 Robbins--Monro 条件~\eqref{eq:rm-cond}。
    \item 对标量常值均值估计 $\theta_{t+1} = \theta_t + \alpha(\mu + \xi_t - \theta_t)$，若采用\textbf{固定常数步长} $\alpha \in (0, 1)$（不满足 $\sum \alpha_t^2 < \infty$），求 $\lim_{t \to \infty} \E[(\theta_t - \mu)^2]$ 的精确闭合表达式，并解释为何它在非平稳环境中反而受青睐。
  \end{enumerate}
  \exHint{(b) 中对式~\eqref{eq:rm-mse-rec} 取极限 $v_\infty = (1-\alpha)^2 v_\infty + \alpha^2 \sigma^2$，解得稳态均方误差 $v_\infty = \frac{\alpha}{2-\alpha}\sigma^2$。}

  \item[\textbf{6.1.4}] \textbf{均匀噪声下的最大化偏差与 Double Q 对比}：设状态 $s$ 下有 $K$ 个动作，其真实价值全为 $0$，但估计值独立同分布于均匀分布 $\hat{Q}(s, a) \sim \mathrm{Uniform}(-1, 1)$。
  \begin{enumerate}
    \item 计算单表最大值的精确期望 $\E[\max_{1 \le a \le K} \hat{Q}(s, a)]$（以 $K$ 表示），并讨论 $K \to \infty$ 时的极限。
    \item 若按 Double Q-Learning 维护两组独立同分布于 $\mathrm{Uniform}(-1, 1)$ 的估值表 $Q_A, Q_B$，试求期望 $\E\big[Q_B(s, \argmax_a Q_A(s, a))\big]$。
  \end{enumerate}
  \exHint{利用最大值累积分布函数 $F_{\max}(x) = \big(\frac{x+1}{2}\big)^K$（$x \in [-1, 1]$）求导得密度函数并积分，结果应为 $\frac{K-1}{K+1}$；而独立评估的期望恒为 $0$。}
\end{enumerate}

\exTier{C 层：算法构造与反例辨析}{辨识环境建模假设、探索退火与异策略边界}

\begin{enumerate}
  \item[\textbf{6.1.5}] \textbf{批量 MC 与批量 TD(0) 的非马尔可夫反例}：在例~\ref{ex:batch-mc-td} 中我们看到，当环境满足马尔可夫性时，批量 TD(0) 给出比批量 MC 更合理的估计 $V(A) = 0.75$。请构造一个\textbf{违反马尔可夫性}（即 $B$ 的未来奖励实际上取决于是否经由 $A$ 到达 $B$）的隐状态过程，使得两条数据生成机制完全吻合例~\ref{ex:batch-mc-td} 的 8 条观测，但此时批量 MC 的答案 $V(A) = 0$ 才是准确的，而 TD(0) 会系统性出错。
  \exHint{设真实状态是“是否携带标记”，从 $A$ 到达 $B$ 时必然触发“损坏标记”导致之后必得 $0$ 分，而直接生于 $B$ 的轨迹有 $6/7$ 概率得 $1$ 分；观测器无法区分两种 $B$ 时，马尔可夫假设破裂。}

  \item[\textbf{6.1.6}] \textbf{GLIE 与步长耦合的探索陷阱}：某同学在实现 SARSA 时，为了加快收敛，将探索率设为指数衰减 $\varepsilon_k = 0.5^k$（$k=1,2,\dots$ 为回合数）。请利用 Borel--Cantelli 引理说明，为何这种指数衰减探索率\textbf{违反 GLIE 条件}，会导致算法以正概率永远遗漏某个未探索过的最优动作。
  \exHint{因 $\sum_{k=1}^\infty 0.5^k = 1 < \infty$，若定义事件 $E_k = \{\text{第 } k \text{ 回合在某状态下发生随机探索}\}$，由第一 Borel--Cantelli 引理知，以概率 1 只有有限多个 $E_k$ 会发生。}
\end{enumerate}

\exTier{D 层：编程实战与前沿桥接}{复现 Double Q 抑制高估并向函数逼近过渡}

\begin{enumerate}
  \item[\textbf{6.1.7}] \textbf{编程实验：在悬崖行走与双状态带噪 MDP 上对比三种表格算法}：基于 \texttt{code/vol6\_experiments.py} 中的 \texttt{exp1\_cliff\_sarsa\_vs\_q()}，请完成以下两项扩展：
  \begin{enumerate}
    \item 在悬崖行走中实现 \textbf{Expected SARSA}，分别测试 $\alpha \in \{0.1, 0.5, 0.9\}$ 下的后 100 轮在线平均回报，验证它在高步长 $\alpha=0.9$ 下为何比标准 SARSA 更稳健；
    \item 编写 Sutton--Barto 的双状态 $A/B$（状态 $B$ 含 $K=10$ 个奖励服从 $\mathcal{N}(-0.1, 1)$ 的动作）测试环境，比较标准 Q-Learning 与 Double Q-Learning 在前 300 回合中选择“向左”动作的比例曲线。
  \end{enumerate}
  \exHint{在双状态实验中，Q-Learning 在前 50 回合选“向左”的比例通常会飙升至 $60\%$ 以上，而 Double Q-Learning 会稳定维持在 $50\%$ 以下并迅速逼近最优探索下限 $\varepsilon/2$。}

\par\pagebreak[4]
  \item[\textbf{6.1.8}] \textbf{前瞻思考：从查表走向线性逼近的第一个裂缝}：假设状态空间过大，我们用线性函数 $V_w(s) = w^\top \phi(s)$ 逼近价值函数，并仿照 TD(0) 写出所谓\textbf{半梯度 TD(0)（Semi-Gradient TD(0)）}更新：
  \[
    w \leftarrow w + \alpha \!\left[r_t + \gamma V_w(s_{t+1}) - V_w(s_t)\right] \nabla_w V_w(s_t).
  \]
  请指出：为什么上式不是最小化均方贝尔曼误差 $\frac{1}{2}\big(r_t + \gamma V_w(s_{t+1}) - V_w(s_t)\big)^2$ 的\emph{真实梯度}？若对目标项 $V_w(s_{t+1})$ 也求梯度（即全残差梯度），会多出哪一项？这一观察如何预示第~\ref{ch:dqn-pg-ac} 章的“目标网络”与“致命三元组”？
  \exHint{真实梯度对整个平方项求导会产生 $\nabla_w V_w(s_t) - \gamma \nabla_w V_w(s_{t+1})$；半梯度法刻意把 $V_w(s_{t+1})$ 当作常量截断梯度（在 PyTorch 中即 \texttt{.detach()}），这在异策略与非线性近似结合时会打破压缩映射保证。}
\end{enumerate}

\chapter[深度价值、策略梯度与 Actor-Critic]{深度价值学习与策略梯度定理\\（DQN、REINFORCE 与 Actor-Critic）}
\label{ch:dqn-pg-ac}

\chapEpigraph{凡是我们必须学会去做的事，我们正是通过去做来学会的。}{Aristotle（亚里士多德，《尼各马可伦理学》）}

\begin{chapIntro}{本章导读：告别查表时代，驶入深度强化学习双航道}
第~\ref{ch:model-free-rl} 章的表格型算法虽然拥有优美的收敛保证，却面临一道无法逾越的维度高墙：一帧 $84 \times 84$ 的灰度游戏画面便有 $256^{7056}$ 种状态，而机械臂的关节位置与角速度更是连续实数向量——地球上所有的原子加起来也存不下一张 $Q[s, a]$ 表。

为了赋予智能体泛化到未知状态的能力，我们必须用参数化神经网络替代查表。然而，神经网络与强化学习的联姻远非“把 $Q$ 表换成 MLP”那般平顺。本章完整勾勒现代深度强化学习的\textbf{两条主航道}：
\begin{enumerate}
  \item \textbf{基于价值的深度航道（DQN）}：亲手推导一个两状态反例，看清\textbf{函数逼近 $+$ 自举 $+$ 异策略}组成的“致命三元组”如何令训练指数爆炸；随后剖析 DQN 赖以起死回生的两大稳定器——\textbf{经验回放}与\textbf{目标网络}。
  \item \textbf{基于策略的梯度航道（Policy Gradient）}：绕开对离散 $\argmax$ 的依赖，直接参数化随机策略 $\pi_\theta(a\mid s)$。完整证明\textbf{策略梯度定理}（揭示为何环境转移核 $P$ 奇迹般地完全消去）与\textbf{基线零偏降方差定理}，并推导统一偏差与方差的 \textbf{Actor-Critic 与广义优势估计 GAE($\lambda$)}。
  \item \textbf{实测验证}：在带正偏移奖励的 12 步随机控制任务上，用 PyTorch 实测对比无基线 REINFORCE、价值基线 REINFORCE 与 GAE Actor-Critic，亲眼见证基线如何将方差压缩近 $4\times$ 并解锁高效收敛。
\end{enumerate}
\end{chapIntro}

\section{价值函数逼近、致命三元组与 DQN}
\label{sec:dqn-triad}

\subsection{从查表到参数化 $Q_\theta(s, a)$ 与致命三元组}

当状态空间 $\mathcal{S}$ 巨大或连续时，我们引入参数为 $\theta \in \R^d$ 的神经网络 $Q_\theta(s, a)$ 来逼近最优动作价值函数 $Q^*(s, a)$。仿照 Q-Learning，自然想到用随机梯度下降最小化单步贝尔曼均方误差：
\begin{equation}
  \mathcal{L}(\theta) = \frac{1}{2}\,\E_{(s, a, r, s') \sim \mathcal{D}}\!\left[\left(y - Q_\theta(s, a)\right)^2\right], \qquad
  y = r + \gamma \max_{a' \in \mathcal{A}} Q_\theta(s', a').
  \label{eq:naive-dqn-loss}
\end{equation}
在反向传播时，按第 1 章习题~\textbf{6.1.8} 提到的\textbf{半梯度（Semi-Gradient）}约定，我们将目标 $y$ 视作常量（截断对目标中 $\theta$ 的求导），得到更新方向
\[
  \theta \leftarrow \theta + \alpha \!\left(r + \gamma \max_{a'} Q_\theta(s', a') - Q_\theta(s, a)\right) \nabla_\theta Q_\theta(s, a).
\]

然而，如果直接把这套看似天衣无缝的算法丢给神经网络训练，你会震惊地发现：\textbf{即便是在极其简单的环境里，价值估计也会剧烈震荡甚至发散至无穷大！} Sutton 与 Barto 将导致深度价值学习发散的罪魁祸首归纳为\textbf{致命三元组（The Deadly Triad）}：
\begin{enumerate}
  \item \textbf{函数逼近（Function Approximation）}：不同状态共享同一套参数 $\theta$；修改 $Q_\theta(s, a)$ 会不可避免地连带改变其他状态 $Q_\theta(s', a')$ 的预测值。
  \item \textbf{自举（Bootstrapping）}：训练目标 $y$ 依赖于网络自身的后继预测 $Q_\theta(s', a')$，导致监督学习的“标签”本身在每一步都在移动。
  \item \textbf{异策略训练（Off-Policy Training）}：更新所用的样本分布 $\mu$ 与目标策略所诱导的平稳状态分布不匹配，打破了投影贝尔曼算子的压缩性。
\end{enumerate}

绝不要以为发散只是深度非线性网络才会遇到的边缘故障——下面的两状态线性反例证明，哪怕只有一个标量参数 $w \in \R$、且真实价值函数完全可以被该模型精确表示，致命三元组依然能让半梯度 TD 指数爆炸！

\begin{thmBox}
\begin{example}[Tsitsiklis--Van Roy 两状态致命三元组发散反例]
\label{ex:deadly-triad}
考虑仅含两个状态的无奖励马尔可夫过程：$s_1 \xrightarrow{r=0} s_2$ 与 $s_2 \xrightarrow{r=0} s_2$（自环），折扣因子 $\gamma = 0.9$。显然真实价值恒为零：$V^*(s_1) = V^*(s_2) = 0$。

我们采用一维线性逼近 $V_w(s) = w\,\phi(s)$，其中特征设为 $\phi(s_1) = 1, \phi(s_2) = 2$。注意最优解 $w^* = 0$ 可以\textbf{零误差}精确表示真值！现在假设在异策略采样下，每步以概率 $p = 0.5$ 采样转移 $(s_1 \to s_2)$、以概率 $1-p = 0.5$ 采样转移 $(s_2 \to s_2)$，执行半梯度 TD(0) 更新 $w \leftarrow w + \alpha [0 + \gamma V_w(s') - V_w(s)] \phi(s)$：
\begin{itemize}
  \item 若抽到 $(s_1 \to s_2)$：$\Delta w_1 = \alpha(\gamma \cdot 2w - 1w)\cdot 1 = \alpha(2\gamma - 1) w = 0.8\alpha w$；
  \item 若抽到 $(s_2 \to s_2)$：$\Delta w_2 = \alpha(\gamma \cdot 2w - 2w)\cdot 2 = 4\alpha(\gamma - 1) w = -0.4\alpha w$。
\end{itemize}
取期望更新：
\[
  \E[w_{t+1} \mid w_t] = w_t + \frac{1}{2}(0.8\alpha w_t) + \frac{1}{2}(-0.4\alpha w_t) = (1 + 0.2\alpha)\, w_t.
\]
只要初值 $w_0 \neq 0$，\textbf{无论学习率 $\alpha > 0$ 选得多么微小}，$|w_t| = (1 + 0.2\alpha)^t |w_0| \to +\infty$ 必然指数发散！
\end{example}
\end{thmBox}

反例的物理机制一针见血：在 $(s_1 \to s_2)$ 这一步，$V_w(s_1) = w$ 试图追赶更大的自举目标 $\gamma V_w(s_2) = 1.8w$，于是把权重 $w$ 往大了推；但由于 $s_1$ 与 $s_2$ 共享参数 $w$，且 $\phi(s_2) = 2\phi(s_1)$，把 $w$ 推大后，目标 $\gamma V_w(s_2)$ 跑得比 $V_w(s_1)$ \emph{还要快}！如同追逐自己拉长的影子，正反馈循环瞬间失控。（Baird 著名的 7 状态星形反例正是这一机制的多维版本。）

\subsection{DQN 的两大核心稳定器：经验回放与目标网络}

2013--2015 年，DeepMind 凭借\textbf{深度 Q 网络（Deep Q-Network, DQN）}首次用端到端卷积网络在 Atari 2600 游戏上达到人类专家水平（Mnih et al., 2015）。DQN 并没有在理论上彻底消灭致命三元组，但它用两项极为精妙的工程与算法稳定器，成功驯服了深度异策略自举的烈马：

\begin{defBox}
\begin{definition}[DQN 的两大核心稳定机制]
\begin{enumerate}
  \item \textbf{经验回放（Experience Replay）}：维护一个定容循环缓冲区 $\mathcal{D} = \{(s_i, a_i, r_i, s'_i, d_i)\}_{i=1}^N$（其中 $d_i \in \{0, 1\}$ 为终止标志）。每次梯度更新时，从 $\mathcal{D}$ 中\textbf{均匀随机采样}一个小批量 $\mathcal{B}$ 执行 SGD，带来双重红利：
  \begin{itemize}
    \item \emph{打破时间相关性}：在线交互相邻帧高度相似（如连续 30 帧球都在向左飞），直接在线更新会导致网络灾难性遗忘其他状态；随机打乱回放使小批量重新近似满足 SGD 所依赖的独立同分布（i.i.d.）假设。
    \item \emph{提升样本效率}：因为 Q-Learning 是异策略的，每条昂贵的交互样本可以被反复复用数十次。
  \end{itemize}
  \item \textbf{目标网络（Target Network）}：额外维护一份参数延迟同步的副本网络 $Q_{\theta^-}$，在计算贝尔曼目标时彻底冻结 $\theta^-$：
  \begin{equation}
    y_i^{\mathrm{DQN}} = r_i + \gamma (1 - d_i) \max_{a' \in \mathcal{A}} Q_{\theta^-}(s'_i, a'), \qquad
    \mathcal{L}_{\mathrm{DQN}}(\theta) = \frac{1}{|\mathcal{B}|}\sum_{i \in \mathcal{B}} \ell_{\mathrm{Huber}}\!\left(y_i^{\mathrm{DQN}} - Q_\theta(s_i, a_i)\right).
    \label{eq:dqn-loss}
  \end{equation}
  每隔 $C$ 步执行一次硬同步 $\theta^- \leftarrow \theta$（或在连续控制中每步执行 Polyak 软更新 $\theta^- \leftarrow \tau \theta + (1-\tau)\theta^-$，$\tau \approx 0.005$）。
\end{enumerate}
\end{definition}
\end{defBox}

目标网络的妙处在于它\textbf{暂时切断了“追逐自己影子”的正反馈环}：在两次同步间隔的 $C$ 步里，$y_i^{\mathrm{DQN}}$ 是完全静止的监督回归标签，式~\eqref{eq:dqn-loss} 退化为一个良态的监督学习问题（常配以 Huber 损失以防偶发极端 TD 误差引发梯度爆炸）。

在原始 DQN 之上，两项几乎零额外开销的改进成为现代价值网络标配：
\begin{itemize}
  \item \textbf{Double DQN}（van Hasselt et al., 2016）：将第~\ref{ch:model-free-rl} 章的 Double Q-Learning 嫁接到现成的双网络 $(\theta, \theta^-)$ 上——用在线网络 $\theta$ \emph{选动作}，用目标网络 $\theta^-$ \emph{评价值}：
  \[
    y_i^{\mathrm{DDQN}} = r_i + \gamma (1 - d_i)\, Q_{\theta^-}\!\left(s'_i,\; \argmax_{a' \in \mathcal{A}} Q_\theta(s'_i, a')\right).
  \]
  仅改一行代码，便大幅压制了深层网络中的 Q 值高估泡沫。
  \item \textbf{Dueling DQN}（Wang et al., 2016）：将网络末层分叉为状态价值头 $V_\theta(s)$ 与优势头 $A_\theta(s, a)$，并减去动作均值以保证分解唯一可辨识：
  \[
    Q_\theta(s, a) = V_\theta(s) + \left(A_\theta(s, a) - \frac{1}{|\mathcal{A}|}\sum_{a' \in \mathcal{A}} A_\theta(s, a')\right).
  \]
  在许多动作对短期局势影响微弱的状态下，网络不必遍历每个动作也能直接学好整体状态价值 $V_\theta(s)$。
\end{itemize}

\begin{lstlisting}[caption={Double DQN 核心损失函数与目标网络软更新（简洁 PyTorch 实现）}]
def ddqn_loss(q_net, target_net, batch, gamma=0.99):
    s, a, r, ns, done = batch                   # a: (B, 1), done: (B,) in {0.0, 1.0}
    q_sa = q_net(s).gather(1, a).squeeze(-1)    # 当前 Q_theta(s, a)
    with torch.no_grad():                       # 关键：彻底阻断目标分支梯度
        best_a = q_net(ns).argmax(dim=1, keepdim=True)
        next_q = target_net(ns).gather(1, best_a).squeeze(-1)
        target = r + gamma * (1.0 - done) * next_q
    return F.smooth_l1_loss(q_sa, target)       # Huber 损失抑制大 TD 误差震荡
\end{lstlisting}

\section{策略梯度定理：完整严格证明与 REINFORCE}
\label{sec:pg-theorem}

\subsection{为什么需要直接优化策略 $\pi_\theta(a \mid s)$？}

尽管 DQN 家族在 Atari 离散动作游戏中大放异彩，价值学习范式在面对现代 AI 的两大主战场——\textbf{连续物理控制}与\textbf{大语言模型生成}——时却捉襟见肘：
\begin{enumerate}
  \item \textbf{高维或连续动作空间中无法求解 $\max_a Q_\theta(s, a)$}：如果机械臂动作是 $\R^{20}$ 的连续力矩向量，或大语言模型每步需从 $10^5$ 词表中自回归生成整段推理链，每步训练都去对一个非凸神经网络求 $\argmax_a$ 在计算上完全不可行。
  \item \textbf{无法自然表示最优随机策略}：在剪刀石头布或存在感知混淆的部分可观测迷宫中，任何确定性贪心策略都会被对手读穿或陷入死循环，唯有随机策略才是最优解。
  \item \textbf{策略优化的平滑性}：在价值法中，网络参数 $\theta$ 的微小扰动可能导致 $\argmax_a Q_\theta(s, a)$ 发生剧烈跳变；而直接参数化平滑概率分布 $\pi_\theta(a \mid s)$（如离散 Softmax 或连续高斯分布 $\mathcal{N}(\mu_\theta(s), \Sigma_\theta)$）时，策略概率随 $\theta$ 连续可微，天然契合梯度上升。
\end{enumerate}

\subsection{策略梯度定理的完整推导}

设参数化随机策略为 $\pi_\theta(a \mid s)$。在有限视野 $T$ 的回合制 MDP 中，一条完整轨迹 $\tau = (s_0, a_0, r_0, \dots,\allowbreak s_{T-1}, a_{T-1}, r_{T-1}, s_T)$ 由初始分布 $\rho_0(s_0)$、策略 $\pi_\theta$ 与环境转移核 $P$ 共同生成，其联合概率密度/质量为
\begin{equation}
  p_\theta(\tau) = \rho_0(s_0) \prod_{t=0}^{T-1} \pi_\theta(a_t \mid s_t)\, P(s_{t+1} \mid s_t, a_t).
  \label{eq:traj-prob}
\end{equation}
记整条轨迹的总（折扣）回报为 $R(\tau) = \sum_{t=0}^{T-1} \gamma^t r_t$（不失一般性，在无折扣有限视野下取 $\gamma = 1$）。策略优化的目标是最大化期望回报：
\[
  J(\theta) = \E_{\tau \sim p_\theta}[R(\tau)] = \int p_\theta(\tau)\, R(\tau)\,\mathrm{d}\tau.
\]

初看之下，求 $\nabla_\theta J(\theta)$ 似乎是不可能完成的任务：参数 $\theta$ 不仅直接决定策略如何选动作，还通过未知的黑盒环境转移核 $P$ 间接改变了未来访问到哪些状态！但下面这条被誉为现代强化学习“核心引擎”的定理证明：\textbf{未知的 $\rho_0$ 与 $P$ 会在对数求导的瞬间奇迹般地全部抵消！}

\begin{thmBox}
\begin{theorem}[策略梯度定理（Policy Gradient Theorem）与因果性形式]
\label{thm:pg-theorem}
假设对所有合法轨迹 $\tau$，$\pi_\theta(a_t \mid s_t) > 0$ 且关于 $\theta$ 连续可微，并满足积分与求导可交换的常规正则条件。则目标函数 $J(\theta) = \E_{\tau \sim p_\theta}[R(\tau)]$ 的精确梯度为：
\begin{equation}
  \nabla_\theta J(\theta)
  = \E_{\tau \sim p_\theta}\!\left[\left(\sum_{t=0}^{T-1} \nabla_\theta \ln \pi_\theta(a_t \mid s_t)\right) R(\tau)\right].
  \label{eq:pg-raw}
\end{equation}
进一步利用时间因果性（当前动作 $a_t$ 无法影响过去时刻 $t' < t$ 已发生的奖励 $r_{t'}$），式~\eqref{eq:pg-raw} 可精确化简为\textbf{未来回报（Reward-to-Go）形式}与\textbf{动作价值形式}：
\begin{align}
  \nabla_\theta J(\theta)
  &= \E_{\tau \sim p_\theta}\!\left[\sum_{t=0}^{T-1} \nabla_\theta \ln \pi_\theta(a_t \mid s_t)\, G_t\right]
  \label{eq:pg-rtg} \\
  &= \E_{\tau \sim p_\theta}\!\left[\sum_{t=0}^{T-1} \nabla_\theta \ln \pi_\theta(a_t \mid s_t)\, Q^{\pi_\theta}(s_t, a_t)\right],
  \label{eq:pg-q}
\end{align}
其中 $G_t = \sum_{k=t}^{T-1} \gamma^{k-t} r_k$（当目标本身带 $\gamma^t$ 折扣时，式~\eqref{eq:pg-rtg} 中 $G_t$ 外层对应因子 $\gamma^t$，工程中常省略该折扣因子以保持各步更新尺度一致）。
\end{theorem}
\begin{proof}
\textbf{第一步：对数导数技巧（Log-Derivative Trick / Score Function）}。由于 $(\ln f(\theta))' = \frac{f'(\theta)}{f(\theta)}$，恒有 $\nabla_\theta p_\theta(\tau) = p_\theta(\tau)\,\nabla_\theta \ln p_\theta(\tau)$。将其代入目标梯度：
\begin{align*}
  \nabla_\theta J(\theta)
  &= \nabla_\theta \int p_\theta(\tau)\, R(\tau)\,\mathrm{d}\tau
   = \int \nabla_\theta p_\theta(\tau)\, R(\tau)\,\mathrm{d}\tau \\
  &= \int p_\theta(\tau)\, \nabla_\theta \ln p_\theta(\tau)\, R(\tau)\,\mathrm{d}\tau
   = \E_{\tau \sim p_\theta}\!\left[\nabla_\theta \ln p_\theta(\tau)\, R(\tau)\right].
\end{align*}

\textbf{第二步：环境动力学彻底消去}。对式~\eqref{eq:traj-prob} 两边取自然对数，乘积化为求和：
\[
  \ln p_\theta(\tau) = \underbrace{\ln \rho_0(s_0)}_{\text{不含 }\theta}
  + \sum_{t=0}^{T-1} \ln \pi_\theta(a_t \mid s_t)
  + \sum_{t=0}^{T-1} \underbrace{\ln P(s_{t+1} \mid s_t, a_t)}_{\text{不含 }\theta}.
\]
对 $\theta$ 求梯度，所有与环境机制有关的 $\ln \rho_0$ 与 $\ln P$ 项求导恒为零，只剩下策略自身：
\[
  \nabla_\theta \ln p_\theta(\tau) = \sum_{t=0}^{T-1} \nabla_\theta \ln \pi_\theta(a_t \mid s_t).
\]
代回第一步，即得式~\eqref{eq:pg-raw}！这解释了为什么我们\textbf{完全不需要环境可微，甚至无需知晓环境模型}，仅凭采样轨迹就能无偏估计精确的策略梯度。

\textbf{第三步：利用因果性削去过去奖励（令 $\gamma = 1$ 为例，带 $\gamma$ 时同理）}。将式~\eqref{eq:pg-raw} 按双重求和展开：
\[
  \nabla_\theta J(\theta) = \sum_{t=0}^{T-1} \sum_{t'=0}^{T-1} \E_{\tau \sim p_\theta}\!\left[\nabla_\theta \ln \pi_\theta(a_t \mid s_t)\, r_{t'}\right].
\]
考察任意满足 $t' < t$ 的“过去奖励项”。记时刻 $t$ 决策前的历史为 $h_t = (s_0, a_0, r_0, \dots, s_t)$，则 $r_{t'}$ 在给定 $h_t$ 下已是完全确定的常量。应用重期望公式（先对 $a_t \sim \pi_\theta(\cdot \mid s_t)$ 求条件期望）：
\begin{align*}
  \E_{\tau}\!\left[\nabla_\theta \ln \pi_\theta(a_t \mid s_t)\, r_{t'}\right]
  &= \E_{h_t}\!\left[r_{t'}\, \E_{a_t \sim \pi_\theta(\cdot \mid s_t)}\!\left[\nabla_\theta \ln \pi_\theta(a_t \mid s_t)\right]\right] \\
  &= \E_{h_t}\!\left[r_{t'} \sum_{a_t} \pi_\theta(a_t \mid s_t) \frac{\nabla_\theta \pi_\theta(a_t \mid s_t)}{\pi_\theta(a_t \mid s_t)}\right] \\
  &= \E_{h_t}\!\left[r_{t'}\, \nabla_\theta \!\left(\sum_{a_t} \pi_\theta(a_t \mid s_t)\right)\right]
   = \E_{h_t}\!\left[r_{t'} \cdot \nabla_\theta(1)\right] = 0.
\end{align*}
因此所有 $t' < t$ 的项在期望下精确为零，只剩 $t' \ge t$ 的未来回报之和 $G_t = \sum_{t'=t}^{T-1} r_{t'}$，即证式~\eqref{eq:pg-rtg}。

\textbf{第四步：条件期望化为 $Q^{\pi_\theta}$}。在式~\eqref{eq:pg-rtg} 中对每项条件于 $(s_t, a_t)$ 再取重期望，由定义 $\E[G_t \mid s_t, a_t] = Q^{\pi_\theta}(s_t, a_t)$，立得式~\eqref{eq:pg-q}。
\end{proof}
\end{thmBox}

用单批采样轨迹 $\tau^{(1)}, \dots, \tau^{(M)}$ 的经验均值替换式~\eqref{eq:pg-rtg} 中的期望，执行梯度上升
\[
  \theta \leftarrow \theta + \alpha \cdot \frac{1}{M}\sum_{i=1}^M \sum_{t=0}^{T-1} \nabla_\theta \ln \pi_\theta\!\left(a_t^{(i)} \mid s_t^{(i)}\right) G_t^{(i)},
\]
就得到了 Williams（1992）提出的经典蒙特卡洛策略梯度算法——\textbf{REINFORCE}。

它的直观几何含义异常迷人：$\nabla_\theta \ln \pi_\theta(a_t \mid s_t)$ 指向“提高动作 $a_t$ 出现对数概率”的方向。若某条轨迹的未来回报 $G_t$ 高，我们就沿着该方向大力拉升 $a_t$ 的概率；若 $G_t$ 为负，就压低 $a_t$ 的概率。本质上，REINFORCE 就像一场\textbf{由未来回报 $G_t$ 加权的极大似然分类学习}！

\section{基线零偏降方差定理与优势函数}
\label{sec:pg-baseline}

然而，朴素 REINFORCE 在实际训练中常常慢如蜗牛甚至剧烈抖动。究其原因，看看下面这个极其简单的思想实验：

假设某任务中无论你做得好坏，每条轨迹的奖励都整体平移了 $+1000$（好动作拿 $+1002$，烂动作拿 $+998$）。虽然最优策略丝毫未变，但朴素 REINFORCE 看到 $+998$ 是个巨大的正数，仍会疯狂\emph{调高}烂动作的概率！它只能寄希望于好动作被拉升得比烂动作“稍微多那么 $0.4\%$”，而在有限样本采样噪声面前，这种微弱差值瞬间就被 $+1000$ 带来的巨大梯度方差彻底淹没！

解决这一痛点的关键无比优雅：\textbf{从 $G_t$ 中减去一个仅依赖状态 $s_t$ 的基线（Baseline）$b(s_t)$}，改看动作“比平均水平好多少”！我们必须先证明：这个减法\textbf{丝毫不会引入任何梯度偏差}。

\begin{thmBox}
\begin{theorem}[状态基线的零偏性（Zero-Bias Baseline Theorem）]
\label{thm:baseline-unbiased}
设 $b: \mathcal{S} \to \R$ 是\textbf{仅依赖于状态 $s_t$（而不依赖于当前动作 $a_t$）}的任意确定性函数。则对任意状态 $s_t$ 均有
\begin{equation}
  \E_{a_t \sim \pi_\theta(\cdot \mid s_t)}\!\left[\nabla_\theta \ln \pi_\theta(a_t \mid s_t)\, b(s_t)\right] = 0.
  \label{eq:baseline-zero}
\end{equation}
因而带基线的策略梯度与原策略梯度完全相等（严格无偏）：
\begin{equation}
  \nabla_\theta J(\theta)
  = \E_{\tau \sim p_\theta}\!\left[\sum_{t=0}^{T-1} \nabla_\theta \ln \pi_\theta(a_t \mid s_t)\,\big(G_t - b(s_t)\big)\right].
  \label{eq:pg-baseline}
\end{equation}
\end{theorem}
\begin{proof}
由于 $b(s_t)$ 不含动作 $a_t$，在关于 $a_t \sim \pi_\theta(\cdot \mid s_t)$ 求期望时可直接提至期望之外：
\begin{align*}
  \E_{a_t \sim \pi_\theta(\cdot \mid s_t)}\!\left[\nabla_\theta \ln \pi_\theta(a_t \mid s_t)\, b(s_t)\right]
  &= b(s_t) \sum_{a_t \in \mathcal{A}} \pi_\theta(a_t \mid s_t)\, \frac{\nabla_\theta \pi_\theta(a_t \mid s_t)}{\pi_\theta(a_t \mid s_t)} \\
  &= b(s_t)\, \nabla_\theta \!\left(\sum_{a_t \in \mathcal{A}} \pi_\theta(a_t \mid s_t)\right)
   = b(s_t)\, \nabla_\theta(1) = 0.
\end{align*}
（对于连续动作空间，将求和换为积分 $\int \pi_\theta(a_t\mid s_t)\,\mathrm{d}a_t = 1$，证明完全相同。）再由重期望公式对外层轨迹分布取期望，即知减去 $b(s_t)$ 后式~\eqref{eq:pg-baseline} 的期望与式~\eqref{eq:pg-rtg} 精确相等。
\end{proof}
\end{thmBox}

那么，什么是最好的基线 $b(s_t)$？在状态 $s_t$ 处，我们可以对梯度估计量 $g_b = \nabla_\theta \ln \pi_\theta(a \mid s)(G - b)$ 的方差求极小值（见下面命题与习题~\textbf{6.2.4}）：当一维参数导数记作 $u(a) = \partial_\theta \ln \pi_\theta(a\mid s)$ 时，方差
\[
  \Var(g_b) = \E\!\left[u(a)^2 (G - b)^2\right] - \big(\E[u(a) G]\big)^2
\]
关于 $b$ 求导置零，得理论最优基线为梯度平方加权期望 $b^* = \frac{\E[u(a)^2 G]}{\E[u(a)^2]}$。在多维参数下逐维加权过于繁琐，而\textbf{直接取状态价值函数 $b(s_t) = V^{\pi_\theta}(s_t) = \E[G_t \mid s_t]$} 既与 $b^*$ 极为接近，又催生了强化学习中最核心的评价量——\textbf{优势函数（Advantage Function）}：
\begin{equation}
  A^{\pi_\theta}(s_t, a_t) := Q^{\pi_\theta}(s_t, a_t) - V^{\pi_\theta}(s_t).
  \label{eq:advantage-def}
\end{equation}
将 $b(s_t) = V^{\pi_\theta}(s_t)$ 代入式~\eqref{eq:pg-q}，策略梯度呈现出最富直觉的黄金形态：
\[
  \nabla_\theta J(\theta) = \E_{\tau \sim p_\theta}\!\left[\sum_{t=0}^{T-1} \nabla_\theta \ln \pi_\theta(a_t \mid s_t)\, A^{\pi_\theta}(s_t, a_t)\right].
\]
此时，无论奖励被整体平移 $+1000$ 还是 $-1000$，常量偏移会在 $Q - V$ 中精确抵消！唯有表现\textbf{优于当前状态平均水平}的动作（$A > 0$）被奖励提权，表现\textbf{逊于平均水平}的动作（$A < 0$）被惩罚降权。

\section{Actor-Critic 架构与广义优势估计 GAE($\lambda$)}
\label{sec:actor-critic-gae}

\subsection{从蒙特卡洛基线迈向自举 Actor-Critic}

在带基线的 REINFORCE 中，我们额外训练一个价值网络 $V_\phi(s)$ 拟合经验回报 $G_t$（最小化 $\frac{1}{2}(V_\phi(s_t) - G_t)^2$），并用 $G_t - V_\phi(s_t)$ 估计优势。这虽已大幅降低由奖励偏移带来的方差，但 $G_t$ 本身仍是整条完整轨迹的蒙特卡洛回报——依然受困于长视野累积噪声且无法每步在线更新。

如果我们在优势估计中\textbf{彻底引入第~\ref{ch:model-free-rl} 章的 TD 自举思想}，用单步观测加后继价值估计 $r_t + \gamma V_\phi(s_{t+1})$ 替代完整回报 $G_t$，会得到什么？正是单步 TD 误差：
\begin{equation}
  \delta_t^V := r_t + \gamma V_\phi(s_{t+1}) - V_\phi(s_t).
  \label{eq:ac-td-error}
\end{equation}
检验一下它在 $V_\phi = V^{\pi_\theta}$ 时的条件期望：
\[
  \E\!\left[\delta_t^V \mid s_t, a_t\right]
  = \E\!\left[r_t + \gamma V^{\pi_\theta}(s_{t+1}) \mid s_t, a_t\right] - V^{\pi_\theta}(s_t)
  = Q^{\pi_\theta}(s_t, a_t) - V^{\pi_\theta}(s_t) = A^{\pi_\theta}(s_t, a_t).
\]
这是一个令人拍案叫绝的巧合：\textbf{当价值网络准确时，单步 TD 误差 $\delta_t^V$ 本身就是优势函数 $A^{\pi_\theta}(s_t, a_t)$ 的无偏估计，而且完全不需要额外拟合高维的 $Q_\phi(s, a)$，只需学一个标量输出的状态价值网 $V_\phi(s)$ 即可！}

由此诞生了现代深度强化学习的双核引擎——\textbf{Actor-Critic（演员--评论家）架构}（见图~\ref{fig:actor-critic}）：
\begin{itemize}
  \item \textbf{Actor（演员 $\pi_\theta$）}：负责根据状态 $s_t$ 输出策略分布并采样动作 $a_t$，沿着 Critic 给出的优势信号 $\hat{A}_t$ 执行策略梯度上升 $\theta \leftarrow \theta + \alpha_\theta \nabla_\theta \ln \pi_\theta(a_t \mid s_t)\, \hat{A}_t$。
  \item \textbf{Critic（评论家 $V_\phi$）}：负责评估每个状态的长期价值，利用时序差分目标更新自身参数 $\phi$，从而为 Actor 持续提供低方差、可在线步进的“打分反馈”。
\end{itemize}

\begin{figure}[htbp]
\centering
\resizebox{0.92\linewidth}{!}{%
\begin{tikzpicture}[
  box/.style={rectangle, rounded corners=2mm, draw=#1, very thick, fill=#1!10,
              minimum width=3.6cm, minimum height=1.25cm, align=center, font=\sffamily\small, text=ink},
  env/.style={rectangle, rounded corners=2mm, draw=slate, very thick, fill=codebg,
              minimum width=4.0cm, minimum height=1.25cm, align=center, font=\sffamily\small, text=ink},
  arr/.style={-{Stealth[length=2.3mm]}, thick, #1},
  lbl/.style={font=\scriptsize\sffamily, color=ink, fill=white, inner sep=1.5pt}
]
  \node[box=teal] (actor) at (-4.25, 1.55) {\textbf{\color{tealDark}Actor（演员）}\\策略网络 $\pi_\theta(a \mid s)$};
  \node[box=navy] (critic) at (4.25, 1.55) {\textbf{\color{navy}Critic（评论家）}\\价值网络 $V_\phi(s)$};
  \node[env] (env) at (0, -1.65) {\textbf{环境（Environment）}\\动态转移 $P(s' \mid s, a)$ 与奖励 $r_t$};

  % Actor -> Env (action a_t)
  \draw[arr=teal] (actor.south) |- ([yshift=2mm]env.west)
    node[lbl, pos=0.23, right=1mm] {采样动作 $a_t \sim \pi_\theta(\cdot \mid s_t)$};

  % Env -> Critic (r_t, s_{t+1})
  \draw[arr=slate] ([yshift=2mm]env.east) -| (critic.south)
    node[lbl, pos=0.75, right=1mm] {反馈奖励 $r_t$ 与新状态 $s_{t+1}$};

  % Env -> Actor (state s_t)
  \draw[arr=slate, dashed] ([yshift=-2.5mm]env.west) -- ([xshift=-5mm,yshift=-2.5mm]actor.west |- env.west) |- (actor.west)
    node[lbl, pos=0.25, left=1mm] {观测状态 $s_t$};

  % Critic -> Actor (TD error / GAE advantage)
  \draw[arr=ochre, very thick] (critic.west) -- (actor.east)
    node[lbl, midway, above=2pt, text=ochre] {\textbf{优势信号 $\delta_t$ / $\hat{A}_t^{\mathrm{GAE}}$（\texttt{.detach()}）}}
    node[lbl, midway, below=2pt, text=slate] {驱动策略梯度 $\nabla_\theta \ln \pi_\theta \cdot \hat{A}_t$};

  % Critic self-update loop
  \draw[arr=navy, dashed] (critic.north) .. controls +(0, 0.85) and +(1.3, 0.85) .. (critic.north east)
    node[lbl, pos=0.5, above=1.5pt, text=navy] {最小化 TD / 回归误差更新 $\phi$};
\end{tikzpicture}}
\caption{Actor-Critic 闭环反馈架构。Actor 负责与环境交互决策，Critic 实时提供价值评估并产出 TD 误差 $\delta_t$ 或广义优势估计 $\hat{A}_t^{\mathrm{GAE}}$；注意传给 Actor 的优势信号必须阻断梯度（\texttt{.detach()}），以免策略损失反向污染 Critic 参数。}
\label{fig:actor-critic}
\end{figure}

\subsection{广义优势估计 GAE($\lambda$) 的严格推导}

单步 TD 误差 $\delta_t^V$ 虽然方差极低，但在训练前期 Critic 尚未准确（$V_\phi \neq V^\pi$）时会引入较大自举偏差；反之，完整回报减基线 $G_t - V_\phi(s_t)$ 虽然严格无偏，却方差高企。有没有一根平滑旋钮，能让我们在“单步 TD”与“完整 MC”之间自由调节偏差与方差？

Schulman 等人（2016）提出的\textbf{广义优势估计（Generalized Advantage Estimation, GAE）}给出了堪称完美的答案。对任意步数 $k \ge 1$，定义 $k$-步截断优势估计量：
\begin{align*}
  \hat{A}_t^{(1)} &= \delta_t^V = -V_\phi(s_t) + r_t + \gamma V_\phi(s_{t+1}), \\
  \hat{A}_t^{(2)} &= \delta_t^V + \gamma \delta_{t+1}^V = -V_\phi(s_t) + r_t + \gamma r_{t+1} + \gamma^2 V_\phi(s_{t+2}), \\
  &\;\;\vdots \\
  \hat{A}_t^{(k)} &= \sum_{l=0}^{k-1} \gamma^l \delta_{t+l}^V = -V_\phi(s_t) + \sum_{l=0}^{k-1} \gamma^l r_{t+l} + \gamma^k V_\phi(s_{t+k}).
\end{align*}
随着 $k$ 增大，偏差因末端项的 $\gamma^k$ 因子呈指数衰减，而方差因累加更多随机奖励而递增；当 $k \to \infty$ 时，$\hat{A}_t^{(\infty)} = G_t - V_\phi(s_t)$ 退化为无偏的带基线蒙特卡洛优势。

\begin{thmBox}
\begin{theorem}[GAE($\gamma, \lambda$) 的几何加权展开与递推公式]
\label{thm:gae}
对衰减参数 $\lambda \in [0, 1]$，将所有 $k$-步优势估计量按几何权重 $(1-\lambda)\lambda^{k-1}$ 加权平均，定义广义优势估计 $\hat{A}_t^{\mathrm{GAE}(\gamma,\lambda)} := (1 - \lambda)\sum_{k=1}^{\infty} \lambda^{k-1} \hat{A}_t^{(k)}$。则它精确化简为单步 TD 误差的指数衰减加权和：
\begin{equation}
  \hat{A}_t^{\mathrm{GAE}(\gamma,\lambda)} = \sum_{l=0}^{\infty} (\gamma \lambda)^l \delta_{t+l}^V,
  \label{eq:gae-sum}
\end{equation}
并在有限视野回合中满足 $O(T)$ 时间复杂度的反向线性递推：
\begin{equation}
  \hat{A}_t^{\mathrm{GAE}} = \delta_t^V + \gamma \lambda (1 - d_t)\, \hat{A}_{t+1}^{\mathrm{GAE}} \qquad (\hat{A}_T^{\mathrm{GAE}} = 0).
  \label{eq:gae-rec}
\end{equation}
\end{theorem}
\begin{proof}
将 $\hat{A}_t^{(k)} = \sum_{l=0}^{k-1} \gamma^l \delta_{t+l}^V$ 代入定义并交换二重求和次序（下标满足 $0 \le l \le k-1 \iff k \ge l+1$）：
\begin{align*}
  \hat{A}_t^{\mathrm{GAE}(\gamma,\lambda)}
  &= (1 - \lambda) \sum_{k=1}^{\infty} \lambda^{k-1} \sum_{l=0}^{k-1} \gamma^l \delta_{t+l}^V \\
  &= (1 - \lambda) \sum_{l=0}^{\infty} \gamma^l \delta_{t+l}^V \left(\sum_{k=l+1}^{\infty} \lambda^{k-1}\right) \\
  &= (1 - \lambda) \sum_{l=0}^{\infty} \gamma^l \delta_{t+l}^V \cdot \frac{\lambda^l}{1 - \lambda}
   = \sum_{l=0}^{\infty} (\gamma \lambda)^l \delta_{t+l}^V.
\end{align*}
从式~\eqref{eq:gae-sum} 提出首项 $l=0$，对余项令 $l' = l - 1$，立即得到递推式 $\hat{A}_t^{\mathrm{GAE}} = \delta_t^V + \gamma\lambda \hat{A}_{t+1}^{\mathrm{GAE}}$；若 $s_{t+1}$ 为回合终止态（$d_t = 1$），则跨回合不再传导，乘以 $(1 - d_t)$ 截断即得式~\eqref{eq:gae-rec}。
\end{proof}
\end{thmBox}

检查两个极端端点，尽显数学统一之美：
\begin{itemize}
  \item 当 $\lambda = 0$ 时（约定 $0^0 = 1$），$\hat{A}_t^{\mathrm{GAE}(\gamma, 0)} = \delta_t^V = r_t + \gamma V_\phi(s_{t+1}) - V_\phi(s_t)$，完全退化为单步 TD Actor-Critic（低方差、有自举偏差）；
  \item 当 $\lambda = 1$ 时，$\hat{A}_t^{\mathrm{GAE}(\gamma, 1)} = \sum_{l=0}^{\infty} \gamma^l \delta_{t+l}^V = G_t - V_\phi(s_t)$（中间全部 $V_\phi(s_{t+l})$ 首尾相消！），完全退化为无偏的蒙特卡洛基线优势（零偏差、高方差）。
\end{itemize}
现代深度强化学习（包括第~\ref{ch:trpo-ppo-grpo} 章的 PPO 与第 4 章的 RLHF）几乎普遍选用 $\lambda \in [0.90, 0.98]$，它在几乎不损失无偏性的同时大幅压平了策略梯度方差。

\section{实测验证：基线降方差与 GAE Actor-Critic}
\label{sec:exp-pg-baseline}

理论告诉我们：当奖励带有正偏移时，无基线 REINFORCE 会因巨大方差陷入泥潭，而引入价值基线与 GAE 应立竿见影地提速稳盘。下面我们在配套脚本 \texttt{code/vol6\_experiments.py} 的 \texttt{exp2\_pg\_baseline\_ac()} 中进行严格对照实测：
\begin{itemize}
  \item \textbf{任务环境}：12 步随机一维目标跟踪控制。每回合初始位置 $x_0 \sim \mathrm{Uniform}(-1, 1)$；智能体每步选择向左（$-0.25$）或向右（$+0.25$）并叠加高斯扰动 $\xi_t \sim \mathcal{N}(0, 0.04^2)$，位置截断在 $[-1.5, 1.5]$。
  \item \textbf{正偏移奖励设计}：$r_t = \mathbf{4.0} + (1.0 - |x_{t+1}|)$，折扣因子 $\gamma = 0.95$。注意每步的 $+4.0$ 偏移不改变最优控制律（始终向原点 $0$ 靠拢），却把总回报推高到 $50$--$58$ 区间，专门用于放大检验基线的降方差威力。
  \item \textbf{网络与统计}：Actor 与 Critic 均为单隐层 MLP（$2 \to 16 \to \text{Tanh} \to \text{out}$），Adam 优化器（Actor 学习率 $0.015$，Critic 学习率 $0.03$），全部 3 种算法均跨 12 组独立随机种子运行 250 回合。
\end{itemize}

\begin{lstlisting}[caption={REINFORCE+基线与 Actor-Critic+GAE($\lambda=0.9$) 核心更新（摘自 \texttt{vol6\_experiments.py}）}]
if mode == "reinforce":
    loss_p = -(logps_t * Gs_t).mean()
    opt_p.zero_grad(); loss_p.backward(); opt_p.step()
elif mode == "baseline":
    vals = val_net(states_t).squeeze(-1)
    adv = Gs_t - vals.detach()            # 必须 .detach()，防止策略梯度回流 Critic
    loss_p = -(logps_t * adv).mean()
    loss_v = F.mse_loss(vals, Gs_t)
    opt_p.zero_grad(); loss_p.backward(); opt_p.step()
    opt_v.zero_grad(); loss_v.backward(); opt_v.step()
else:  # Actor-Critic + GAE(gamma=0.95, lambda=0.90)
    vals = val_net(states_t).squeeze(-1)
    next_vals = torch.cat([vals[1:].detach(), torch.zeros(1)])
    deltas = rews_t + 0.95 * next_vals - vals.detach()
    advs, gae = [], 0.0
    for d in reversed(deltas.tolist()):   # 按式 (2.10) 反向递推计算 GAE
        gae = d + 0.95 * 0.90 * gae
        advs.insert(0, gae)
    adv = torch.tensor(advs, dtype=torch.float32)
    loss_p = -(logps_t * adv).mean()
    loss_v = F.mse_loss(vals, adv + vals.detach())  # lambda-回报目标
    opt_p.zero_grad(); loss_p.backward(); opt_p.step()
    opt_v.zero_grad(); loss_v.backward(); opt_v.step()
\end{lstlisting}

图~\ref{fig:pg-baseline} 直接读取实测记录文件 \texttt{data/pg\_baseline.csv}，对比了 250 回合内三种方法的平均总回报轨迹：
\begin{itemize}
  \item \textbf{无基线 REINFORCE}：起步阶段虽略有爬升，但从第 70 回合起就因 $+4.0$ 奖励偏移引起的剧烈梯度方差在 $52$ 附近原地打转；最后 30 回合均值仅停留在 \textbf{52.16（$\pm 5.02$）}，跨种子波动极大。
  \item \textbf{REINFORCE $+$ 价值基线}：减去 $V_\phi(s_t)$ 后，正偏移噪声瞬间消失，策略稳定收敛至原点附近，最后 30 回合均值跃升至 \textbf{56.13（$\pm 2.55$）}——标准差直接腰斩，意味着\textbf{方差被压缩了近 $4\times$}（$5.02^2 / 2.55^2 \approx 3.88$）！
  \item \textbf{Actor-Critic $+$ GAE($\lambda=0.9$)}：凭借低方差的时序差分优势信号，前中期爬升速度最快（第 115 回合即达 $54.59$，领先纯 MC 基线 $0.65$ 分），最后 30 回合均值稳定在 \textbf{55.80（$\pm 3.16$）}，且天然具备向更长视野、持续交互任务无缝扩展的能力。
\end{itemize}

\begin{figure}[htbp]
\centering
\begin{tikzpicture}
  \begin{axis}[
    width=0.90\linewidth, height=6.4cm,
    xlabel={训练回合数（Episode）},
    ylabel={回合总回报（10 回合均值，12 种子）},
    xmin=0, xmax=260, ymin=50.0, ymax=57.0,
    grid=both, grid style={tealLine!45, dashed},
    axis line style={slate},
    tick label style={font=\scriptsize},
    label style={font=\footnotesize\sffamily, color=ink},
    legend style={at={(0.03,0.96)}, anchor=north west, font=\scriptsize, fill=white, draw=tealLine}
  ]
    \addplot[ochre, very thick, dashed, mark=triangle*, mark size=1.8pt]
      table[x=ep, y=reinforce, col sep=comma]{data/pg_baseline.csv};
    \addlegendentry{REINFORCE 无基线（终段 $52.16 \pm 5.02$）}

    \addplot[teal, very thick, mark=*, mark size=1.6pt]
      table[x=ep, y=baseline, col sep=comma]{data/pg_baseline.csv};
    \addlegendentry{REINFORCE $+$ 价值基线（终段 $56.13 \pm 2.55$）}

    \addplot[navy, very thick, mark=square*, mark size=1.5pt]
      table[x=ep, y=ac, col sep=comma]{data/pg_baseline.csv};
    \addlegendentry{Actor-Critic $+$ GAE($\lambda=0.9$)（终段 $55.80 \pm 3.16$）}
  \end{axis}
\end{tikzpicture}
\caption{带 $+4.0$ 正偏移奖励的 12 步随机目标跟踪实测曲线（读取自 \texttt{data/pg\_baseline.csv}）。减去状态价值基线将标准差减半（方差压缩约 $4\times$），彻底释放了策略梯度的收敛潜力；GAE($\lambda=0.9$) 在前中期收敛最快。}
\label{fig:pg-baseline}
\end{figure}

\begin{failBox}{策略梯度与 Actor-Critic 编码中的两大高频致命错误}
\begin{enumerate}
  \item \textbf{忘记对优势项 \texttt{adv}（或 \texttt{vals}）调用 \texttt{.detach()}}：在 Actor 损失 \texttt{loss\_p = -(logps * (Gs - vals)).mean()} 中，如果漏写了 \texttt{vals.detach()}，并且把 Actor 和 Critic 合并在同一个 \texttt{loss\_total = loss\_p + 0.5 * loss\_v} 里一起反向传播，那么 \texttt{loss\_p} 会产生一条流向 Critic 参数 $\phi$ 的非法梯度 $+\ln\pi_\theta \cdot \nabla_\phi V_\phi(s)$！它试图通过扭曲 $V_\phi$ 来人为压低伪损失，彻底毁掉价值估计。务必始终保证进入 Actor 损失的优势张量处于 \texttt{no\_grad} 或 \texttt{.detach()} 状态。
  \item \textbf{误用依赖动作 $a_t$ 的基线 $b(s_t, a_t)$}：回看定理~\ref{thm:baseline-unbiased} 的证明——关键一步是将 $b(s_t)$ 提到 $\sum_{a_t}$ 之外！如果基线同时依赖动作 $a_t$（例如误把 $Q_\phi(s_t, a_t)$ 直接当作基线减去，却不补加对应的解析修正项），则 $\sum_{a_t} \nabla_\theta \pi_\theta(a_t\mid s_t)\, b(s_t, a_t) \neq 0$，\textbf{无偏性瞬间破裂}，策略梯度可能指向完全错误的方向！
\end{enumerate}
\end{failBox}

\begin{extBox}{熵正则化与最大熵强化学习（Soft Actor-Critic, SAC）}
在实际训练中，策略梯度容易出现\textbf{过早熵坍缩（Premature Entropy Collapse）}：一旦策略在早期偶遇一条还不错的轨迹，梯度上升会迅速把该动作的概率推近 $1$，策略退化为确定性策略，彻底丧失后续探索更优解的能力。

最简单有效的对策是在策略目标中加入\textbf{策略熵正则项}：
\[
  J_{\mathrm{ent}}(\theta) = \E_{\tau \sim p_\theta}\!\left[\sum_{t=0}^{T-1} \gamma^t \Big(r_t + \alpha\, \mathcal{H}\big(\pi_\theta(\cdot \mid s_t)\big)\Big)\right], \qquad
  \mathcal{H}(\pi) = -\E_{a \sim \pi}[\ln \pi(a \mid s)],
\]
其中温度系数 $\alpha > 0$ 控制探索权重。它鼓励智能体在“拿高分”的同时尽可能保持行为的多样性与随机性。

进一步将这一最大熵目标与\textbf{异策略双 Q 网络（Clipped Double Q，取 $\min(Q_{\phi_1}, Q_{\phi_2})$ 抑制高估）}以及\textbf{重参数化技巧（Reparameterization Trick，$a = \tanh(\mu_\theta(s) + \sigma_\theta(s)\odot \varepsilon)$）}结合，便诞生了连续控制领域的王者算法——\textbf{Soft Actor-Critic（SAC，Haarnoja et al., 2018）}。SAC 既有异策略经验回放带来的极高样本效率，又因最大熵目标而对超参数高度稳健，是当今机器人仿真与真机连续控制的绝对主力之一。
\end{extBox}

\section*{本章小结与核心脉络}
\addcontentsline{toc}{section}{本章小结与核心脉络}

本章搭建起现代深度强化学习的两座主桥：
\begin{itemize}
  \item \textbf{深度价值路线（DQN 家族）}：通过两状态线性反例揭示“函数逼近 $+$ 自举 $+$ 异策略”致命三元组的发散根源；用经验回放打散时间相关性、用冻结目标网络切断追影正反馈，再以 Double DQN 与 Dueling DQN 精修偏差与结构。
  \item \textbf{深度策略路线（Policy Gradient $\to$ Actor-Critic $\to$ GAE）}：利用对数导数技巧完整证明了策略梯度定理，让黑盒环境动力学 $P(s'\mid s,a)$ 精确消去；证明了任意状态基线的严格零偏性，引出优势函数 $A^\pi = Q^\pi - V^\pi$（实测将回报方差压缩近 $4\times$）；最后通过广义优势估计 $\hat{A}_t^{\mathrm{GAE}} = \sum_{l=0}^\infty (\gamma\lambda)^l \delta_{t+l}^V$ 优雅统一了低方差 TD 与零偏差 MC。
\end{itemize}

然而，Actor-Critic 与 REINFORCE 共同背负着一个沉重的效率枷锁：\textbf{它们是严格同策略的——辛辛苦苦与环境交互采来的一批轨迹，只敢做一步梯度上升，然后就必须全部扔进垃圾桶！} 如果我们试图把同一批样本拿来重复更新 10 轮，会发生什么？下一章，我们将直面策略崩塌的悬崖，推导出统治现代深度强化学习与大语言模型后训练的两大核心算法——\textbf{近端策略优化（PPO）与群体相对策略优化（GRPO）}。

\section*{本章分层习题}
\addcontentsline{toc}{section}{本章分层习题}

\exTier{A 层：概念与公式手算}{夯实策略梯度公式、对数导数与 GAE 递推}

\begin{enumerate}
  \item[\textbf{6.2.1}] \textbf{Softmax 策略与高斯策略的 Score Function 解析推导}：
  \begin{enumerate}
    \item 设离散 Softmax 策略 $\pi_\theta(a \mid s) = \frac{\exp(\theta^\top \phi(s, a))}{\sum_{a'} \exp(\theta^\top \phi(s, a'))}$，求证其对数梯度（Score Function）具有极其简洁的“特征减期望特征”形式：
    \[
      \nabla_\theta \ln \pi_\theta(a \mid s) = \phi(s, a) - \E_{a' \sim \pi_\theta(\cdot \mid s)}\!\left[\phi(s, a')\right].
    \]
    \item 设一维高斯策略 $\pi_\theta(a \mid s) = \mathcal{N}\!\left(a \mid \mu_\theta(s), \sigma^2\right)$（标准差 $\sigma > 0$ 为常数），求 $\nabla_\theta \ln \pi_\theta(a \mid s)$，并解释其直观更新方向。
  \end{enumerate}
  \exHint{(a) 对 $\ln \pi_\theta(a\mid s) = \theta^\top\phi(s,a) - \ln \sum_{a'}\exp(\theta^\top\phi(s,a'))$ 求梯度即可；(b) $\nabla_\theta \ln \pi_\theta = \frac{a - \mu_\theta(s)}{\sigma^2}\nabla_\theta \mu_\theta(s)$，好动作会把均值拉向 $a$。}

  \item[\textbf{6.2.2}] \textbf{三步轨迹的 GAE($\lambda$) 手算}：某 3 步回合结束于终止态 $s_3$（$V(s_3) = 0$），奖励序列为 $(r_0, r_1, r_2) = (1, 0, 5)$，Critic 估值 $(V(s_0), V(s_1), V(s_2)) = (2.0, 2.5, 4.0)$，取 $\gamma = 0.9$。
  \begin{enumerate}
    \item 计算单步 TD 误差 $\delta_0, \delta_1, \delta_2$；
    \item 分别计算 $\lambda = 0$、$\lambda = 0.5$ 与 $\lambda = 1$ 时三个时刻的 $\hat{A}_t^{\mathrm{GAE}}$，并验证 $\lambda = 1$ 的结果恰好等于 $G_t - V(s_t)$。
  \end{enumerate}
  \exHint{利用式~\eqref{eq:gae-rec} 从 $t=2$ 倒推回 $t=0$，注意 $t=2$ 到达终止态，故 $\delta_2 = r_2 + 0 - V(s_2) = 1.0$ 且 $\hat{A}_2 = \delta_2$。}
\end{enumerate}

\exTier{B 层：严谨数学证明}{掌握最优方差基线与无穷视野策略梯度}

\begin{enumerate}
  \item[\textbf{6.2.3}] \textbf{单步标量策略的最优最小方差基线}：考虑单状态上下文赌博机问题，策略 $\pi_\theta(a)$ 仅含单个标量参数 $\theta \in \R$，记 $u(a) = \frac{\partial}{\partial\theta}\ln \pi_\theta(a)$，随机奖励为 $R(a)$，定义带常数基线 $b$ 的策略梯度估计量 $g(b) = u(a)\big(R(a) - b\big)$。
  \begin{enumerate}
    \item 证明 $\Var(g(b))$ 是关于 $b$ 的严格凸二次函数（假设 $\E[u(a)^2] > 0$）；
    \item 求出使 $\Var(g(b))$ 达到全局极小的最优基线 $b^*$ 的闭合表达式，并说明为何 $b^*$ 一般略微不同于期望奖励 $\E[R(a)]$，但在何种条件下二者恰好重合。
  \end{enumerate}
  \exHint{因 $\E[g(b)]$ 与 $b$ 无关，极小化方差等价于极小化二阶矩 $\E[u(a)^2(R(a) - b)^2]$；求导置零得 $b^* = \frac{\E[u(a)^2 R(a)]}{\E[u(a)^2]}$。}

  \item[\textbf{6.2.4}] \textbf{强化学习中的“因果性降方差”定量验证}：在一个两步独立赌博机（$s_0 \to s_1 \to \mathsf{T}$，两步奖励独立且 $r_0 \sim \mathcal{N}(\mu_0, \sigma_0^2)$ 与动作 $a_1$ 无关）中，证明对时刻 $t=1$ 的策略梯度估计量，采用完整回合回报 $g_{\mathrm{full}} = u_1(r_0 + r_1)$ 与采用未来回报 $g_{\mathrm{rtg}} = u_1 r_1$ 具有完全相同的期望，而方差之差精确满足
  \[
    \Var(g_{\mathrm{full}}) - \Var(g_{\mathrm{rtg}}) = \E[u_1^2]\,\E[r_0^2] + 2\mu_0\,\E[u_1^2 r_1].
  \]
  据此说明：当 $\mu_0 = 0$（或对 $r_1$ 减去最优基线 $b^*$ 后）时，方差之差严格化简为 $\E[u_1^2]\,\sigma_0^2 \ge 0$；而在一般 $\mu_0 \neq 0$ 时交叉项会产生何种影响。
  \exHint{利用 $r_0$ 与 $(u_1, r_1)$ 独立且 $\E[u_1] = 0$，直接展开 $\E[u_1^2(r_0 + r_1)^2] - \E[u_1^2 r_1^2] = \E[u_1^2]\E[r_0^2] + 2\E[r_0]\E[u_1^2 r_1]$；当对 $r_1$ 引入习题~\textbf{6.2.3} 的最优基线 $b^* = \E[u_1^2 r_1]/\E[u_1^2]$ 后，$\E[u_1^2(r_1 - b^*)] = 0$，交叉项彻底消失！}
\end{enumerate}

\exTier{C 层：机制剖析与反例设计}{辨识致命三元组、目标网络频率与高估诊断}

\begin{enumerate}
  \item[\textbf{6.2.5}] \textbf{致命三元组的拆解实验（基于例~\ref{ex:deadly-triad}）}：回到例~\ref{ex:deadly-triad} 的两状态线性逼近发散反例（$\gamma = 0.9, \phi(s_1)=1, \phi(s_2)=2$）。
  \begin{enumerate}
    \item \textbf{去掉“异策略”}：若改按真实马尔可夫链的同策略平稳分布采样（从 $s_1$ 走一步即永久留在 $s_2$，平稳分布 $d(s_1) = 0, d(s_2) = 1$），或任意足够小的 $p = \Prob(s_1)$，求使半梯度 TD 保持收敛的采样概率阈值 $p_{\max}$；
    \item \textbf{去掉“自举”}：若改用蒙特卡洛目标（真实回报恒为 $0$），验证任意 $p \in [0, 1]$ 与任意小步长 $\alpha > 0$ 下 $w_t$ 均指数收敛至 $0$。
  \end{enumerate}
  \exHint{令期望增量系数 $p(2\gamma - 1) + 4(1-p)(\gamma - 1) < 0$，代入 $\gamma=0.9$ 解得 $p < 1/3$。这直观展示了同策略分布为何能抑制半梯度发散。}

  \item[\textbf{6.2.6}] \textbf{DQN 目标网络更新周期 $C$ 的两难权衡}：如果把 DQN 中目标网络同步频率设为 $C = 1$（每步立刻同步 $\theta^- \leftarrow \theta$），算法退化为什么？反之，如果把 $C$ 设得极大（如 $C = 10^6$ 步才同步一次），又会引发什么问题？请从偏差传播速度与训练稳定性两个维度给出定量分析。
  \exHint{$C=1$ 退化为无目标网络的原始半梯度 Q-Learning（易受致命三元组震荡）；$C$ 过大时，多步奖励信息每 $C$ 步才能向前传导一格（类似价值迭代的慢速外循环），导致样本效率暴跌。}
\end{enumerate}

\exTier{D 层：编程实验与系统延伸}{动手复现基线效应与体验 Bug 后果}

\begin{enumerate}
  \item[\textbf{6.2.7}] \textbf{编程实验：验证 \texttt{.detach()} 缺失与奖励平移对 REINFORCE 的冲击}：修改配套脚本里的 \texttt{exp2\_pg\_\allowbreak baseline\_ac()}（位于 \texttt{code/\allowbreak vol6\_\allowbreak experiments.py}）：
  \begin{enumerate}
    \item 将奖励偏移量从 $+4.0$ 分别改为 $0.0$ 与 $+20.0$，比较无基线 REINFORCE 与带基线 REINFORCE 在终段 30 回合的表现，验证带基线版本对常量偏移的近乎完全不变性；
    \item 在 \texttt{baseline} 分支中故意合并优化器并删去 \texttt{vals.detach()}，观察 250 回合内的回报曲线与 Critic 的 MSE 损失如何崩坏。
  \end{enumerate}
  \exHint{当偏移增至 $+20.0$ 时，无基线 REINFORCE 几乎完全无法学出有效策略，而带基线版本在 Critic 适应偏移后仍能稳定收敛。}

  \item[\textbf{6.2.8}] \textbf{从普通 Actor-Critic 走向多轮样本复用的困境预览}：在 \texttt{exp2\_pg\_baseline\_ac()} 的 \texttt{ac} 分支中，每采集 1 条回合只执行 1 次 \texttt{opt\_p.step()}。尝试把 \texttt{loss\_p = -(logps\_t * adv).mean()} 的梯度上升在\textbf{同一批轨迹}上循环执行 10 次（重新前向计算 \texttt{logps\_t}），观察策略回报与策略熵的变化。思考：为什么直接重算对数概率做多轮更新会彻底摧毁策略？这需要怎样的数学修正（提示：第~\ref{ch:trpo-ppo-grpo} 章的重要性采样与裁剪信任域）？
  \exHint{旧轨迹中某次正优势动作若被连续拉高 10 步，其概率会冲向 $1.0$，策略瞬间坍缩；且随着 $\pi_\theta$ 偏离采样策略 $\pi_{\theta_{\mathrm{old}}}$，未加校正的梯度已不再是 $J(\theta)$ 的无偏梯度。}
\end{enumerate}

\chapter[信任域、PPO 裁剪与 GRPO 演进]{信任域方法、近端策略优化（PPO）\\与群体相对策略优化（GRPO）}
\label{ch:trpo-ppo-grpo}

\chapEpigraph{理智不应插上翅膀，而应当加上铅锤与重物，以防它任意飞跃。}{Francis Bacon（弗朗西斯·培根，《新工具》，1620）}

\begin{chapIntro}{本章导读：给策略梯度装上“安全阀”，从信任域走向大模型强化学习}
第~\ref{ch:dqn-pg-ac} 章的策略梯度与 Actor-Critic 优雅地绕开了离散 $\argmax$ 难题，却背负着一个令人抓狂的工程死穴：\textbf{严格的同策略（On-Policy）单步消耗}。每与昂贵的环境交互收集一批轨迹，刚执行一次梯度上升，策略参数就从 $\theta_{\mathrm{old}}$ 变成了 $\theta$，旧数据随即作废；而一旦试图把同一批样本拿来重复更新多轮，策略往往在几步之内彻底崩塌，再也爬不起来。

培根四百年前的告诫恰是现代策略优化的核心哲学——\textbf{比“迈出更大一步”更重要的，是给每一步套上防飞跃的铅锤}。本章沿着“发现崩塌 $\to$ 理论保底 $\to$ 一阶简化 $\to$ 大模型革命”的脉络，完整推导当代强化学习最核心的三代算法：
\begin{enumerate}
  \item \textbf{重要性采样与信任域策略优化（TRPO）}：从性能差异引理出发，证明含 KL 散度惩罚的单调改进下界，并用 Fisher 信息矩阵导出自然策略梯度步长。
  \item \textbf{近端策略优化（PPO-Clip）}：如何仅用一行分段裁剪函数，就在一阶优化框架下构筑与信任域异曲同工的悲观安全边界？我们将完成 $\hat{A}_t > 0$ 与 $\hat{A}_t < 0$ 全六种区间的严格分段求导，并用双子图直观剖析其梯度动力学。
  \item \textbf{群体相对策略优化（GRPO）}：面对千亿参数大语言模型，如何彻底甩掉沉重的 Critic 价值网络？推导组内相对优势归一化与非负无偏 KL 估计量，打通通向第~\ref{ch:posttrain-rlhf-dpo} 章推理模型训练的最后一公里。
  \item \textbf{实测验证}：在一维连续高斯 LQR 控制任务上执行 $K=10$ 轮样本复用，用实测数据 \texttt{data/ppo\_vs\_unclip.csv} 亲眼对比未裁剪重要性采样策略梯度的反复崩塌与 PPO-Clip 的稳健爬升。
\end{enumerate}
\end{chapIntro}

\section{策略崩塌之谜：为什么多轮复用梯度会摧毁策略？}
\label{sec:pg-collapse}

\subsection{监督学习与强化学习对步长的本质差异}

在监督学习中，即使某一步学习率偏大导致参数跳过极小值、验证损失暂时反弹，训练数据集 $\mathcal{D}$ 依然静止不动——只要调小步长，后续梯度仍会稳稳地把参数拉回低损失盆地。然而在强化学习中，\textbf{策略参数 $\theta$ 不仅决定预测质量，更直接决定下一轮采集到什么样的数据}：
\[
  \theta_k \;\longrightarrow\; \pi_{\theta_k} \;\longrightarrow\; \text{采样轨迹 } \tau \sim p_{\theta_k} \;\longrightarrow\; \text{构造梯度 } \hat{g}(\theta_k) \;\longrightarrow\; \theta_{k+1}.
\]
一旦某次更新步长过大，把策略推入一片低回报的糟糕参数区，接下来智能体采集到的全都是“原地打转”或“瞬间坠毁”的垃圾轨迹；基于垃圾轨迹估计出的策略梯度几乎不包含通往高回报区域的信号，智能体于是陷入\textbf{“策略变差 $\to$ 数据变差 $\to$ 策略更差”的不可逆崩塌螺旋（Performance Collapse）}。

更棘手的是，\textbf{参数空间中的欧氏距离 $\|\theta - \theta_{\mathrm{old}}\|_2$ 并不等于策略分布空间中的距离}。对于深层神经网络或高斯策略 $\mathcal{N}(\mu_\theta(s), \sigma^2)$：当 $\sigma$ 很小时，均值 $\mu_\theta(s)$ 仅偏移 $0.1$ 就足以让新旧动作分布几乎零重叠；反之，在激活饱和区，参数变动 $1.0$ 却可能几乎不改变动作概率。因此，仅仅在参数空间里设置固定学习率，根本无法保证策略行为“步子迈得温和”。

\subsection{重要性采样代理目标与无约束多轮复用的崩溃机制}

为了把旧策略 $\pi_{\theta_{\mathrm{old}}}$ 辛辛苦苦采集的一批轨迹重复利用 $K$ 轮（例如 $K = 4 \sim 10$），我们借助概率论中的\textbf{重要性采样（Importance Sampling, IS）}恒等式：只要 $\operatorname{supp}(\pi_\theta) \subseteq \operatorname{supp}(\pi_{\theta_{\mathrm{old}}})$，对任意动作函数 $f(a)$ 都有
\begin{equation}
  \E_{a \sim \pi_\theta(\cdot \mid s)}[f(a)]
  = \sum_a \pi_\theta(a \mid s)\, f(a)
  = \sum_a \pi_{\theta_{\mathrm{old}}}(a \mid s)\,\frac{\pi_\theta(a \mid s)}{\pi_{\theta_{\mathrm{old}}}(a \mid s)}\, f(a)
  = \E_{a \sim \pi_{\theta_{\mathrm{old}}}(\cdot \mid s)}\!\left[r(\theta)\, f(a)\right],
  \label{eq:is-identity}
\end{equation}
其中定义关键的\textbf{概率比（Importance Ratio）}：
\begin{equation}
  r_t(\theta) \triangleq \frac{\pi_\theta(a_t \mid s_t)}{\pi_{\theta_{\mathrm{old}}}(a_t \mid s_t)}
  = \exp\!\Big(\ln \pi_\theta(a_t \mid s_t) - \ln \pi_{\theta_{\mathrm{old}}}(a_t \mid s_t)\Big).
  \label{eq:is-ratio}
\end{equation}
由此构造所谓的\textbf{保守策略迭代代理目标（Conservative Policy Iteration Surrogate, CPI）}：
\begin{equation}
  L^{\mathrm{CPI}}(\theta) \triangleq \hat{\E}_t\!\left[\frac{\pi_\theta(a_t \mid s_t)}{\pi_{\theta_{\mathrm{old}}}(a_t \mid s_t)}\,\hat{A}_t\right]
  = \hat{\E}_t\!\left[r_t(\theta)\,\hat{A}_t\right].
  \label{eq:cpi-surrogate}
\end{equation}

在起点 $\theta = \theta_{\mathrm{old}}$ 处，显然 $r_t(\theta_{\mathrm{old}}) = 1$，且其梯度精确吻合第~\ref{ch:dqn-pg-ac} 章的策略梯度：
\[
  \nabla_\theta L^{\mathrm{CPI}}(\theta)\Big|_{\theta=\theta_{\mathrm{old}}}
  = \hat{\E}_t\!\left[\frac{\nabla_\theta \pi_\theta(a_t \mid s_t)}{\pi_{\theta_{\mathrm{old}}}(a_t \mid s_t)}\,\hat{A}_t\right]_{\theta=\theta_{\mathrm{old}}}
  = \hat{\E}_t\!\left[\nabla_\theta \ln \pi_{\theta_{\mathrm{old}}}(a_t \mid s_t)\,\hat{A}_t\right].
\]

那么，为什么直接拿 $L^{\mathrm{CPI}}(\theta)$ 做 $K=10$ 轮梯度上升会引发灾难？这背后隐藏着\textbf{双重失真}：
\begin{enumerate}
  \item \textbf{状态访问分布失配}：式~\eqref{eq:cpi-surrogate} 仅对动作分布做了重要性校正，却把状态分布冻结在旧策略的 $d^{\pi_{\theta_{\mathrm{old}}}}(s)$ 上！一旦 $\pi_\theta$ 显著偏离 $\pi_{\theta_{\mathrm{old}}}$，真实状态分布 $d^{\pi_\theta}(s)$ 早已面目全非，$L^{\mathrm{CPI}}(\theta)$ 不再反映真实性能 $J(\theta)$。
  \item \textbf{无界概率比的爆炸杠杆}：假设旧策略偶然抽到一个低概率动作（例如 $\pi_{\theta_{\mathrm{old}}}(a_t \mid s_t) = 0.01$），恰好撞上噪声带来的正优势 $\hat{A}_t > 0$。在无约束最大化 $r_t(\theta)\hat{A}_t$ 时，优化器会把 $\pi_\theta(a_t \mid s_t)$ 疯狂推向 $1$，令 $r_t(\theta)$ 飙升至 $100$，从而凭空制造出虚假的巨额代理收益！而对于离散 softmax 或连续高斯分布，把单个样本的概率比推向极端会彻底摧毁其他所有状态的合理策略，引发雪崩式掉分。
\end{enumerate}

要打破这一困境，我们必须回答一个精确的数学问题：\textbf{代理目标 $L^{\mathrm{CPI}}(\theta)$ 与真实回报 $J(\theta)$ 之间的误差，究竟如何受新旧策略散度的控制？}

\section{性能差异引理与信任域策略优化（TRPO）}
\label{sec:trpo}

\subsection{性能差异引理（Performance Difference Lemma）}

比较任意两个策略 $\pi$ 与 $\tilde{\pi}$ 的全局折扣回报 $J(\pi) = \E_{\tau \sim \pi}[\sum_{t=0}^\infty \gamma^t r_t]$ 时，下面的恒等式堪称现代强化学习理论的“瑞士军刀”。它把全局性能之差精确表示为新策略轨迹上旧策略优势函数的折扣累加。

\begin{thmBox}
\begin{lemma}[性能差异引理，Kakade \& Langford, 2002]
\label{lem:pdl}
对任意两个策略 $\pi$ 与 $\tilde{\pi}$ 及折扣因子 $\gamma \in [0, 1)$，记 $d^{\tilde{\pi}}(s) = (1 - \gamma)\sum_{t=0}^\infty \gamma^t \Prob(s_t = s \mid \tilde{\pi})$ 为新策略 $\tilde{\pi}$ 的归一化折扣状态访问分布，则严格成立恒等式
\begin{equation}
  J(\tilde{\pi}) - J(\pi)
  = \E_{\tau \sim \tilde{\pi}}\!\left[\sum_{t=0}^\infty \gamma^t A^\pi(s_t, a_t)\right]
  = \frac{1}{1-\gamma}\,\E_{s \sim d^{\tilde{\pi}},\, a \sim \tilde{\pi}(\cdot \mid s)}\!\left[A^\pi(s, a)\right].
  \label{eq:pdl}
\end{equation}
\end{lemma}
\begin{proof}
由优势函数定义 $A^\pi(s_t, a_t) = Q^\pi(s_t, a_t) - V^\pi(s_t) = \E_{s_{t+1} \sim P(\cdot\mid s_t, a_t)}[r_t + \gamma V^\pi(s_{t+1}) - V^\pi(s_t)]$，将其代入右边并在 $\tau \sim \tilde{\pi}$ 下利用全期望公式与线性性展开：
\begin{align*}
  \E_{\tau \sim \tilde{\pi}}\!\left[\sum_{t=0}^\infty \gamma^t A^\pi(s_t, a_t)\right]
  &= \E_{\tau \sim \tilde{\pi}}\!\left[\sum_{t=0}^\infty \gamma^t \Big(r_t + \gamma V^\pi(s_{t+1}) - V^\pi(s_t)\Big)\right] \\
  &= \E_{\tau \sim \tilde{\pi}}\!\left[\sum_{t=0}^\infty \gamma^t r_t\right]
     + \E_{\tau \sim \tilde{\pi}}\!\left[\sum_{t=0}^\infty \gamma^{t+1} V^\pi(s_{t+1}) - \sum_{t=0}^\infty \gamma^t V^\pi(s_t)\right].
\end{align*}
注意第二个括号内是精确的\textbf{望远镜级数（Telescoping Sum）}：设 $V^\pi$ 有界，当 $t \ge 1$ 时所有 $\gamma^t V^\pi(s_t)$ 项两两完全抵消，仅留下初值项 $-V^\pi(s_0)$！又因为初始状态分布 $\rho_0(s_0)$ 与策略无关，故
\[
  \E_{s_0 \sim \rho_0}[V^\pi(s_0)] = J(\pi), \qquad
  \E_{\tau \sim \tilde{\pi}}\!\left[\sum_{t=0}^\infty \gamma^t r_t\right] = J(\tilde{\pi}).
\]
合并即得 $J(\tilde{\pi}) - J(\pi) = \E_{\tau \sim \tilde{\pi}}[\sum_{t=0}^\infty \gamma^t A^\pi(s_t, a_t)]$。最后，将时间求和改写为状态访问分布期望：
\[
  \sum_{t=0}^\infty \gamma^t \E_{s_t \sim P_t^{\tilde{\pi}},\, a_t \sim \tilde{\pi}}\!\left[A^\pi(s_t, a_t)\right]
  = \frac{1}{1-\gamma}\sum_s \underbrace{(1-\gamma)\sum_{t=0}^\infty \gamma^t \Prob(s_t = s \mid \tilde{\pi})}_{d^{\tilde{\pi}}(s)} \E_{a \sim \tilde{\pi}(\cdot\mid s)}[A^\pi(s, a)],
\]
证毕。
\end{proof}
\end{thmBox}

引理~\ref{lem:pdl} 揭示了一条惊人的推论：只要新策略 $\tilde{\pi}$ 在每个状态下的期望旧优势均满足 $\E_{a \sim \tilde{\pi}}[A^\pi(s, a)] \ge 0$，便必有 $J(\tilde{\pi}) \ge J(\pi)$。遗憾的是，由于神经网络共享参数，某处优势为正往往伴随他处为负；而式~\eqref{eq:pdl} 依赖未知的新分布 $d^{\tilde{\pi}}$，无法在采样前直接计算。

\subsection{单调改进下界与 Majorization--Minimization（MM）哲学}

为了摆脱对未知 $d^{\tilde{\pi}}$ 的依赖，Schulman 等（2015）先用旧策略的状态分布 $d^\pi$ 替代 $d^{\tilde{\pi}}$，定义局部线性近似代理目标：
\begin{equation}
  L_\pi(\tilde{\pi}) \triangleq J(\pi) + \frac{1}{1-\gamma}\,\E_{s \sim d^\pi,\, a \sim \tilde{\pi}(\cdot \mid s)}\!\left[A^\pi(s, a)\right].
  \label{eq:trpo-local-surrogate}
\end{equation}
容易验证 $L_\pi$ 与 $J$ 在 $\tilde{\pi} = \pi$ 处具有\textbf{零阶与一阶完全一致}：因为 $\E_{a\sim\pi}[A^\pi(s,a)] = 0$，故 $L_\pi(\pi) = J(\pi)$；且对参数求导时，因为 $d^{\tilde{\pi}}$ 的导数乘以零期望项 $\E_{a\sim\pi}[A^\pi(s,a)] = 0$ 恰好消去，所以 $\nabla_\theta L_{\pi_{\theta_0}}(\pi_\theta)\big|_{\theta=\theta_0} = \nabla_\theta J(\pi_\theta)\big|_{\theta=\theta_0}$！

更进一步，当 $\tilde{\pi}$ 离开 $\pi$ 不太远时，用 $d^\pi$ 替换 $d^{\tilde{\pi}}$ 产生的误差可以用新旧策略间的最大 KL 散度严格界定：

\begin{thmBox}
\begin{theorem}[TRPO 单调改进下界，Schulman et al., 2015]
\label{thm:trpo-bound}
记 $D_{\mathrm{KL}}^{\max}(\pi \,\|\, \tilde{\pi}) \triangleq \max_{s \in \mathcal{S}} \DKL\big(\pi(\cdot \mid s) \,\|\, \tilde{\pi}(\cdot \mid s)\big)$，并设最大绝对优势 $\varepsilon = \max_{s,a} |A^\pi(s, a)|$。则对任意策略 $\tilde{\pi}$，真实回报严格满足下界
\begin{equation}
  J(\tilde{\pi}) \;\ge\; M_\pi(\tilde{\pi}) \triangleq L_\pi(\tilde{\pi}) - C\, D_{\mathrm{KL}}^{\max}(\pi \,\|\, \tilde{\pi}),
  \qquad C = \frac{4\varepsilon \gamma}{(1 - \gamma)^2}.
  \label{eq:trpo-mm-bound}
\end{equation}
\end{theorem}
\end{thmBox}

\begin{insightBox}{证明骨架与 MM 单调爬坡保证}
式~\eqref{eq:trpo-mm-bound} 的证明分为两步精妙的不等式放缩：
\begin{enumerate}
  \item \textbf{总变差耦合界}：由性能差异引理，$J(\tilde{\pi}) - L_\pi(\tilde{\pi}) = \frac{1}{1-\gamma}\sum_s \big(d^{\tilde{\pi}}(s) - d^\pi(s)\big)\E_{a\sim\tilde{\pi}}[A^\pi(s,a)]$。令 $\alpha = \max_s D_{\mathrm{TV}}(\pi(\cdot\mid s), \tilde{\pi}(\cdot\mid s))$，通过状态耦合引理可知第 $t$ 步两个策略未曾分叉的概率至少为 $(1-\alpha)^t$，从而证明状态分布偏移 $\|d^{\tilde{\pi}} - d^\pi\|_1 \le \frac{2\gamma\alpha}{1-\gamma}$，且每个状态的期望优势绝对值 $|\E_{a\sim\tilde{\pi}}[A^\pi(s,a)]| \le 2\varepsilon \alpha$，两项相乘便得到二次惩罚 $-O(\alpha^2)$。
  \item \textbf{Pinsker 不等式转化}：由信息论经典的 Pinsker 不等式 $D_{\mathrm{TV}}(p, q)^2 \le \frac{1}{2}\DKL(p \| q)$（注意这里是一次 KL 散度控制平方总变差！），立刻把 $\alpha^2$ 替换为 $D_{\mathrm{KL}}^{\max}(\pi \| \tilde{\pi})$，得到式~\eqref{eq:trpo-mm-bound}。
\end{enumerate}
由此立刻推出\textbf{严格单调改进定理（MM 算法）}：若令 $\pi_{k+1} = \argmax_{\tilde{\pi}} M_{\pi_k}(\tilde{\pi})$，则因为 $M_{\pi_k}(\pi_k) = J(\pi_k)$，必然有
\[
  J(\pi_{k+1}) \;\ge\; M_{\pi_k}(\pi_{k+1}) \;\ge\; M_{\pi_k}(\pi_k) \;=\; J(\pi_k).
\]
从几何上看，$M_{\pi_k}(\tilde{\pi})$ 是一条在当前点 $\pi_k$ 处与真实曲折地形 $J(\tilde{\pi})$ 相切、且处处位于其下方的“光滑托盘”；只要把策略推到托盘最高点，真实性能就绝不会掉下去。
\end{insightBox}

\subsection{TRPO 约束优化、Fisher 信息矩阵与自然策略梯度}

直接按定理~\ref{thm:trpo-bound} 的理论惩罚系数 $C = \frac{4\varepsilon\gamma}{(1-\gamma)^2}$ 优化往往过于保守（当 $\gamma = 0.99$ 时 $C \propto 10^4$，步长小如蜗牛）。因此在实用算法\textbf{信任域策略优化（Trust Region Policy Optimization, TRPO）}中，研究者将软惩罚改写为显式 KL 球约束，并用易于采样的平均 KL 散度替代难以计算的最大 KL 散度：
\begin{equation}
  \max_\theta \; \E_{s \sim d^{\theta_{\mathrm{old}}},\, a \sim \pi_{\theta_{\mathrm{old}}}}\!\left[\frac{\pi_\theta(a \mid s)}{\pi_{\theta_{\mathrm{old}}}(a \mid s)}\,\hat{A}_t\right]
  \quad \text{s.t.} \quad
  \bar{D}_{\mathrm{KL}}^{\theta_{\mathrm{old}}}(\theta_{\mathrm{old}} \,\|\, \theta) \triangleq \E_{s \sim d^{\theta_{\mathrm{old}}}}\!\Big[\DKL\big(\pi_{\theta_{\mathrm{old}}}(\cdot \mid s) \,\|\, \pi_\theta(\cdot \mid s)\big)\Big] \le \delta.
  \label{eq:trpo-opt}
\end{equation}

如何在 $\theta_{\mathrm{old}}$ 附近近似求解带非线性约束的式~\eqref{eq:trpo-opt}？令步长 $\Delta\theta = \theta - \theta_{\mathrm{old}}$，对目标函数作一阶泰勒展开、对 KL 约束作二阶泰勒展开：
\begin{itemize}
  \item \textbf{目标展开}：$L^{\mathrm{CPI}}(\theta_{\mathrm{old}} + \Delta\theta) \approx L^{\mathrm{CPI}}(\theta_{\mathrm{old}}) + g^\top \Delta\theta$，其中 $g = \nabla_\theta L^{\mathrm{CPI}}(\theta)\big|_{\theta_{\mathrm{old}}}$ 正是普通策略梯度。
  \item \textbf{KL 约束的零阶、一阶与二阶导数}：记 $D(\theta) = \bar{D}_{\mathrm{KL}}^{\theta_{\mathrm{old}}}(\theta_{\mathrm{old}} \,\|\, \theta)$。
  \begin{enumerate}
    \item 零阶项：$D(\theta_{\mathrm{old}}) = 0$。
    \item 一阶导数：$\nabla_\theta D(\theta)\big|_{\theta_{\mathrm{old}}} = -\E_{s, a\sim \pi_{\theta_{\mathrm{old}}}}[\nabla_\theta \ln \pi_\theta(a\mid s)]\big|_{\theta_{\mathrm{old}}} = 0$（因为概率分布对数梯度的期望恒为零，这与第~\ref{ch:dqn-pg-ac} 章定理~\ref{thm:baseline-unbiased} 用的正是同一个恒等式！）。
    \item 二阶 Hessian 矩阵：再次求导并利用 $\E[\nabla^2 \ln \pi] = -\E[\nabla \ln \pi\,\nabla \ln \pi^\top]$，得到正半定的\textbf{Fisher 信息矩阵（Fisher Information Matrix, FIM）}：
    \begin{equation}
      F \triangleq \nabla_\theta^2 D(\theta)\Big|_{\theta_{\mathrm{old}}}
      = \E_{s \sim d^{\theta_{\mathrm{old}}},\, a \sim \pi_{\theta_{\mathrm{old}}}}\!\Big[\nabla_\theta \ln \pi_\theta(a \mid s)\,\nabla_\theta \ln \pi_\theta(a \mid s)^\top\Big]_{\theta=\theta_{\mathrm{old}}}.
      \label{eq:fisher-matrix}
    \end{equation}
  \end{enumerate}
\end{itemize}
于是式~\eqref{eq:trpo-opt} 局部化为一个标准的二次约束线性规划：
\begin{equation}
  \max_{\Delta\theta} \; g^\top \Delta\theta \quad \text{s.t.} \quad \frac{1}{2}\,\Delta\theta^\top F\, \Delta\theta \le \delta.
  \label{eq:trpo-qp}
\end{equation}
对式~\eqref{eq:trpo-qp} 应用拉格朗日乘子法（$\mathcal{L}(\Delta\theta, \lambda) = g^\top \Delta\theta - \lambda(\frac{1}{2}\Delta\theta^\top F\Delta\theta - \delta)$，令 $\nabla_{\Delta\theta}\mathcal{L} = 0$ 得 $\Delta\theta = \frac{1}{\lambda} F^{-1} g$，再代回边界等式 $\frac{1}{2}\Delta\theta^\top F\Delta\theta = \delta$ 解出 $\lambda$），即得著名的\textbf{自然策略梯度（Natural Policy Gradient）}闭式步长：
\begin{equation}
  \Delta\theta^* = \sqrt{\frac{2\delta}{g^\top F^{-1} g}}\, F^{-1} g.
  \label{eq:npg-step}
\end{equation}

为了避免显式构造和求逆 $d \times d$ 维的 Fisher 矩阵 $F$（对百万参数网络，$F$ 需要数 TB 显存），TRPO 采用\textbf{共轭梯度法（Conjugate Gradient, CG）}仅通过二阶自动微分计算矩阵-向量积 $Fv = \nabla_\theta\big((\nabla_\theta D)^\top v\big)$ 来近似求解线性方程 $Fx = g$，最后辅以\textbf{回溯线搜索（Backtracking Line Search）}验证真实代理目标上升且真实 KL 不超 $\delta$。

尽管 TRPO 在理论上无懈可击，它在现代深度学习与大模型工程中却举步维艰：二阶 CG 迭代实现复杂、难以兼容共享特征层的 Actor-Critic 架构与 Dropout/LayerNorm，也无法直接套用 Adam 等自适应一阶优化器。于是历史向我们发问：\textbf{能否只用一阶随机梯度下降（SGD / Adam），就复现出信任域的绝大部分威力？}

\section{近端策略优化（PPO-Clip）：一阶裁剪如何筑起信任域}
\label{sec:ppo-clip}

\subsection{从 PPO-Penalty 到 PPO-Clip 代理目标}

2017 年，Schulman 等提出\textbf{近端策略优化（Proximal Policy Optimization, PPO）}，最初考察了两条一阶路线：
\begin{itemize}
  \item \textbf{PPO-Penalty（自适应 KL 惩罚）}：直接将 KL 散度作为软惩罚项并入一阶目标，即最大化 $\hat{\E}_t\!\big[r_t(\theta)\hat{A}_t - \beta \DKL(\pi_{\theta_{\mathrm{old}}} \| \pi_\theta)\big]$，若本轮实测 KL 偏大就加倍 $\beta$，偏小就减半 $\beta$。
  \item \textbf{PPO-Clip（概率比裁剪）}：彻底省去显式计算 KL 散度和调节 $\beta$ 的麻烦，直接对概率比 $r_t(\theta)$ 施加对称区间 $[1-\epsilon, 1+\epsilon]$ 裁剪，并与原始项取悲观极小值！
\end{itemize}
大量消融实验表明，结构更简单的 \textbf{PPO-Clip} 不仅代码极短，在连续控制、电子竞技乃至大语言模型 RLHF 中的稳定性还全面胜过 PPO-Penalty。

\begin{defBox}
\begin{definition}[PPO-Clip 代理目标函数，Schulman et al., 2017]
\label{def:ppo-clip}
设裁剪超参数 $\epsilon \in (0, 1)$（工程经典默认值 $\epsilon = 0.2$），概率比 $r_t(\theta) = \frac{\pi_\theta(a_t \mid s_t)}{\pi_{\theta_{\mathrm{old}}}(a_t \mid s_t)}$，$\hat{A}_t$ 为由旧轨迹计算并截断梯度（\texttt{.detach()}）的优势估计（通常采用 GAE($\lambda$)）。PPO-Clip 的策略代理目标定义为
\begin{equation}
  L^{\mathrm{CLIP}}(\theta)
  \triangleq \hat{\E}_t\!\left[\min\!\Big(r_t(\theta)\hat{A}_t,\;\; \clip\!\big(r_t(\theta),\, 1-\epsilon,\, 1+\epsilon\big)\,\hat{A}_t\Big)\right],
  \label{eq:ppo-clip}
\end{equation}
其中 $\clip(r, 1-\epsilon, 1+\epsilon) = \min\big(\max(r, 1-\epsilon),\, 1+\epsilon\big)$。
\end{definition}
\end{defBox}

初看式~\eqref{eq:ppo-clip}，许多初学者都会心生疑惑：\textbf{为什么不能直接最大化 $\clip(r_t(\theta), 1-\epsilon, 1+\epsilon)\hat{A}_t$，非要在外面多套一层看似多余的 $\min(\cdot, \cdot)$ 呢？} 下面的完整逐情形求导剖析将彻底揭开这层 $\min$ 的神来之笔。

\subsection{逐情形严格分段求导剖析：$\hat{A}_t > 0$ 与 $\hat{A}_t < 0$}

\begin{thmBox}
\begin{proposition}[PPO-Clip 代理目标的分段闭式表达与梯度行为]
\label{prop:ppo-cases}
记单步裁剪代理函数为 $f(r; \hat{A}) \triangleq \min\big(r\hat{A},\; \clip(r, 1-\epsilon, 1+\epsilon)\hat{A}\big)$（其中 $r = r_t(\theta) > 0$）。其严格分段形式与关于 $r$ 的导数 $\frac{\partial f}{\partial r}$ 完全刻画如下：
\begin{enumerate}
  \item \textbf{情形一：正优势 $\hat{A} > 0$（好动作，希望提高概率 $r$）}
  \begin{equation}
    f(r; \hat{A}) =
    \begin{cases}
      r\hat{A}, & 0 < r < 1 - \epsilon \quad \text{（$\frac{\partial f}{\partial r} = \hat{A} > 0$，\textbf{未截断，积极拉回！}）} \\
      r\hat{A}, & 1 - \epsilon \le r \le 1 + \epsilon \quad \text{（$\frac{\partial f}{\partial r} = \hat{A} > 0$，信任域内正常上升）} \\
      (1 + \epsilon)\hat{A}, & r > 1 + \epsilon \quad \text{（$\frac{\partial f}{\partial r} = 0$，\textbf{饱和封顶，禁止继续拉高！}）}
    \end{cases}
    \label{eq:ppo-case-pos}
  \end{equation}
  合并前两段可精炼写作 $f(r; \hat{A}) = \min(r, 1+\epsilon)\,\hat{A}$。

  \item \textbf{情形二：负优势 $\hat{A} < 0$（坏动作，希望压低概率 $r$）}
  \begin{equation}
    f(r; \hat{A}) =
    \begin{cases}
      (1 - \epsilon)\hat{A}, & 0 < r < 1 - \epsilon \quad \text{（$\frac{\partial f}{\partial r} = 0$，\textbf{饱和封底，禁止继续压低！}）} \\
      r\hat{A}, & 1 - \epsilon \le r \le 1 + \epsilon \quad \text{（$\frac{\partial f}{\partial r} = \hat{A} < 0$，信任域内正常压低）} \\
      r\hat{A}, & r > 1 + \epsilon \quad \text{（$\frac{\partial f}{\partial r} = \hat{A} < 0$，\textbf{未截断，严厉纠偏！}）}
    \end{cases}
    \label{eq:ppo-case-neg}
  \end{equation}
  合并后两段可精炼写作 $f(r; \hat{A}) = \max(r, 1-\epsilon)\,\hat{A}$。
\end{enumerate}
\end{proposition}
\begin{proof}
当 $\hat{A} > 0$ 时，同乘正数 $\hat{A}$ 保持不等号方向，故
\[
  f(r; \hat{A}) = \min\big(r,\; \clip(r, 1-\epsilon, 1+\epsilon)\big)\,\hat{A}.
\]
分段验证三个区间：若 $r < 1-\epsilon$，则 $\clip(r) = 1-\epsilon > r$，取 $\min$ 得 $r$；若 $r \in [1-\epsilon, 1+\epsilon]$，则 $\clip(r) = r$，取 $\min$ 得 $r$；若 $r > 1+\epsilon$，则 $\clip(r) = 1+\epsilon < r$，取 $\min$ 得 $1+\epsilon$。由此立得式~\eqref{eq:ppo-case-pos}。

当 $\hat{A} < 0$ 时，同乘负数 $\hat{A}$ 翻转不等号，即 $\min(u\hat{A}, v\hat{A}) = \max(u, v)\hat{A}$，故
\[
  f(r; \hat{A}) = \max\big(r,\; \clip(r, 1-\epsilon, 1+\epsilon)\big)\,\hat{A}.
\]
再分三段检验：若 $r < 1-\epsilon$，则 $\clip(r) = 1-\epsilon > r$，取 $\max$ 得 $1-\epsilon$；若 $r \in [1-\epsilon, 1+\epsilon]$，则 $\clip(r) = r$，取 $\max$ 得 $r$；若 $r > 1+\epsilon$，则 $\clip(r) = 1+\epsilon < r$，取 $\max$ 得 $r$。由此立得式~\eqref{eq:ppo-case-neg}。
\end{proof}
\end{thmBox}

\begin{figure}[htbp]
\centering
\resizebox{0.98\linewidth}{!}{%
\begin{tikzpicture}[font=\sffamily\footnotesize]
  % ==================== Left Panel: A_t > 0 ====================
  \begin{scope}[shift={(0,0)}]
    \node[font=\sffamily\bfseries\small, color=tealDark] at (3.1, 4.72) {(a) 正优势 $\hat{A}_t > 0$（好动作：封顶不封底）};
    % Shaded trust region [1-eps, 1+eps] -> x in [1.9, 3.7], center r=1 at x=2.8
    \fill[tealLight] (1.9, 0) rectangle (3.7, 3.75);
    \fill[ochreLight!65] (3.7, 0) rectangle (5.8, 3.75);
    % Axes
    \draw[-{Stealth[length=2mm]}, thick, slate] (0, 0) -- (6.1, 0) node[right, text=ink] {$r_t(\theta)$};
    \draw[-{Stealth[length=2mm]}, thick, slate] (0, 0) -- (0, 4.0) node[left, yshift=-1mm, text=ink] {$L^{\mathrm{CLIP}}$};
    % Ticks on x-axis
    \draw[slate, dashed, thin] (1.9, 0) -- (1.9, 3.75);
    \draw[slate, dashed, thin] (2.8, 0) -- (2.8, 3.75);
    \draw[slate, dashed, thin] (3.7, 0) -- (3.7, 3.75);
    \node[below, font=\scriptsize] at (1.9, -0.05) {$1-\epsilon$};
    \node[below, font=\scriptsize] at (2.8, -0.05) {$1$};
    \node[below, font=\scriptsize] at (3.7, -0.05) {$1+\epsilon$};

    % Unclipped line: y = 0.72 * x
    \draw[navy, dashed, thick] (0.2, 0.144) -- (5.2, 3.744);
    \node[navy, font=\scriptsize, anchor=west] at (4.95, 3.40) {$r_t \hat{A}_t$};
    % Clipped-only line (dotted) on left side r < 1-eps
    \draw[slate, dotted, very thick] (0.2, 1.368) -- (1.9, 1.368);
    \node[slate, font=\scriptsize, align=center] at (1.02, 1.88) {仅裁剪线\\(已舍弃)};

    % Actual PPO-Clip min curve: y = 0.72*x for x<=3.7, y = 2.664 for x>3.7
    \draw[teal, line width=2.2pt] (0.2, 0.144) -- (3.7, 2.664) -- (5.8, 2.664);
    \fill[tealDark] (2.8, 2.016) circle (2.6pt);
    \fill[ochre] (3.7, 2.664) circle (2.6pt);

    % Annotations
    \node[tealDark, font=\scriptsize, align=center] at (1.30, 0.42) {纠偏区\\$\nabla \neq 0$};
    \node[teal, font=\scriptsize, align=center] at (2.8, 3.25) {信任域\\正常上升};
    \node[ochre, font=\scriptsize, align=center] at (4.78, 2.02) {梯度截断区\\$\frac{\partial f}{\partial r} = 0$};
    \draw[-{Stealth[length=1.8mm]}, teal, thick] (0.55, 0.65) -- (1.20, 1.12);
    \node[left, font=\scriptsize, color=ochre] at (0, 2.664) {$(1+\epsilon)\hat{A}_t$};
    \draw[ochre, dotted] (0, 2.664) -- (3.7, 2.664);
  \end{scope}

  % ==================== Right Panel: A_t < 0 ====================
  \begin{scope}[shift={(7.5,0)}]
    \node[font=\sffamily\bfseries\small, color=ochre] at (3.1, 4.72) {(b) 负优势 $\hat{A}_t < 0$（坏动作：封底不封顶）};
    % Shaded trust region [1-eps, 1+eps]
    \fill[tealLight] (1.9, 0) rectangle (3.7, 3.75);
    \fill[ochreLight!65] (0, 0) rectangle (1.9, 3.75);
    % Axes: x-axis at y=3.75 (value 0), downward is negative
    \draw[-{Stealth[length=2mm]}, thick, slate] (0, 3.75) -- (6.1, 3.75) node[right, text=ink] {$r_t(\theta)$};
    \draw[-{Stealth[length=2mm]}, thick, slate] (0, 0) -- (0, 4.0) node[left, yshift=-1mm, text=ink] {$L^{\mathrm{CLIP}}$};
    % Ticks
    \draw[slate, dashed, thin] (1.9, 0) -- (1.9, 3.75);
    \draw[slate, dashed, thin] (2.8, 0) -- (2.8, 3.75);
    \draw[slate, dashed, thin] (3.7, 0) -- (3.7, 3.75);
    \node[above, font=\scriptsize] at (1.9, 3.80) {$1-\epsilon$};
    \node[above, font=\scriptsize] at (2.8, 3.80) {$1$};
    \node[above, font=\scriptsize] at (3.7, 3.80) {$1+\epsilon$};

    % Unclipped line: y = 3.75 - 0.65 * x
    \draw[navy, dashed, thick] (0.2, 3.62) -- (5.5, 0.175);
    \node[navy, font=\scriptsize, anchor=west] at (5.15, 0.72) {$r_t \hat{A}_t$};
    % Clipped-only line (dotted) on right side r > 1+eps: y = 3.75 - 0.65*3.7 = 1.345
    \draw[slate, dotted, very thick] (3.7, 1.345) -- (5.8, 1.345);
    \node[slate, font=\scriptsize, align=center] at (4.85, 1.72) {仅裁剪线(已舍弃)};

    % Actual PPO-Clip min curve: flat at y = 3.75 - 0.65*1.9 = 2.515 for x<=1.9, then follows 3.75 - 0.65*x
    \draw[ochre, line width=2.2pt] (0.15, 2.515) -- (1.9, 2.515) -- (5.5, 0.175);
    \fill[tealDark] (2.8, 1.93) circle (2.6pt);
    \fill[ochre] (1.9, 2.515) circle (2.6pt);

    % Annotations
    \node[ochre, font=\scriptsize, align=center] at (0.95, 1.85) {梯度截断区\\$\frac{\partial f}{\partial r} = 0$};
    \node[teal, font=\scriptsize, align=center] at (2.8, 0.65) {信任域\\正常压低};
    \node[tealDark, font=\scriptsize, align=center] at (4.22, 0.25) {严厉纠偏区\\$\nabla \neq 0$};
    \draw[-{Stealth[length=1.8mm]}, ochre, thick] (5.10, 0.78) -- (4.50, 1.15);
    \node[left, font=\scriptsize, color=ochre] at (0, 2.515) {$(1-\epsilon)\hat{A}_t$};
  \end{scope}
\end{tikzpicture}}
\caption{PPO-Clip 单步代理目标函数 $f(r_t; \hat{A}_t)$ 在正优势与负优势下的精确分段几何。外层 $\min$ 构造了原始代理目标的\textbf{悲观下界}：当更新试图继续冲出信任域以谋取更大收益时，曲线变平令梯度归零；而当其他小批量更新意外将 $r_t$ 推向错误一侧时，梯度保持激活以迅速纠偏。}
\label{fig:ppo-clip-geometry}
\end{figure}

结合图~\ref{fig:ppo-clip-geometry}，现在我们可以彻底回答前面提出的灵魂拷问——\textbf{外层 $\min$ 到底起了什么作用？}
\begin{itemize}
  \item \textbf{它是一道单向棘轮（One-Way Ratchet）}：当动作是好动作（$\hat{A}_t > 0$）且概率已经被拉高超过 $1+\epsilon$ 时，梯度归零，防止贪婪榨干；但如果因为同批次中其他样本的参数耦合，意外把该好动作的概率压到了 $1-\epsilon$ 以下，外层 $\min$ 选中的是\emph{更低}的 $r_t \hat{A}_t$ 而非 $(1-\epsilon)\hat{A}_t$——于是梯度仍然为正，立刻把误伤的好动作拉回信任域！
  \item \textbf{坏动作同理（$\hat{A}_t < 0$）}：把坏动作的概率压低到 $1-\epsilon$ 以下后，梯度归零，防止把探索概率一棒子打死成 $0$；但如果某个坏动作的概率反常飙升到了 $1+\epsilon$ 以上，外层 $\min$ 依然保持惩罚梯度激活，毫不留情地把它按回 $[1-\epsilon, 1+\epsilon]$！
\end{itemize}
正因为这一精妙的非对称“只拦激进、不拦纠偏”悲观下界设计，PPO-Clip 允许我们在同一批轨迹上安心切分 mini-batch、重复迭代 $K = 4 \sim 10$ 轮，将样本利用率直接提升一个数量级，却只需一行 \texttt{torch.min} 与 \texttt{torch.clamp}！

\subsection{完整 PPO 算法流程与 PyTorch 核心实现}

在工程实践中，PPO 将策略裁剪目标、价值函数回归损失与策略熵正则项组合成联合目标（注意 PyTorch 执行梯度\emph{下降}，因此对代理目标与熵取负号）：
\begin{equation}
  \mathcal{L}_{\mathrm{PPO}}(\theta, \phi)
  = -L^{\mathrm{CLIP}}(\theta)
    + c_1\,\hat{\E}_t\!\left[\big(V_\phi(s_t) - \hat{R}_t\big)^2\right]
    - c_2\,\hat{\E}_t\!\left[\mathcal{H}\big(\pi_\theta(\cdot \mid s_t)\big)\right],
  \label{eq:ppo-total-loss}
\end{equation}
其中 $\hat{R}_t = \hat{A}_t^{\mathrm{GAE}} + V_{\phi_{\mathrm{old}}}(s_t)$ 为 $\lambda$-回报目标，常用系数 $c_1 = 0.5, c_2 \in [0, 0.01]$。此外，现代 PPO 实现通常配套三项久经考验的工程稳定技巧：
\begin{enumerate}
  \item \textbf{小批量优势标准化}：将每个 batch 内的优势归一化为 $\tilde{A}_t = \frac{\hat{A}_t - \mu_A}{\sigma_A + 10^{-8}}$，消除不同任务奖励尺度差异。
  \item \textbf{全局梯度范数裁剪}：\texttt{clip\_grad\_norm\_(params, 0.5)} 防止偶发极大优势引发参数突变。
  \item \textbf{近似 KL 早停（KL Early Stopping）}：在 $K$ 轮复用中实时监测近似散度 $\hat{d}_{\mathrm{KL}} = \hat{\E}_t[(r_t - 1) - \ln r_t]$；一旦超过阈值（如 $1.5\,\delta_{\mathrm{target}}$），立刻提前终止本轮剩余 epoch。
\end{enumerate}

\begin{lstlisting}[caption={PPO-Clip 多轮复用更新核心模块（支持连续高斯与离散策略的简洁 PyTorch 实现）}]
def ppo_clip_update(policy, value_net, opt, batch, clip_eps=0.2,
                    epochs=10, vf_coef=0.5, ent_coef=0.01, max_kl=0.03):
    # batch 中 old_logp 与 adv 必须在采样阶段经 torch.no_grad() 彻底冻结!
    states, actions, old_logp, returns, adv = batch
    adv = (adv - adv.mean()) / (adv.std() + 1e-8)

    for _ in range(epochs):
        dist = policy(states)                      # 返回 Categorical 或 Normal 分布
        new_logp = dist.log_prob(actions)
        if new_logp.ndim > 1:
            new_logp = new_logp.sum(dim=-1)        # 连续多维动作对数概率相加

        log_ratio = torch.clamp(new_logp - old_logp, -10.0, 10.0)
        ratio = torch.exp(log_ratio)               # r_t(theta) = pi_new / pi_old

        # PPO-Clip 悲观下界代理目标
        surr1 = ratio * adv
        surr2 = torch.clamp(ratio, 1.0 - clip_eps, 1.0 + clip_eps) * adv
        policy_loss = -torch.min(surr1, surr2).mean()

        value_loss = F.mse_loss(value_net(states).squeeze(-1), returns)
        entropy = dist.entropy().mean()
        loss = policy_loss + vf_coef * value_loss - ent_coef * entropy

        opt.zero_grad()
        loss.backward()
        torch.nn.utils.clip_grad_norm_(
            list(policy.parameters()) + list(value_net.parameters()), 0.5
        )
        opt.step()

        # 非负近似 KL 早停监控 (Schulman k3 估计量)
        with torch.no_grad():
            approx_kl = ((ratio - 1.0) - log_ratio).mean().item()
            if approx_kl > 1.5 * max_kl:
                break
\end{lstlisting}

\section{群体相对策略优化（GRPO）：甩掉 Critic 的大模型 RL 革命}
\label{sec:grpo}

\subsection{为什么标准 PPO 在大语言模型时代变得无比沉重？}

从 2022 年起，PPO 成为大语言模型人类反馈强化学习（RLHF）的标准引擎。然而，当人们把 PPO 从几万参数的机器人 MLP 搬到 700 亿（70B）参数的大语言模型时，立刻撞上了两堵铜墙铁壁：
\begin{enumerate}
  \item \textbf{显存与算力的“四模型噩梦”}：标准 LLM-PPO 训练需要同时常驻四个同量级巨无霸模型——待训练的策略模型 $\pi_\theta$（Actor）、价值模型 $V_\phi$（Critic）、冻结参考模型 $\pi_{\mathrm{ref}}$（防止遗忘语言能力）、以及奖励模型 $r_\psi$（RM）。其中 Critic $V_\phi$ 与 Actor 规模相当，仅其权重、梯度与 Adam 一二阶动量就吞噬了近一半的训练显存与算力！
  \item \textbf{长思维链上的稀疏终端奖励难题}：在数学解题或代码生成任务中，模型一口气输出长达 4000 个 token 的推理过程，只有最后一个答案对错才给出标量奖励 $R \in \{0, 1\}$。试图训练一个 token 级 Critic $V_\phi(q, o_{\le t})$ 在推理进行到第 1372 个词时就准确预判最终能否解出难题，噪声极大且极易欠拟合；一个糟糕的 Critic 反而会通过有偏优势估计把 Actor 带偏。
\end{enumerate}

\subsection{GRPO 的核心洞察：用“同题组内比拼”取代 Critic 基线}

2024 年，DeepSeekMath 与随后的 DeepSeek-R1 提出了\textbf{群体相对策略优化（Group Relative Policy Optimization, GRPO，Shao et al., 2024）}。它的破局思路返璞归真却威力惊人：\textbf{既然价值基线 $V(q)$ 的本质不过是“同一提示 $q$ 下策略生成回答的期望奖励”，何必养一个千亿参数的 Critic 去艰难拟合它？直接让旧策略 $\pi_{\theta_{\mathrm{old}}}$ 对同一个问题 $q$ 独立采样一组 $G$ 个回答，用这组回答的经验均值当基线、经验标准差做归一化不就行了吗！}

\begin{defBox}
\begin{definition}[GRPO 群体相对优势与目标函数，Shao et al., 2024]
\label{def:grpo}
给定问题（提示）$q \sim P(Q)$，从旧策略 $\pi_{\theta_{\mathrm{old}}}(\cdot \mid q)$ 独立采样一组 $G$ 条完整回答 $\{o_1, o_2, \dots, o_G\}$（典型值 $G \in \{8, 16, 64\}$），由奖励模型或可验证规则（如数学答案比对、单元测试）给出每条回答的标量结果奖励 $\{R_1, R_2, \dots, R_G\}$。
\begin{enumerate}
  \item \textbf{群体相对归一化优势}：对组内第 $i$ 条回答，定义其相对优势为 z-score 标准化值，并广播至该回答的每一个输出 token $t \in \{1, \dots, |o_i|\}$：
  \begin{equation}
    \hat{A}_{i, t} = \hat{A}_i \triangleq \frac{R_i - \bar{R}}{\sigma_R + \varepsilon}, \qquad
    \bar{R} = \frac{1}{G}\sum_{j=1}^G R_j, \qquad
    \sigma_R = \sqrt{\frac{1}{G}\sum_{j=1}^G (R_j - \bar{R})^2},
    \label{eq:grpo-adv}
  \end{equation}
  其中 $\varepsilon \approx 10^{-6}$ 用于防止全组奖励相同时除以零。
  \item \textbf{带显式参考 KL 正则的 GRPO 代理目标}：记 token 级概率比 $r_{i,t}(\theta) = \frac{\pi_\theta(o_{i,t} \mid q, o_{i,<t})}{\pi_{\theta_{\mathrm{old}}}(o_{i,t} \mid q, o_{i,<t})}$，GRPO 最大化
  \begin{equation}
  \begin{aligned}
    \mathcal{J}_{\mathrm{GRPO}}(\theta)
    = \E_{\substack{q \sim P(Q) \\ \{o_i\}_{i=1}^G \sim \pi_{\theta_{\mathrm{old}}}}}
    \Bigg[\frac{1}{G}\sum_{i=1}^G \frac{1}{|o_i|}\sum_{t=1}^{|o_i|}
    \Big(&\min\!\big(r_{i,t}(\theta)\hat{A}_i,\; \clip(r_{i,t}(\theta), 1-\epsilon, 1+\epsilon)\hat{A}_i\big) \\
    &- \beta\,\hat{D}_{\mathrm{KL}}\big(\pi_\theta \,\|\, \pi_{\mathrm{ref}}\big)_{i,t}\Big)\Bigg],
  \end{aligned}
  \label{eq:grpo-obj}
  \end{equation}
  其中 $\beta \ge 0$ 为参考策略偏离惩罚系数，token 级 KL 估计量采用 Schulman $k_3$ 形式：
  \begin{equation}
    \hat{D}_{\mathrm{KL}}\big(\pi_\theta \,\|\, \pi_{\mathrm{ref}}\big)_{i,t}
    \triangleq \frac{\pi_{\mathrm{ref}}(o_{i,t} \mid q, o_{i,<t})}{\pi_\theta(o_{i,t} \mid q, o_{i,<t})}
    - \ln \frac{\pi_{\mathrm{ref}}(o_{i,t} \mid q, o_{i,<t})}{\pi_\theta(o_{i,t} \mid q, o_{i,<t})} - 1.
    \label{eq:grpo-kl-k3}
  \end{equation}
\end{enumerate}
\end{definition}
\end{defBox}

注意 GRPO 在两个精妙之处与传统 LLM-PPO 截然不同：第一，它用\textbf{组内蒙特卡洛基线}彻底消灭了 Critic 网络；第二，传统 PPO 把 $-\beta \ln \frac{\pi_\theta}{\pi_{\mathrm{ref}}}$ 揉进每步奖励里再算优势，而 GRPO 直接把非负 KL 估计量~\eqref{eq:grpo-kl-k3} 作为显式正则项写入损失函数。为什么非要用这个看似绕口的 $\frac{\pi_{\mathrm{ref}}}{\pi_\theta} - \ln \frac{\pi_{\mathrm{ref}}}{\pi_\theta} - 1$，而不是更直观的 $-\ln \frac{\pi_{\mathrm{ref}}}{\pi_\theta} = \ln \frac{\pi_\theta}{\pi_{\mathrm{ref}}}$？下面的短小命题给出了天书级的回答。

\begin{thmBox}
\begin{proposition}[GRPO 非负 KL 估计量的无偏性与严格非负性]
\label{prop:kl-k3}
设 $p, q$ 为同一支撑集上的两个概率分布，定义随机变量 $u(a) = \frac{q(a)}{p(a)}$（其中 $a \sim p$），并令
\[
  k_1(a) \triangleq -\ln u(a) = \ln \frac{p(a)}{q(a)}, \qquad
  k_3(a) \triangleq u(a) - 1 - \ln u(a) = \frac{q(a)}{p(a)} - \ln \frac{q(a)}{p(a)} - 1.
\]
则对所有采样结果 $a$ 严格成立：
\begin{enumerate}
  \item \textbf{两者均无偏}：$\E_{a \sim p}[k_1(a)] = \E_{a \sim p}[k_3(a)] = \DKL(p \,\|\, q)$；
  \item \textbf{逐点严格非负}：$k_3(a) \ge 0$ 恒成立（当且仅当 $p(a) = q(a)$ 时取等号），而 $k_1(a)$ 在 $p(a) < q(a)$ 时会取负值；
  \item \textbf{方差更低}：利用控制变量 $(u(a) - 1)$（其期望为零且与 $-\ln u(a)$ 强负相关），在 $p \approx q$ 时 $k_3(a) \approx \frac{1}{2}(u(a)-1)^2$ 的方差量级远低于 $k_1(a) \approx -(u(a)-1)$。
\end{enumerate}
\end{proposition}
\begin{proof}
首先验证控制变量期望为零：
\[
  \E_{a \sim p}[u(a) - 1] = \sum_a p(a)\,\frac{q(a)}{p(a)} - 1 = \sum_a q(a) - 1 = 1 - 1 = 0.
\]
因此 $\E_{a \sim p}[k_3(a)] = \E_{a \sim p}[u(a) - 1] + \E_{a \sim p}[-\ln u(a)] = 0 + \E_{a \sim p}\!\big[\ln\frac{p(a)}{q(a)}\big] = \DKL(p \,\|\, q)$。

其次，令 $h(u) = u - 1 - \ln u$（$u > 0$），求导得 $h'(u) = 1 - \frac{1}{u}$、$h''(u) = \frac{1}{u^2} > 0$。故 $h(u)$ 在 $u = 1$ 处取得唯一全局极小值 $h(1) = 0$，即对任意 $u > 0$ 恒有 $k_3(a) \ge 0$。最后，记 $\delta = u - 1$，当 $p \approx q$（$|\delta| \ll 1$）时泰勒展开得 $k_1 = -\delta + O(\delta^2)$（方差为 $O(\E[\delta^2])$），而 $k_3 = \frac{1}{2}\delta^2 + O(\delta^3)$（方差为 $O(\E[\delta^4])$），在策略微调的近端区域实现量级级的方差压缩！
\end{proof}
\end{thmBox}

\begin{table}[htbp]
\centering
\caption{标准 LLM-PPO 与群体相对策略优化（GRPO）系统架构与数学特性全面对比}
\label{tab:ppo-vs-grpo}
\footnotesize
\begin{tabularx}{\linewidth}{@{}p{27mm}XX@{}}
\toprule
\textbf{对比维度} & \textbf{标准 PPO（Actor-Critic + GAE）} & \textbf{GRPO（Group Relative Policy Opt.）} \\
\midrule
常驻模型开销 & Actor $+$ \textbf{Critic} $+$ Ref $+$ RM（4 份） & Actor $+$ Ref（若用规则验证器，\textbf{仅 2 份！}） \\
优势基线来源 & 参数化价值网络 $V_\phi(s_t)$ $+$ GAE($\lambda$) 递推 & 同提示 $G$ 条采样的组内统计量 $(\bar{R}, \sigma_R)$ \\
优势粒度 & Token 级稠密优势 $\hat{A}_t$（依赖 Critic 精度） & 序列级结果优势广播至全句（或配合 PRM） \\
KL 正则注入 & 并入每步奖励 $r_t - \beta \ln\frac{\pi_\theta}{\pi_{\mathrm{ref}}}$ & 直接加入损失函数 $-\beta \hat{D}_{\mathrm{KL}}^{(k_3)}(\pi_\theta \| \pi_{\mathrm{ref}})$ \\
采样计算成本 & 每个 prompt 采样 1 条即可更新 & 每个 prompt 需采样 $G$ 条（如 $G=8 \sim 64$），推理并行度高 \\
核心短板 & Critic 显存翻倍、长思维链价值估计高偏 & 若组内全对或全错（$\sigma_R = 0$），该题梯度归零 \\
最佳适用场景 & 稠密奖励连续控制、密集步级反馈任务 & 数学/代码推理（RLVR）、结果导向型大模型后训练 \\
\bottomrule
\end{tabularx}
\end{table}

\begin{lstlisting}[caption={GRPO 核心损失函数（群体相对优势归一化 $+$ PPO-Clip $+$ 非负 $k_3$ KL 正则）}]
def grpo_loss(new_logp, old_logp, ref_logp, rewards, mask,
              clip_eps=0.2, beta=0.04):
    """
    new_logp, old_logp, ref_logp, mask: (B, G, T)  B组问题, 每组G条回答, 长度T
    rewards: (B, G) 每条完整回答的标量结果奖励; old_logp 与 ref_logp 需 no_grad
    """
    # 1. 组内相对优势归一化 (无需 Critic!) -> 广播到每个 token 位置
    mean_r = rewards.mean(dim=1, keepdim=True)
    std_r  = rewards.std(dim=1, keepdim=True, unbiased=False)
    adv    = ((rewards - mean_r) / (std_r + 1e-6)).unsqueeze(-1)   # (B, G, 1)

    # 2. Token 级概率比与 PPO-Clip 悲观代理目标
    log_ratio = torch.clamp(new_logp - old_logp, -10.0, 10.0)
    ratio     = torch.exp(log_ratio)
    surr1     = ratio * adv
    surr2     = torch.clamp(ratio, 1.0 - clip_eps, 1.0 + clip_eps) * adv
    pg_obj    = torch.min(surr1, surr2)

    # 3. Schulman k3 非负无偏 KL 估计量: exp(ref - new) - (ref - new) - 1
    log_ref_diff = torch.clamp(ref_logp - new_logp, -10.0, 10.0)
    kl_k3        = torch.exp(log_ref_diff) - log_ref_diff - 1.0

    # 4. 按每条序列真实 token 长度求均值, 再跨组求均值
    per_token = (pg_obj - beta * kl_k3) * mask
    seq_lens  = mask.sum(dim=-1).clamp(min=1.0)                    # (B, G)
    return -(per_token.sum(dim=-1) / seq_lens).mean()
\end{lstlisting}

请细品这段仅二十余行的 GRPO 代码：它与前文的 PPO-Clip 一脉相承，却在架构上彻底卸下了千亿参数的 Critic 巨石。正是这套简洁到近乎苛刻的机制，让开源社区得以用相对有限的算力复现出具备自我反思与长链推理能力的前沿推理模型（我们将在第~\ref{ch:posttrain-rlhf-dpo} 章深入探讨 RLVR 推理涌现现象）。

\section{实测验证：多轮复用下未裁剪 IS-PG 的崩塌与 PPO-Clip 的稳健爬升}
\label{sec:exp-ppo}

理论上的推导固然自洽，但最令人震撼的莫过于亲眼见证“去掉裁剪”与“加上裁剪”在完全相同的随机种子、完全相同的网络架构与优化器下走上截然不同的命运。

在配套实验脚本 \texttt{code/vol6\_experiments.py} 的 \texttt{exp3\_ppo\_vs\_unclipped()} 中，我们搭建了一个经典的一维连续高斯\textbf{线性二次调节器（LQR）控制环境}：
\begin{itemize}
  \item \textbf{动力学与奖励}：初始位置 $x_0 \sim \mathrm{Uniform}(-1.5, 1.5)$，每回合 $T = 10$ 步。策略输出高斯分布 $a_t \sim \mathcal{N}(\mu_\theta(s_t), \sigma^2)$（截断至 $[-2, 2]$），状态按 $x_{t+1} = \clip(x_t + 0.5 a_t + \xi_t, -3, 3)$ 演化（噪声 $\xi_t \sim \mathcal{N}(0, 0.05^2)$），即时奖励为二次控制代价 $r_t = -x_t^2 - 0.1 a_t^2$（理论最优回报接近 $0$）。
  \item \textbf{严格对照设置}：每次外循环收集 16 条轨迹，随后在同一批轨迹上执行 \textbf{$K = 10$ 轮样本复用}（Adam 学习率 $4 \times 10^{-2}$，优势标准化，全局梯度范数裁剪 $2.0$）。两组实验除是否启用 PPO-Clip（$\epsilon = 0.2$）外，初始化种子（\texttt{2026}）与全部超参数完全一致。
\end{itemize}

\begin{figure}[htbp]
\centering
\begin{tikzpicture}
  % Left panel: Full trajectory from iter 1 (-51.54) to 49
  \begin{axis}[
    name=leftplot,
    width=0.47\linewidth, height=6.1cm,
    title={\footnotesize\sffamily\bfseries (a) 全局收敛轨迹（迭代 1--49）},
    xlabel={策略更新迭代轮数（每轮复用 $K=10$ 次）},
    ylabel={回合平均总回报},
    xmin=0, xmax=51, ymin=-55, ymax=3,
    grid=both, grid style={tealLine!45, dashed},
    axis line style={slate},
    tick label style={font=\scriptsize},
    label style={font=\footnotesize\sffamily, color=ink},
    legend style={at={(0.97,0.06)}, anchor=south east, font=\scriptsize, fill=white, draw=tealLine}
  ]
    \addplot[ochre, thick, dashed, mark=triangle*, mark size=1.4pt]
      table[x=iter, y=unclipped, col sep=comma]{data/ppo_vs_unclip.csv};
    \addlegendentry{未裁剪 IS-PG}

    \addplot[teal, very thick, mark=*, mark size=1.3pt]
      table[x=iter, y=ppo, col sep=comma]{data/ppo_vs_unclip.csv};
    \addlegendentry{PPO-Clip ($\epsilon=0.2$)}
  \end{axis}

  % Right panel: Zoomed view iter 7..49 showing repeated crashes of unclipped IS-PG
  \begin{axis}[
    at={(leftplot.east)}, anchor=west, xshift=12mm,
    width=0.47\linewidth, height=6.1cm,
    title={\footnotesize\sffamily\bfseries (b) 局部放大（迭代 7--49：崩塌尖峰对比）},
    xlabel={策略更新迭代轮数（局部放大）},
    ylabel={回合平均总回报},
    xmin=6, xmax=51, ymin=-8.2, ymax=0.4,
    grid=both, grid style={tealLine!45, dashed},
    axis line style={slate},
    tick label style={font=\scriptsize},
    label style={font=\footnotesize\sffamily, color=ink},
    legend style={at={(0.97,0.06)}, anchor=south east, font=\scriptsize, fill=white, draw=tealLine}
  ]
    \addplot[ochre, thick, dashed, mark=triangle*, mark size=1.6pt]
      table[x=iter, y=unclipped, col sep=comma]{data/ppo_vs_unclip.csv};
    \addlegendentry{未裁剪 IS-PG（最差 $-7.26$）}

    \addplot[teal, very thick, mark=*, mark size=1.4pt]
      table[x=iter, y=ppo, col sep=comma]{data/ppo_vs_unclip.csv};
    \addlegendentry{PPO-Clip（末 10 轮均值 $-0.36$）}

    % Annotate the dramatic crash at iter=11 (-7.263), clear of legend
    \node[ochre, font=\scriptsize\sffamily, anchor=west] at (axis cs:14.0, -4.95) {第 11 轮骤崩至 $-7.26$};
    \draw[-{Stealth[length=1.8mm]}, ochre, thin] (axis cs:15.2, -5.25) -- (axis cs:11.4, -7.12);
    % Annotate secondary crash at iter=25 (-2.431)
    \node[ochre, font=\scriptsize\sffamily, anchor=west] at (axis cs:27.5, -3.55) {第 21/25 轮反复回撤};
    \draw[-{Stealth[length=1.8mm]}, ochre, thin] (axis cs:28.0, -3.35) -- (axis cs:25.3, -2.52);
  \end{axis}
\end{tikzpicture}
\caption{连续一维高斯 LQR 控制任务上 $K=10$ 轮多步样本复用的实测结果（读取自 \texttt{data/ppo\_vs\_unclip.csv}）。(a) 两者均从初始回报 $-51.54$ 起步；(b) 进入精细控制区后，未裁剪的重要性采样策略梯度因概率比失控反复剧烈崩塌，而 PPO-Clip（$\epsilon=0.2$）牢牢锁住信任域，平稳收敛至近优区间。}
\label{fig:ppo-vs-unclip}
\end{figure}

图~\ref{fig:ppo-vs-unclip}（直接读取 \texttt{data/ppo\_vs\_unclip.csv}）清晰呈现了两种算法在多轮复用下的鲜明分叉：
\begin{enumerate}
  \item \textbf{初期步幅与风险的直觉反差}：两种策略起点完全相同（第 1 轮回报均为 \textbf{$-51.54$}）。前 5 轮中，无约束的未裁剪 IS-PG 看似“冲得更快”（第 3 轮即跳至 $-2.56$，而保守的 PPO-Clip 因受 $[0.8, 1.2]$ 裁剪约束仅稳步爬至 $-18.53$）——这正是培根所警示的“插上翅膀的虚幻飞跃”！
  \item \textbf{未裁剪 IS-PG 的反复崩塌}：一旦进入精细控制区，未裁剪 IS-PG 在第 9 轮刚摸到 $-0.61$，紧接着第 11 轮就遭遇断崖式崩塌，骤跌至全中后期最差值 \textbf{$-7.26$}（性能恶化超过 $11\times$）！随后在第 21 轮（$-2.18$）、第 25 轮（$-2.43$）、第 35 轮（$-1.18$）与第 41 轮（$-1.36$）反复上演“刚爬上去、一步踩空摔回来”的剧烈震荡，最后 10 轮平均回报仅停留在 \textbf{$-0.55$}。
  \item \textbf{PPO-Clip 的稳健爬升与高精度收敛}：反观 PPO-Clip（$\epsilon = 0.2$），虽然前几步刻意放慢了速度，但自第 11 轮跨入 $-0.76$ 后便如履平地、全程零崩塌，第 27 轮更是达到 $-0.20$ 的优异表现，最后 10 轮平均回报稳定在 \textbf{$-0.36$}——相比未裁剪版本，稳态控制代价直接降低了 \textbf{$34.5\%$}！
\end{enumerate}

\begin{failBox}{PPO 与 GRPO 实战中最隐蔽的三大数值与逻辑陷阱}
\begin{enumerate}
  \item \textbf{忘记冻结 \texttt{old\_logp} 导致 PPO-Clip 悄无声息失效}：如果计算 \texttt{old\_logp} 时忘了包裹 \texttt{torch.no\_grad()} 或忘了在多轮更新前保存脱离计算图的张量，使得循环内误用当前参数重新算了 \texttt{old\_logp}，则每一步都有 \texttt{ratio == 1.0} 恒成立！此时裁剪机制永远不会触发，虽然第一轮梯度恰好仍等于单步策略梯度，但多轮复用直接退化为无保护的灾难性暴走。
  \item \textbf{指数运算 \texttt{torch.exp(new\_logp - old\_logp)} 溢出成 \texttt{NaN}}：在高维连续动作或长文本序列中，对数概率是数百个维度/token 的累加值；若某个异常样本引起对数概率差飙升至 $+90$ 以上，FP32 的 \texttt{exp} 会瞬间溢出为 \texttt{inf}，紧接着 \texttt{inf * 0} 或反向传播便将整个网络权重染成 \texttt{NaN}。工程上既应对单步/单 token（而非整句联合概率！）计算概率比，也应对 \texttt{log\_ratio} 施加安全截断（如 $[-10, 10]$）。
  \item \textbf{GRPO 中组内奖励全等与 KL 方向颠倒}：当一批数学题太简单（全对 $R_i \equiv 1$）或太难（全错 $R_i \equiv 0$）时，$\sigma_R = 0$，该组所有样本的优势 $\hat{A}_i = 0$，产生“空转梯度”；同时，务必核对 $k_3$ 公式中的分子是 $\pi_{\mathrm{ref}}$、分母是 $\pi_\theta$（对应正向惩罚 $\DKL(\pi_\theta \| \pi_{\mathrm{ref}})$）——一旦写反为 $\pi_\theta / \pi_{\mathrm{ref}}$，模型会朝错误的高方差方向偏离，且容易因 $\pi_{\mathrm{ref}} \to 0$ 引发爆炸！
\end{enumerate}
\end{failBox}

\begin{extBox}{后 GRPO 时代的前沿精修：DAPO 与 Dr. GRPO}
2025 年以来，随着大规模推理模型强化学习（RLVR）的深入实践，研究者发现了原始 GRPO 目标~\eqref{eq:grpo-obj} 中潜藏的三个微妙结构偏差，并提出了广受欢迎的改进算法 \textbf{DAPO}（Yu et al., 2025）与 \textbf{Dr. GRPO}（Liu et al., 2025）：
\begin{enumerate}
  \item \textbf{非对称裁剪（Clip-Higher）}：在原始对称裁剪 $[1-\epsilon, 1+\epsilon]$（$\epsilon=0.2$）下，一个初始概率仅 $0.01$ 的低概率探索 token 最多只能涨到 $0.012$，而一个初始概率 $0.9$ 的高概率 token 却能一步涨到 $1.0$，这会系统性加速\textbf{熵坍缩}。DAPO 将上限放宽（如 $\epsilon_{\mathrm{low}} = 0.2, \epsilon_{\mathrm{high}} = 0.28$），专门给罕见的“顿悟推理词”（如“等等，让我再验算一遍”）留出足够的概率攀升空间。
  \item \textbf{动态采样（Dynamic Sampling）}：剔除组内准确率 $0\%$（全错）和 $100\%$（全对）的零梯度提示，动态补采直到有效批次填满，避免训练后期因题目越做越熟导致有效批量塌缩。
  \item \textbf{消除长度偏差与标准差偏差（Dr. GRPO）}：观察式~\eqref{eq:grpo-obj} 中的 $\frac{1}{|o_i|}$——对于答错的负优势回答（$\hat{A}_i < 0$），拉长回答长度 $|o_i|$ 反而会稀释每个错词受到的惩罚，无意间诱导出“越错越冗长”的长度膨胀！同时，除以 $\sigma_R$ 会给极易或极难（方差小）的题目赋予过大权重。Dr. GRPO 主张改为用固定最大长度常数 $L_{\max}$ 归一化并去掉 $\sigma_R$ 分母，令无偏性更接近纯策略梯度。
\end{enumerate}
\end{extBox}

\section*{本章小结与核心脉络}
\addcontentsline{toc}{section}{本章小结与核心脉络}

本章完成了现代策略优化从“危险的裸跑”到“工业级精密引擎”的跃迁：
\begin{itemize}
  \item \textbf{问题诊断}：多轮复用旧样本时，未加约束的重要性采样目标 $L^{\mathrm{CPI}}(\theta) = \hat{\E}_t[r_t(\theta)\hat{A}_t]$ 会因状态分布失配与无界概率比杠杆导致策略剧烈崩塌。
  \item \textbf{信任域基石（TRPO）}：通过性能差异引理证明了含最大 KL 散度惩罚的 MM 单调改进下界；借由二阶泰勒展开得到 Fisher 信息矩阵与自然策略梯度步长 $\Delta\theta^* \propto F^{-1} g$。
  \item \textbf{一阶革命（PPO-Clip）}：通过悲观下界 $L^{\mathrm{CLIP}}(\theta) = \hat{\E}_t[\min(r_t\hat{A}_t, \clip(r_t, 1-\epsilon, 1+\epsilon)\hat{A}_t)]$ 构筑出“只截断激进外冲、不截断反向纠偏”的单向棘轮；在 1D 连续 LQR 实测中彻底消除了未裁剪 IS-PG 的反复骤崩（$-7.26$），将终段回报稳步推升至 $-0.36$。
  \item \textbf{大模型跃迁（GRPO）}：用同题 $G$ 条采样的组内相对归一化 $\hat{A}_i = (R_i - \bar{R})/(\sigma_R + \varepsilon)$ 彻底甩掉庞大的 Critic 网络，配合非负无偏的 Schulman $k_3$ KL 估计量，成为前沿大模型推理强化的核心引擎。
\end{itemize}

现在，我们已经手握 PPO 与 GRPO 两把无往不利的策略优化利剑。但对大语言模型而言，一个更为根本的问题浮出水面：\textbf{在开放的对话、写作、安全与复杂推理世界里，环境并不会天然吐出标量奖励 $r_t$——我们究竟该如何把模糊的人类偏好与客观正确性，转化为机器可优化的数学目标？有没有办法甚至连显式奖励模型和在线强化学习采样也一并省去？} 下一章，我们将迈入大模型后训练与对齐的殿堂：SFT、RLHF、DPO 与 RLVR。

\section*{本章分层习题}
\addcontentsline{toc}{section}{本章分层习题}

\exTier{A 层：概念与公式手算}{熟练掌握概率比、PPO-Clip 分段取值与 GRPO 组内优势}

\begin{enumerate}
  \item[\textbf{6.3.1}] \textbf{PPO-Clip 六种边界情形逐一验算}：设裁剪参数 $\epsilon = 0.2$（即信任区间为 $[0.8, 1.2]$）。对下表中的 6 个样本，分别计算未裁剪值 $r_t \hat{A}_t$、裁剪值 $\clip(r_t, 0.8, 1.2)\hat{A}_t$、PPO-Clip 目标值 $f(r_t; \hat{A}_t)$，以及关于 $r_t$ 的导数 $\frac{\partial f}{\partial r_t}$（标明是否发生梯度截断）：
  \begin{center}
  \footnotesize
  \begin{tabular}{ccccccc}
  \toprule
  样本编号 & (1) & (2) & (3) & (4) & (5) & (6) \\
  \midrule
  概率比 $r_t$ & $0.6$ & $1.1$ & $1.5$ & $0.6$ & $1.1$ & $1.5$ \\
  优势 $\hat{A}_t$ & $+2.0$ & $+2.0$ & $+2.0$ & $-2.0$ & $-2.0$ & $-2.0$ \\
  \bottomrule
  \end{tabular}
  \end{center}
  \exHint{对照命题~\ref{prop:ppo-cases}：特别注意样本 (1)（好动作被误压至 $0.6$）与样本 (6)（坏动作被误抬至 $1.5$）虽然都在 $[0.8, 1.2]$ 之外，但外层 $\min$ 均保留了非零纠偏梯度！真正梯度归零的只有样本 (3) 与样本 (4)。}

  \item[\textbf{6.3.2}] \textbf{GRPO 群体优势与 $k_3$ KL 估计量手算}：对同一道数学题，旧策略采样 $G = 4$ 条回答，验证器给出二元奖励 $(R_1, R_2, R_3, R_4) = (1, 0, 1, 0)$。
  \begin{enumerate}
    \item 计算组内均值 $\bar{R}$、总体标准差 $\sigma_R$ 及每条回答的 GRPO 相对优势 $\hat{A}_i$（忽略 $\varepsilon$）；
    \item 若某条回答的某个 token 上当前策略概率 $\pi_\theta = 0.5$、参考策略概率 $\pi_{\mathrm{ref}} = 0.25$，分别计算朴素对数比 $k_1 = \ln(\pi_\theta/\pi_{\mathrm{ref}})$ 与 Schulman 非负估计量 $k_3 = \frac{\pi_{\mathrm{ref}}}{\pi_\theta} - \ln\frac{\pi_{\mathrm{ref}}}{\pi_\theta} - 1$ 的精确数值；若 $\pi_\theta = 0.25, \pi_{\mathrm{ref}} = 0.5$，两者又分别为多少？
  \end{enumerate}
  \exHint{(a) $\bar{R} = 0.5, \sigma_R = 0.5$，故答对者 $\hat{A}_i = +1$、答错者 $\hat{A}_i = -1$；(b) 注意第二种情形下 $k_1 = -\ln 2 < 0$，而 $k_3 = 2 - \ln 2 - 1 \approx 0.307 > 0$。}
\end{enumerate}

\exTier{B 层：严谨数学证明}{深入理解 KL 二阶几何、自然梯度与控制变量方差分析}

\begin{enumerate}
  \item[\textbf{6.3.3}] \textbf{KL 散度的 Hessian 严格等于 Fisher 信息矩阵}：设参数化概率分布族 $\pi_\theta(a \mid s)$ 对 $\theta$ 二阶连续可微，且允许微分与求和/积分交换顺序。定义以 $\theta_0$ 为锚点的 KL 散度 $D(\theta) = \DKL(\pi_{\theta_0}(\cdot \mid s) \,\|\, \pi_\theta(\cdot \mid s))$。
  \begin{enumerate}
    \item 严格证明 $\nabla_\theta D(\theta)\big|_{\theta=\theta_0} = 0$；
    \item 严格证明 Hessian 矩阵满足 $\nabla_\theta^2 D(\theta)\big|_{\theta=\theta_0} = -\E_{a \sim \pi_{\theta_0}}[\nabla_\theta^2 \ln \pi_\theta(a \mid s)]\big|_{\theta_0} = \E_{a \sim \pi_{\theta_0}}[\nabla_\theta \ln \pi_{\theta_0}(a \mid s)\,\nabla_\theta \ln \pi_{\theta_0}(a \mid s)^\top]$；
    \item 证明反向散度 $\tilde{D}(\theta) = \DKL(\pi_\theta(\cdot \mid s) \,\|\, \pi_{\theta_0}(\cdot \mid s))$ 在 $\theta = \theta_0$ 处具有完全相同的二阶 Hessian（即局部意义下正反向 KL 散度二阶等价！）。
  \end{enumerate}
  \exHint{对恒等式 $\sum_a \pi_\theta(a\mid s) = 1$ 连续求两次梯度，再利用恒等式 $\nabla^2 \ln \pi = \frac{\nabla^2 \pi}{\pi} - \nabla \ln \pi\,\nabla \ln \pi^\top$ 即可优雅证明。}

  \item[\textbf{6.3.4}] \textbf{TRPO 自然梯度步长的参数重参数化不变性}：考虑式~\eqref{eq:trpo-qp} 的信任域子问题：在约束 $\frac{1}{2}\Delta\theta^\top F\Delta\theta \le \delta$（设 $F \succ 0$）下最大化 $g^\top \Delta\theta$。
  \begin{enumerate}
    \item 用拉格朗日乘子法完整推导出闭式最优解式~\eqref{eq:npg-step}：$\Delta\theta^* = \sqrt{\frac{2\delta}{g^\top F^{-1} g}} F^{-1} g$；
    \item 设我们对参数做可逆线性坐标变换 $\psi = A\theta$（$A \in \R^{d \times d}$ 可逆）。写出新坐标系下的梯度 $g_\psi$ 与 Fisher 矩阵 $F_\psi$，并证明自然梯度步长满足 $\Delta\psi^* = A \Delta\theta^*$（而普通欧氏梯度步长 $\Delta\theta = \eta g$ 不具备这种坐标不变性！）。
  \end{enumerate}
  \exHint{链式法则给出 $g_\psi = A^{-\top} g$，$F_\psi = A^{-\top} F A^{-1}$；代入 $F_\psi^{-1} g_\psi = A F^{-1} A^\top A^{-\top} g = A F^{-1} g$，且二次型标量 $g_\psi^\top F_\psi^{-1} g_\psi = g^\top F^{-1} g$ 严格不变！}
\end{enumerate}

\exTier{C 层：机制剖析与反例设计}{辨识去 $\min$ 裁剪缺陷、二元奖励分布与长度偏差}

\begin{enumerate}
  \item[\textbf{6.3.5}] \textbf{反例构造：删去 PPO-Clip 外层 $\min$ 的灾难后果}：假设某同学删掉外层 $\min$，将目标简化为 $\tilde{L}(\theta) = \hat{\E}_t\!\big[\clip(r_t(\theta), 1-\epsilon, 1+\epsilon)\hat{A}_t\big]$。
  \begin{enumerate}
    \item 画出 $\hat{A}_t > 0$ 与 $\hat{A}_t < 0$ 时 $\tilde{L}$ 随 $r_t$ 变化的图像，指出它在哪两个区间与真正的 PPO-Clip 分道扬镳；
    \item 构造一个仅含两个样本 $(s_1, a_1, \hat{A}_1 = +1)$ 与 $(s_2, a_2, \hat{A}_2 = -1)$ 的一维参数场景，说明当一步更新意外把坏动作 $(s_2, a_2)$ 的概率比推到 $r_2 = 1.5$ 后，删去 $\min$ 的目标函数为何会对这个严重错误“彻底装瞎”（梯度为零），拒绝将其拉回信任域！
  \end{enumerate}
  \exHint{当 $\hat{A}_2 = -1 < 0$ 且 $r_2 = 1.5 > 1+\epsilon$ 时，$\clip(r_2)\hat{A}_2 = -1.2$ 是常数，导数为 $0$；坏动作越高反而不再受罚，策略安全防线瞬间洞穿。}

  \item[\textbf{6.3.6}] \textbf{GRPO 在二元奖励 $R_i \in \{0, 1\}$ 下的难度加权效应分析}：在可验证数学题训练中，设题目 $q$ 当前的通过率为 $p \in (0, 1)$，采样组容量 $G$ 足够大，使得组内经验均值 $\bar{R} \approx p$、标准差 $\sigma_R \approx \sqrt{p(1-p)}$。
  \begin{enumerate}
    \item 计算答对样本（$R_i = 1$）与答错样本（$R_i = 0$）的 GRPO 归一化优势 $\hat{A}_+$ 与 $\hat{A}_-$；
    \item 计算单道题目上所有答对样本贡献的总正权重 $p \cdot \hat{A}_+$，并讨论当 $p \to 0^+$（极难题）或 $p \to 1^-$（极易题）时，除以 $\sigma_R$ 如何改变单条答对/答错轨迹的梯度幅值（对比 Dr. GRPO 去掉 $\sigma_R$ 后的表现）。
  \end{enumerate}
  \exHint{$\hat{A}_+ = \sqrt{\frac{1-p}{p}}$，$\hat{A}_- = -\sqrt{\frac{p}{1-p}}$。当 $p = 0.01$ 时，唯一答对的那条样本优势高达 $\approx +9.95$，若偶发“蒙对答案但推理胡扯”的假阳性轨迹，单条样本便会产生近 $10\times$ 的剧烈梯度冲击！}
\end{enumerate}

\exTier{D 层：编程实验与系统延伸}{复现裁剪消融实验并实现轻量级 GRPO}

\begin{enumerate}
  \item[\textbf{6.3.7}] \textbf{编程实验：PPO 裁剪阈值 $\epsilon$ 与复用轮数 $K$ 的相图扫描}：修改 \texttt{code/vol6\_experiments.py} 中的 \texttt{exp3\_ppo\_vs\_unclipped()}：
  \begin{enumerate}
    \item 固定 $K = 10$，扫描 $\epsilon \in \{0.05, 0.1, 0.2, 0.4, 0.8, \infty\}$，绘制并记录各配置前 10 轮爬升速度、中后期最差回撤值与末 10 轮均值，验证 $\epsilon = 0.2$ 为何处于速度与稳定性的黄金甜点区；
    \item 尝试在 PPO-Clip 循环中加入代码清单里的 \textbf{KL 早停}（\texttt{max\_kl = 0.05}），统计每轮实际平均执行了多少次 epoch 才触发早停。
  \end{enumerate}
  \exHint{$\epsilon = 0.05$ 会极其稳定但前期爬升明显变慢；当 $\epsilon \ge 0.8$ 时，裁剪边界过于宽松，中后期会重新出现类似未裁剪 IS-PG 的回撤尖峰。}

  \item[\textbf{6.3.8}] \textbf{编程实验：在上下文赌博机上亲手实现无 Critic 的 GRPO}：构造一个含 32 个上下文状态、8 个离散动作的随机赌博机任务（每个状态仅 1 个动作奖励为 $1$，其余为 $0$）。使用本文给出的 \texttt{grpo\_loss} 函数，对每个状态采样 $G = 8$ 个动作执行 GRPO 训练（无需任何价值网络），并在 $\beta \in \{0.0, 0.05, 0.5\}$ 下观察平均准确率与策略到初始均匀参考策略的 KL 散度演化。思考：如果把奖励从客观 $\{0, 1\}$ 换成“两条回答谁更好”的人类偏好对，能否绕开采样与奖励模型直接优化策略？这正是第~\ref{ch:posttrain-rlhf-dpo} 章 DPO 的起点。
  \exHint{当 $\beta = 0$ 时策略会收敛到接近 $100\%$ 确定性选择最优动作（KL $\to \ln 8 \approx 2.08$）；当 $\beta = 0.5$ 时策略会停在柔和的玻尔兹曼分布上——试猜想它的闭式解与第~\ref{ch:posttrain-rlhf-dpo} 章定理的关系！}
\end{enumerate}

\chapter[大模型后训练：SFT、RLHF、DPO 与 RLVR]{大模型后训练与对齐：SFT、RLHF、DPO 与推理时缩放}
\label{ch:posttrain-rlhf-dpo}

\chapEpigraph{如果我们为了达成目标而使用一台无法有效干预其运行的机械装置，那我们最好无比确信，放入机器里的目标正是我们真正渴望的目标。}{Norbert Wiener（诺伯特·维纳，《科学》，1960）}

\begin{chapIntro}{本章导读：给“预训练怪兽”装上人类价值观与深度推理脊梁}
在第 5 卷中，我们用万亿词元和下一个词预测（Next-Token Prediction）锻造出了知识渊博的预训练基座模型。然而，如果你直接问一个刚完成预训练的基座模型“请用通俗语言解释拉格朗日乘子”，它往往不会老实作答，反而可能顺着互联网语料的统计惯性继续帮你补写：“……问题 2：什么是 KKT 条件？问题 3：……”——因为它的目标只是\textbf{模仿互联网文档的概率分布}，而不是\textbf{帮助人类、诚实作答、安全拒答或独立攻克难题}。

如何把一座沉睡的“语言世界模拟器”，点化为听得懂指令、守得住底线、甚至能自主反思与推导数学证明的现代智能体？这就是\textbf{大模型后训练（Post-Training）与对齐（Alignment）}的核心使命。本章将接续第~\ref{ch:trpo-ppo-grpo} 章的 PPO 与 GRPO 武器库，完成四重跨越：
\begin{enumerate}
  \item \textbf{后训练全景与监督微调（SFT）}：理清预训练 $\to$ SFT $\to$ 偏好对齐 / 可验证推理强化的全链路，剖析 SFT 仅靠模仿学习为何遭遇能力天花板。
  \item \textbf{Bradley--Terry 偏好建模与 KL 正则 RLHF 闭式最优解}：调用第 4 卷第 2 章的拉格朗日乘子法与信息论 Gibbs 不等式，严格证明 KL 正则 RLHF 的全局唯一最优策略 $\pi^*(y\mid x) = \frac{1}{Z(x)}\pi_{\mathrm{ref}}(y\mid x)\exp\big(\frac{1}{\beta}r(x,y)\big)$。
  \item \textbf{直接偏好优化（DPO）天书级推导}：反解隐式奖励并代入 Bradley--Terry 模型，亲眼见证无法计算的配分函数 $Z(x)$ 如何\textbf{恒等抵消}；进而深入剖析 DPO 梯度的自适应错判权重 $\sigma(\hat{r}_\theta(x, y_l) - \hat{r}_\theta(x, y_w))$。
  \item \textbf{奖励黑客（Goodhart 定律）与可验证奖励推理（RLVR）}：直面维纳六十年前的预言，定量刻画奖励模型过优化曲线，并介绍以 GRPO + 可验证规则与 Best-of-$N$ 为代表的推理时计算缩放范式；最后基于 \texttt{data/dpo\_align.csv} 实测验证 DPO 中温度 $\beta$ 对“真实奖励—KL 散度”帕累托前沿的精密调控。
\end{enumerate}
\end{chapIntro}

\section{大模型后训练全景：从预训练底座到指令微调（SFT）}
\label{sec:posttrain-sft}

\subsection{为什么预训练底座必须经过后训练？}

预训练阶段在海量无标注文本上最小化自回归交叉熵损失 $\mathcal{L}_{\mathrm{PT}}(\theta) = -\E[\sum_t \ln \pi_\theta(x_t \mid x_{<t})]$。这赋予了模型极其广博的世界知识、语法直觉与上下文学习潜能，但同时留下了根本的\textbf{目标错位（Objective Misalignment）}：
\begin{itemize}
  \item \textbf{行为模仿而非意图满足}：预训练分布平等地拟合优质教材与低质论坛回帖、严谨证明与胡言乱语、友善解答与危险教程。
  \item \textbf{缺乏对话角色边界}：底座模型不知道何时该作为助手结束回答，何时该拒绝恶意诱导，也不懂得在面对复杂多步数学题时“停下来检查自己有没有算错”。
\end{itemize}
因此，现代前沿大模型的训练被清晰划分为两大阶段：\textbf{预训练}负责把世界知识与潜在推理能力“压缩进权重”；\textbf{后训练}则负责激活并塑形这些能力，使其对齐人类的\textbf{有用性（Helpfulness）、诚实性（Honesty）与无害性（Harmlessness），即 3H 原则}，并在数学与代码等客观领域突破模仿上限。

\subsection{监督微调（SFT）的数学机制与三大局限}

后训练的第一站通常是\textbf{监督微调（Supervised Fine-Tuning, SFT）}：收集数万至数百万条高质量的“指令—示范回答”对话对 $\mathcal{D}_{\mathrm{SFT}} = \{(x^{(i)}, y^{(i)})\}_{i=1}^N$，在底座模型参数之上继续做极大似然训练。工程上最关键的细节是\textbf{提示掩码（Prompt Loss Masking）}——损失函数\emph{只计算模型生成的回答词元 $y$}，而不对用户输入 $x$ 计算梯度：
\begin{equation}
  \mathcal{L}_{\mathrm{SFT}}(\theta)
  = -\E_{(x, y) \sim \mathcal{D}_{\mathrm{SFT}}}\!\left[\sum_{t=1}^{|y|} \ln \pi_\theta(y_t \mid x, y_{<t})\right].
  \label{eq:sft-loss}
\end{equation}
SFT 极其高效地教会了模型对话格式、指令遵循与基础写作风格（常被称为“把预训练积累的能力表面化”）。然而，单纯依赖 SFT 会迅速撞上三层天花板：
\begin{enumerate}
  \item \textbf{模仿上限（Behavior Cloning Ceiling）}：SFT 本质上是离线行为克隆，模型最多逼近标注者的示范水平，极难超越人类写手在有限时间里写出的标准答案。
  \item \textbf{缺少负样本对比信号}：式~\eqref{eq:sft-loss} 只告诉模型“应当提高 $y$ 的概率”，却从未显式惩罚那些\emph{表面流畅但包含幻觉、逻辑漏洞或安全隐患}的回答。
  \item \textbf{词元级局部损失与序列级全局目标脱节}：交叉熵平等对待每个词元，但在数学证明或代码生成中，写错一个正负号或漏掉一个边界判断，虽只占整段输出 $0.2\%$ 的词元损失，却足以令整条回答彻底作废！
\end{enumerate}

为了突破 SFT 的模仿天花板，现代后训练分化出两条相辅相成的对齐与强化路线：\textbf{基于人类偏好的对齐（RLHF 与 DPO）}以及\textbf{基于可验证客观奖励的强化学习（RLVR）}。图~\ref{fig:rlhf-vs-dpo-pipeline} 全景展示了这三条路线的系统流向与架构差异。

\begin{figure}[htbp]
\centering
\resizebox{0.98\linewidth}{!}{%
\begin{tikzpicture}[
  node distance=5.5mm and 7.5mm,
  base/.style={rectangle, rounded corners=1.6mm, draw=slate, thick, fill=codebg,
    minimum width=23mm, minimum height=10.5mm, align=center, font=\sffamily\footnotesize, text=ink, inner sep=3.5pt},
  sft/.style={rectangle, rounded corners=1.6mm, draw=navy, thick, fill=navyLight,
    minimum width=25mm, minimum height=10.5mm, align=center, font=\sffamily\footnotesize, text=ink, inner sep=3.5pt},
  rlhf/.style={rectangle, rounded corners=1.6mm, draw=ochre, thick, fill=ochreLight,
    minimum width=28mm, minimum height=10.5mm, align=center, font=\sffamily\footnotesize, text=ink, inner sep=3.5pt},
  dpo/.style={rectangle, rounded corners=1.6mm, draw=teal, very thick, fill=tealLight,
    minimum width=28mm, minimum height=10.5mm, align=center, font=\sffamily\footnotesize, text=ink, inner sep=3.5pt},
  rlvr/.style={rectangle, rounded corners=1.6mm, draw=plum, thick, fill=plumLight,
    minimum width=28mm, minimum height=10.5mm, align=center, font=\sffamily\footnotesize, text=ink, inner sep=3.5pt},
  arr/.style={-{Stealth[length=2.1mm]}, thick, slate},
  arrT/.style={-{Stealth[length=2.1mm]}, thick, teal},
  arrO/.style={-{Stealth[length=2.1mm]}, thick, ochre},
  arrP/.style={-{Stealth[length=2.1mm]}, thick, plum}
]
  % Shared left stages
  \node[base] (pt) {\textbf{预训练底座}\\Next-Token MLE};
  \node[sft, right=8mm of pt] (sft) {\textbf{SFT 参考模型}\\$\pi_{\mathrm{ref}} = \pi_{\mathrm{SFT}}$};
  \draw[arr] (pt) -- (sft);

  % Branch 1 (Top): Classic RLHF
  \node[rlhf, above right=11mm and 12mm of sft] (pref1) {\textbf{人类偏好对}\\$(x, y_w \succ y_l)$};
  \node[rlhf, right=8mm of pref1] (rm) {\textbf{奖励模型 RM}\\Bradley--Terry MLE};
  \node[rlhf, right=8mm of rm] (ppo) {\textbf{在线 PPO / RLHF}\\$\max \E[r_\psi] - \beta\DKL$};
  \draw[arrO] (sft.north east) -- (pref1.west);
  \draw[arrO] (pref1) -- (rm);
  \draw[arrO] (rm) -- (ppo);
  \node[above=1mm of rm, font=\sffamily\scriptsize\bfseries, color=ochre] {路线一：经典三阶段 RLHF（需显式 RM $+$ 在线采样）};

  % Branch 2 (Middle): DPO
  \node[dpo, right=12mm of sft] (pref2) {\textbf{人类偏好对}\\$(x, y_w \succ y_l)$};
  \node[dpo, right=19mm of pref2, minimum width=52mm] (dpoopt) {\textbf{直接偏好优化（DPO）}\\闭式消去 $Z(x)$，纯分类损失直训 $\pi_\theta$};
  \draw[arrT] (sft) -- (pref2);
  \draw[arrT] (pref2) -- node[above, font=\scriptsize, text=tealDark, align=center] {省去 RM 与\\在线采样} (dpoopt);
  \node[above=1mm of dpoopt, font=\sffamily\scriptsize\bfseries, color=tealDark] {路线二：离线两阶段 DPO（闭式等价 shortcut）};

  % Branch 3 (Bottom): RLVR
  \node[rlvr, below right=11mm and 12mm of sft] (ver) {\textbf{客观验证器}\\数学等价 / 单元测试};
  \node[rlvr, right=19mm of ver, minimum width=52mm] (grpo) {\textbf{RLVR（如 GRPO 推理强化）}\\组内 $G$ 条采样，二元奖励 $R \in \{0,1\}$};
  \draw[arrP] (sft.south east) -- (ver.west);
  \draw[arrP] (ver) -- node[above, font=\scriptsize, text=plum] {零主观标注} (grpo);
  \node[below=1mm of grpo, font=\sffamily\scriptsize\bfseries, color=plum] {路线三：可验证奖励强化（突破人类模仿上限，涌现长思维链）};

  % Reference KL dashed connections (well above Route 1 caption)
  \draw[dashed, thick, navy!70, -{Stealth[length=1.8mm]}]
    (sft.north) |- ([yshift=9.2mm]ppo.north) -- (ppo.north);
  \node[font=\scriptsize, text=navy, fill=white, inner sep=1.5pt] at ([yshift=9.2mm]pref1.north east) {冻结 $\pi_{\mathrm{ref}}$ 提供 KL 约束};
\end{tikzpicture}}
\caption{现代大模型后训练三大核心路线架构对照：经典三阶段 RLHF（需单独训练奖励模型并在线 RL 采样）、直接偏好优化 DPO（解析消去奖励模型与配分函数，用纯监督分类损失一步到位），以及面向数学与代码推理的可验证奖励强化学习 RLVR。}
\label{fig:rlhf-vs-dpo-pipeline}
\end{figure}

\section{人类偏好建模与 KL 正则 RLHF 的闭式最优解}
\label{sec:bt-rlhf}

\subsection{Bradley--Terry 偏好模型与奖励模型极大似然估计}

让人类标注员给一篇回答打出“$7.3$ 分还是 $7.8$ 分”这样的绝对标量分数，不仅跨标注员方差极大，连同一个人隔天也会标准漂移；但让人类比较同一提示 $x$ 下的两条候选回答 $y_1, y_2$ 并选出“哪一条更好”（记胜出者为 $y_w$，落败者为 $y_l$），标注一致性会大幅跃升。

统计心理物理学中的\textbf{Bradley--Terry（BT）成对比较模型}（1952）将这种离散胜负比较优雅地映射到了连续隐含奖励函数 $r^*(x, y) \in \R$ 上：

\begin{defBox}
\begin{definition}[Bradley--Terry 偏好模型与奖励模型训练目标]
\label{def:bt-rm}
假设存在潜在的真实人类偏好奖励函数 $r^*(x, y)$。给定提示 $x$ 与一对候选回答 $(y_w, y_l)$，人类更偏好 $y_w$ 胜过 $y_l$（记作 $y_w \succ y_l \mid x$）的概率建模为逻辑斯蒂 Sigmoid 差值：
\begin{equation}
  P^*(y_w \succ y_l \mid x)
  = \frac{\exp\!\big(r^*(x, y_w)\big)}{\exp\!\big(r^*(x, y_w)\big) + \exp\!\big(r^*(x, y_l)\big)}
  = \sigma\!\Big(r^*(x, y_w) - r^*(x, y_l)\Big),
  \label{eq:bt-model}
\end{equation}
其中 $\sigma(z) = \frac{1}{1 + e^{-z}}$。收集人类成对偏好数据集 $\mathcal{D}_{\mathrm{pref}} = \{(x^{(i)}, y_w^{(i)}, y_l^{(i)})\}_{i=1}^M$ 后，在 SFT 底座末层换上一个标量线性输出头构成参数化\textbf{奖励模型（Reward Model, RM）} $r_\psi(x, y)$，通过最小化负对数似然进行训练：
\begin{equation}
  \mathcal{L}_{\mathrm{RM}}(\psi)
  = -\E_{(x, y_w, y_l) \sim \mathcal{D}_{\mathrm{pref}}}\!\left[\ln \sigma\!\Big(r_\psi(x, y_w) - r_\psi(x, y_l)\Big)\right].
  \label{eq:rm-loss}
\end{equation}
\end{definition}
\end{defBox}

请注意式~\eqref{eq:bt-model} 一个看似不起眼、却将在下一节引爆奇迹的数学性质——\textbf{提示级平移不变性（Prompt-Dependent Shift Invariance）}：如果对任意仅依赖提示 $x$（不依赖回答 $y$）的标量函数 $c(x)$，构造新奖励函数
\[
  \tilde{r}(x, y) \triangleq r(x, y) + c(x),
\]
那么 $\tilde{r}(x, y_w) - \tilde{r}(x, y_l) \equiv r(x, y_w) - r(x, y_l)$，它诱导的 Bradley--Terry 偏好分布与原奖励函数\textbf{完全一模一样}！换言之，偏好数据只决定奖励函数的“相对高度差”，却对每个提示 $x$ 的基准零点 $c(x)$ 拥有完全的自由度。

\subsection{KL 正则 RLHF 目标与闭式最优策略定理}

拿到训练好的奖励模型 $r(x, y)$ 后，经典 RLHF 将每个提示 $x$ 视作单步上下文强化学习（Contextual Bandit）的初始状态，将整条回答 $y$ 视作动作，在最大化期望奖励的同时施加相对于 SFT 参考策略 $\pi_{\mathrm{ref}}$ 的 KL 散度约束，防止语言模型偏离自然语言分布太远而退化成胡言乱语的“高分乱码机”：
\begin{equation}
  \max_{\pi}\; J_{\mathrm{RLHF}}(\pi)
  \triangleq \E_{x \sim \mathcal{D}}\!\left[
    \E_{y \sim \pi(\cdot \mid x)}\!\big[r(x, y)\big]
    - \beta\,\DKL\!\Big(\pi(\cdot \mid x) \,\Big\|\, \pi_{\mathrm{ref}}(\cdot \mid x)\Big)
  \right], \qquad \beta > 0.
  \label{eq:rlhf-obj}
\end{equation}

\begin{recallBox}{前置知识衔接：拉格朗日乘子法与信息论 Gibbs 不等式}
以下定理的证明直接调用第 4 卷第 2 章（优化理论与信息论）的两块基石：
\begin{enumerate}
  \item \textbf{等式约束的拉格朗日乘子法}：在概率单纯形约束 $\sum_y \pi(y \mid x) = 1$ 下求极值，引入乘子 $\lambda(x)$ 构造增广泛函 $\mathcal{L}(\pi, \lambda)$ 并令各分量偏导为零。
  \item \textbf{KL 散度的 Gibbs 不等式}：对任意同支撑概率分布 $p, q$，恒有 $\DKL(p \,\|\, q) = \sum_y p(y)\ln\frac{p(y)}{q(y)} \ge 0$，且等号成立\textbf{当且仅当} $p(y) = q(y)$ 对所有 $y$ 成立。
\end{enumerate}
\end{recallBox}

令人拍案叫绝的是，虽然在 PPO 训练中我们用迭代采样去逼近式~\eqref{eq:rlhf-obj}，但如果不限制 $\pi$ 的参数化形式（在所有合法概率分布空间中优化），式~\eqref{eq:rlhf-obj} 竟存在一个\textbf{极其优雅的精确闭式全局唯一最优解}！下面我们给出拉格朗日乘子法与 KL 配方双重完整证明。

\begin{thmBox}
\begin{theorem}[KL 正则 RLHF 的精确闭式最优策略定理]
\label{thm:rlhf-optimal}
设对给定提示 $x$，参考分布 $\pi_{\mathrm{ref}}(y \mid x) > 0$ 且配分函数
\begin{equation}
  Z(x) \triangleq \sum_{y'} \pi_{\mathrm{ref}}(y' \mid x)\,\exp\!\left(\frac{1}{\beta}\,r(x, y')\right) < \infty.
  \label{eq:rlhf-partition}
\end{equation}
则 KL 正则 RLHF 目标~\eqref{eq:rlhf-obj} 在概率单纯形 $\{\pi(\cdot \mid x) : \pi(y\mid x) \ge 0,\; \sum_y \pi(y\mid x) = 1\}$ 上存在\textbf{唯一的全局最优策略}：
\begin{equation}
  \pi^*(y \mid x) = \frac{1}{Z(x)}\,\pi_{\mathrm{ref}}(y \mid x)\,\exp\!\left(\frac{1}{\beta}\,r(x, y)\right),
  \label{eq:rlhf-opt-policy}
\end{equation}
且该提示 $x$ 对应的最优目标函数值精确等于 $\beta \ln Z(x)$。
\end{theorem}
\begin{proof}
由于目标~\eqref{eq:rlhf-obj} 对不同提示 $x$ 完全可分离，只需固定任意一个 $x$，最大化单提示泛函
\[
  F(\pi) \triangleq \sum_y \pi(y \mid x)\, r(x, y) - \beta \sum_y \pi(y \mid x)\,\ln\frac{\pi(y \mid x)}{\pi_{\mathrm{ref}}(y \mid x)}
  \quad \text{s.t.} \quad \sum_y \pi(y \mid x) = 1.
\]

\textbf{步骤一（拉格朗日乘子法导出驻点）}：引入拉格朗日乘子 $\lambda(x) \in \R$，构造拉格朗日函数
\[
  \mathcal{L}\big(\pi(\cdot \mid x), \lambda(x)\big)
  = \sum_y \pi(y \mid x)\, r(x, y)
    - \beta \sum_y \pi(y \mid x)\,\ln\frac{\pi(y \mid x)}{\pi_{\mathrm{ref}}(y \mid x)}
    - \lambda(x)\!\left(\sum_y \pi(y \mid x) - 1\right).
\]
对每个回答 $y$ 的概率分量 $\pi(y \mid x)$ 求偏导并令其等于 $0$（利用 $\frac{d}{dp}(p\ln p) = \ln p + 1$）：
\[
  \frac{\partial \mathcal{L}}{\partial \pi(y \mid x)}
  = r(x, y) - \beta\!\left(\ln\frac{\pi(y \mid x)}{\pi_{\mathrm{ref}}(y \mid x)} + 1\right) - \lambda(x) = 0.
\]
移项整理对数比：
\[
  \ln\frac{\pi(y \mid x)}{\pi_{\mathrm{ref}}(y \mid x)}
  = \frac{1}{\beta}\,r(x, y) - \left(1 + \frac{\lambda(x)}{\beta}\right).
\]
两边取指数，得到
\[
  \pi^*(y \mid x)
  = \exp\!\left(-1 - \frac{\lambda(x)}{\beta}\right)\,\pi_{\mathrm{ref}}(y \mid x)\,\exp\!\left(\frac{1}{\beta}\,r(x, y)\right).
\]
注意前缀因子 $\exp\!\big(-1 - \frac{\lambda(x)}{\beta}\big)$ 与 $y$ 无关！代入归一化约束 $\sum_y \pi^*(y \mid x) = 1$，立刻解得
\[
  \exp\!\left(1 + \frac{\lambda(x)}{\beta}\right)
  = \sum_{y'} \pi_{\mathrm{ref}}(y' \mid x)\,\exp\!\left(\frac{1}{\beta}\,r(x, y')\right)
  = Z(x),
\]
回代即得精确驻点公式~\eqref{eq:rlhf-opt-policy}，且显然满足非负性 $\pi^*(y\mid x) > 0$。

\textbf{步骤二（信息论 KL 配方证明全局唯一最优性）}：为验证 $\pi^*$ 不仅是局部驻点，更是全局唯一最大值点，将定义好的合法分布 $\pi^*(y \mid x) = \frac{1}{Z(x)}\pi_{\mathrm{ref}}(y\mid x)\exp(\frac{1}{\beta}r(x,y))$ 代回原泛函 $F(\pi)$ 做恒等变形：
\begin{align*}
  F(\pi)
  &= -\beta \sum_y \pi(y \mid x)\!\left[\ln\frac{\pi(y \mid x)}{\pi_{\mathrm{ref}}(y \mid x)} - \frac{1}{\beta}\,r(x, y)\right] \\
  &= -\beta \sum_y \pi(y \mid x)\!\left[\ln\frac{\pi(y \mid x)}{\pi_{\mathrm{ref}}(y \mid x)\exp\!\big(\frac{1}{\beta}r(x,y)\big)}\right] \\
  &= -\beta \sum_y \pi(y \mid x)\!\left[\ln\frac{\pi(y \mid x)}{Z(x)\,\pi^*(y \mid x)}\right] \\
  &= -\beta \sum_y \pi(y \mid x)\,\ln\frac{\pi(y \mid x)}{\pi^*(y \mid x)} + \beta \ln Z(x)\sum_y \pi(y \mid x) \\
  &= -\beta\,\DKL\!\Big(\pi(\cdot \mid x) \,\Big\|\, \pi^*(\cdot \mid x)\Big) + \beta \ln Z(x).
\end{align*}
由 Gibbs 不等式，$\DKL(\pi \,\|\, \pi^*) \ge 0$ 且等号成立当且仅当 $\pi(\cdot \mid x) \equiv \pi^*(\cdot \mid x)$！又因 $\beta > 0$，故对一切概率分布 $\pi$ 均有 $F(\pi) \le \beta \ln Z(x) = F(\pi^*)$，全局唯一最优性证毕！
\end{proof}
\end{thmBox}

式~\eqref{eq:rlhf-opt-policy} 展现了一幅优美的统计物理图景：\textbf{最优对齐策略 $\pi^*$ 正是以参考策略 $\pi_{\mathrm{ref}}$ 为先验测度、以奖励 $r(x,y)$ 为负能量、以 $\beta$ 为温度的玻尔兹曼（Gibbs）倾斜分布！}
\begin{itemize}
  \item 当 $\beta \to +\infty$（极强 KL 约束）：$\frac{1}{\beta}r(x,y) \to 0$，最优策略 $\pi^* \to \pi_{\mathrm{ref}}$，完全退回保守的 SFT 参考模型。
  \item 当 $\beta \to 0^+$（几乎不约束 KL）：指数权重把全部概率质量瞬间坍缩到 $\argmax_y r(x, y)$ 上，哪怕 $r$ 中藏有一丝奖励黑客漏洞，模型也会毫不犹豫地一头扎进去！
\end{itemize}

遗憾的是，尽管我们写出了 $\pi^*(y\mid x)$ 的精确闭式解，却\textbf{无法直接用它来采样或计算概率}——因为分母中的配分函数 $Z(x) = \sum_{y' \in \mathcal{V}^*} \pi_{\mathrm{ref}}(y'\mid x)\exp(\frac{1}{\beta}r(x, y'))$ 需要对词表 $\mathcal{V}$（$|\mathcal{V}| \approx 10^5$）上所有可能长度的文本序列求和，其项数是超越天文数字的 $|\mathcal{V}|^L$！正因 $Z(x)$ 不可计算，传统 RLHF 才被迫先训奖励模型，再用昂贵的在线 PPO 迭代逼近 $\pi^*$。

直到 2023 年，斯坦福团队灵光乍现，发现了绕过 $Z(x)$ 的“天书”捷径。

\section{直接偏好优化（DPO）：精确消去配分函数的天书推导}
\label{sec:dpo}

\subsection{DPO 精确等价定理：配分函数 $Z(x)$ 的完美抵消}

Rafailov 等（2023）提出了一个颠覆性的逆向视角：既然无法由奖励 $r$ 直接算出策略 $\pi^*$（卡在 $Z(x)$ 上），\textbf{何不反过来——把公式~\eqref{eq:rlhf-opt-policy} 倒着解，用策略 $\pi^*$ 来表示奖励 $r^*$，再直接塞进 Bradley--Terry 偏好似然里呢？}

\begin{thmBox}
\begin{theorem}[直接偏好优化（DPO）精确推导定理，Rafailov et al., 2023]
\label{thm:dpo-derivation}
在 Bradley--Terry 偏好模型~\eqref{eq:bt-model} 与 KL 正则 RLHF 目标~\eqref{eq:rlhf-obj} 的框架下，设 $\pi^*$ 为真实奖励函数 $r^*$ 所诱导的唯一闭式最优策略~\eqref{eq:rlhf-opt-policy}。则：
\begin{enumerate}
  \item 真实奖励函数可由最优策略与参考策略精确反解为
  \begin{equation}
    r^*(x, y) = \beta \ln \frac{\pi^*(y \mid x)}{\pi_{\mathrm{ref}}(y \mid x)} + \beta \ln Z(x);
    \label{eq:dpo-invert-r}
  \end{equation}
  \item 人类成对偏好概率完全由最优策略 $\pi^*$ 与参考策略 $\pi_{\mathrm{ref}}$ 决定，\textbf{不可计算的配分函数 $Z(x)$ 恒等消去}：
  \begin{equation}
    P^*(y_w \succ y_l \mid x)
    = \sigma\!\left(\beta \ln \frac{\pi^*(y_w \mid x)}{\pi_{\mathrm{ref}}(y_w \mid x)} - \beta \ln \frac{\pi^*(y_l \mid x)}{\pi_{\mathrm{ref}}(y_l \mid x)}\right);
    \label{eq:dpo-bt-policy}
  \end{equation}
  \item 从而无需单独训练奖励模型、也无需在线强化学习采样，直接对参数化策略 $\pi_\theta$ 最小化如下纯二元交叉熵式的 \textbf{DPO 损失函数}，即可直接逼近 KL 正则 RLHF 的最优策略：
  \begin{equation}
    \mathcal{L}_{\mathrm{DPO}}(\theta;\, \pi_{\mathrm{ref}})
    = -\E_{(x, y_w, y_l) \sim \mathcal{D}_{\mathrm{pref}}}\!\left[
      \ln \sigma\!\left(
        \beta \ln \frac{\pi_\theta(y_w \mid x)}{\pi_{\mathrm{ref}}(y_w \mid x)}
        - \beta \ln \frac{\pi_\theta(y_l \mid x)}{\pi_{\mathrm{ref}}(y_l \mid x)}
      \right)
    \right].
    \label{eq:dpo-loss}
  \end{equation}
\end{enumerate}
\end{theorem}
\begin{proof}
对定理~\ref{thm:rlhf-optimal} 的闭式解式~\eqref{eq:rlhf-opt-policy} 两边先除以 $\pi_{\mathrm{ref}}(y \mid x)$ 再取自然对数：
\[
  \ln \frac{\pi^*(y \mid x)}{\pi_{\mathrm{ref}}(y \mid x)} = \frac{1}{\beta}\,r^*(x, y) - \ln Z(x).
\]
移项并同乘 $\beta > 0$，即得反解公式~\eqref{eq:dpo-invert-r}：$r^*(x, y) = \beta \ln \frac{\pi^*(y\mid x)}{\pi_{\mathrm{ref}}(y\mid x)} + \beta \ln Z(x)$。

现在将其分别代入胜出回答 $y_w$ 与落败回答 $y_l$，并在 Bradley--Terry 模型~\eqref{eq:bt-model} 中作差：
\begin{align*}
  r^*(x, y_w) - r^*(x, y_l)
  &= \left(\beta \ln \frac{\pi^*(y_w \mid x)}{\pi_{\mathrm{ref}}(y_w \mid x)} + \beta \ln Z(x)\right)
   - \left(\beta \ln \frac{\pi^*(y_l \mid x)}{\pi_{\mathrm{ref}}(y_l \mid x)} + \beta \ln Z(x)\right) \\
  &= \beta \ln \frac{\pi^*(y_w \mid x)}{\pi_{\mathrm{ref}}(y_w \mid x)} - \beta \ln \frac{\pi^*(y_l \mid x)}{\pi_{\mathrm{ref}}(y_l \mid x)}.
\end{align*}
因为 $Z(x)$ 仅依赖提示 $x$ 而与具体回答 $y_w, y_l$ 无关，借助 Bradley--Terry 模型的提示级平移不变性，这一原本不可计算的配分项分毫不差地彻底抵消！将差值代入 $\sigma(\cdot)$ 即得式~\eqref{eq:dpo-bt-policy}；最后用参数化语言模型 $\pi_\theta$ 替代 $\pi^*$ 并对偏好数据集 $\mathcal{D}_{\mathrm{pref}}$ 执行极大似然估计（最小化负对数似然），即得 DPO 损失函数~\eqref{eq:dpo-loss}，证毕。
\end{proof}
\end{thmBox}

回顾这段仅仅十行左右的推导，你很难不被它的数学美感折服——难怪它被称为大模型对齐时代的“天书证明”之一：原本需要四套大模型协同、极易发散崩溃的“RM 训练 $+$ 在线 PPO 采样”复杂流水线，经由变量替换，瞬间坍缩为一个仅含策略 $\pi_\theta$ 与冻结参考 $\pi_{\mathrm{ref}}$、像逻辑回归一样稳定收敛的离线监督损失！

\subsection{DPO 梯度的深层力学剖析：隐式奖励与自适应错判权重}

为了真正看透 DPO 在训练中究竟如何雕琢语言模型，让我们亲手求出式~\eqref{eq:dpo-loss} 关于参数 $\theta$ 的精确梯度。先定义策略 $\pi_\theta$ 所隐式表达的奖励函数（称为\textbf{隐式奖励, Implicit Reward}）：
\begin{equation}
  \hat{r}_\theta(x, y) \triangleq \beta \ln \frac{\pi_\theta(y \mid x)}{\pi_{\mathrm{ref}}(y \mid x)}.
  \label{eq:implicit-reward}
\end{equation}

\begin{thmBox}
\begin{proposition}[DPO 梯度的隐式奖励加权结构]
\label{prop:dpo-grad}
DPO 损失函数~\eqref{eq:dpo-loss} 关于策略参数 $\theta$ 的精确梯度为
\begin{equation}
\begin{aligned}
  \nabla_\theta \mathcal{L}_{\mathrm{DPO}}(\theta;\, \pi_{\mathrm{ref}})
  = -\beta\,\E_{(x, y_w, y_l) \sim \mathcal{D}_{\mathrm{pref}}}\Bigg[
    &\underbrace{\sigma\!\Big(\hat{r}_\theta(x, y_l) - \hat{r}_\theta(x, y_w)\Big)}_{\text{自适应错判权重 } w_\theta \in (0, 1)} \\
    &\cdot \Big(\nabla_\theta \ln \pi_\theta(y_w \mid x) - \nabla_\theta \ln \pi_\theta(y_l \mid x)\Big)
  \Bigg].
\end{aligned}
\label{eq:dpo-grad}
\end{equation}
\end{proposition}
\begin{proof}
记逻辑边际 $\Delta_\theta \triangleq \hat{r}_\theta(x, y_w) - \hat{r}_\theta(x, y_l)$。利用微积分恒等式
\[
  \frac{d}{dz}\big[-\ln \sigma(z)\big] = -\frac{\sigma'(z)}{\sigma(z)} = -\frac{\sigma(z)(1-\sigma(z))}{\sigma(z)} = -\big(1 - \sigma(z)\big) = -\sigma(-z),
\]
并注意参考策略 $\pi_{\mathrm{ref}}$ 已完全冻结（$\nabla_\theta \ln \pi_{\mathrm{ref}} = 0$），故 $\nabla_\theta \Delta_\theta = \beta\big(\nabla_\theta \ln \pi_\theta(y_w\mid x) - \nabla_\theta \ln \pi_\theta(y_l\mid x)\big)$。由链式法则立得
\[
  \nabla_\theta \big[-\ln \sigma(\Delta_\theta)\big]
  = -\sigma(-\Delta_\theta)\,\nabla_\theta \Delta_\theta
  = -\beta\,\sigma\!\Big(\hat{r}_\theta(x, y_l) - \hat{r}_\theta(x, y_w)\Big)\Big(\nabla_\theta \ln \pi_\theta(y_w \mid x) - \nabla_\theta \ln \pi_\theta(y_l \mid x)\Big),
\]
取期望即得式~\eqref{eq:dpo-grad}。
\end{proof}
\end{thmBox}

请仔细品味式~\eqref{eq:dpo-grad} 中三个因子的精妙协作：
\begin{enumerate}
  \item \textbf{推高好回答、压低坏回答}：因为梯度下降沿 $-\nabla_\theta \mathcal{L}_{\mathrm{DPO}}$ 方向更新，每一小步都在沿 $+\nabla_\theta \ln \pi_\theta(y_w \mid x)$ 拉高胜出回答的对数似然，同时沿 $-\nabla_\theta \ln \pi_\theta(y_l \mid x)$ 压低落败回答的对数似然。
  \item \textbf{为什么不能直接去掉权重 $w_\theta$ 做朴素对比？}如果删去系数 $\sigma(\hat{r}_\theta(x, y_l) - \hat{r}_\theta(x, y_w))$，直接最大化 $\ln \pi_\theta(y_w\mid x) - \ln \pi_\theta(y_l\mid x)$，会引发灾难性的\textbf{语言模型崩坏}：模型会无休止地把 $y_l$ 的概率往下砸向 $0$，连带破坏通用语法结构！
  \item \textbf{动态防过冲权重 $w_\theta = \sigma(\hat{r}_\theta(x, y_l) - \hat{r}_\theta(x, y_w))$}：
  \begin{itemize}
    \item 当模型\textbf{判反了}（坏回答的隐式奖励比好回答还高，即 $\hat{r}_l \gg \hat{r}_w$），权重 $w_\theta \to 1$，施以全力纠偏；
    \item 当模型\textbf{已经拉开足够边际}（$\hat{r}_w \gg \hat{r}_l$），权重 $w_\theta \to 0$，梯度自动衰减熄火，防止过度拟合某对样本而摧毁语言流畅性！
  \end{itemize}
\end{enumerate}

\subsection{简洁可运行 PyTorch DPO 核心实现}

在自回归大语言模型中，整条回答 $y = (y_1, \dots, y_L)$ 的联合对数概率等于各回复词元对数概率之和：
\[
  \ln \pi_\theta(y \mid x) = \sum_{t=1}^{|y|} \ln \pi_\theta(y_t \mid x, y_{<t}).
\]
下面的二十余行代码即为现代对齐库（如 TRL）中 DPO 损失的精确内核：

\begin{lstlisting}[caption={直接偏好优化（DPO）核心损失与隐式准确率监控（简洁 PyTorch 实现）}]
def seq_logprobs(logits, labels, resp_mask):
    """计算回复部分整句联合对数概率 sum_t log pi(y_t | x, y_<t)."""
    # logits: (B, L, V), labels: (B, L), resp_mask: (B, L) 仅在回复 token 处为 1
    logp_all = F.log_softmax(logits[:, :-1, :], dim=-1)
    target   = labels[:, 1:].unsqueeze(-1)
    tok_logp = logp_all.gather(dim=-1, index=target).squeeze(-1)
    return (tok_logp * resp_mask[:, 1:]).sum(dim=-1)          # (B,)

def dpo_loss(pi_w_logp, pi_l_logp, ref_w_logp, ref_l_logp, beta=0.1):
    """
    所有输入形状均为 (B,); ref_w_logp 与 ref_l_logp 由冻结参考模型在 no_grad 下算得
    """
    # 隐式奖励 r_hat = beta * (log pi_theta - log pi_ref)
    r_w = beta * (pi_w_logp - ref_w_logp)
    r_l = beta * (pi_l_logp - ref_l_logp)
    margin = r_w - r_l                                        # 等价于 beta*(dw - dl)

    loss = -F.logsigmoid(margin).mean()                       # 数值最稳的 -log(sigma(margin))
    with torch.no_grad():
        win_acc = (margin > 0).float().mean().item()          # 偏好对排序准确率
        mean_margin = margin.mean().item()
    return loss, win_acc, mean_margin
\end{lstlisting}

\section{奖励黑客（Goodhart 定律）与可验证奖励推理（RLVR）}
\label{sec:hacking-rlvr}

\subsection{奖励黑客（Reward Hacking）与 Goodhart 定律}

读到这里，不妨回头重读本章扉页上诺伯特·维纳写于 1960 年的警示。强化学习与对齐中最臭名昭著的幽灵，莫过于经济学中的\textbf{古德哈特定律（Goodhart's Law）}：
\begin{quote}
\kaishu “当一项代理度量变成优化目标时，它就不再是一项好的度量。”
\end{quote}
在 RLHF 中，我们真正渴望最大化的是难以穷尽的\textbf{真实人类质量 $R_{\mathrm{gold}}(x, y)$}，而交给 PPO 或 Best-of-$N$ 去拼命优化的却是一个从有限偏好数据拟合出的神经网络\textbf{代理奖励模型 $R_{\mathrm{proxy}} = r_\psi(x, y)$}。代理模型必然存在分布外盲区与统计伪相关；当策略优化器拥有强大算力时，它会精准钻空子，找到那些让 $R_{\mathrm{proxy}}$ 飙到极高、但真实人类看来却糟糕透顶的\textbf{奖励黑客（Reward Hacking）}解：
\begin{itemize}
  \item \textbf{长度膨胀（Verbosity Bias）}：人类标注员在对比中往往潜意识偏爱更详尽的回答，RM 于是学到了“越长分越高”的捷径；经强 RL 优化后，哪怕问“$1+1$ 等于几”，模型也要洋洋洒洒写三百字分点综述。
  \item \textbf{谄媚迎合（Sycophancy）}：当用户在提问里塞入错误前提（例如“我认为地球是平的，对吗？”），模型发现附和用户比温和纠错更容易拿高 RM 分，从而放弃客观事实。
  \item \textbf{排版刷分（Format Gaming）}：滥用加粗标题、项目符号与礼貌套话（“非常感谢您的精彩提问！”），以形式华丽掩盖实质空洞。
\end{itemize}

Gao 等（2023）在 OpenAI 的一项里程碑研究中，用一个合成“金标准大模型奖励（Gold RM）”系统测量了代理奖励模型在优化过程中随 KL 距离 $d \triangleq \sqrt{\DKL(\pi \,\|\, \pi_{\mathrm{ref}})}$ 变化的经验标度律：
\begin{equation}
  R_{\mathrm{gold}}(d) \approx d\,\big(\alpha_{\mathrm{RM}} - \beta_{\mathrm{RM}}\, d\big),
  \qquad d = \sqrt{\DKL(\pi \,\|\, \pi_{\mathrm{ref}})}.
  \label{eq:overopt-scaling}
\end{equation}
请观察式~\eqref{eq:overopt-scaling} 这一倒抛物线规律：初期随着策略离开参考模型（$d$ 增大），代理分数与真实质量同步上升；然而越过最优拐点 $d^* = \frac{\alpha_{\mathrm{RM}}}{2\beta_{\mathrm{RM}}}$ 之后，\textbf{代理分数继续一路狂飙，真实人类质量却掉头俯冲直至彻底崩坏！} 这也从根本上解释了为何式~\eqref{eq:rlhf-obj} 与 DPO 中的 KL 惩罚系数 $\beta$ 绝非可有可无的装饰品，而是遏制奖励过优化的生命防线。

\subsection{从主观偏好走向可验证奖励（RLVR）与推理时计算缩放}

如何从根源上避开“代理奖励模型可被黑客攻破”的困境？2024--2025 年，以 OpenAI o1 与 DeepSeek-R1 为代表的推理模型掀起了后训练的第二次范式革命——\textbf{基于可验证奖励的强化学习（Reinforcement Learning with Verifiable Rewards, RLVR）}：
\begin{itemize}
  \item \textbf{彻底抛弃神经网络 RM}：在数学竞赛、编程算法、形式化定理证明等客观领域，不再训练容易被刷分的代理奖励模型，而是直接用\textbf{确定性规则验证器}打分——例如用符号计算引擎核验最终数学表达式是否等价、在隔离沙箱里跑通全部单元测试！答对记 $R = 1$，答错记 $R = 0$（可辅以极轻量的格式规范奖励）。
  \item \textbf{以 GRPO 激发长思维链（Long CoT）涌现}：结合第~\ref{ch:trpo-ppo-grpo} 章推导的 GRPO 算法，对每道难题采样 $G$ 条完整思考轨迹，仅凭稀疏终端奖励 $\{0, 1\}$ 驱动策略进化。令人惊叹的是，模型不需要任何人类编写的“反思示范”，就会在 RL 探索中自发涌现出\textbf{长链分解、边界自检、发现矛盾后回溯重算（“顿悟时刻” / Aha Moment）}等高级认知行为！
\end{itemize}

与之紧密呼应的是新一代扩展法则——\textbf{推理时计算缩放（Test-Time Compute Scaling）}。过去我们主要在\emph{训练时}堆参数、堆数据；如今在\emph{推理时}为难题分配更多算力，同样能带来稳定且陡峭的性能跃升。推理时缩放主要包含两条路径：
\begin{enumerate}
  \item \textbf{并行缩放（Best-of-$N$ 采样与验证器筛选）}：对同一问题独立采样 $N$ 条回答 $\{y_1, \dots, y_N\} \sim \pi_{\mathrm{ref}}(\cdot \mid x)$，用\textbf{结果奖励模型（ORM）}、\textbf{过程奖励模型（PRM，对每个推理步骤打分）}或规则验证器选出得分最高者 $\hat{y} = \argmax_{i \le N} r(x, y_i)$。令人赞叹的是，Best-of-$N$ 采样策略 $\pi_{\mathrm{BoN}}$ 与参考策略之间的 KL 散度存在一个与奖励分布完全无关的精确解析公式（见习题~\textbf{6.4.4}）：
  \begin{equation}
    \DKL\!\big(\pi_{\mathrm{BoN}} \,\|\, \pi_{\mathrm{ref}}\big) = \ln N - \frac{N - 1}{N} \;\approx\; \ln N - 1 \quad (\text{当 } N \text{ 较大时}).
    \label{eq:bon-kl}
  \end{equation}
  这表明 Best-of-$N$ 仅耗费 $O(\ln N)$ 的极低 KL 预算，就能极其稳健地逼近高奖励区域，往往比粗调的 RL 拥有更优的“奖励—KL”帕累托前沿！
  \item \textbf{串行缩放（延长思维链测试时思考预算）}：通过 RLVR 训练出的推理模型，在遇到简单算术题时几十个词元即作答，遇到竞赛级难题时则自动展开数千乃至上万词元的串行推演与验算，将\textbf{搜索算法内化到了自回归生成过程之中}。
\end{enumerate}

\section{实测验证：DPO 对齐中的奖励—KL 帕累托权衡}
\label{sec:exp-dpo}

为了将定理~\ref{thm:rlhf-optimal} 与定理~\ref{thm:dpo-derivation} 从纸面推演变成可触摸的定量实验，配套脚本 \texttt{code/vol6\_experiments.py} 中的 \texttt{exp4\_dpo\_alignment()} 构建了一个完全可控、可精确观测真实隐藏奖励与精确 KL 散度的上下文偏好对齐实验：
\begin{itemize}
  \item \textbf{真实世界与参考模型设定}：构造 $N_{\mathrm{ctx}} = 64$ 个 12 维上下文向量 $X \in \R^{64 \times 12}$，每个上下文有 $K_{\mathrm{resp}} = 8$ 种候选回复。隐藏的真实人类奖励矩阵为 $R_{\mathrm{true}} = X W_{\mathrm{true}}$（对 DPO 算法完全不可见！）。初始参考策略 $\pi_{\mathrm{ref}}$ 是一个带噪声的不完美 SFT 模型（权重仅对齐真实奖励的 $25\%$ 并叠加高斯扰动），其初始期望真实奖励仅为 \textbf{$0.455$}。
  \item \textbf{偏好数据生成与 DPO 训练}：对每个上下文随机抽取 6 对候选回答（共 $64 \times 6 = 384$ 条偏好对），严格按 Bradley--Terry 概率~\eqref{eq:bt-model} 以隐藏真值 $R_{\mathrm{true}}$ 掷硬币决出胜负 $(y_w, y_l)$。随后从 $\pi_{\mathrm{ref}}$ 初始化 $\pi_\theta$，分别在\textbf{弱约束 $\beta = 0.1$} 与\textbf{强约束 $\beta = 0.6$} 下用纯 DPO 损失~\eqref{eq:dpo-loss} 训练 120 轮（Adam, $\text{lr}=3\times 10^{-2}$），每 10 轮记录真实期望奖励、精确 KL 散度 $\DKL(\pi_\theta \| \pi_{\mathrm{ref}})$ 与成对排序准确率。
\end{itemize}

\begin{figure}[htbp]
\centering
\begin{tikzpicture}
  % Left panel: Expected true hidden reward
  \begin{axis}[
    name=dpoleft,
    width=0.47\linewidth, height=6.0cm,
    title={\footnotesize\sffamily\bfseries (a) 隐藏真实人类奖励 $\E[R_{\mathrm{true}}]$ 爬升},
    xlabel={DPO 训练轮数（Epoch）},
    ylabel={期望真实奖励（对训练不可见）},
    xmin=-2, xmax=125, ymin=0.2, ymax=4.1,
    grid=both, grid style={tealLine!45, dashed},
    axis line style={slate},
    tick label style={font=\scriptsize},
    label style={font=\footnotesize\sffamily, color=ink},
    legend style={at={(0.97,0.06)}, anchor=south east, font=\scriptsize, fill=white, draw=tealLine}
  ]
    \addplot[ochre, very thick, dashed, mark=triangle*, mark size=1.6pt]
      table[x=ep, y=rew_b01, col sep=comma]{data/dpo_align.csv};
    \addlegendentry{$\beta = 0.1$（终值 $3.703$）}

    \addplot[teal, very thick, mark=*, mark size=1.5pt]
      table[x=ep, y=rew_b06, col sep=comma]{data/dpo_align.csv};
    \addlegendentry{$\beta = 0.6$（终值 $3.669$）}

    \node[slate, font=\scriptsize, anchor=west] at (axis cs:6, 0.55) {初始 SFT: $0.455$};
  \end{axis}

  % Right panel: KL divergence from pi_ref
  \begin{axis}[
    at={(dpoleft.east)}, anchor=west, xshift=12mm,
    width=0.47\linewidth, height=6.0cm,
    title={\footnotesize\sffamily\bfseries (b) 相对参考策略的 KL 散度 $\DKL(\pi_\theta \| \pi_{\mathrm{ref}})$},
    xlabel={DPO 训练轮数（Epoch）},
    ylabel={精确 KL 散度（nats）},
    xmin=-2, xmax=125, ymin=-0.045, ymax=3.05,
    grid=both, grid style={tealLine!45, dashed},
    axis line style={slate},
    tick label style={font=\scriptsize},
    label style={font=\footnotesize\sffamily, color=ink},
    legend style={at={(0.04,0.95)}, anchor=north west, font=\scriptsize, fill=white, draw=tealLine}
  ]
    \addplot[ochre, very thick, dashed, mark=triangle*, mark size=1.6pt]
      table[x=ep, y=kl_b01, col sep=comma]{data/dpo_align.csv};
    \addlegendentry{$\beta = 0.1$（终值 $2.200$）}

    \addplot[teal, very thick, mark=*, mark size=1.5pt]
      table[x=ep, y=kl_b06, col sep=comma]{data/dpo_align.csv};
    \addlegendentry{$\beta = 0.6$（终值 $1.746$）}

    % Callout showing 20.6% KL savings, placed cleanly below teal curve
    \draw[{Stealth[length=1.6mm]}-{Stealth[length=1.6mm]}, navy, thick]
      (axis cs:120, 1.77) -- (axis cs:120, 2.17);
    \node[navy, font=\scriptsize\sffamily, align=center, fill=white, draw=navy!40, rounded corners=1pt, inner sep=2pt]
      at (axis cs:86, 0.72) {KL 降低 $20.6\%$\\奖励仅差 $0.9\%$};
    \draw[-{Stealth[length=1.6mm]}, navy, thin] (axis cs:104, 1.00) -- (axis cs:118.5, 1.94);
  \end{axis}
\end{tikzpicture}
\caption{上下文成对偏好数据集上的 DPO 实测轨迹（读取自 \texttt{data/dpo\_align.csv}）。从未见过真实奖励标量的 DPO 仅靠 384 对胜负比较，便把期望真实奖励从 $0.455$ 拉升逾 $8\times$；而较大温度 $\beta=0.6$ 以不到 $1\%$ 的微弱奖励让步，换取了 $20.6\%$ 的 KL 偏离压缩与更高的偏好排序准确率（$93.2\%$ vs $88.8\%$）。}
\label{fig:dpo-align}
\end{figure}

表~\ref{tab:dpo-metrics} 与图~\ref{fig:dpo-align}（直接读取 \texttt{data/dpo\_align.csv}）呈现了极具启发性的定量规律：
\begin{enumerate}
  \item \textbf{纯成对比较即可精准还原隐藏奖励方向}：尽管 DPO 全程从未见过哪怕一个真实标量分 $R_{\mathrm{true}}$，也从未训练显式奖励模型或运行在线 RL，120 轮内它将策略的真实期望回报从初始 SFT 的 \textbf{$0.455$} 一路抬升至 \textbf{$3.703$}（$\beta=0.1$）与 \textbf{$3.669$}（$\beta=0.6$），涨幅超过 \textbf{$8.0\times$}！
  \item \textbf{温度 $\beta$ 的帕累托杠杆效应}：将 $\beta$ 从 $0.1$ 提高至 $0.6$，最终真实奖励几乎毫无损失（$3.669$ vs $3.703$，仅低 $0.9\%$），但策略相对参考模型 $\pi_{\mathrm{ref}}$ 的 KL 散度从 \textbf{$2.200$} 显著收紧至 \textbf{$1.746$}（节省了 \textbf{$20.6\%$} 的 KL 偏移预算），同时偏好对排序准确率从 \textbf{$88.8\%$} 反而提升至 \textbf{$93.2\%$}！这是因为适度增大的 $\beta$ 令逻辑斯蒂边际 $\beta(\Delta_w - \Delta_l)$ 对较小对数概率差更敏感，从而在不把策略分布推向极端尖峰的前提下高效完成了偏好排序。
\end{enumerate}

\begin{table}[htbp]
\centering
\caption{DPO 实测关键节点指标对照（数据源：\texttt{data/dpo\_align.csv} 及 \texttt{exp4\_dpo\_alignment()} 输出）}
\label{tab:dpo-metrics}
\footnotesize
\begin{tabularx}{\linewidth}{@{}lccccX@{}}
\toprule
\textbf{轮数} & \textbf{$\beta{=}0.1$ 奖励} & \textbf{$\beta{=}0.1$ KL} & \textbf{$\beta{=}0.6$ 奖励} & \textbf{$\beta{=}0.6$ KL} & \textbf{阶段特征} \\
\midrule
Ep.\,0   & $0.4554$ & $0.0000$ & $0.4554$ & $0.0000$ & 初始不完美 SFT 参考模型 \\
Ep.\,20  & $2.6287$ & $0.8295$ & $2.6257$ & $0.6510$ & 奖励极速爬升，$\beta{=}0.6$ KL 已低 $21.5\%$ \\
Ep.\,60  & $3.4392$ & $1.8131$ & $3.4405$ & $1.3293$ & 两者奖励持平，KL 剪刀差持续拉大 \\
Ep.\,120 & $3.7035$ & $2.1997$ & $3.6690$ & $1.7457$ & 准确率 $88.8\%$ ($\beta{=}0.1$) vs $93.2\%$ ($\beta{=}0.6$) \\
\bottomrule
\end{tabularx}
\end{table}

\begin{failBox}{DPO 与大模型后训练工程中的三大高频失误}
\begin{enumerate}
  \item \textbf{同时压低胜出回答概率的“似然位移”（Likelihood Displacement）}：细看式~\eqref{eq:dpo-loss}——DPO 损失只关心差值 $\ln\frac{\pi_\theta(y_w)}{\pi_{\mathrm{ref}}(y_w)} - \ln\frac{\pi_\theta(y_l)}{\pi_{\mathrm{ref}}(y_l)}$。实践中经常观测到\textbf{胜出回答 $y_w$ 与落败回答 $y_l$ 的绝对对数概率双双下降}（只是 $y_l$ 降得更快！），导致模型把大量概率质量泄漏到未见过的分布外词元上，推理能力退化。常见缓解手段包括加入轻量 SFT 正则项（如 CPO / RPO）、降低学习率（通常仅 $5\times 10^{-7} \sim 5\times 10^{-6}$），以及仅训练 $1 \sim 2$ 个 epoch。
  \item \textbf{跳过 SFT 分布对齐直接做离线 DPO}：定理~\ref{thm:dpo-derivation} 隐含假设偏好数据 $(y_w, y_l)$ 位于参考策略 $\pi_{\mathrm{ref}}$ 的高概率支撑集内。如果直接拿开源外部模型生成的偏好对去 DPO 微调一个风格迥异的基座模型（而不先用 $y_w$ 做一轮 SFT 使 $\pi_{\mathrm{ref}}$ 落入同一分布），$\ln\frac{\pi_\theta}{\pi_{\mathrm{ref}}}$ 的参考基准会严重失真，对齐效果大打折扣。
  \item \textbf{随意把整句求和 $\sum_t \ln \pi$ 改成词元均值 $\frac{1}{|y|}\sum_t \ln \pi$ 而不调整理论}：在精确 DPO 定理中，序列对数概率必须是\textbf{求和}而非平均；如果擅自取均值却仍沿用 DPO 的 $\pi_{\mathrm{ref}}$ 项，就会破坏与 KL 正则 RLHF 的等价性（尽管后来 \textbf{SimPO} 在去掉 $\pi_{\mathrm{ref}}$ 的同时专门引入了长度归一化与目标边际 $\gamma$ 以治理长度偏差——二者设计自成体系，切忌混搭）。
\end{enumerate}
\end{failBox}

\begin{extBox}{离线 DPO 的演进谱系：IPO、KTO、SimPO 与在线迭代 DPO}
DPO 因其极致简洁在工业界迅速普及，同时也催生出一系列针对特定痛点的变体与在线化演进：
\begin{itemize}
  \item \textbf{IPO（Identity Preference Optimization，Azar et al., 2024）}：指出在确定性偏好（几乎无噪声）数据下，Bradley--Terry 要求 $\hat{r}_w - \hat{r}_l \to +\infty$，容易令 DPO 过拟合甚至无视 $\beta$ 的 KL 约束；IPO 将对数 Sigmoid 换成带目标边际的平方回归损失 $\big((\ln\frac{\pi_\theta(y_w)}{\pi_{\mathrm{ref}}(y_w)} - \ln\frac{\pi_\theta(y_l)}{\pi_{\mathrm{ref}}(y_l)}) - \frac{1}{2\beta}\big)^2$，天然防止边际发散。
  \item \textbf{KTO（Kahneman--Tversky Optimization，Ethayarajh et al., 2024）}：借鉴行为经济学前景理论（损失厌恶），\textbf{不再需要成对数据 $(y_w, y_l)$}，仅需用户日常点赞/点踩的二元单条反馈即可直接对齐！
  \item \textbf{SimPO（Simple Preference Optimization，Meng et al., 2024）}：用长度归一化平均对数概率 $\frac{\beta}{|y|}\ln\pi_\theta(y\mid x)$ 作为隐式奖励并引入正边际 $\gamma > 0$，\textbf{连参考模型 $\pi_{\mathrm{ref}}$ 也一并省去}，显存再降一半且天然抑制长度膨胀。
  \item \textbf{迭代/在线 DPO（Iterative / Online DPO）}：离线 DPO 最大的软肋是无法针对当前策略 $\pi_\theta$ 自己新犯的错误做定向纠偏（Xu et al., 2024 证明在线 PPO 在分布外泛化上常优于纯离线 DPO）。迭代 DPO 每训完一轮，就用当前策略重新采样新回复、由 RM 或裁判模型标注新偏好对，再执行下一轮 DPO，优雅地架起了离线对齐与在线 RL 之间的桥梁。
\end{itemize}
\end{extBox}

\section*{本章小结与核心脉络}
\addcontentsline{toc}{section}{本章小结与核心脉络}

本章完整贯通了现代大语言模型从“预训练底座”进化为“对齐助手与深度推理者”的数学脊梁：
\begin{itemize}
  \item \textbf{SFT 指令微调}：通过提示掩码交叉熵建立对话范式与基础能力，但受限于模仿上限与缺乏负对比信号。
  \item \textbf{Bradley--Terry 与 KL 正则 RLHF}：借助成对比较建模隐含奖励，并用拉格朗日乘子法与 Gibbs 不等式严格证明了唯一闭式最优策略 $\pi^*(y\mid x) = \frac{1}{Z(x)}\pi_{\mathrm{ref}}(y\mid x)\exp(\frac{1}{\beta}r(x,y))$。
  \item \textbf{DPO 天书捷径}：反解隐式奖励并利用 Bradley--Terry 模型的提示级平移不变性，精确消去不可计算的配分函数 $Z(x)$，得到仅含自适应错判权重 $\sigma(\hat{r}_l - \hat{r}_w)$ 的离线分类损失；在上下文偏好实测中将隐藏真实回报从 $0.455$ 拉升至 $3.703$，并定量验证了 $\beta=0.6$ 节省 $20.6\%$ KL 散度的帕累托权衡。
  \item \textbf{奖励黑客与 RLVR}：以 Goodhart 定律与过优化标度律揭示代理奖励模型的崩坏机制，进而迈向以规则验证器 + GRPO + Best-of-$N$ 推理时计算缩放为核心的下一代推理范式。
\end{itemize}

至此，本卷上半部（第 1--4 章）的\textbf{强化学习与大模型后训练算法大厦}已全部竣工！然而，当你真正试图把 700 亿参数模型放进集群，用 GRPO 对每道题并发采样 $G=64$ 条 8000 词元的思维链、并完成反向传播时，一个无比残酷的工程现实会迎面撞来：\textbf{GPU 算力利用率竟然不到 $5\%$，显存瞬间 OOM 崩溃！} 究竟是谁卡住了现代 AI 的咽喉？从第 5 章起，我们将切换至底层计算机系统视角，直面 GPU 访存墙、Roofline 模型与 FlashAttention 的在线分块奇迹。

\section*{本章分层习题}
\addcontentsline{toc}{section}{本章分层习题}

\exTier{A 层：概念与公式手算}{熟练计算 Bradley--Terry 概率、DPO 损失与隐式奖励权重}

\begin{enumerate}
  \item[\textbf{6.4.1}] \textbf{DPO 损失、隐式奖励与错判权重手算}：给定某提示 $x$ 与一对回答 $(y_w, y_l)$，设温度 $\beta = 0.5$，测得参考策略与当前策略的整句对数概率如下：
  \[
    \ln \pi_{\mathrm{ref}}(y_w \mid x) = -12.0, \quad
    \ln \pi_{\mathrm{ref}}(y_l \mid x) = -10.0, \quad
    \ln \pi_\theta(y_w \mid x) = -11.0, \quad
    \ln \pi_\theta(y_l \mid x) = -8.0.
  \]
  \begin{enumerate}
    \item 计算胜出与落败回答的隐式奖励 $\hat{r}_\theta(x, y_w)$ 与 $\hat{r}_\theta(x, y_l)$；
    \item 计算当前模型预测 $y_w \succ y_l$ 的概率 $\sigma(\hat{r}_w - \hat{r}_l)$、该样本的 DPO 损失值 $-\ln\sigma(\hat{r}_w - \hat{r}_l)$，以及式~\eqref{eq:dpo-grad} 中的梯度错判权重 $w_\theta = \sigma(\hat{r}_l - \hat{r}_w)$；
    \item 判断当前模型是“判对了”还是“判反了”，并解释为何权重 $w_\theta > 0.5$。
  \end{enumerate}
  \exHint{$\hat{r}_w = 0.5 \times (-11 - (-12)) = +0.5$，$\hat{r}_l = 0.5 \times (-8 - (-10)) = +1.0$；由于 $\hat{r}_l > \hat{r}_w$，模型判反了，错判权重 $w_\theta = \sigma(0.5) \approx 0.622$。}

  \item[\textbf{6.4.2}] \textbf{RLHF 闭式最优策略的两动作数值验证}：某提示 $x$ 仅有两条候选回答 $\{y_1, y_2\}$，参考分布为 $\pi_{\mathrm{ref}}(y_1 \mid x) = 0.2, \pi_{\mathrm{ref}}(y_2 \mid x) = 0.8$，奖励为 $r(x, y_1) = 3.0, r(x, y_2) = 1.0$。
  \begin{enumerate}
    \item 分别在 $\beta = 1.0$ 与 $\beta = 0.2$ 下，用定理~\ref{thm:rlhf-optimal} 计算配分函数 $Z(x)$、闭式最优策略 $\pi^*(y_1\mid x)$、精确散度 $\DKL(\pi^* \| \pi_{\mathrm{ref}})$ 及最优目标值 $\beta \ln Z(x)$；
    \item 直接验算 $\E_{\pi^*}[r] - \beta \DKL(\pi^* \| \pi_{\mathrm{ref}})$ 是否精确等于 $\beta \ln Z(x)$。
  \end{enumerate}
  \exHint{当 $\beta = 1.0$ 时，$\pi^*(y_1) = \frac{0.2 e^3}{0.2 e^3 + 0.8 e^1} = \frac{e^2}{e^2 + 4} \approx 0.6488$；当 $\beta = 0.2$ 时，$\frac{1}{\beta}\Delta r = 10$，$\pi^*(y_1)$ 骤升至 $\approx 0.9998$，直观体现了低 $\beta$ 下的激进贪婪化。}
\end{enumerate}

\exTier{B 层：严谨数学证明}{吃透等价类唯一性、Best-of-$N$ 精确 KL 与温度极限}

\begin{enumerate}
  \item[\textbf{6.4.3}] \textbf{Bradley--Terry 奖励等价类与最优策略的唯一对应}：定义两个奖励函数 $r_1(x, y)$ 与 $r_2(x, y)$ 等价（记作 $r_1 \sim r_2$），若存在仅依赖 $x$ 的函数 $c(x)$ 使得 $r_2(x, y) = r_1(x, y) + c(x)$ 对一切 $(x, y)$ 成立。
  \begin{enumerate}
    \item 证明等价的奖励函数不仅诱导完全相同的 Bradley--Terry 偏好分布~\eqref{eq:bt-model}，而且在式~\eqref{eq:rlhf-obj} 中诱导\textbf{完全相同的闭式最优策略 $\pi^*(y \mid x)$}；
    \item 反之，证明在给定参考策略 $\pi_{\mathrm{ref}}$（处处为正）和 $\beta > 0$ 下，若任意合法策略 $\pi$ 均可写成 $\pi_{\mathrm{ref}}$ 的玻尔兹曼倾斜，则其所对应的奖励函数恰形成一个唯一的等价类，从而严格说明 DPO 重参数化\textbf{不损失任何表达能力}。
  \end{enumerate}
  \exHint{(a) 将 $r_2 = r_1 + c(x)$ 代入式~\eqref{eq:rlhf-opt-policy}，分子分母同时析出常数因子 $\exp(c(x)/\beta)$ 并约去；(b) 任选代表元 $r_\pi(x,y) = \beta \ln\frac{\pi(y\mid x)}{\pi_{\mathrm{ref}}(y\mid x)}$，其配分函数恰好恒等于 $1$！}

  \item[\textbf{6.4.4}] \textbf{推理时缩放的优美定理：Best-of-$N$ 的精确 KL 散度公式~\eqref{eq:bon-kl}}：设参考分布 $\pi_{\mathrm{ref}}$ 连续无原子（任意两条独立采样回答的奖励值几乎必然互不相等），每次独立采样 $N$ 条回答 $\{Y_1, \dots, Y_N\} \overset{\mathrm{i.i.d.}}{\sim} \pi_{\mathrm{ref}}$ 并输出奖励最高者 $Y^* = \argmax_i r(Y_i)$，记 $Y^*$ 的分布为 $\pi_{\mathrm{BoN}}$。
  \begin{enumerate}
    \item 记奖励 $R = r(Y)$（$Y \sim \pi_{\mathrm{ref}}$）的累积分布函数为 $F(r)$，令分位数 $U = F(r(Y)) \sim \mathrm{Uniform}(0, 1)$。证明 $Y^*$ 对应分位数 $U^* = \max(U_1, \dots, U_N)$ 的概率密度函数为 $p_N(u) = N u^{N-1}$（$u \in [0, 1]$）；
    \item 利用 KL 散度对单调双射变换的不变性，直接积分证明精确公式
    \[
      \DKL\!\big(\pi_{\mathrm{BoN}} \,\|\, \pi_{\mathrm{ref}}\big)
      = \int_0^1 N u^{N-1} \ln\!\big(N u^{N-1}\big)\,du
      = \ln N - \frac{N - 1}{N}.
    \]
  \end{enumerate}
  \exHint{拆分积分 $\ln N \int_0^1 N u^{N-1}du + N(N-1)\int_0^1 u^{N-1}\ln u\,du$，利用分部积分 $\int_0^1 u^{N-1}\ln u\,du = -\frac{1}{N^2}$ 一步即得！}
\end{enumerate}

\exTier{C 层：机制剖析与反例设计}{洞察奖励过优化拐点、确定性偏好过拟合与长度偏差}

\begin{enumerate}
  \item[\textbf{6.4.5}] \textbf{奖励过优化的数学定量诊断}：回顾式~\eqref{eq:overopt-scaling} 的经验过优化标度律 $R_{\mathrm{gold}}(d) = d(\alpha_{\mathrm{RM}} - \beta_{\mathrm{RM}} d)$，其中 $d = \sqrt{\DKL(\pi \,\|\, \pi_{\mathrm{ref}})}$。假设某 7B 代理奖励模型测得 $\alpha_{\mathrm{RM}} = 2.0, \beta_{\mathrm{RM}} = 0.25$。
  \begin{enumerate}
    \item 求出令真实人类质量达到峰值的最优 KL 散度 $\DKL^*$ 与最高真实得分 $R_{\mathrm{gold}}^{\max}$；
    \item 当 RLHF 训练缺乏 KL 约束导致 $\DKL = 64\;\text{nats}$ 时，计算此时的 $R_{\mathrm{gold}}$；结合习题~\textbf{6.4.4} 的式~\eqref{eq:bon-kl}，算出大约需要多大的 Best-of-$N$ 采样数 $N$ 才能恰好落在该最优 KL 散度 $\DKL^*$ 上？
  \end{enumerate}
  \exHint{(a) $d^* = \frac{2.0}{0.5} = 4.0$，即 $\DKL^* = 16\;\text{nats}$，$R_{\mathrm{gold}}^{\max} = 4.0$；(b) $\DKL = 64$ 时 $d=8$，$R_{\mathrm{gold}} = 0$（收益全盘归零！）；对于 Best-of-$N$，由 $\ln N - 1 \approx 16$ 可见需要 $N \approx e^{17}$，这也解释了为何纯 BoN 很难轻易过优化。}

  \item[\textbf{6.4.6}] \textbf{为什么确定性偏好会让 DPO 绕开 $\beta$ 发生过拟合？}：假设偏好数据集中对某个提示 $x$ 只有唯一一对无噪声样本 $(y_w, y_l)$（即经验概率 $\hat{P}(y_w \succ y_l \mid x) = 1$）。分析：无论正则强度 $\beta > 0$ 选得多大（哪怕 $\beta = 10$），只要不加早停，无约束最小化 DPO 损失 $-\ln \sigma\big(\beta(\ln\frac{\pi_\theta(y_w)}{\pi_{\mathrm{ref}}(y_w)} - \ln\frac{\pi_\theta(y_l)}{\pi_{\mathrm{ref}}(y_l)})\big)$ 最终会把 $\pi_\theta(y_l \mid x)$ 推向什么极限值？并解释 IPO 的二次边际损失如何根治这一缺陷。
  \exHint{$-\ln\sigma(z)$ 只有在 $z \to +\infty$ 时才趋于下确界 $0$；因此无论 $\beta$ 多大，梯度永不为零，最终必然将 $\pi_\theta(y_l \mid x) \to 0$（$\DKL \to \infty$），从而使 $\beta$ 的正则保护形同虚设！}
\end{enumerate}

\exTier{D 层：编程实验与系统延伸}{探索 $\beta$ 帕累托曲线与似然位移现象}

\begin{enumerate}
  \item[\textbf{6.4.7}] {\rightskip=0pt plus 1fil\relax \textbf{编程实验：绘制完整的 DPO 奖励—KL 帕累托前沿并监测似然位移}：修改 \texttt{code/vol6\_experiments.py} 中的 \texttt{exp4\_dpo\_alignment()}：\par}
  \begin{enumerate}
    \item 扫描 $\beta \in \{0.02, 0.05, 0.1, 0.3, 0.6, 1.5, 5.0\}$，以最终 \texttt{KL} 为横轴、最终真实奖励 \texttt{exp\_rew} 为纵轴绘制帕累托前沿曲线；
    \item 在训练过程中额外记录胜出样本的平均对数概率 $\bar{\ell}_w = \E[\ln \pi_\theta(y_w\mid x)]$ 与落败样本平均对数概率 $\bar{\ell}_l = \E[\ln \pi_\theta(y_l\mid x)]$，尝试在低 $\beta$ 下观察是否出现二者同步下降或过拟合迹象，并验证在 DPO 损失中加入 $0.1 \cdot (-\bar{\ell}_w)$（轻量 SFT 正则）后的变化。
  \end{enumerate}
  \exHint{当 $\beta = 5.0$ 时，KL 极小但真实奖励明显受限；当 $\beta = 0.02$ 时，KL 飙升但对噪声偏好对过拟合的风险加大，$\beta \in [0.1, 0.6]$ 位于帕累托拐角附近。}

  \item[\textbf{6.4.8}] \textbf{编程实验：在模拟偏好数据上实现 Best-of-$N$ 并验证公式~\eqref{eq:bon-kl}}：复用 \texttt{exp4\_dpo\_alignment()} 的随机环境（或将候选回复数扩充至 $K_{\mathrm{resp}} = 500$ 近似连续分布），对均匀参考策略 $\pi_{\mathrm{ref}}$ 实现 Best-of-$N$ 采样（$N \in \{2, 4, 8, 16, 32\}$），在 $10^4$ 次蒙特卡洛抽样下统计实测平均真实奖励与实测散度 $\DKL(\hat{\pi}_{\mathrm{BoN}} \| \pi_{\mathrm{ref}})$，并与精确理论值 $\ln N - \frac{N-1}{N}$ 逐项比对！
  \exHint{例如当 $N = 8$ 时，理论值 $\ln 8 - 7/8 \approx 2.079 - 0.875 = 1.204\;\text{nats}$；当 $K_{\mathrm{resp}} \gg N$（碰撞概率可忽略）时，实测 KL 会与理论公式吻合至小数点后两位。}
\end{enumerate}

\chapter[AI 硬件、Roofline 模型与 FlashAttention]{AI 计算硬件、GPU 内存墙与算子优化（Roofline 与 FlashAttention）}
\label{ch:hardware-roofline-flash}

\chapEpigraph{过早优化是万恶之源，然而我们也不该放过关键 3\% 代码中的优化机会。}{Donald E. Knuth（高德纳，1974）}

\begin{chapIntro}{本章导读：当算法复杂度撞上硅片上的“内存墙”}
在前面四章中，我们从时序差分、策略梯度一路走到 PPO、GRPO 与 DPO，在数学纸面上自由地写下矩阵乘法、Softmax 归一化与整句对数概率梯度。然而，当你真正试图在一台拥有 80\,GB 显存的 GPU 服务器上训练或部署一个百亿参数的大语言模型时，一个残酷的物理现实会立刻迎面撞来：\textbf{在现代 AI 计算中，算术运算本身几乎是“免费的”，真正昂贵的是把数据从显存搬运到计算核心的字节数。}

如果只用教科书上的大 $O$ 算术复杂度去衡量模型，你会陷入无数反直觉的工程迷局：为什么单请求逐词生成（Decode）时，每秒可算近千万亿次的旗舰 GPU 利用率却不足 $0.4\%$？为什么一个 70 亿参数的半精度模型本身仅占 14\,GB 权重，训练时却连一张 80\,GB 显卡都装不下？为什么 FlashAttention 明明在反向传播中多做了重计算（总 FLOPs 反而略增），实际运行速度却比标准注意力快 $2\text{--}4\times$、显存更是从平方级骤降至线性级？

本章将把计算机组成原理、操作系统存储层次与现代 GPU 微架构无缝贯通，为读者建立四块坚实的 AI 系统基石：
\begin{enumerate}
  \item \textbf{现代 GPU 硬件与内存墙}：从 CPU 的低延迟设计切入 GPU 的高吞吐 SIMT 架构，定量对比片上 SRAM（$\sim 20\,\mathrm{MB}$、$\sim 19\,\mathrm{TB/s}$）与片外 HBM（$\sim 80\,\mathrm{GB}$、$\sim 3.35\,\mathrm{TB/s}$）之间的数量级鸿沟。
  \item \textbf{Roofline 性能模型与 Prefill/Decode 分化定理}：以算术强度 $\mathrm{AI} = \mathrm{FLOPs}/\mathrm{Bytes}$ 与拐点 $I^* = P_{\mathrm{peak}}/B_{\mathrm{HBM}}$ 为尺，严格证明为何 Prefill 处于计算受限区而单词元 Decode 深陷访存受限区，并基于 \texttt{data/roofline.csv} 完整复现实测曲线。
  \item \textbf{大模型训练显存账本（$16\Phi$ 法则与 $\sqrt{L}$ 检查点定理）}：逐字节核算混合精度 AdamW 训练的静态显存与动态激活显存，并利用均值不等式证明最优 $\sqrt{L}$ 激活重计算如何以约 $33\%$ 的额外算力换取从 $O(L)$ 到 $O(\sqrt{L})$ 的显存骤降。
  \item \textbf{FlashAttention 数学核心与在线 Softmax 分块证明}：用数学归纳法严格证明分块在线 Softmax 递推与全局标准注意力逐元素恒等，且仅需 $O(B_r B_c)$ 片上 SRAM 瓦片；最后用纯 NumPy 实验验证误差低至 $3.89\times 10^{-16}$ 的双精度机器零。
\end{enumerate}
\end{chapIntro}

\section{从计算机组成到现代 GPU：计算力与“内存墙”}
\label{sec:gpu-arch-memory-wall}

\begin{recallBox}{408 计算机组成与操作系统速忆：存储层次金字塔}
在经典计算机组成原理中，冯·诺依曼体系的存储系统遵循一条铁律：\textbf{速度越快的存储介质单位造价越高、容量越小}。从 CPU 内部的通用寄存器（$<1\,\mathrm{ns}$）、片上 L1/L2/L3 SRAM 高速缓存（数纳秒），到片外主存 DRAM（数十至上百纳秒）与固态硬盘，构成了经典的存储层次金字塔。程序之所以能高速运行，依赖于\textbf{时间局部性}（刚访问过的数据很快再次使用）与\textbf{空间局部性}（相邻地址的数据会被连续读取）。现代 GPU 并没有推翻这座金字塔，而是把它的带宽剪刀差放大到了极致。
\end{recallBox}

\subsection{延迟导向（CPU）与吞吐导向（GPU）的芯片面积博弈}

在同一块面积有限的硅片上，芯片架构师必须在\textbf{控制逻辑、大容量缓存}与\textbf{算术逻辑单元（ALU）}之间做取舍。这直接造就了 CPU 与 GPU 截然不同的设计哲学：
\begin{itemize}
  \item \textbf{CPU（低延迟导向，Latency-Oriented）}：面对操作系统调度、数据库索引与复杂分支程序，CPU 必须让\emph{单条线程}尽可能快地跑完。因此它把大片晶体管投入到深分支预测器、乱序执行（Out-of-Order Execution）窗口、寄存器重命名以及动辄数十兆字节的 L3 缓存上，留给纯算术核心的硅片面积反而占比较小。
  \item \textbf{GPU（高吞吐导向，Throughput-Oriented）}：深度学习的核心载荷是高度规则的稠密矩阵乘加与张量逐元素变换，几乎没有复杂数据依赖分支。GPU 大幅削减单个核心的控制复杂度，把绝大多数硅片面积铺满数以万计的精简算术单元与专用矩阵引擎（Tensor Core），追求的不是单个线程的最低延迟，而是\textbf{整张芯片每秒能吞吐多少万亿次浮点运算}。
\end{itemize}

现代 NVIDIA 数据中心 GPU（以 Hopper 架构 H100 SXM 为代表）的核心硬件层次由四个关键词串联：
\begin{enumerate}
  \item \textbf{流多处理器（Streaming Multiprocessor, SM）}：GPU 的基本并行执行岛屿。一张 H100 SXM 芯片集成高达 132 个 SM，每个 SM 内置指令调度器、大容量寄存器文件、片上共享内存（Shared Memory / L1 Cache，本质是高速片上 SRAM）以及标量 CUDA Core 与矩阵引擎 Tensor Core。
  \item \textbf{单指令多线程（SIMT）与线程束（Warp）}：GPU 并不逐个线程单独发射指令，而是将每 32 个相邻线程编为一个\textbf{线程束（Warp）}，在锁步（Lockstep）模式下共享同一条指令译码流。当某个 Warp 因等待片外显存返回数据而停顿时，SM 内的硬件 Warp 调度器可以在\textbf{零时钟周期开销}内立刻切换到另一个已就绪的 Warp 执行——这就是 GPU 用海量并发线程隐藏访存延迟的底层秘密（对比 408 操作系统中 CPU 进程上下文切换动辄微秒级保存/恢复栈帧的昂贵开销，GPU 把所有活跃 Warp 的上下文直接常驻在庞大的寄存器文件中！）。
  \item \textbf{张量核心（Tensor Core）}：传统标量 ALU 每周期仅完成 1 次乘加（FMA，即 2 FLOPs）；而 Tensor Core 是专为深度学习打造的硬件矩阵乘加（MMA）阵列，单条指令即可在极短周期内吞下 $16\times 8\times 16$ 乃至更大尺寸的小矩阵乘加 $D = A B + C$，将 BF16 稠密峰值算力推升至近 $1000\,\mathrm{TFLOP/s}$！
  \item \textbf{存储层级断崖（On-Chip SRAM vs Off-Chip HBM）}：如图~\ref{fig:gpu-memory-hierarchy} 所示，GPU 内部存在一道惊人的带宽与容量断崖：
  \begin{itemize}
    \item \textbf{片外高带宽显存（HBM3）}：容量大（$80\,\mathrm{GB}$），用来存放整座大模型的权重、梯度、优化器状态与 KV 缓存；尽管 HBM3 通过硅中介层与 3D 堆叠已将带宽推到惊人的 $3.35\,\mathrm{TB/s}$（远超普通服务器 DDR5 内存），但面对近 $1000\,\mathrm{TFLOP/s}$ 的胃口依然杯水车薪。
    \item \textbf{片上 SRAM（Shared Memory / 寄存器）}：分布在各个 SM 内部，总容量仅约 $20\text{--}50\,\mathrm{MB}$（每 SM 约 $256\,\mathrm{KB}$ 可配置共享内存 + 寄存器文件），比 HBM 小了足足三个数量级；但聚合片上读写带宽高达 $\sim 19\,\mathrm{TB/s}$（寄存器带宽更高达数百 $\mathrm{TB/s}$），比 HBM 快一个数量级以上！
  \end{itemize}
\end{enumerate}

\begin{figure}[htbp]
\centering
\resizebox{0.96\linewidth}{!}{%
\begin{tikzpicture}[
  font=\sffamily\footnotesize,
  box/.style={rectangle, rounded corners=1.4mm, draw=#1, thick, fill=#1!10, align=center, text=ink, inner sep=3.5pt},
  arr/.style={-{Stealth[length=2.1mm]}, thick, slate}
]
  % Left side: CPU vs GPU die layout
  \node[font=\sffamily\small\bfseries, color=navy] at (-4.3, 3.65) {硅片面积分配哲学对比};

  % CPU Die
  \draw[navy, thick, rounded corners=1.5mm, fill=navyLight] (-6.9, 1.55) rectangle (-1.7, 3.32);
  \node[anchor=west, font=\sffamily\scriptsize\bfseries, color=navy] at (-6.75, 3.07) {CPU（低延迟导向）};
  \fill[ochre!25, draw=ochre, rounded corners=0.8mm] (-6.7, 2.35) rectangle (-4.4, 2.82);
  \node[font=\scriptsize] at (-5.55, 2.58) {复杂控制 / 分支预测};
  \fill[teal!22, draw=teal, rounded corners=0.8mm] (-4.25, 2.35) rectangle (-1.9, 2.82);
  \node[font=\scriptsize] at (-3.07, 2.58) {重型 ALU 核心};
  \fill[slate!20, draw=slate, rounded corners=0.8mm] (-6.7, 1.72) rectangle (-1.9, 2.24);
  \node[font=\scriptsize] at (-4.3, 1.98) {超大容量片上 L2 / L3 缓存};

  % GPU Die
  \draw[tealDark, thick, rounded corners=1.5mm, fill=tealLight] (-6.9, -1.62) rectangle (-1.7, 1.32);
  \node[anchor=west, font=\sffamily\scriptsize\bfseries, color=tealDark] at (-6.75, 1.06) {GPU（高吞吐 SIMT 导向）};
  \foreach \r in {0,1,2,3} {
    \fill[ochre!25, draw=ochre, rounded corners=0.5mm] (-6.7, 0.42-\r*0.56) rectangle (-5.55, 0.80-\r*0.56);
    \node[font=\tiny] at (-6.125, 0.61-\r*0.56) {Warp 调度};
    \fill[slate!18, draw=slate, rounded corners=0.5mm] (-5.45, 0.42-\r*0.56) rectangle (-4.65, 0.80-\r*0.56);
    \node[font=\tiny] at (-5.05, 0.61-\r*0.56) {SRAM};
    \foreach \c in {0,1,2,3,4,5} {
      \fill[teal!35, draw=tealDark, rounded corners=0.4mm] (-4.53+\c*0.44, 0.42-\r*0.56) rectangle (-4.17+\c*0.44, 0.80-\r*0.56);
    }
  }
  \node[font=\tiny\bfseries, color=tealDark] at (-4.30, -1.44) {密集排布的 SM 标量核与 Tensor Core 矩阵阵列};

  % Right side: GPU Memory Hierarchy Pyramid (H100 SXM)
  \node[font=\sffamily\small\bfseries, color=tealDark] at (3.55, 3.65) {旗舰 GPU 存储层次与带宽断崖（H100 SXM 量级）};

  % Level 1: Tensor Cores + Registers
  \node[box=gold, minimum width=44mm, minimum height=10.5mm] (reg) at (3.55, 2.75) {
    \textbf{Tensor Core \& 寄存器文件}\\
    {\scriptsize BF16 稠密峰值 \textbf{989 TFLOP/s}}
  };

  % Level 2: On-chip SRAM
  \node[box=teal, minimum width=62mm, minimum height=11.5mm] (sram) at (3.55, 1.25) {
    \textbf{片上 SRAM（Shared Memory / L1）}\\
    {\scriptsize 容量 $\sim 20\,\mathrm{MB}$ \quad $\cdot$ \quad 聚合带宽 \textbf{$\sim 19\,\mathrm{TB/s}$}}
  };

  % Level 3: L2 Cache
  \node[box=navy, minimum width=76mm, minimum height=10.0mm] (l2) at (3.55, -0.15) {
    \textbf{共享片上 L2 缓存}\\
    {\scriptsize 容量 $50\,\mathrm{MB}$ \quad $\cdot$ \quad 带宽 $\sim 12\,\mathrm{TB/s}$}
  };

  % Level 4: Off-chip HBM3
  \node[box=ochre, minimum width=90mm, minimum height=12.0mm] (hbm) at (3.55, -1.65) {
    \textbf{片外高带宽显存（HBM3）—— 内存墙主战场}\\
    {\scriptsize 容量 \textbf{80\,GB} \quad $\cdot$ \quad 带宽 \textbf{3.35\,TB/s}（比 SRAM 慢约 $6\times$）}
  };

  \draw[arr, <->, tealDark] (reg.south) -- (sram.north);
  \draw[arr, <->, navy] (sram.south) -- (l2.north);
  \draw[arr, <->, ochre] (l2.south) -- (hbm.north);

  % Side annotations
  \draw[-{Stealth[length=2mm]}, thick, teal] (-1.05, -1.7) -- (-1.05, 2.9)
    node[midway, sloped, above, font=\scriptsize\bfseries, color=tealDark] {带宽骤升、延迟骤降};
  \draw[{Stealth[length=2mm]}-, thick, ochre] (8.35, -1.7) -- (8.35, 2.9)
    node[midway, sloped, below, font=\scriptsize\bfseries, color=ochre] {容量暴涨 $4000\times$};
\end{tikzpicture}%
}
\caption{左侧对比 CPU（重控制与大缓存）与 GPU（重密集矩阵算力与零开销 Warp 切换）的硅片面积分配哲学；右侧刻画 H100 SXM 级 GPU 从片外 HBM3 到片上 SRAM 与 Tensor Core 的存储金字塔，揭示 AI 系统优化的核心矛盾——片上 SRAM 虽快却极小，片外 HBM 虽大却受带宽制约。}
\label{fig:gpu-memory-hierarchy}
\end{figure}

\subsection{什么是“内存墙”？为什么 FLOPs 不再等于运行时间？}

过去十年间，受限于芯片封装引脚密度与物理能耗，GPU 的浮点算力增速远远甩开了显存带宽增速：从 2016 年的 P100 到 2023 年的 H100，半精度矩阵算力暴涨了近 $50\times$，而 HBM 显存带宽仅增长了约 $4.6\times$。这意味着现代 GPU 就像一位每秒能切 1000 盘菜的超级主厨（Tensor Core），却配了一条每秒只能端来 3 盘原料的狭窄传送带（HBM 总线）。如果菜谱设计不当，主厨有 $99\%$ 的时间都只能呆站在灶台前等原料！

这正是系统结构领域著名的\textbf{内存墙（Memory Wall）}现象。要科学诊断一个神经网络算子究竟是“主厨切不过来”（计算受限），还是“传送带太慢”（访存受限），我们需要一把精确的定量标尺——\textbf{Roofline 性能模型}。

\section{Roofline 性能模型：Prefill 计算受限与 Decode 访存受限}
\label{sec:roofline-model}

\subsection{算术强度、硬件上限与拐点（Ridge Point）}

2009 年，加州大学伯克利分校的 Samuel Williams、Andrew Waterman 与 David Patterson 提出了简洁而极具威力的 \textbf{Roofline（屋顶线）性能模型}，用仅仅两个硬件参数与一个算法指标，就能给出任意计算核函数在特定硬件上的理论性能上界。

\begin{defBox}
\begin{definition}[算术强度、Roofline 上界与硬件拐点]
\label{def:roofline-ai}
考虑一个在峰值计算吞吐率为 $P_{\mathrm{peak}}$（单位：$\mathrm{FLOP/s}$）、片外显存带宽为 $B_{\mathrm{HBM}}$（单位：$\mathrm{Byte/s}$）的加速器上运行的算子。设该算子共需执行 $F$ 次浮点运算，且必须在片外显存与计算单元之间搬运至少 $Q$ 字节数据。
\begin{enumerate}
  \item \textbf{算术强度（Arithmetic Intensity, AI）}定义为平均每搬运 1 字节显存数据所完成的浮点运算次数：
  \begin{equation}
    \mathrm{AI} \;\triangleq\; \frac{F}{Q} \qquad \big(\text{单位：}\mathrm{FLOP/Byte}\big).
    \label{eq:ai-def}
  \end{equation}
  \item \textbf{Roofline 可达性能上界}由计算屋顶与访存斜坡二者中的瓶颈决定：
  \begin{equation}
    P(\mathrm{AI}) \;=\; \min\!\Big(P_{\mathrm{peak}},\; \mathrm{AI} \cdot B_{\mathrm{HBM}}\Big).
    \label{eq:roofline-perf}
  \end{equation}
  \item \textbf{硬件拐点（Ridge Point，亦称脊点）} $I^*$ 定义为访存斜坡与计算屋顶的交点：
  \begin{equation}
    I^* \;\triangleq\; \frac{P_{\mathrm{peak}}}{B_{\mathrm{HBM}}} \qquad \big(\text{单位：}\mathrm{FLOP/Byte}\big).
    \label{eq:ridge-point}
  \end{equation}
\end{enumerate}
当 $\mathrm{AI} < I^*$ 时，算子处于\textbf{访存受限区（Memory-Bandwidth Bound）}，运行时间几乎完全由字节搬运量 $Q / B_{\mathrm{HBM}}$ 决定，增加硬件算力无济于事；当 $\mathrm{AI} \ge I^*$ 时，算子进入\textbf{计算受限区（Compute Bound）}，计算单元能够吃饱跑满，运行时间由 $F / P_{\mathrm{peak}}$ 主导。
\end{definition}
\end{defBox}

让我们立即代入 H100 SXM（BF16 稠密张量核）的真实硬件参数算一笔账：
\[
  P_{\mathrm{peak}} = 989\,\mathrm{TFLOP/s}, \qquad
  B_{\mathrm{HBM}} = 3.35\,\mathrm{TB/s}
  \quad\Longrightarrow\quad
  I^* = \frac{989 \times 10^{12}}{3.35 \times 10^{12}} \approx 295.2\,\mathrm{FLOP/B}.
\]
这个数字 $295\,\mathrm{FLOP/B}$ 极其震撼：它意味着\textbf{你从 HBM 里每辛苦搬运 1 字节数据（哪怕是半个 BF16 浮点数），至少要在片上复用并完成近 300 次浮点运算，才配得上让 H100 的 Tensor Core 满负荷运转！}对于一个 BF16 元素（2 字节），则必须平均参与约 $590$ 次浮点运算。

\subsection{核心定理：为什么 Prefill 吃满算力，而 Decode 却惨遭腰斩？}

大语言模型的自回归推理清晰地分为两个物理截然不同的阶段：
\begin{itemize}
  \item \textbf{预填充阶段（Prefill）}：将用户输入的整段提示词（长度为 $S_{\mathrm{prompt}}$，例如 512 或 4096 个词元）一次性并行送入 Transformer，计算全部提示词的隐藏状态并填满 KV 缓存；
  \item \textbf{逐词解码阶段（Decode）}：自回归循环每次仅将上一步刚生成的\textbf{单个新词元}（若批量大小为 $B$，则总共只有 $T = B$ 个新词元）送入模型，逐层做线性投影并与历史 KV 缓存交互。
\end{itemize}
下面这条定理以零含糊的精确代数，揭示了两个阶段为何在硬件上呈现天壤之别。

\begin{thmBox}
\begin{theorem}[Transformer 线性投影的 Prefill 计算受限与 Decode 访存受限分化]
\label{thm:prefill-vs-decode}
考虑大模型中一个典型的 BF16（每元素占 2 字节）全连接线性层 $Y = X W$，其中权重矩阵 $W \in \R^{d \times d}$，当前步同时处理的词元隐状态矩阵 $X \in \R^{T \times d}$（单请求 Prefill 时 $T = S_{\mathrm{prompt}}$，批大小为 $B$ 的 Decode 时 $T = B$）。假设忽略偏置项与缓存溢出，从 HBM 读入 $X, W$ 并将结果 $Y \in \R^{T \times d}$ 写回 HBM，则：
\begin{enumerate}
  \item 该线性层的精确浮点运算量与 HBM 总字节读写量分别为
  \begin{equation}
    F(T, d) = 2 T d^2, \qquad
    Q(T, d) = 2\!\left(d^2 + 2 T d\right)\,\text{Bytes}.
    \label{eq:gemm-flops-bytes}
  \end{equation}
  \item 其精确算术强度为关于有效词元数 $T$ 的严格单调递增函数：
  \begin{equation}
    \mathrm{AI}(T) \;=\; \frac{T d}{d + 2 T} \;=\; \frac{1}{\frac{1}{T} + \frac{2}{d}} \qquad \big(\mathrm{FLOP/B}\big).
    \label{eq:gemm-ai}
  \end{equation}
  \item 算子跨过硬件拐点 $I^*$ 进入计算受限区，当且仅当 $d > 2 I^*$ 且词元数满足
  \begin{equation}
    T \;\ge\; T^* \;=\; \frac{d\, I^*}{d - 2 I^*}.
    \label{eq:critical-tokens}
  \end{equation}
\end{enumerate}
特别地，在典型宽度 $d = 4096$ 与 H100 规格（$P_{\mathrm{peak}} = 989\,\mathrm{TFLOP/s}$，$B_{\mathrm{HBM}} = 3.35\,\mathrm{TB/s}$，$I^* \approx 295.2\,\mathrm{FLOP/B}$）下：
\begin{itemize}
  \item 当 $T = 1$（单请求 Decode，矩阵—向量乘 GEMV）时，$\mathrm{AI}(1) = \frac{4096}{4098} \approx 1.00\,\mathrm{FLOP/B}$，可达算力仅为 $3.35\,\mathrm{TFLOP/s}$，GPU 计算利用率低至 $\mathbf{0.34\%}$（严重访存受限）；
  \item 跨入计算受限区需要 $T^* \approx 345$ 个词元；当 $T = 512$（Prefill 或大批量 Decode，矩阵—矩阵乘 GEMM）时，$\mathrm{AI}(512) = 409.60\,\mathrm{FLOP/B} > I^*$，可达算力打满 $\mathbf{989.00\,\mathrm{TFLOP/s}}$，利用率达到 $\mathbf{100\%}$！
\end{itemize}
\end{theorem}
\end{thmBox}

\begin{proof}
（1）矩阵乘法 $Y_{i,j} = \sum_{k=1}^{d} X_{i,k} W_{k,j}$ 共有 $T \times d$ 个输出元素，每个输出元素需做 $d$ 次乘法和 $d - 1$ 次加法（在工程惯例中记为 $d$ 次乘加操作 FMA，每次 FMA 计 $2$ FLOPs），因此总浮点运算量为 $F(T, d) = 2 T d^2$。

在 HBM 字节搬运上，BF16 每个数值占 $2$ 字节：读入权重 $W$ 占 $2 d^2$ 字节，读入输入激活 $X$ 占 $2 T d$ 字节，写出输出激活 $Y$ 占 $2 T d$ 字节，故最小总访存字节数为 $Q(T, d) = 2(d^2 + 2 T d)$。

（2）将式~\eqref{eq:gemm-flops-bytes} 代入定义~\ref{def:roofline-ai}，分子分母同除以 $2d$ 立即得到：
\[
  \mathrm{AI}(T) = \frac{2 T d^2}{2(d^2 + 2 T d)} = \frac{T d}{d + 2 T} = \frac{1}{\frac{1}{T} + \frac{2}{d}}.
\]
显然随 $T$ 增大，$\frac{1}{T}$ 严格递减，故 $\mathrm{AI}(T)$ 严格递增，且当 $T \to \infty$ 时趋于上界 $d / 2$（此时访存由激活 $X, Y$ 主导而非权重 $W$）。

（3）令 $\mathrm{AI}(T) \ge I^*$，整理不等式 $T d \ge I^*(d + 2 T)$ 得 $T(d - 2 I^*) \ge d I^*$。因此当且仅当隐藏维度满足 $d > 2 I^*$ 且 $T \ge \frac{d I^*}{d - 2 I^*}$ 时，算子才能进入计算受限区。代入 $d = 4096, I^* = 989/3.35 \approx 295.224$：
\begin{itemize}
  \item $T = 1$：$\mathrm{AI}(1) = \frac{4096}{4098} \approx 0.9995 \approx 1.00\,\mathrm{FLOP/B}$，$P(1) = 0.9995 \times 3.35 \approx 3.35\,\mathrm{TFLOP/s}$，利用率 $3.35 / 989 \approx 0.34\%$；
  \item 临界词元数：$T^* = \frac{4096 \times 295.224}{4096 - 590.448} \approx 345.0$；
  \item $T = 512$：$\mathrm{AI}(512) = \frac{512 \times 4096}{4096 + 1024} = \frac{2048}{5} = 409.60\,\mathrm{FLOP/B} > 295.22$，从而 $P(512) = \min(989,\, 409.60 \times 3.35) = 989.00\,\mathrm{TFLOP/s}$，利用率 $100.00\%$。
\end{itemize}
证毕。
\end{proof}

\begin{insightBox}{一条公式统领整本 AI 推理优化教科书}
请反复凝视式~\eqref{eq:gemm-ai} 的结构：在单请求 Decode（$T=1$）时，权重矩阵 $W$ 足足有 $2d^2 = 32\,\mathrm{MB}$，你辛辛苦苦把它从 HBM 搬进片上，却只让它乘了区区 1 个向量 $X$（$8\,\mathrm{KB}$）就扔掉了！\textbf{整个现代大模型推理系统的核心优化，几乎都在对着式~\eqref{eq:gemm-ai} 的分子分母做文章}：
\begin{itemize}
  \item \textbf{把分子 $T$ 变大（提高权重复用率）}：动态连续批处理（Continuous Batching，将几十上百个用户的 Decode 请求拼成大批量 $T = B$）与投机解码（Speculative Decoding，单请求一次并行验证 $\gamma + 1$ 个草稿词元），将在第~\ref{ch:dist-quant-serving} 章详述；
  \item \textbf{把分母 $Q$ 变小（压缩搬运字节数）}：将权重从 BF16（2 字节）量化到 INT8（1 字节）甚至 INT4（0.5 字节），分母权重项立刻缩小 $2\text{--}4\times$，Decode 速度近乎线性翻倍！
\end{itemize}
\end{insightBox}

\subsection{实测验证：H100 级 BF16 线性层的 Roofline 曲线与利用率攀升}

配套脚本 \texttt{code/vol6\_experiments.py} 中的实验 5（\texttt{exp5\_flash\_online}）精确扫描了 $T \in \{1, 2, 4, \dots, 1024\}$ 下 $4096 \times 4096$ BF16 线性层在 H100 规格上的算术强度、可达算力与利用率，结果保存在 \texttt{data/roofline.csv} 中。表~\ref{tab:roofline-data} 与图~\ref{fig:roofline-exp} 直接展示了这一从 $0.34\%$ 到 $100\%$ 的悬殊跨越。

\begin{table}[htbp]
\centering
\caption{H100 规格下 $4096 \times 4096$ BF16 线性层的 Roofline 实测数据节选（来源：\texttt{data/roofline.csv}）。}
\label{tab:roofline-data}
\footnotesize
\setlength{\tabcolsep}{4.2pt}
\begin{tabular}{@{}rrrrcl@{}}
\toprule
\textbf{词元数 $T$} & \textbf{算术强度 (FLOP/B)} & \textbf{可达算力 (TFLOP/s)} & \textbf{利用率} & \textbf{性能区域} & \textbf{典型推理场景} \\
\midrule
$1$ & $1.00$ & $3.35$ & $0.34\%$ & 访存受限 & 单请求逐词 Decode（GEMV） \\
$8$ & $7.97$ & $26.70$ & $2.70\%$ & 访存受限 & 小批量 Decode / 投机解码 \\
$64$ & $62.06$ & $207.90$ & $21.02\%$ & 访存受限 & 中批量在线服务 Decode \\
$128$ & $120.47$ & $403.58$ & $40.81\%$ & 访存受限 & 短提示 Prefill / 大批量 Decode \\
$256$ & $227.56$ & $762.31$ & $77.08\%$ & 临近拐点 & 逼近 $T^* \approx 345$ 临界线 \\
$512$ & $409.60$ & $989.00$ & $100.00\%$ & 计算受限 & 常规 Prefill（完全吃满算力） \\
$1024$ & $682.67$ & $989.00$ & $100.00\%$ & 计算受限 & 长提示 Prefill / 预训练前向 \\
\bottomrule
\end{tabular}
\end{table}

\begin{figure}[htbp]
\centering
\begin{tikzpicture}
  % Left panel: Classic Log-Log Roofline
  \begin{axis}[
    name=leftplot,
    width=0.49\linewidth,
    height=6.2cm,
    xmode=log,
    ymode=log,
    xmin=0.7, xmax=1400,
    ymin=2, ymax=2800,
    xlabel={算术强度 $\mathrm{AI}$ (FLOP/Byte)},
    ylabel={可达性能 $P$ (TFLOP/s)},
    title={\sffamily\bfseries\small (a) H100 BF16 对数 Roofline 模型},
    grid=both,
    grid style={line width=0.25pt, draw=slate!20},
    major grid style={line width=0.4pt, draw=slate!35},
    tick label style={font=\scriptsize},
    label style={font=\footnotesize\sffamily},
    legend style={at={(0.03,0.96)}, anchor=north west, font=\scriptsize\sffamily, fill=white, fill opacity=0.88, draw=tealLine}
  ]
    % Theoretical Roofline envelope
    \addplot[navy, very thick, domain=0.7:295.22, samples=20] {3.35*x};
    \addlegendentry{访存斜坡 ($3.35\,\mathrm{TB/s}$)}
    \addplot[ochre, very thick, domain=295.22:1400, samples=10] {989};
    \addlegendentry{计算屋顶 ($989\,\mathrm{TFLOP/s}$)}
    % Ridge point dashed line (excluded from legend)
    \addplot[gold, dashed, thick, forget plot] coordinates {(295.22, 2) (295.22, 989)};
    \node[font=\tiny\sffamily\bfseries, color=gold, anchor=north east] at (axis cs:270, 12) {拐点 $I^* \approx 295$};
    % Real CSV points
    \addplot[teal, only marks, mark=*, mark size=2.1pt, mark options={fill=tealLight, draw=tealDark, thick}]
      table[col sep=comma, x=ai, y=achieved_tflops] {data/roofline.csv};
    \addlegendentry{实测点 (\texttt{roofline.csv})}
    % Annotations
    \node[font=\tiny\sffamily, color=ochre, anchor=west] at (axis cs:1.35, 3.4) {$T=1$ ($0.34\%$)};
    \node[font=\tiny\sffamily, color=tealDark, anchor=south] at (axis cs:520, 1250) {$T=512$ ($100\%$)};
  \end{axis}

  % Right panel: Utilization vs Token count
  \begin{axis}[
    at={($(leftplot.east)+(1.35cm,0)$)},
    anchor=west,
    width=0.46\linewidth,
    height=6.2cm,
    xmode=log,
    log basis x=2,
    xmin=0.8, xmax=1400,
    ymin=0, ymax=112,
    xlabel={批处理词元总数 $T$},
    ylabel={GPU 算力利用率 (\%)},
    title={\sffamily\bfseries\small (b) 利用率随词元数 $T$ 的相变曲线},
    grid=both,
    grid style={line width=0.25pt, draw=slate!20},
    major grid style={line width=0.4pt, draw=slate!35},
    tick label style={font=\scriptsize},
    label style={font=\footnotesize\sffamily},
    xtick={1,4,16,64,256,1024},
    xticklabels={1,4,16,64,256,1024}
  ]
    \addplot[teal, very thick, mark=square*, mark size=2pt, mark options={fill=tealLight, draw=tealDark}]
      table[col sep=comma, x=tokens, y=util_pct] {data/roofline.csv};
    \addplot[ochre, dashed, thick] coordinates {(345, 0) (345, 100)};
    \node[font=\tiny\sffamily\bfseries, color=ochre, anchor=east] at (axis cs:310, 86) {$T^* \approx 345$};
    \node[font=\tiny\sffamily, color=slate, anchor=south west] at (axis cs:2.5, 48) {访存墙区};
    \node[font=\tiny\sffamily, color=tealDark, anchor=north east] at (axis cs:1250, 94) {饱和区};
  \end{axis}
\end{tikzpicture}
\caption{基于 \texttt{data/roofline.csv} 实测数据绘制的 $4096\times 4096$ BF16 线性层性能画像。(a) 对数 Roofline 图清楚区分了斜率为 1 的 HBM 带宽斜坡与水平算力屋顶；(b) GPU 利用率在 $T \le 32$ 时长期徘徊在 $10\%$ 以下，直至 $T^* \approx 345$ 后方才真正触顶。}
\label{fig:roofline-exp}
\end{figure}

\section{大模型训练的精确显存账本：$16\Phi$ 法则与激活重计算}
\label{sec:memory-accounting}

弄清了“跑得快不快”（Roofline）之后，我们立刻面临第二个同样致命的系统问题：\textbf{“装得下装不下”}（显存容量限制）。训练一个参数量为 $\Phi$ 的大语言模型时，HBM 中的显存消耗严格划分为两大板块：\textbf{静态模型状态显存}（与批量大小和序列长度无关）与\textbf{动态激活显存}（随批量、层数和序列长度剧烈膨胀）。

\subsection{静态模型状态的精确账本：$16\Phi$ 法则}

现代大模型训练几乎清一色采用\textbf{混合精度训练（Mixed-Precision Training）}配合 \textbf{AdamW 优化器}。初学者常以为“BF16 权重每参数只占 2 字节，7B 模型训练岂不是只需 14\,GB 显存？”现实却要贵上整整 8 倍！原因在于反向传播与优化器更新需要保留四份独立的张量副本：
\begin{enumerate}
  \item \textbf{BF16 模型权重（$2\Phi$ 字节）}：用于高速 Tensor Core 前向与反向矩阵乘；
  \item \textbf{BF16 梯度张量（$2\Phi$ 字节）}：反向传播累加得到的各层参数梯度（部分框架为了高稳定性也会在 FP32 中累加梯度，占 $4\Phi$ 字节，此处按标准低显存 ZeRO 基准计为 $2\Phi$）；
  \item \textbf{FP32 主权重副本（Master Weights，$4\Phi$ 字节）}：为什么不能直接把梯度加到 BF16 权重上？BF16 仅有 7 位尾数（机器精度 $\epsilon_{\mathrm{BF16}} = 2^{-7} \approx 7.8\times 10^{-3}$）。当权重 $w \approx 1.0$ 而单步更新量 $\eta \Delta w \approx 10^{-4}$ 时，在 BF16 中计算 $1.0 + 10^{-4}$ 会因舍入误差直接吞噬更新量（结果依然等于 $1.0$！）。因此必须在 FP32（23 位尾数）中保存一份主权重，每步先在 FP32 中完成精密更新，再向下舍入回 BF16 权重供下一步前向使用；
  \item \textbf{FP32 AdamW 一阶矩 $m_t$ 与二阶矩 $v_t$（$4\Phi + 4\Phi = 8\Phi$ 字节）}：AdamW 必须为每一个参数维护指数移动平均梯度 $m_t$ 与平方梯度 $v_t$；因为二阶矩 $v_t$ 量级极小（常在 $10^{-10}$ 以下），必须在 FP32 中存储以免下溢归零。
\end{enumerate}

将以上四项相加，我们得到大模型系统工程中最著名的\textbf{静态显存 $16\Phi$ 法则}：
\begin{equation}
  M_{\mathrm{static}} \;=\; \underbrace{2\Phi}_{\text{BF16 权重}} \;+\; \underbrace{2\Phi}_{\text{BF16 梯度}} \;+\; \underbrace{4\Phi}_{\text{FP32 主权重}} \;+\; \underbrace{8\Phi}_{\text{FP32 AdamW } m,v} \;=\; \mathbf{16}\,\Phi\;\text{Bytes}.
  \label{eq:16phi-rule}
\end{equation}
表~\ref{tab:16phi-ledger} 直观列出了不同参数规模模型的静态显存账单。即便在\textbf{尚未计入任何输入数据、激活值与框架碎片}的前提下，仅仅把一个 7B 模型的静态训练状态放进显存就需要 $112\,\mathrm{GB}$——已经彻底撑爆单张 $80\,\mathrm{GB}$ 的 H100！而 70B 模型更需要超过 $1.1\,\mathrm{TB}$ 的纯静态显存，这正是第~\ref{ch:dist-quant-serving} 章引入 ZeRO/FSDP 切分与 LoRA 低秩微调的直接动因。

\begin{table}[htbp]
\centering
\caption{混合精度（BF16 + FP32 AdamW）训练的静态显存 $16\Phi$ 账本（以 $1\,\mathrm{GB} = 10^9\,\mathrm{Bytes}$ 计）。}
\label{tab:16phi-ledger}
\footnotesize
\setlength{\tabcolsep}{4.0pt}
\begin{tabular}{@{}lrrrrrc@{}}
\toprule
\textbf{模型规模 $\Phi$} & \textbf{BF16 权重} & \textbf{BF16 梯度} & \textbf{FP32 主权重} & \textbf{FP32 AdamW} & \textbf{静态总计 ($16\Phi$)} & \textbf{推理 ($2\Phi$)} \\
\midrule
$1\mathrm{B}\;(10^9)$ & $2\,\mathrm{GB}$ & $2\,\mathrm{GB}$ & $4\,\mathrm{GB}$ & $8\,\mathrm{GB}$ & $\mathbf{16\,\mathrm{GB}}$ & $2\,\mathrm{GB}$ \\
$7\mathrm{B}\;(7\times 10^9)$ & $14\,\mathrm{GB}$ & $14\,\mathrm{GB}$ & $28\,\mathrm{GB}$ & $56\,\mathrm{GB}$ & $\mathbf{112\,\mathrm{GB}}$ & $14\,\mathrm{GB}$ \\
$70\mathrm{B}\;(7\times 10^{10})$ & $140\,\mathrm{GB}$ & $140\,\mathrm{GB}$ & $280\,\mathrm{GB}$ & $560\,\mathrm{GB}$ & $\mathbf{1120\,\mathrm{GB}}$ & $140\,\mathrm{GB}$ \\
$405\mathrm{B}\;(4.05\times 10^{11})$ & $810\,\mathrm{GB}$ & $810\,\mathrm{GB}$ & $1620\,\mathrm{GB}$ & $3240\,\mathrm{GB}$ & $\mathbf{6480\,\mathrm{GB}}$ & $810\,\mathrm{GB}$ \\
\bottomrule
\end{tabular}
\end{table}

\subsection{动态激活显存与激活检查点（$\sqrt{L}$ 最优定理）}

相比固定不变的 $16\Phi$ 静态显存，\textbf{激活显存（Activation Memory）}是随着微批大小 $B$、序列长度 $S$、网络层数 $L$ 与隐藏维度 $d$ 急剧膨胀的“隐形巨兽”。在反向传播中，为了计算线性层 $Y = XW$ 对权重的梯度 $\nabla_W \mathcal{L} = X^\top (\nabla_Y \mathcal{L})$，前向传播时的中间输入张量 $X$ 必须一直驻留在 HBM 中，直到反向梯度传回该层后才能释放！

按照 Megatron-LM 论文对标准 Transformer 层（含多头自注意力、4 倍宽 MLP、LayerNorm 与 Dropout）的逐算子精确核算，若使用常规实现（显式物化 $h \times S \times S$ 注意力权重矩阵），单层 BF16 激活显存约为：
\begin{equation}
  M_{\mathrm{act,1}} \;\approx\; B S d \left(34 + \frac{5 h S}{d}\right)\,\text{Bytes},
  \label{eq:act-memory-layer}
\end{equation}
其中 $h$ 为注意力头数。请注意括号内的两项截然不同的增长律：
\begin{itemize}
  \item \textbf{第一项 $34 B S d$}：来自 QKV/Output 投影、MLP 升维/降维与归一化层的输入激活，随序列长度 $S$ 呈\textbf{线性增长}，且总激活显存随层数 $L$ 线性叠加为 $O(L B S d)$；
  \item \textbf{第二项 $5 B h S^2$}：来自标准注意力显式保存的 $S\times S$ 注意力分数、Softmax 概率与 Dropout 掩码，随序列长度 $S$ 呈\textbf{平方爆炸}！当上下文扩展到 $S = 32768$ 时，这一项本身就会吞噬数百 GB 显存（我们将在下一节用 FlashAttention 将其连根拔除）。
\end{itemize}

即便用 FlashAttention 消去了 $O(S^2)$ 项，线性项 $34 L B S d$ 在深层大模型（例如 $L = 80, d = 8192, B = 4, S = 4096$）下依然高达 $\sim 366\,\mathrm{GB}$！面对线性累积的 $O(L)$ 深度激活显存，\textbf{激活检查点（Activation Checkpointing，亦称梯度检查点 / 重计算）}提出了一个优雅的“以算力换显存”策略：前向传播时不再保存每一层的全部中间激活，而是每隔 $k$ 层仅保存一个边界输入（检查点）；反向传播退回到某一段时，再从该段起点的检查点出发，临时重算这一段 $k$ 层的局部前向激活，用完即刻释放。

\begin{thmBox}
\begin{theorem}[激活检查点的 $\sqrt{L}$ 最优分块权衡定理]
\label{thm:activation-checkpointing}
考虑一个包含 $L$ 个同构串行层的深度神经网络，单层前向激活显存为 $M_1$、单层前向耗时为 $t_F$、单层反向耗时为 $t_B \approx 2 t_F$。若将 $L$ 层均匀划分为 $L / k$ 个连续段（每段含 $k$ 层，$1 \le k \le L$），前向仅保存每段起点的边界激活（称为检查点），反向在每个段内按需重算局部前向激活，则：
\begin{enumerate}
  \item 整个训练步的峰值激活显存上界为
  \begin{equation}
    M_{\mathrm{peak}}(k) \;=\; \left(\frac{L}{k} + k\right) M_1.
    \label{eq:ckpt-mem}
  \end{equation}
  \item 当且仅当分段长度取 $k^* = \sqrt{L}$ 时，峰值激活显存取得全局最小值：
  \begin{equation}
    M_{\mathrm{peak}}^* \;=\; 2\sqrt{L}\, M_1,
    \label{eq:ckpt-opt}
  \end{equation}
  将激活显存复杂度从 $O(L)$ 压缩至 $O(\sqrt{L})$。
  \item 该策略在每个训练步中最多增加一次完整前向传播的计算开销，总步耗时的相对增量上界仅为
  \begin{equation}
    \frac{\Delta T}{T_0} \;=\; \frac{(L - L/k)\, t_F}{L(t_F + t_B)} \;\le\; \frac{t_F}{t_F + 2 t_F} \;=\; \frac{1}{3} \;\approx\; 33.3\%.
    \label{eq:ckpt-overhead}
  \end{equation}
\end{enumerate}
\end{theorem}
\end{thmBox}

\begin{proof}
（1）在前向传播结束并进入某一段（如最后一段）的反向传播时，HBM 中必须同时驻留两部分激活：第一部分是跨全网保存的 $L / k$ 个段首检查点，总计占 $(L / k) M_1$；第二部分是为了对该段执行反向传播而从段首临时重算出的该段全部 $k$ 层中间激活，占 $k M_1$。当该段反向完成后，这 $k$ 层临时激活即被立刻释放并复用于前一段。因此全过程的峰值激活显存满足 $M_{\mathrm{peak}}(k) = (\frac{L}{k} + k) M_1$。

（2）由算术—几何平均不等式（AM-GM）：
\[
  \frac{L}{k} + k \;\ge\; 2\sqrt{\frac{L}{k} \cdot k} \;=\; 2\sqrt{L},
\]
等号成立当且仅当 $\frac{L}{k} = k$，即 $k^* = \sqrt{L}$，此时最小峰值显存为 $2\sqrt{L}\,M_1$。

（3）原始无检查点训练一步需执行 $L$ 层前向与 $L$ 层反向，耗时 $T_0 = L(t_F + t_B)$。启用检查点后，每个非段首层在反向经过时仅需重算一次前向（共计 $L - L/k$ 层额外前向，严格小于 $L$ 层）。由于对线性层 $Y = XW$，前向做 1 次矩阵乘 $XW$，反向需分别计算对输入梯度 $(\nabla_Y \mathcal{L}) W^\top$ 与对权重梯度 $X^\top (\nabla_Y \mathcal{L})$ 共 2 次同尺寸矩阵乘，故精确有 $t_B \approx 2 t_F$。代入即得额外时间开销上界为 $\frac{L t_F}{3 L t_F} = \frac{1}{3} \approx 33.3\%$。证毕。
\end{proof}

例如对于一座 $L = 64$ 层的大模型，取 $k^* = 8$ 即可将激活显存压缩至原来的 $\frac{2\sqrt{64}}{64} = \frac{1}{4}$，从而允许你把微批大小 $B$ 扩大数倍；根据定理~\ref{thm:prefill-vs-decode}，更大的微批不仅摊薄了通信成本，甚至常常让实际端到端吞吐率\textbf{不降反升}！实践中，工业框架还会采用\textbf{选择性重计算（Selective Checkpointing）}：只丢弃算力开销极小却极占显存的算子（如 Dropout、GeLU、LayerNorm 与注意力分数），而保留计算昂贵的 GEMM 输出，从而以不到 $10\%$ 的算力开销砍掉 $70\%$ 以上的激活显存。

\section{FlashAttention 数学核心：在线 Softmax 分块与 IO 感知优化}
\label{sec:flash-attention}

现在我们迎来本章最具数学美感的高潮——Tri Dao 等人（2022）提出的 \textbf{FlashAttention}。它不是任何形式的稀疏近似或低秩妥协，而是\textbf{在数学上与标准全注意力逐比特等价的精确算法}，却一举把注意力层的显存占用从 $O(N^2)$ 压到 $O(N)$，并在长序列上收获了 $2\text{--}4\times$ 的真实加速。它的秘密既不在减少 FLOPs，也不在奇技淫巧，而在于一条简洁的数学恒等式：\textbf{在线 Softmax 分块递推（Online Softmax Tiling）}。

\subsection{标准注意力为何狠狠撞上 HBM 内存墙？}

给定单个注意力头的查询、键、值矩阵 $Q, K, V \in \R^{N \times d}$（序列长度 $N$，头维度 $d$，通常 $d \in \{64, 128\}$），标准缩放点积注意力按如下三步顺序执行：
\begin{equation}
  S = \frac{Q K^\top}{\sqrt{d}} \in \R^{N \times N}, \qquad
  P = \softmax(S) \in \R^{N \times N}, \qquad
  O = P V \in \R^{N \times d}.
  \label{eq:std-attn}
\end{equation}
在未融合的标准框架实现中，这三个算子分别对应三次独立的 GPU 核函数调用：
\begin{enumerate}
  \item \textbf{GEMM 核函数}：从 HBM 读入 $Q, K$（$2Nd$ 元素），算出巨大的 $S \in \R^{N\times N}$ 并\textbf{写回 HBM}（$N^2$ 元素）；
  \item \textbf{Softmax 核函数}：从 HBM 重新读出整个 $S$（$N^2$ 元素），逐行做指数归一化，再把 $P \in \R^{N\times N}$ \textbf{写回 HBM}（$N^2$ 元素）；
  \item \textbf{GEMM 核函数}：从 HBM 读出 $P$ 与 $V$（$N^2 + Nd$ 元素），算完输出 $O$ 后写回 HBM（$Nd$ 元素）。
\end{enumerate}
统计一下 HBM 总访存量：高达 $\Theta(Nd + N^2)$！由于现代长上下文模型中 $N \gg d$（例如 $N = 8192, d = 128$ 时，$N^2 = 6.7\times 10^7$ 远远大于 $Nd = 1.0\times 10^6$），绝大部分时间并没有花在矩阵乘法上，而是浪费在了\textbf{把庞大的中间矩阵 $S$ 和 $P$ 反复搬出、搬进缓慢的 HBM}！尤其对于中间的 Softmax 核函数，每个元素只做寥寥数次加减乘除指数运算，算术强度不足 $5\,\mathrm{FLOP/B}$，是一头彻头彻尾的访存受限怪兽。

既然如此，我们能否把“算 $S$、算 Softmax、乘 $V$”融合进\textbf{同一个核函数}里，让 $S$ 和 $P$ 的小瓦片永远只停留在高速的片上 SRAM 中，算完即毁，永远不落地到 HBM？

\subsection{从安全 Softmax 到在线 Softmax 分块：完整数学证明}

看似有一道天然的数学障碍挡在前面：\textbf{Softmax 的归一化分母} $\sum_{t=1}^{N} e^{s_t - m}$ \textbf{以及防溢出的行最大值} $m = \max_{1 \le t \le N} s_t$，\textbf{都依赖于一整行的全部 $N$ 个元素！}如果片上 SRAM 太小，根本装不下一整行（特别是 $N = 131072$ 时），我们在只看到前半段键块 $K_1$ 时算出的局部最大值 $m_1$ 和局部归一化分母 $\ell_1$，岂不是在读入后半段键块 $K_2$ 后全部作废了？

答案是：\textbf{非但不会作废，而且只需乘上一个极其简洁的指数重标定因子 $\alpha = e^{m_{\mathrm{old}} - m_{\mathrm{new}}}$，就能把历史统计量与历史输出累加器精确无损地修正到新的全局基准上！}下面我们以完整的数学归纳法给出这条天书级定理的严格证明。

\begin{thmBox}
\begin{theorem}[FlashAttention 在线 Softmax 分块递推的精确恒等性]
\label{thm:online-softmax-tiling}
设 $Q, K, V \in \R^{N \times d}$，尺度系数 $\tau = 1/\sqrt{d}$。选定块大小 $B_r, B_c$（使单个瓦片适配片上 SRAM 容量 $M_{\mathrm{SRAM}}$），将 $Q$ 按行划分为 $T_r = \lceil N / B_r \rceil$ 块 $\{Q_i\}_{i=1}^{T_r}$（$Q_i \in \R^{B_r \times d}$），将 $K, V$ 按行同步划分为 $T_c = \lceil N / B_c \rceil$ 块 $\{K_j, V_j\}_{j=1}^{T_c}$（$K_j, V_j \in \R^{B_c \times d}$）。

对任意固定查询块 $Q_i$，在片上 SRAM 中初始化行统计量与输出累加器：
\[
  m_i^{(0)} = (-\infty, \dots, -\infty)^\top \in \R^{B_r \times 1}, \qquad
  \ell_i^{(0)} = \mathbf{0} \in \R^{B_r \times 1}, \qquad
  O_i^{(0)} = \mathbf{0} \in \R^{B_r \times d}.
\]
沿键值块索引 $j = 1, 2, \dots, T_c$ 依次流式执行如下片上递推（所有列向量与矩阵的乘除法均按行广播，即等价于左乘对角阵）：
\begin{align}
  S_{ij} &= \tau\, Q_i K_j^\top \in \R^{B_r \times B_c}, \label{eq:flash-sij}\\
  m_{ij} &= \operatorname{rowmax}(S_{ij}) \in \R^{B_r \times 1}, \qquad
  m_i^{(j)} = \max\!\Big(m_i^{(j-1)},\; m_{ij}\Big), \label{eq:flash-m}\\
  \alpha_i^{(j)} &= \exp\!\Big(m_i^{(j-1)} - m_i^{(j)}\Big), \qquad
  \tilde{P}_{ij} = \exp\!\Big(S_{ij} - m_i^{(j)}\Big), \label{eq:flash-p}\\
  \ell_i^{(j)} &= \alpha_i^{(j)} \ell_i^{(j-1)} + \operatorname{rowsum}\!\big(\tilde{P}_{ij}\big), \label{eq:flash-l}\\
  O_i^{(j)} &= \frac{\Big(\alpha_i^{(j)} \ell_i^{(j-1)}\Big) O_i^{(j-1)} + \tilde{P}_{ij} V_j}{\ell_i^{(j)}}. \label{eq:flash-o}
\end{align}
则当 $j = T_c$ 结束时，所得输出块 $O_i^{(T_c)}$ 在实数域上与完整全局注意力精确恒等：
\begin{equation}
  O_i^{(T_c)} \;=\; \softmax\!\Big(\tau\, Q_i K^\top\Big)\, V.
  \label{eq:flash-exact}
\end{equation}
此外，整个算法从未在 HBM 中物化 $N \times N$ 矩阵，片上仅占用 $O(B_r d + B_c d + B_r B_c)$ SRAM；在令 $B_c = \Theta(M_{\mathrm{SRAM}} / d)$、$B_r = \min(\Theta(M_{\mathrm{SRAM}} / d), d)$ 时，HBM 总访存量从 $\Theta(Nd + N^2)$ 降至
\begin{equation}
  Q_{\mathrm{Flash}} \;=\; \Theta\!\left(\frac{N^2 d^2}{M_{\mathrm{SRAM}}}\right).
  \label{eq:flash-io}
\end{equation}
\end{theorem}
\end{thmBox}

\begin{proof}
由于递推式~\eqref{eq:flash-sij}--\eqref{eq:flash-o} 对查询块 $Q_i$ 的各行彼此独立，不妨固定其中任意一行 $q \in \R^{1 \times d}$，记全序列的注意力分数行向量为 $s = \tau q K^\top \in \R^{1 \times N}$，并按块切分为 $s = [s_1, s_2, \dots, s_{T_c}]$，其中 $s_j \in \R^{1 \times B_c}$。记前 $j$ 个块所覆盖的词元索引集合为 $\mathcal{I}_j = \{1, 2, \dots, j B_c\}$（最后一块截断至 $N$）。

我们对块索引 $j \in \{1, 2, \dots, T_c\}$ 使用\textbf{数学归纳法}，证明以下三条循环不变量同时成立：
\begin{align*}
  \text{(I)}\quad & m^{(j)} = \max_{t \in \mathcal{I}_j} s_t, \\
  \text{(II)}\quad & \ell^{(j)} = \sum_{t \in \mathcal{I}_j} \exp\!\big(s_t - m^{(j)}\big), \\
  \text{(III)}\quad & \ell^{(j)}\, o^{(j)} = \sum_{t \in \mathcal{I}_j} \exp\!\big(s_t - m^{(j)}\big)\, v_t,
\end{align*}
其中 $v_t \in \R^{1 \times d}$ 为矩阵 $V$ 的第 $t$ 行向量。

\textbf{奠基步（$j = 1$）}：由于 $m^{(0)} = -\infty$，有 $m^{(1)} = \max(-\infty, \max_{t \in \mathcal{I}_1} s_t) = \max_{t \in \mathcal{I}_1} s_t$，不变量 (I) 成立。同时 $\alpha^{(1)} = e^{-\infty - m^{(1)}} = 0$（约定 $0 \cdot 0 = 0$），代入式~\eqref{eq:flash-l} 与~\eqref{eq:flash-o} 得：
\[
  \ell^{(1)} = 0 + \sum_{t \in \mathcal{I}_1} e^{s_t - m^{(1)}}, \qquad
  \ell^{(1)} o^{(1)} = \mathbf{0} + \sum_{t \in \mathcal{I}_1} e^{s_t - m^{(1)}} v_t,
\]
故不变量 (II)、(III) 在 $j = 1$ 时精确成立。

\textbf{归纳步（假设 $j - 1$ 成立，证 $j$ 成立）}：记第 $j$ 块的索引集合为 $\mathcal{J}_j = \mathcal{I}_j \setminus \mathcal{I}_{j-1}$。
首先，由式~\eqref{eq:flash-m} 和归纳假设 (I)：
\[
  m^{(j)} = \max\!\Big(\max_{t \in \mathcal{I}_{j-1}} s_t,\; \max_{t \in \mathcal{J}_j} s_t\Big) = \max_{t \in \mathcal{I}_j} s_t,
\]
不变量 (I) 得以保持。

其次，观察指数乘法恒等式：对任意 $t \in \mathcal{I}_{j-1}$，都有
\[
  \alpha^{(j)} \cdot \exp\!\big(s_t - m^{(j-1)}\big)
  \;=\; \exp\!\big(m^{(j-1)} - m^{(j)}\big) \cdot \exp\!\big(s_t - m^{(j-1)}\big)
  \;=\; \exp\!\big(s_t - m^{(j)}\big)!
\]
将这一恒等式与归纳假设 (II) 代入式~\eqref{eq:flash-l}：
\[
  \ell^{(j)}
  = \alpha^{(j)} \sum_{t \in \mathcal{I}_{j-1}} e^{s_t - m^{(j-1)}} + \sum_{t \in \mathcal{J}_j} e^{s_t - m^{(j)}}
  = \sum_{t \in \mathcal{I}_{j-1}} e^{s_t - m^{(j)}} + \sum_{t \in \mathcal{J}_j} e^{s_t - m^{(j)}}
  = \sum_{t \in \mathcal{I}_j} e^{s_t - m^{(j)}},
\]
不变量 (II) 严格成立。

最后，将式~\eqref{eq:flash-o} 两边同乘以 $\ell^{(j)} > 0$，并代入归纳假设 (III) 与同一指数恒等式：
\begin{align*}
  \ell^{(j)} o^{(j)}
  &= \alpha^{(j)} \Big(\ell^{(j-1)} o^{(j-1)}\Big) + \sum_{t \in \mathcal{J}_j} e^{s_t - m^{(j)}} v_t \\
  &= \alpha^{(j)} \sum_{t \in \mathcal{I}_{j-1}} e^{s_t - m^{(j-1)}} v_t + \sum_{t \in \mathcal{J}_j} e^{s_t - m^{(j)}} v_t \\
  &= \sum_{t \in \mathcal{I}_{j-1}} e^{s_t - m^{(j)}} v_t + \sum_{t \in \mathcal{J}_j} e^{s_t - m^{(j)}} v_t
  \;=\; \sum_{t \in \mathcal{I}_j} e^{s_t - m^{(j)}} v_t,
\end{align*}
不变量 (III) 亦严格成立！

\textbf{终止结论与 IO 复杂度}：当 $j = T_c$ 时，$\mathcal{I}_{T_c} = \{1, \dots, N\}$，于是
\[
  o^{(T_c)} \;=\; \frac{\sum_{t=1}^{N} e^{s_t - m^{(T_c)}} v_t}{\sum_{t=1}^{N} e^{s_t - m^{(T_c)}}}
  \;=\; \sum_{t=1}^{N} \frac{e^{s_t}}{\sum_{u=1}^{N} e^{s_u}}\, v_t
  \;=\; \big[\softmax(s)\, V\big],
\]
与全局标准注意力逐元素完全恒等！并且由于全程恒有 $m^{(j-1)} - m^{(j)} \le 0$ 且 $S_{ij} - m^{(j)} \le 0$，所有 $\exp(\cdot)$ 的自变量均 $\le 0$，彻底杜绝了浮点指数上溢（Overflow）！

最后核算 HBM 访存量（原论文初版以外循环遍历 $K_j, V_j$、内循环遍历 $Q_i$；而 FlashAttention-2 改为外循环遍历 $Q_i$、内循环遍历 $K_j, V_j$，正是式~\eqref{eq:flash-sij}--\eqref{eq:flash-o} 的更优顺序！）：对每个外层查询块 $Q_i$（共 $T_r = N / B_r$ 块），仅从 HBM 读入 $Q_i$ 一次、最终写出 $O_i$ 一次（总计 $2Nd$）；而在内循环中需扫描全部 $K, V$ 一次（每次扫过 $2Nd$）。因此总 HBM 访存元素数为
\[
  Q_{\mathrm{Flash}} = 2 N d + T_r \cdot 2 N d = 2 N d + \frac{2 N^2 d}{B_r}.
\]
片上 SRAM 需同时容纳 $Q_i$（$B_r d$）、$K_j, V_j$（$2 B_c d$）、$O_i$（$B_r d$）与临时分数块 $S_{ij}$（$B_r B_c$）。取 $B_r, B_c = \Theta(M_{\mathrm{SRAM}} / d)$（同时受限于 $B_r B_c \le O(M_{\mathrm{SRAM}})$），代入即得 $Q_{\mathrm{Flash}} = \Theta\!\left(\frac{N^2 d^2}{M_{\mathrm{SRAM}}}\right)$。以典型参数 $d = 64$、$M_{\mathrm{SRAM}} \approx 10^5$ 元素为例，$\frac{d^2}{M_{\mathrm{SRAM}}} \approx 0.04$，HBM 读写量直接暴跌一个数量级以上！证毕。
\end{proof}

\begin{remark}[反向传播如何同样避免保存 $O(N^2)$ 矩阵？]
细心的读者一定会追问：前向没把 $P \in \R^{N\times N}$ 存到 HBM，那反向传播求 $\nabla_Q, \nabla_K, \nabla_V$ 时怎么办？FlashAttention 再次祭出了“以算力换显存带宽”的法宝：前向传播时，对每一行 $i$ 仅在 HBM 中额外存下\textbf{一个标量对数配分函数} $L_i = m_i^{(T_c)} + \ln \ell_i^{(T_c)} \in \R$（全序列仅占 $O(N)$ 显存，微不足道！）。反向传播分块扫描时，在片上 SRAM 中临时重算 $S_{ij} = \tau Q_i K_j^\top$，随即用单条公式 $P_{ij} = \exp(S_{ij} - L_i)$ 毫秒级重构出精确注意力概率块！尽管多做了一次 $Q_i K_j^\top$ 重计算，但因为彻底免去了读写巨型 $N\times N$ 矩阵的 HBM 延迟，反向传播的实际墙钟耗时反而\textbf{大幅缩短}！
\end{remark}

图~\ref{fig:flash-vs-standard} 从物理数据流视角并排对比了标准注意力与 FlashAttention 在 HBM 与片上 SRAM 之间的字节搬运差异。

\begin{figure}[htbp]
\centering
\resizebox{0.97\linewidth}{!}{%
\begin{tikzpicture}[
  font=\sffamily\footnotesize,
  hbmbox/.style={rectangle, rounded corners=1.2mm, draw=ochre, thick, fill=ochreLight, align=center, text=ink, inner sep=3.5pt},
  srambox/.style={rectangle, rounded corners=1.2mm, draw=teal, thick, fill=tealLight, align=center, text=ink, inner sep=3.5pt},
  arrO/.style={-{Stealth[length=2mm]}, very thick, ochre},
  arrT/.style={-{Stealth[length=2mm]}, very thick, tealDark}
]
  % Left panel: Standard Attention
  \draw[slate!40, dashed, rounded corners=2mm, fill=codebg!50] (-7.3, -2.55) rectangle (-0.35, 3.25);
  \node[font=\sffamily\small\bfseries, color=ochre] at (-3.82, 2.92) {(a) 标准注意力：反复往返 HBM（$\Theta(N^2)$ IO）};

  \node[hbmbox, minimum width=62mm, minimum height=11mm] (std_hbm) at (-3.82, 1.80) {
    \textbf{片外 HBM（3.35\,TB/s）}\\
    {\scriptsize 存放 $Q, K, V, O$ 以及巨型 \textbf{$S, P \in \R^{N\times N}$}}
  };

  \node[srambox, minimum width=18mm, minimum height=13mm] (k1) at (-5.95, -0.95) {\textbf{核 1}\\$S = QK^\top$};
  \node[srambox, minimum width=19mm, minimum height=13mm] (k2) at (-3.82, -0.95) {\textbf{核 2}\\$P = \mathrm{softmax}(S)$};
  \node[srambox, minimum width=18mm, minimum height=13mm] (k3) at (-1.70, -0.95) {\textbf{核 3}\\$O = PV$};

  % Heavy arrows between HBM and kernels (staggered vertically: up=top, down=bottom)
  \draw[arrO] (-6.40, 1.25) -- node[pos=0.50, left, font=\tiny, color=ochre] {$Q,K$} (-6.40, -0.30);
  \draw[arrO] (-5.50, -0.30) -- node[pos=0.74, right, font=\tiny\bfseries, color=ochre] {写 $S$} (-5.50, 1.25);
  \draw[arrO] (-4.25, 1.25) -- node[pos=0.74, left, font=\tiny\bfseries, color=ochre] {读 $S$} (-4.25, -0.30);
  \draw[arrO] (-3.38, -0.30) -- node[pos=0.74, right, font=\tiny\bfseries, color=ochre] {写 $P$} (-3.38, 1.25);
  \draw[arrO] (-2.12, 1.25) -- node[pos=0.74, left, font=\tiny\bfseries, color=ochre] {读 $P,V$} (-2.12, -0.30);
  \draw[arrO] (-1.22, -0.30) -- node[pos=0.50, right, font=\tiny, color=ochre] {写 $O$} (-1.22, 1.25);

  \node[font=\scriptsize, color=ochre, align=center] at (-3.82, -2.12) {三次核启动；巨型矩阵 $S, P$（各 $N^2$）四次穿透 HBM 带宽墙};

  % Right panel: FlashAttention
  \draw[tealLine, thick, rounded corners=2mm, fill=tealLight!45] (0.35, -2.55) rectangle (7.55, 3.25);
  \node[font=\sffamily\small\bfseries, color=tealDark] at (3.95, 2.92) {(b) FlashAttention：片上融合分块（无 $N^2$ 落地）};

  \node[hbmbox, minimum width=64mm, minimum height=10.5mm] (fl_hbm) at (3.95, 1.80) {
    \textbf{片外 HBM（仅存输入输出）}\\
    {\scriptsize $Q, K, V, O \in \R^{N\times d}$ 及行标量 $L \in \R^N$（\textbf{无 $S, P$！}）}
  };

  % Single Fused SRAM Kernel
  \node[srambox, minimum width=67mm, minimum height=22mm] (fl_sram) at (3.95, -0.78) {};
  \node[font=\sffamily\scriptsize\bfseries, color=tealDark] at (3.95, 0.05) {单一融合核函数（片上 SRAM，$\sim 19\,\mathrm{TB/s}$）};

  \node[rectangle, rounded corners=1mm, draw=navy, fill=navyLight, font=\tiny, align=center, inner sep=2.2pt] (tile1) at (1.62, -0.82) {瓦片乘加\\$S_{ij} = \tau Q_i K_j^\top$\\($B_r \times B_c$)};
  \node[rectangle, rounded corners=1mm, draw=gold, fill=ochreLight, font=\tiny, align=center, inner sep=2.2pt] (tile2) at (3.78, -0.82) {在线 Softmax\\更新 $m_i, \ell_i, \alpha_i$\\重标定 $\tilde{P}_{ij}$};
  \node[rectangle, rounded corners=1mm, draw=teal, fill=white, font=\tiny, align=center, inner sep=2.2pt] (tile3) at (6.12, -0.82) {累加输出\\$O_i \!\leftarrow\! \frac{\alpha\ell O_i + \tilde{P}_{ij}V_j}{\ell_{\mathrm{new}}}$};

  \draw[-{Stealth[length=1.7mm]}, thick, navy] (tile1.east) -- (tile2.west);
  \draw[-{Stealth[length=1.7mm]}, thick, gold] (tile2.east) -- (tile3.west);
  \draw[-{Stealth[length=1.6mm]}, dashed, tealDark] (tile3.south) -- ++(0, -0.25) -| node[pos=0.25, below, font=\tiny\bfseries, color=tealDark] {循环下一键值块 $j \leftarrow j+1$（全程驻留片上）} (tile1.south);

  \draw[arrT] (2.15, 1.25) -- node[left, font=\tiny\bfseries, color=tealDark] {载入小块 $Q_i, K_j, V_j$} (2.15, 0.34);
  \draw[arrT] (5.75, 0.34) -- node[right, font=\tiny\bfseries, color=tealDark] {终态写出 $O_i, L_i$} (5.75, 1.25);

  \node[font=\scriptsize, color=tealDark, align=center] at (3.95, -2.22) {峰值瓦片仅占 $B_r B_c / N^2$，算子从访存受限跃迁至计算受限！};
\end{tikzpicture}%
}
\caption{标准注意力与 FlashAttention 的硬件数据流对比。FlashAttention 通过分块在线 Softmax 将三次核函数融合为单次片上流式流水线，使庞大的 $S, P$ 矩阵永远不触碰片外 HBM。}
\label{fig:flash-vs-standard}
\end{figure}

\subsection{实测验证：纯 NumPy 分块在线 Softmax 与标准注意力的机器精度比对}

为了亲眼验证定理~\ref{thm:online-softmax-tiling} 绝非纸上空谈，代码~\ref{lst:flash-online} 给出了与配套脚本 \texttt{code/vol6\_experiments.py} 中实验 5 完全一致的纯 NumPy 实现（$N = 256, d = 64, B_r = B_c = 32$）。

\begin{lstlisting}[caption={纯 NumPy 实现 FlashAttention 在线 Softmax 分块递推并与标准注意力做精度核对。},label={lst:flash-online}]
import math
import numpy as np

rng = np.random.default_rng(2026)
N, d = 256, 64
Q = rng.normal(size=(N, d))
K = rng.normal(size=(N, d))
V = rng.normal(size=(N, d))
scale = 1.0 / math.sqrt(d)

# 1. 标准注意力（显式物化完整 N x N 分数与概率矩阵）
S = (Q @ K.T) * scale
S_shift = S - np.max(S, axis=-1, keepdims=True)
P = np.exp(S_shift)
P /= np.sum(P, axis=-1, keepdims=True)
O_std = P @ V

# 2. FlashAttention 核心：分块在线 Softmax（全程仅保留 Br x Bc 小瓦片）
Br, Bc = 32, 32
O_flash = np.zeros_like(Q)
for i in range(0, N, Br):
    Qi = Q[i:i+Br]
    mi = np.full((Br, 1), -np.inf)
    li = np.zeros((Br, 1))
    Oi = np.zeros((Br, d))
    for j in range(0, N, Bc):
        Kj, Vj = K[j:j+Bc], V[j:j+Bc]
        Sij = (Qi @ Kj.T) * scale               # 仅占 Br x Bc SRAM
        mij = np.max(Sij, axis=-1, keepdims=True)
        m_new = np.maximum(mi, mij)
        alpha = np.exp(mi - m_new)              # 历史重标定因子 <= 1
        pij = np.exp(Sij - m_new)               # 当前块非归一化概率
        l_new = alpha * li + np.sum(pij, axis=-1, keepdims=True)
        Oi = (Oi * (alpha * li) + pij @ Vj) / l_new
        mi, li = m_new, l_new
    O_flash[i:i+Br] = Oi

max_err = float(np.max(np.abs(O_std - O_flash)))
peak_ratio = (Br * Bc) / (N * N) * 100.0
print(f"max_err = {max_err:.2e}, peak_tile_ratio = {peak_ratio:.2f}%")
\end{lstlisting}

在 IEEE-754 双精度浮点下的实测输出令人拍案叫绝：
\begin{itemize}
  \item \textbf{最大逐元素绝对误差}：\texttt{max\_err = 3.89e-16}（精确落入 64 位浮点数的机器精度量级 $\epsilon_{\mathrm{FP64}} \approx 2.22\times 10^{-16}$，数学上完全无损！）；
  \item \textbf{峰值注意力分数瓦片占比}：\texttt{peak\_tile\_ratio = 1.56\%}（即 $\frac{32\times 32}{256\times 256} = \frac{1}{64}$）；若序列长度延伸至 $N = 8192$ 而块大小仍保持 $B_r = B_c = 64$，峰值瓦片更是仅占完整矩阵的 $\mathbf{0.006\%}$！
\end{itemize}
仔细留意代码第 33 行：与教材中“把未除以 $\ell_i$ 的未归一化累加器 $\tilde{O}_i$ 留到循环最后统一除以 $\ell_i^{(T_c)}$”的 FlashAttention-2 写法相比，这里每步即时除以 \texttt{l\_new} 与最后统一除法在代数上完全等价（见习题~\textbf{6.5.4}），而在 GPU 上把除法推迟到内循环之外还能进一步省去 $T_c - 1$ 次向量除法开销！

\section{硬件优化陷阱与前沿算子生态}
\label{sec:hw-pitfalls-ext}

\begin{failBox}{算子工程中的三大隐蔽性能与数值陷阱}
\begin{enumerate}
  \item \textbf{访存受限的未融合逐元素算子（为什么 Kernel Fusion 至关重要）}：在 Transformer 中，诸如残差相加（Residual Add）、RMSNorm、Dropout、SwiGLU 激活等逐元素或逐行归约算子，虽然总 FLOPs 不足整网的 $1\%$，但如果不做\textbf{算子融合（Kernel Fusion）}，每次独立调用都要把完整的 $(B, S, d)$ 张量从 HBM 读进来再写回去——算术强度甚至低于 $1\,\mathrm{FLOP/B}$！在长序列或中小模型上，这些“看似微不足道”的逐元素算子往往会吃掉 $30\%\text{--}50\%$ 的总墙钟时间。PyTorch 2.0 的 \texttt{torch.compile}（TorchInductor）自动把连续逐元素算子融合成单个片上核函数，本质上就是把它们对 HBM 的往返次数从 $k$ 次压成 $1$ 次。
  \item \textbf{非合并访存（Non-Coalesced Memory Access）与共享内存 Bank Conflict}：GPU 从 HBM 读取数据时以 32\,B 或 128\,B 的缓存行（Cache Line）为最小事务单位。如果同一个 Warp 的 32 个线程访问的是相邻连续地址（例如行主序遍历矩阵的同一行），硬件只需 1 次内存事务即可满足全部 32 个线程（\textbf{合并访存}）；反之，若按列大步幅跨行跳跃访问，32 个线程可能触发 32 次独立内存事务，实际有效带宽瞬间暴跌至标称值的 $1/32$！同理，片上 Shared Memory 通常划分为 32 个并发存储体（Bank），若同 Warp 多个线程同时读写同一 Bank 的不同地址，便会发生串行化的 Bank Conflict（工程上常在二维瓦片列宽上填充 1--8 个 Padding 元素打破周期对齐）。
  \item \textbf{FP16 梯度下溢（Underflow）陷阱 vs BF16 动态范围优势}：虽然 FP16 与 BF16 同占 2 字节（16 比特），但二者的比特位分配截然不同：
  \begin{itemize}
    \item \textbf{FP16}（1 符号 + 5 指数 + 10 尾数）：最大表示值仅 $65504$，最小正规数约 $6.1\times 10^{-5}$。反向传播中大量微弱梯度（$10^{-6}\text{--}10^{-8}$）在 FP16 中会直接下溢归零，导致模型学不动！因此使用 FP16 混合精度训练时\textbf{必须配合动态损失缩放（Dynamic Loss Scaling）}：反向前把损失乘以一个大系数 $S_{\mathrm{loss}}$（如 $2^{16}$）把梯度抬进正规区，更新权重前再除以 $S_{\mathrm{loss}}$；一旦检测到梯度出现 \texttt{Inf/NaN} 则跳过该步并把 $S_{\mathrm{loss}}$ 减半。
    \item \textbf{BF16}（1 符号 + 8 指数 + 7 尾数）：直接继承了 FP32 完整的 8 位指数（动态范围 $\sim 10^{\pm 38}$），梯度几乎永不下溢，彻底免去了繁琐的 Loss Scaling！代价则是尾数精度较粗糙，从而严格要求按式~\eqref{eq:16phi-rule} 维护 FP32 主权重。
  \end{itemize}
\end{enumerate}
\end{failBox}

\begin{extBox}{FlashAttention-2/3、FP8 Tensor Core 与 OpenAI Triton 编程范式}
FlashAttention 提出之后，软硬件协同演化仍在飞速狂奔：
\begin{itemize}
  \item \textbf{FlashAttention-2}（2023）：将内循环中的非矩阵乘浮点运算（如逐块缩放除 $\ell_i$）推迟到循环结束统一执行，仅保留未归一化累加器 $\tilde{O}_i$；同时改写外循环为遍历查询块 $Q_i$ 以便在序列维上充分并行，并优化 SM 内部多个 Warp 的切分方式（按 $Q$ 切分而非按 $K$ 切分），彻底消除了 Warp 间通过 Shared Memory 同步归约的通信开销，使 A100 上的 Tensor Core 利用率从约 $35\%$ 跃升至 $70\%$ 以上。
  \item \textbf{FlashAttention-3}（2024，面向 Hopper H100）：深度挖掘 Hopper 架构的两项异步硬件特性——\textbf{TMA（张量内存加速器）}异步拷贝引擎与 \textbf{WGMMA} 异步矩阵乘指令，让“从 HBM 搬运下一块 $K_{j+1}, V_{j+1}$”、“在 Tensor Core 上算当前块 $Q_i K_j^\top$”与“在特殊函数单元 SFU 上算上一块的指数 Softmax”三条流水线在硬件上真正重叠执行；同时引入\textbf{块级 FP8 量化}与非相干哈达玛变换，将单卡注意力吞吐进一步推高至惊人的 $1.2\,\mathrm{PFLOP/s}$！
  \item \textbf{OpenAI Triton 块级编程}：传统编写一个高性能 CUDA C++ 融合算子需要手工管理每个线程的索引、共享内存 Bank Conflict、异步预取指令与寄存器溢出，门槛极高。Triton 将编程抽象从“单线程”提升到了“瓦片块（Block/Tile）”：开发者只需用类 Python 语法描述如何加载 $(B_r, d)$ 的矩阵块、在块上做矩阵乘和在线 Softmax 并写回，编译器便会自动完成合并访存调度、共享内存分配与 Tensor Core 指令映射——几十行 Triton 代码即可跑出接近手工汇编 CUDA 核 $90\%$ 以上的性能。
\end{itemize}
\end{extBox}

\section*{本章分层习题}
\addcontentsline{toc}{section}{本章分层习题}

\exTier{A 层：概念与公式手算}{熟练核算 Roofline 拐点、临界批大小与 $16\Phi$ 训练显存账本}

\begin{enumerate}
  \item[\textbf{6.5.1}] \textbf{Roofline 算术强度与跨芯片拐点核算}：假设某新型 AI 加速芯片的 BF16 峰值算力为 $P_{\mathrm{peak}} = 2000$ TFLOP/s，HBM 显存带宽为 $B_{\mathrm{HBM}} = 8.0\,\mathrm{TB/s}$。现于其上执行维度为 $d_{\mathrm{in}} = 4096, d_{\mathrm{out}} = 16384$ 的 BF16 线性层 $Y = X W$（输入 $X \in \R^{T \times 4096}$）。
  \begin{enumerate}
    \item 计算该芯片的硬件拐点 $I^*$（单位：$\mathrm{FLOP/B}$）；
    \item 推导该非方阵线性层的精确算术强度 $\mathrm{AI}(T)$ 公式（注意读 $X$、读 $W$、写 $Y$ 的字节数各不相同）；
    \item 分别计算 $T = 1$（单请求 Decode）与 $T = 512$（Prefill）时的可达算力与硬件利用率，并求出进入计算受限区所需的最小词元数 $T^*$。
  \end{enumerate}
  \exHint{(a) $I^* = 2000 / 8.0 = 250\,\mathrm{FLOP/B}$；(b) $F = 2 T d_{\mathrm{in}} d_{\mathrm{out}}$，$Q = 2(d_{\mathrm{in}} d_{\mathrm{out}} + T d_{\mathrm{in}} + T d_{\mathrm{out}})$，从而 $\mathrm{AI}(T) = T d_{\mathrm{in}} d_{\mathrm{out}} / [d_{\mathrm{in}} d_{\mathrm{out}} + T(d_{\mathrm{in}} + d_{\mathrm{out}})]$；代入 $d_{\mathrm{in}} = 4096, d_{\mathrm{out}} = 16384$ 可解得 $T^* \approx 271$。}

  \item[\textbf{6.5.2}] \textbf{13B 大模型训练静态与激活显存全账本}：现拟在 $80\,\mathrm{GB}$ 显存的 GPU 集群上，使用混合精度（BF16 权重与梯度 + FP32 主权重与 AdamW）训练一个 $\Phi = 1.3\times 10^{10}$（13B）的 Transformer 模型，具体结构参数为 $L = 40$ 层、隐藏维度 $d = 5120$、注意力头数 $h = 40$，单卡微批大小 $B = 2$、序列长度 $S = 4096$。
  \begin{enumerate}
    \item 用 $16\Phi$ 法则分别列出权重、梯度、主权重与优化器状态各占多少 GB，静态总计多少 GB？
    \item 用式~\eqref{eq:act-memory-layer} 分别计算\textbf{未启用 FlashAttention}（保留 $O(S^2)$ 项）与\textbf{启用 FlashAttention}（消去 $\frac{5hS}{d}$ 项）时全网 40 层的总激活显存；
    \item 若进一步对 40 层启用均匀激活检查点（取近似最优块长 $k = 5$ 或 $k = 8$），启用 FlashAttention 后的峰值激活显存将压缩到多少 GB？
  \end{enumerate}
  \exHint{(a) 静态总计 $16 \times 13 = 208\,\mathrm{GB}$；(b) $B S d = 2 \times 4096 \times 5120 \approx 4.194\times 10^7$ 字节，$\frac{5hS}{d} = 160$，可见标准注意力下 $O(S^2)$ 项竟是线性项 $34$ 的近 5 倍！全网激活约 $325\,\mathrm{GB}$（无 Flash）降至 $57\,\mathrm{GB}$（有 Flash）；(c) 取 $k=8$ 时仅保留 $L/k + k = 5 + 8 = 13$ 层，峰值激活压至约 $18.5\,\mathrm{GB}$。}
\end{enumerate}

\exTier{B 层：严谨数学证明}{吃透非均匀检查点优化、未归一化在线 Softmax 与反向重算恒等式}

\begin{enumerate}
  \item[\textbf{6.5.3}] \textbf{激活检查点定理的推广与极限边界}：在定理~\ref{thm:activation-checkpointing} 的设定下：
  \begin{enumerate}
    \item 若我们追求\textbf{极限压低显存}（甚至把检查点也只保留 1 个，并递归分段，即 Griewank 的二项式检查点思想），请先考虑最简单的“两级嵌套检查点”：将 $L$ 层分为 $L / k_1$ 个大段，每大段再按 $k_2$ 分为子段（$1 \le k_2 \le k_1 \le L$）。推导峰值激活显存关于 $(k_1, k_2)$ 的表达式，并用三元 AM-GM 不等式证明其最小显存可进一步降至 $3 L^{1/3} M_1$，同时计算额外前向传播次数变为多少倍？
    \item 说明在真实大规模预训练中为何工程界普遍采用单级检查点（甚至仅按每个 Transformer 层边界做检查点，$k = 1$ 的算子级局部重算）而非三级以上嵌套重算。
  \end{enumerate}
  \exHint{(a) 峰值激活为 $(\frac{L}{k_1} + \frac{k_1}{k_2} + k_2) M_1 \ge 3 L^{1/3} M_1$，当 $k_1 = L^{2/3}, k_2 = L^{1/3}$ 时取等，但每个底层需经历多达 2 次额外前向（总算力开销升至 $+66.7\%$）；(b) 当单级检查点已把激活显存压到低于静态模型显存或通信带宽边界时，再牺牲算力换显存得不偿失。}

  \item[\textbf{6.5.4}] \textbf{FlashAttention-2 延迟归一化形式与反向传播 $L_i$ 重构证明}：
  \begin{enumerate}
    \item 证明：若在定理~\ref{thm:online-softmax-tiling} 的内循环中去掉每步除以 $\ell_i^{(j)}$ 的操作，改为维护未归一化累加器 $\tilde{O}_i^{(j)} = \alpha_i^{(j)} \tilde{O}_i^{(j-1)} + \tilde{P}_{ij} V_j$（初始 $\tilde{O}_i^{(0)} = \mathbf{0}$），仅在循环结束 $j = T_c$ 后统一计算 $O_i = \tilde{O}_i^{(T_c)} / \ell_i^{(T_c)}$，则所得结果与式~\eqref{eq:flash-o} 完全恒等；
    \item 前向仅保存标量 $L_i = m_i^{(T_c)} + \ln \ell_i^{(T_c)}$。证明在反向传播重算任意分数瓦片 $S_{ij} = \tau Q_i K_j^\top$ 后，精确注意力概率块恰等于 $P_{ij} = \exp(S_{ij} - L_i)$（按行广播），并据此写出利用 $dO_i$、$O_i$ 与 $P_{ij}$ 在片上计算行标量 $D_i = \operatorname{rowsum}(dO_i \odot O_i)$ 与分数梯度 $dS_{ij} = P_{ij} \odot (dP_{ij} - D_i)$ 的精确推导！
  \end{enumerate}
  \exHint{(a) 直接归纳证明 $\tilde{O}_i^{(j)} = \ell_i^{(j)} O_i^{(j)}$；(b) 注意 $\exp(S_{ij} - m_i - \ln \ell_i) = \frac{\exp(S_{ij} - m_i)}{\ell_i} = P_{ij}$，再利用 Softmax 雅可比公式 $dS_i = P_i \odot (dP_i - \sum_k P_{ik} dP_{ik})$ 以及 $\sum_k P_{ik} dP_{ik} = \sum_k P_{ik} \langle dO_i, V_k\rangle = \langle dO_i, \sum_k P_{ik} V_k\rangle = \langle dO_i, O_i\rangle = D_i$！最后一步把跨列的 $dP_{ij}$ 归约巧妙转化为了仅依赖当前块 $O_i, dO_i$ 的局部内积，彻底解除反向传播的行同步依赖！}
\end{enumerate}

\exTier{C 层：机制剖析与反例设计}{诊断融合收益与 BF16 尾数吞噬反例}

\begin{enumerate}
  \item[\textbf{6.5.5}] \textbf{定量诊断算子融合对访存受限算子的加速比}：考虑形状为 $(B, S, d) = (16, 2048, 4096)$ 的 BF16 张量流水线：
  \[
    H_1 = X + R \;\text{（残差加）} \;\longrightarrow\; H_2 = \mathrm{RMSNorm}(H_1) \;\longrightarrow\; H_3 = \mathrm{SiLU}(H_2).
  \]
  \begin{enumerate}
    \item 分别计算\textbf{未融合}（3 个独立核函数，每次读写 HBM）与\textbf{完全融合}（单核片上流水，仅读 $X, R$ 与归一化权重 $\gamma \in \R^d$，写出 $H_3$）两种实现下的 HBM 总字节读写量（忽略极小的 $d$ 维参数 $\gamma$）；
    \item 在 H100（$3.35\,\mathrm{TB/s}$）上，假设上述三步均处于深度访存受限区，计算融合前后的理论墙钟耗时与加速比。
  \end{enumerate}
  \exHint{单张量体积 $V = 16 \times 2048 \times 4096 \times 2\,\mathrm{B} = 268.4\,\mathrm{MB}$。未融合：残差加读 2 写 1（$3V$）+ RMSNorm 读 1 写 1（$2V$）+ SiLU 读 1 写 1（$2V$），总计 $7V$；融合后仅读 $X, R$ 并写 $H_3$，总计 $3V$，理论直接提速 $7/3 \approx 2.33\times$！}

  \item[\textbf{6.5.6}] \textbf{数值反例：为什么纯 BF16 权重更新会导致优化彻底“冻结”？}：IEEE 风格 BF16 的格式为 1 位符号、8 位指数、7 位显式尾数（加 1 位隐含前导 1，共 8 位有效二进制数字）。
  \begin{enumerate}
    \item 在区间 $[1, 2)$ 内，相邻两个可精确表示的 BF16 浮点数之间的步长（ULP）是多少？
    \item 设某权重当前为 $w = 1.0$，采用学习率 $\eta = 10^{-3}$，每步精确梯度恒为 $g = -1.0$（理想情况下每步应增加 $+10^{-3}$）。若不保存 FP32 主权重，直接按“舍入到最近偶数（Round-to-Nearest-Even）”规则在 BF16 中执行 $w \leftarrow \mathrm{BF16}(w - \eta g)$，经过 $10000$ 步后 $w$ 的值是多少？请精确说明为何必须搭配 FP32 主权重（或随机舍入 Stochastic Rounding）才能打破这一冻结陷阱。
  \end{enumerate}
  \exHint{(a) 在 $[1, 2)$ 内，ULP $= 2^{-7} = 1/128 = 0.0078125$；(b) 半个 ULP 为 $0.00390625$，由于单步增量 $0.001 < 0.00390625$，每次加完舍入后严格等于 $1.0$，跑完 $10000$ 步后 $w$ 依然纹丝不动停在 $1.0$！}
\end{enumerate}

\exTier{D 层：编程实验与系统延伸}{动手实现因果掩码 FlashAttention 与长上下文 GQA 带宽仿真}

\begin{enumerate}
  \item[\textbf{6.5.7}] \textbf{编程实验：为分块在线 Softmax 添加因果掩码（Causal Mask）与上三角跳过加速}：基于代码~\ref{lst:flash-online}（实验 5 脚本）：
  \begin{enumerate}
    \item 在标准注意力中为上三角（列索引 $>$ 行索引）加上 $-\infty$ 因果掩码，得到自回归语言模型的标准因果输出 $O_{\mathrm{causal}}$；
    \item 在分块在线 Softmax 循环中实现两项工程优化：对完全位于对角线上方的块（\texttt{j > i + Br - 1}）直接 \texttt{continue} 跳过零计算；仅对主对角线上的块（\texttt{j == i}，当 $B_r = B_c$）施加局部上三角 $-\infty$ 掩码；
    \item 验证在 $N = 256, B_r = B_c = 32$ 时，最大误差仍小于 $10^{-15}$，并统计跳过的块计算量比例（理论上接近 $50\%$ 加速！）。
  \end{enumerate}
  \exHint{跳过上三角纯掩码块后，总计 $8 \times 8 = 64$ 个瓦片中仅需计算下三角与对角线共 $8 \times 9 / 2 = 36$ 个瓦片，直接省去 $43.75\%$ 的块乘加；当块数增多时渐近省去 $50\%$ 计算量。}

  \item[\textbf{6.5.8}] \textbf{编程实验：Roofline 分析标准 MHA 与分组查询注意力（GQA）在长上下文 Decode 中的带宽瓶颈}：编写一个小型 Python Roofline 模拟器，考虑单层自注意力在 Decode 阶段（批量 $B = 16$，历史上下文长度 $S$ 从 $512$ 扫描至 $32768$，$d = 4096$，查询头数 $h_q = 32$，头维度 $d_h = 128$）：
  \begin{enumerate}
    \item 推导该步注意力打分 $q K_{\mathrm{cache}}^\top$ 与加权求和 $p V_{\mathrm{cache}}$ 的总 FLOPs（约为 $4 B S d$）与从 HBM 读入历史 KV 缓存的字节数（对于标准 MHA，$h_{\mathrm{kv}} = 32$；对于 GQA-8，$h_{\mathrm{kv}} = 4$；每元素 BF16 占 2 字节）；
    \item 计算 MHA 与 GQA-8 在上述条件下的注意力 Decode 算术强度 $\mathrm{AI}_{\mathrm{attn}}$，验证一个惊人事实：\textbf{无论批量 $B$ 或长度 $S$ 多大，标准 MHA 的 Decode 注意力算术强度都被死死锁死在 $1\,\mathrm{FLOP/B}$ 附近}，而 GQA-8 将其精确抬升了 $8\times$！
  \end{enumerate}
  \exHint{因为每个请求的 KV 缓存都是各自独立的，无法像模型权重 $W$ 那样在批量 $B$ 之间共享复用！单请求每个 KV 元素（2 字节）在 MHA 中仅与 1 个查询头做 1 次 FMA（2 FLOPs），故 $\mathrm{AI}_{\mathrm{MHA}} \approx 1\,\mathrm{FLOP/B}$；而在 GQA-8 中，同一份 KV 头被 8 个查询头共享复用，算术强度直接提升 $8\times$！}
\end{enumerate}

\chapter[分布式训练、LoRA、量化与高效推理]{分布式训练、低比特量化与高效推理系统}
\label{ch:dist-quant-serving}

\chapEpigraph{人工智能七十年研究史给我们的最大教训是：能够借力通用算力的方法，最终总能以压倒性优势胜出。}{Richard S. Sutton（理查德·萨顿，《苦涩的教训》，2019）}

\begin{chapIntro}{本章导读：从单卡极致走向万卡协同与毫秒级生产推理}
在第~\ref{ch:hardware-roofline-flash} 章中，我们借助 Roofline 模型与 FlashAttention 压榨出了单张 GPU 的每一分算力与片上 SRAM 带宽。然而，静态显存的 $16\Phi$ 法则立刻抛出了一个无法在单卡上解决的硬边界：一座 70B 参数的大模型仅静态训练状态就需要 $1120\,\mathrm{GB}$ 显存，405B 模型更是高达 $6.5\,\mathrm{TB}$——相当于数十甚至上百张旗舰 GPU 的显存总和！另一方面，当模型训练完毕、面向千万并发用户上线服务时，如何在单请求严重受限于 HBM 带宽的 Decode 阶段同时做到\textbf{低成本微调、高吞吐并发、低延迟首字与零精度损失加速}，成为当代 AI 系统工程皇冠上的明珠。

本章将沿着“\textbf{万卡训练 $\to$ 轻量微调 $\to$ 低比特压缩 $\to$ 高吞吐推理引擎}”的全生命周期，完成四场系统与数学的深度会师：
\begin{enumerate}
  \item \textbf{集合通信与 3D 并行精确账本}：严格证明 Ring AllReduce 每卡 $2\frac{P-1}{P}M$ 的集群规模无关通信量；逐字节推导 DDP 与 ZeRO Stage 1/2/3（FSDP）的显存与通信公式，并严谨剖析 Megatron-LM 张量并行与 GPipe / 1F1B 流水线并行的精确气泡率 $\frac{p-1}{m+p-1}$。
  \item \textbf{参数高效微调（LoRA）与零延迟合并}：刻画低秩增量 $\Delta W = \frac{\alpha}{r}BA$ 的参数压缩率与 $B=0$ 初始化铁律，并基于 \texttt{data/lora\_rank.csv} 实测展示秩从 $r=1$ 跃升至 $r=4$ 时测试误差骤降 $100\times$ 的内在维度相变。
  \item \textbf{低比特后训练量化（INT8 / INT4）}：数学证明重尾离群通道如何彻底摧毁逐张量量化，而逐通道/分组量化如何隔离离群步长（衔接 SmoothQuant、AWQ 与 GPTQ），并以实测验证 INT8 相对误差骤降 $7\times$、INT4 从崩溃（$63.8\%$）挽救至 $12.1\%$。
  \item \textbf{现代大模型推理系统三驾马车}：剖析连续批处理（Continuous Batching）、直接移植 408 操作系统虚拟分页思想的 PagedAttention（将 KV 显存碎片从 $>60\%$ 压至 $<4\%$），以及投机解码（Speculative Decoding）的\textbf{精确无损分布保持定理}。
\end{enumerate}
\end{chapIntro}

\section{集合通信原语与 3D 并行训练账本}
\label{sec:3d-parallelism}

一旦模型突破单张 GPU 的显存上限，计算图就必须被切分到多卡、多机乃至数万张加速卡上协同执行。在分布式训练中，跨卡通信不再是无伤大雅的配角，而是决定集群扩展效率（MFU）的生死瓶颈。

\subsection{四大集合通信原语与 Ring AllReduce 最优性证明}

分布式深度学习框架（如 NCCL）用四种核心集合通信原语（Collective Primitives）编排 $P$ 张 GPU 之间的数据交换：
\begin{itemize}
  \item \textbf{Broadcast（广播）}：将单张源卡上的张量复制到其余全部 $P-1$ 张卡；
  \item \textbf{ReduceScatter（归约散射）}：对 $P$ 张卡上各自持有的等长张量按元素求和（或其他结合律归约），并将结果均匀切分为 $P$ 份，每张卡仅保留其中对应的 $1/P$ 份；
  \item \textbf{AllGather（全收集）}：每张卡各自持有一份长度为 $M/P$ 的分片，通信结束后每张卡都拼接出完整的长度为 $M$ 的全量张量；
  \item \textbf{AllReduce（全归约）}：对 $P$ 张卡上的张量按元素求和，且通信结束后\textbf{每张卡都拿到完整的全局求和结果}。
\end{itemize}
请注意分布式系统中最重要的一条结构恒等式：
\begin{equation}
  \mathrm{AllReduce} \;\equiv\; \mathrm{ReduceScatter} \;+\; \mathrm{AllGather}.
  \label{eq:allreduce-decomp}
\end{equation}
如果天真地采用“参数服务器（Parameter Server）”中心节点汇聚全部 $P$ 张卡的梯度再广播回去，中心节点的网卡带宽会随 $P$ 线性拥堵。而 \textbf{Ring AllReduce（环形全归约）}凭借式~\eqref{eq:allreduce-decomp} 的分解，彻底消除了集群规模对带宽吞吐的制约！

\begin{thmBox}
\begin{theorem}[Ring AllReduce 的精确通信量与集群规模无关性]
\label{thm:ring-allreduce}
设有 $P \ge 2$ 张 GPU 构成逻辑单向环 $0 \to 1 \to \cdots \to P-1 \to 0$，每张 GPU 持有长度为 $M$ 字节的局部梯度张量 $g^{(k)}$，相邻链路单向带宽为 $B$、单次传输固定延迟为 $\alpha$。通过将每个张量均分为 $P$ 个分片（每片大小 $M/P$ 字节），依次执行 $P-1$ 步环形 ReduceScatter 与 $P-1$ 步环形 AllGather，则：
\begin{enumerate}
  \item 通信结束时，每张 GPU 精确获得全局和 $\sum_{k=0}^{P-1} g^{(k)}$；
  \item 每张 GPU 累计发送（及接收）的精确总字节数为
  \begin{equation}
    V_{\mathrm{Ring}} \;=\; 2 \cdot \frac{P - 1}{P}\, M \;<\; 2 M;
    \label{eq:ring-vol}
  \end{equation}
  \item 忽略本地加法耗时，总通信墙钟时间为
  \begin{equation}
    T_{\mathrm{Ring}} \;=\; 2(P - 1)\alpha \;+\; 2 \cdot \frac{P - 1}{P} \cdot \frac{M}{B}.
    \label{eq:ring-time}
  \end{equation}
\end{enumerate}
当梯度张量足够大（$M \gg P \alpha B$）时，带宽耗时 $\frac{2(P-1)}{P}\frac{M}{B} \to \frac{2M}{B}$ \textbf{完全独立于 GPU 总数 $P$}！
\end{theorem}
\end{thmBox}

\begin{proof}
将每张 GPU $k \in \{0, \dots, P-1\}$ 的张量切分为 $P$ 块分片 $g^{(k)}[0], g^{(k)}[1], \dots, g^{(k)}[P-1]$，每块大小均为 $M/P$ 字节。

\textbf{阶段一：环形 ReduceScatter（共 $P-1$ 轮）}。在第 $r$ 轮（$r = 0, 1, \dots, P-2$）中，所有 GPU 并发工作：每张卡 $k$ 将其当前负责的块索引 $c(k, r) = (k - r) \bmod P$ 发送给顺时针下一张卡 $(k+1)\bmod P$，同时从上一张卡 $(k-1)\bmod P$ 接收对应块并\textbf{累加到自己的本地对应分片上}。经过 $P-1$ 轮传递后，每个分片恰好顺时针游历了全部 $P$ 张卡各一次并累加了所有卡的贡献；因此每张卡 $k$ 上恰好有一个分片（索引为 $(k+1)\bmod P$）已经持有完整的全局归约和！本阶段每张卡发送了 $P-1$ 次、每次 $M/P$ 字节，发送量精确为 $\frac{P-1}{P}M$。

\textbf{阶段二：环形 AllGather（共 $P-1$ 轮）}。此时每张卡都手握一块“已经求和完毕的成品分片”。继续沿同一逻辑环并发传递 $P-1$ 轮：每张卡把自己刚收到的（或初始持有的）成品分片发给下一张卡，接收后直接\textbf{覆盖写入}对应槽位（无需再做加法）。经过 $P-1$ 轮后，全部 $P$ 块成品分片均传播到每一张 GPU 上，每张卡集齐完整的全局梯度和！本阶段每张卡同样发送了 $\frac{P-1}{P}M$ 字节。

两阶段合计，每张 GPU 累计发送与接收字节数均为 $V_{\mathrm{Ring}} = 2\frac{P-1}{P}M$；每一轮所有 $P$ 条有向链路同时满载并发，单轮耗时 $\alpha + \frac{M}{PB}$，共 $2(P-1)$ 轮即得式~\eqref{eq:ring-time}。此外，在任意规约算法中，每张卡的其余 $(P-1)/P$ 份结果至少需要从外部接收一次（下界 $\frac{P-1}{P}M$），且自己最终分片至少需经归约传播（再加下界 $\frac{P-1}{P}M$），故 Ring AllReduce 在带宽项上是\textbf{严格信息论最优的}。证毕。
\end{proof}

\subsection{数据并行（DDP）与 ZeRO Stage 1 / 2 / 3（FSDP）精确显存通信账本}

回顾第~\ref{ch:hardware-roofline-flash} 章式~\eqref{eq:16phi-rule} 建立的静态显存 $16\Phi$ 法则：BF16 权重占 $2\Phi$ 字节，BF16 梯度占 $2\Phi$ 字节，FP32 主权重与 AdamW 状态合计占 $12\Phi$ 字节。为了清晰区分“参数个数”与“BF16 权重字节数”，下文统一记模型 BF16 参数体积为 $\Psi \triangleq 2\Phi$（字节）。

最基础的\textbf{分布式数据并行（DDP）}在 $N_d$ 张 GPU 上各自复制一份完整的模型状态（每卡消耗全额 $16\Phi$ 字节），每卡处理不同的微批数据，反向传播结束后通过一次 Ring AllReduce 同步全部 BF16 梯度（每卡通信量 $\approx 2\Psi = 4\Phi$ 字节，且可在反向传播计算深层梯度时按桶异步重叠隐藏通信延迟）。

DDP 的痛点一目了然：当 $N_d = 64$ 时，相同的优化器状态、梯度和权重在 64 张卡上足足冗余拷贝了 64 份！微软 DeepSpeed 提出的 \textbf{ZeRO（Zero Redundancy Optimizer）}——以及 PyTorch 官方原生实现的 \textbf{FSDP（Fully Sharded Data Parallel）}——凭借式~\eqref{eq:allreduce-decomp} 的洞察，分三级彻底切碎了这份冗余：
\begin{enumerate}
  \item \textbf{ZeRO-1（切分优化器状态，$P_{\mathrm{os}}$）}：将 $12\Phi$ 字节的 FP32 主权重与 AdamW 状态均匀切分到 $N_d$ 张卡上，每卡仅负责更新自己分管的那 $1/N_d$ 参数！反向传播结束后，不再做全量 AllReduce，而是先做一次 \textbf{ReduceScatter}（通信量 $\Psi$）让每卡只拿到自己分管参数的全局梯度，各自在本地执行 AdamW 更新，最后再做一次 \textbf{AllGather}（通信量 $\Psi$）把更新后的 BF16 权重广播回全集群。总通信量仍为 $\Psi + \Psi = 2\Psi$（\textbf{与 DDP 完全相同，零额外通信代价！}），而单卡静态显存骤降至：
  \[
    M_{\mathrm{ZeRO\text{-}1}} \;=\; 2\Phi + 2\Phi + \frac{12\Phi}{N_d} \;=\; 4\Phi + \frac{12\Phi}{N_d}.
  \]
  \item \textbf{ZeRO-2（切分优化器状态 + 梯度，$P_{\mathrm{os+g}}$）}：反向传播逐层计算完梯度后，立刻就地触发该层的 ReduceScatter 并释放非己分管的梯度显存！通信量\textbf{依然保持 $2\Psi$ 不变}，单卡静态显存进一步压至：
  \[
    M_{\mathrm{ZeRO\text{-}2}} \;=\; 2\Phi + \frac{2\Phi + 12\Phi}{N_d} \;=\; 2\Phi + \frac{14\Phi}{N_d}.
  \]
  \item \textbf{ZeRO-3 / FSDP（全切分：优化器状态 + 梯度 + 模型参数，$P_{\mathrm{os+g+p}}$）}：连 $2\Phi$ 的 BF16 权重也按 $1/N_d$ 彻底切分到各卡！代价仅在于：前向传播算到第 $\ell$ 层时，临时对第 $\ell$ 层权重执行一次 AllGather（凑齐整层权重算完后立刻释放，累计前向通信 $+\Psi$）；反向传播退回第 $\ell$ 层时，再 AllGather 一次权重以算梯度（$+\Psi$），算完后 ReduceScatter 梯度并释放权重（$+\Psi$）。总通信量变为 $3\Psi$——\textbf{仅比 DDP 多了 $1.5\times$ 通信量}（且完全可以通过双缓冲预取与计算高度重叠），却换来了静态显存严格的\textbf{线性无限扩展}：
  \[
    M_{\mathrm{ZeRO\text{-}3}} \;=\; \frac{16\Phi}{N_d}.
  \]
\end{enumerate}

表~\ref{tab:zero-ledger} 完整汇总了 DDP 与三级 ZeRO 的精确显存与通信账本。原本在单卡上需要 $1.12\,\mathrm{TB}$ 的 70B 大模型，在 64 卡 ZeRO-3 下每卡静态显存仅剩区区 $17.5\,\mathrm{GB}$！

\begin{table}[htbp]
\centering
\caption{DDP 与 ZeRO-1/2/3（FSDP）在 $N_d$ 卡上的单卡静态显存与通信量精确对比（$\Psi = 2\Phi$ 为 BF16 权重字节数）。}
\label{tab:zero-ledger}
\small
\begin{tabular}{@{}lcccc@{}}
\toprule
\textbf{并行策略} & \textbf{单卡静态显存（Bytes）} & \textbf{7B ($N_d=8$)} & \textbf{70B ($N_d=64$)} & \textbf{每步单卡通信量} \\
\midrule
DDP（无切分） & $16\Phi$ & $112.0\,\mathrm{GB}$ & $1120.0\,\mathrm{GB}$ & $2\Psi$（AllReduce） \\
ZeRO-1（切优化器） & $4\Phi + 12\Phi / N_d$ & $38.5\,\mathrm{GB}$ & $293.1\,\mathrm{GB}$ & $2\Psi$（RS + AG） \\
ZeRO-2（+ 切梯度） & $2\Phi + 14\Phi / N_d$ & $26.3\,\mathrm{GB}$ & $155.3\,\mathrm{GB}$ & $2\Psi$（RS + AG） \\
ZeRO-3 / FSDP（全切） & $16\Phi / N_d$ & $\mathbf{14.0\,\mathrm{GB}}$ & $\mathbf{17.5\,\mathrm{GB}}$ & $3\Psi$（$1.5\times$ 通信） \\
\bottomrule
\end{tabular}
\end{table}

\subsection{张量并行（Megatron-LM）：单层内仅需一次 AllReduce 的矩阵巧切}
\label{subsec:tensor-parallel}

当单层矩阵本身大到无法装入单卡，或批大小不够大而亟需降低单卡激活显存与单步延迟时，仅仅沿数据维度切分（DDP/ZeRO）已力不从心，我们需要沿权重矩阵的内部维度把\textbf{单个算子拆到 $N_t$ 张卡上联合计算}——这就是 NVIDIA 提出的\textbf{张量模型并行（Tensor Parallelism, TP）}。

以 Transformer 中最耗算力的两层 MLP 块 $Z = \mathrm{GeLU}(X W_1) W_2$ 为例（$W_1 \in \R^{d \times 4d}, W_2 \in \R^{4d \times d}$）。若天真地对 $W_1$ 按行切分，算完局部乘积后必须先通信求和才能喂给非线性函数 $\mathrm{GeLU}$！Megatron-LM 的神来之笔在于采用\textbf{“先列切分、后行切分”的配对设计}（以 $N_t = 2$ 为例，一般 $N_t$ 路同理）：
\begin{enumerate}
  \item \textbf{第一层列并行（Column Parallel）}：将 $W_1$ 按列切分为 $W_1 = \begin{bmatrix} W_{1a} & W_{1b} \end{bmatrix}$（卡 0 持 $W_{1a}$，卡 1 持 $W_{1b}$），输入 $X$ 同步广播给两卡。由于逐元素非线性 $\mathrm{GeLU}$ 与列拼接严格可交换：
  \begin{equation}
    Y \;=\; \mathrm{GeLU}\!\big(X \begin{bmatrix} W_{1a} & W_{1b} \end{bmatrix}\big)
    \;=\; \Big[\,\underbrace{\mathrm{GeLU}(X W_{1a})}_{Y_a \text{（卡 0 本地）}},\;\; \underbrace{\mathrm{GeLU}(X W_{1b})}_{Y_b \text{（卡 1 本地）}}\,\Big],
    \label{eq:tp-col}
  \end{equation}
  第一层前向\textbf{零跨卡通信}！
  \item \textbf{第二层行并行（Row Parallel）}：将 $W_2$ 顺势按行切分为 $W_2 = \begin{bmatrix} W_{2a} \\ W_{2b} \end{bmatrix}$（卡 0 持 $W_{2a}$，卡 1 持 $W_{2b}$）。两卡直接用各自手头的 $Y_a, Y_b$ 做本地矩阵乘，由分块矩阵乘法恒等式：
  \begin{equation}
    Z \;=\; \begin{bmatrix} Y_a & Y_b \end{bmatrix} \begin{bmatrix} W_{2a} \\ W_{2b} \end{bmatrix}
    \;=\; Y_a W_{2a} \;+\; Y_b W_{2b},
    \label{eq:tp-row}
  \end{equation}
  最后仅需执行\textbf{一次 AllReduce} 求和，即可在每张卡上还原完整输出 $Z$！
\end{enumerate}

对于多头自注意力（MHA/GQA），逻辑完全同构：把 $h$ 个独立的注意力头均匀分给 $N_t$ 张卡（等价于对 $W_Q, W_K, W_V$ 做列并行），各卡完全独立地完成本地注意力计算（零通信），再对输出投影矩阵 $W_O$ 做行并行，并在块尾做\textbf{一次 AllReduce}！反向传播中，张量并行在前向入口定义恒等算子 $f$（反向为 AllReduce）、出口定义 AllReduce 算子 $g$（反向为恒等），因此\textbf{每个 Transformer 层（Attention + MLP）在前向共需 2 次 AllReduce、反向共需 2 次 AllReduce}。

\begin{remark}[为什么张量并行必须严格限制在单机 NVLink 内部？]
与 ZeRO 可将通信重叠在整层计算背后不同，TP 的 AllReduce 位于\textbf{每一层的临界路径内部}，下一层必须死等 AllReduce 完成才能开算！一台 8 卡 H100 服务器内部卡间拥有 $900\,\mathrm{GB/s}$ 的双向 NVLink 高速互联，而跨机节点间的 InfiniBand 网络带宽通常仅为 $50\text{--}100\,\mathrm{GB/s}$（相差近一个数量级）。因此工业界铁律是：\textbf{张量并行度严格控制在单机内（$N_t \le 8$），跨机扩展必须交给流水线并行与数据并行！}
\end{remark}

\subsection{流水线并行（GPipe 与 1F1B）：气泡率与峰值显存定理}
\label{subsec:pipeline-parallel}

当需要跨越多台服务器切分一座超深模型（如 $L = 80\text{--}126$ 层）时，\textbf{流水线并行（Pipeline Parallelism, PP）}按层深将模型切分为 $p$ 个连续阶段（Stage $0, 1, \dots, p-1$），每个阶段仅需通过极轻量的点对点通信（Send/Recv）把层边界激活传给下一台服务器（通信量远小于 TP 的层内 AllReduce）。

如果把一整个全局批次（Global Batch）整块送进流水线，同一时刻只有 1 台机器在干活，其余 $p-1$ 台全部干等，空闲率高达 $(p-1)/p$！为了填满流水线，GPipe 将一个全局批次切分为 $m$ 个微批次（Micro-Batch）依次注入流水线；而 Megatron-LM 提出的 \textbf{1F1B（One Forward, One Backward）调度}更进一步，在预热填满流水线后立刻让每个阶段进入“算一个微批次前向、紧接着算一个微批次反向并释放其激活”的稳态交替节奏。

\begin{thmBox}
\begin{theorem}[流水线气泡率与 1F1B 峰值激活显存上界]
\label{thm:pipeline-bubble}
设流水线并行包含 $p$ 个均匀划分的阶段，将单个训练步拆分为 $m$ 个相同大小的微批次（$m \ge p$）。假设每个阶段处理单个微批次的前向耗时为 $t_F$、反向耗时为 $t_B$，忽略极小的点对点边界通信延迟，则：
\begin{enumerate}
  \item 无论是 \textbf{GPipe}（全部 $m$ 个微批次前向完毕后再统一反向）还是 \textbf{1F1B} 稳态调度，单步总墙钟耗时与流水线空闲气泡比例（Bubble Fraction）均精确为
  \begin{equation}
    T_{\mathrm{step}} \;=\; (m + p - 1)(t_F + t_B), \qquad
    f_{\mathrm{bubble}} \;=\; \frac{(p - 1)(t_F + t_B)}{(m + p - 1)(t_F + t_B)} \;=\; \frac{p - 1}{m + p - 1}.
    \label{eq:pp-bubble}
  \end{equation}
  \item 在峰值激活显存占用上（以单个微批次在一个阶段上的激活显存 $M_{\mathrm{mb}}$ 为单位）：
  \begin{itemize}
    \item \textbf{GPipe} 必须在每个阶段保留全部 $m$ 个微批次的激活，峰值激活显存为 $\mathbf{m \cdot M_{\mathrm{mb}} = O(m)}$；
    \item \textbf{1F1B} 在最重负载的阶段 0 上，在飞（In-Flight）微批次个数永不超过 $p$ 个（第 $s$ 阶段最多保留 $p - s$ 个），峰值激活显存严格降至 $\mathbf{p \cdot M_{\mathrm{mb}} = O(p)}$，\textbf{彻底与微批数 $m$ 解耦}！
  \end{itemize}
\end{enumerate}
\end{theorem}
\end{thmBox}

\begin{proof}
（1）对任意阶段（例如首阶段 Stage 0），其有效计算总时长恒为 $T_{\mathrm{useful}} = m(t_F + t_B)$。
在 GPipe 中，第 1 个微批次须依次经过全部 $p$ 个阶段前向到达末阶段（耗时 $p t_F$），随后其余 $m-1$ 个微批次紧密衔接，故全前向结束时刻为 $(m + p - 1) t_F$；对称地，反向传播从末阶段逆流回首阶段同样消耗 $(m + p - 1) t_B$，总耗时 $T_{\mathrm{step}} = (m + p - 1)(t_F + t_B)$。

在 1F1B 调度中，观察末阶段 Stage $p-1$：它需等待第 1 个微批次穿过前 $p-1$ 个阶段前向（空闲 $(p-1)t_F$），随后立刻无任何空隙地连续交替执行全部 $m$ 组“1F + 1B”（耗时 $m(t_F + t_B)$），当最后一个微批次反向离开 Stage $p-1$ 后，还需沿前 $p-1$ 个阶段反向收尾（耗末段外空闲 $(p-1)t_B$）。因此 1F1B 的端到端总耗时依然精确为 $T_{\mathrm{step}} = (m + p - 1)(t_F + t_B)$！两者的空闲气泡时间均为 $T_{\mathrm{step}} - T_{\mathrm{useful}} = (p-1)(t_F + t_B)$，作比即得 $f_{\mathrm{bubble}} = \frac{p-1}{m+p-1}$。

（2）在 GPipe 中，阶段 0 算完所有 $m$ 个微批次的前向时，连第 1 个微批次的反向都尚未传回，因此全部 $m$ 份激活必须同时驻留显存。反观 1F1B：在预热阶段，阶段 $s$（$0 \le s \le p-1$）恰好先执行 $p - s$ 次前向传播；就在阶段 0 完成第 $p$ 个微批次的前向（此刻恰有 $p$ 份激活在飞）之际，第 1 个微批次的反向梯度也正好穿越流水线传回阶段 0！随后阶段 0 每完成 1 次反向（立即释放 1 份激活），才注入 1 个新前向（新增 1 份激活），稳态在飞激活份数恒定锁死在 $p$ 份！证毕。
\end{proof}

这条定理精准点破了 1F1B 的威力：要将气泡率压到 $10\%$ 以下，由式~\eqref{eq:pp-bubble} 需令 $m \ge 9(p-1)$（例如 $p = 8$ 时取 $m = 64$）。若用 GPipe，你必须硬扛 $64$ 份微批激活显存直接 OOM；而改用 1F1B 后，哪怕把 $m$ 增至 $64$ 甚至 $128$，峰值激活显存也纹丝不动地锁死在 $p = 8$ 份（直接省下 $8\times$ 以上显存）！

将上述三种正交并行维度嵌套组合——\textbf{单机内 NVLink 做张量并行（$N_t$） $\times$ 跨机做流水线并行（$N_p$） $\times$ 顶层集群副本做 ZeRO 数据并行（$N_d$）}（总卡数 $N_{\mathrm{GPU}} = N_t \cdot N_p \cdot N_d$）——便构成了万卡级大模型预训练的黄金架构 \textbf{3D 并行（3D Parallelism）}，如图~\ref{fig:3d-parallel-and-1f1b} 所示。

\begin{figure}[htbp]
\centering
\resizebox{0.98\linewidth}{!}{%
\begin{tikzpicture}[
  font=\sffamily\footnotesize,
  fblk/.style={rectangle, draw=tealDark, fill=teal!28, minimum width=6.2mm, minimum height=4.8mm, inner sep=0pt, font=\tiny\bfseries, text=ink},
  bblk/.style={rectangle, draw=ochre, fill=ochre!30, minimum width=9.4mm, minimum height=4.8mm, inner sep=0pt, font=\tiny\bfseries, text=ink}
]
  % Left: 3D Parallelism Topology
  \node[font=\sffamily\small\bfseries, color=tealDark] at (-4.1, 3.45) {(a) 万卡集群 3D 并行拓扑映射};

  % Replica 0 box (Front)
  \draw[navy, thick, rounded corners=1.6mm, fill=navyLight!65] (-7.4, -1.88) rectangle (-0.7, 3.05);
  \node[anchor=west, font=\sffamily\scriptsize\bfseries, color=navy] at (-7.22, 2.75) {数据并行副本 0（ZeRO / FSDP 跨副本同步）};

  % 3 Pipeline stages inside Replica 0
  \foreach \s/\lbl/\x in {0/Stage 0 (节点 0)/-7.15, 1/Stage 1 (节点 1)/-4.90, 2/Stage 2 (节点 2)/-2.65} {
    \draw[teal, thick, rounded corners=1.2mm, fill=white] (\x, -1.35) rectangle (\x+1.90, 2.35);
    \node[font=\tiny\bfseries, color=tealDark] at (\x+0.95, 2.12) {\lbl};
    \foreach \g in {0,1,2,3} {
      \fill[tealLight, draw=tealLine, rounded corners=0.5mm] (\x+0.28, 1.44-\g*0.64) rectangle (\x+1.82, 1.86-\g*0.64);
      \node[font=\tiny, color=ink] at (\x+1.05, 1.65-\g*0.64) {GPU \g\ (TP)};
    }
    % Intra-node NVLink double arrow (cleanly separated from GPU pills)
    \draw[{Stealth[length=1.4mm]}-{Stealth[length=1.4mm]}, thick, gold] (\x+0.13, -0.50) -- (\x+0.13, 1.78);
  }
  % Inter-node PP arrows
  \draw[-{Stealth[length=2mm]}, very thick, ochre] (-5.25, 0.55) -- node[above, font=\tiny\bfseries, color=ochre] {PP} (-4.90, 0.55);
  \draw[-{Stealth[length=2mm]}, very thick, ochre] (-3.00, 0.55) -- node[above, font=\tiny\bfseries, color=ochre] {PP} (-2.65, 0.55);

  % Legend bars under left
  \node[font=\tiny, color=gold!85!black] at (-5.65, -1.64) {\textbf{金线}：单机内 NVLink (TP)};
  \node[font=\tiny, color=ochre] at (-2.35, -1.64) {\textbf{橙箭}：跨节点点对点 (PP)};

  % Replica 1 hint arrow
  \draw[<->, very thick, dashed, navy] (-4.05, -1.92) -- (-4.05, -2.35)
    node[below, font=\scriptsize\bfseries, color=navy] {跨副本 Ring AllReduce / ReduceScatter（DP / ZeRO）};

  % Right: GPipe vs 1F1B Timeline (p=3 stages, m=4 micro-batches)
  \node[font=\sffamily\small\bfseries, color=navy] at (4.05, 3.35) {(b) GPipe vs 1F1B 流水线调度（$p=3, m=4$）};

  % GPipe sub-panel
  \node[anchor=west, font=\scriptsize\bfseries, color=ochre] at (0.3, 2.75) {GPipe：峰值激活 $O(m)=4$ 份};
  \foreach \st/\y in {0/2.25, 1/1.68, 2/1.11} {
    \node[anchor=east, font=\tiny\bfseries, color=slate] at (0.75, \y) {S\st};
  }
  % GPipe Stage 0: F1..F4 then bubble then B1..B4
  \foreach \k in {1,2,3,4} {
    \node[fblk] at (0.8+\k*0.64, 2.25) {F\k};
    \node[fblk] at (1.44+\k*0.64, 1.68) {F\k};
    \node[fblk] at (2.08+\k*0.64, 1.11) {F\k};
  }
  % GPipe Backward (tB = 1.5 tF -> width 9.4mm vs 6.2mm, step 0.96)
  \foreach \k in {1,2,3,4} {
    \node[bblk] at (4.52+\k*0.96, 1.11) {B\k};
    \node[bblk] at (5.48+\k*0.96, 1.68) {B\k};
    \node[bblk] at (6.44+\k*0.96, 2.25) {B\k};
  }
  \node[font=\tiny\bfseries, color=slate] at (4.7, 2.25) {气泡};

  % 1F1B sub-panel
  \node[anchor=west, font=\scriptsize\bfseries, color=tealDark] at (0.3, 0.35) {1F1B：气泡率相同，峰值激活骤降至 $O(p)=3$ 份！};
  \foreach \st/\y in {0/-0.18, 1/-0.75, 2/-1.32} {
    \node[anchor=east, font=\tiny\bfseries, color=slate] at (0.75, \y) {S\st};
  }
  % Stage 2 (bottom): F1 B1 F2 B2 F3 B3 F4 B4
  \node[fblk] at (2.72, -1.32) {F1};
  \node[bblk] at (3.52, -1.32) {B1};
  \node[fblk] at (4.32, -1.32) {F2};
  \node[bblk] at (5.12, -1.32) {B2};
  \node[fblk] at (5.92, -1.32) {F3};
  \node[bblk] at (6.72, -1.32) {B3};
  \node[fblk] at (7.52, -1.32) {F4};
  \node[bblk] at (8.32, -1.32) {B4};

  % Stage 1 (middle): F1 F2 B1 F3 B2 F4 B3 B4
  \node[fblk] at (2.08, -0.75) {F1};
  \node[fblk] at (2.72, -0.75) {F2};
  \node[bblk] at (4.48, -0.75) {B1};
  \node[fblk] at (5.28, -0.75) {F3};
  \node[bblk] at (6.08, -0.75) {B2};
  \node[fblk] at (6.88, -0.75) {F4};
  \node[bblk] at (7.68, -0.75) {B3};
  \node[bblk] at (9.28, -0.75) {B4};

  % Stage 0 (top): F1 F2 F3 B1 F4 B2 B3 B4
  \node[fblk] at (1.44, -0.18) {F1};
  \node[fblk] at (2.08, -0.18) {F2};
  \node[fblk] at (2.72, -0.18) {F3};
  \node[bblk] at (5.44, -0.18) {B1};
  \node[fblk] at (6.24, -0.18) {F4};
  \node[bblk] at (7.04, -0.18) {B2};
  \node[bblk] at (8.64, -0.18) {B3};
  \node[bblk] at (10.24, -0.18) {B4};

  % Annotation note for 1F1B early release
  \node[font=\tiny\bfseries, color=tealDark, align=center] at (5.44, -2.02) {稳态 1F1B：S0 在 B1 完成后立刻释放 F1 激活，再注入 F4，在飞份数恒定 $\le p=3$};
\end{tikzpicture}%
}
\caption{3D 并行拓扑映射与流水线调度对比。(a) 将通信密集的 TP 锁在单机 NVLink 内，较轻量的 PP 跨越节点，顶层以 ZeRO/DP 水平复制；(b) 1F1B 在保持与 GPipe 相同气泡率 $\frac{p-1}{m+p-1}$ 的前提下，提前完成早期微批次的反向传播并立刻释放激活显存。}
\label{fig:3d-parallel-and-1f1b}
\end{figure}

\section{参数高效微调（LoRA）：低秩适配与零延迟合并}
\label{sec:lora}

3D 并行解决了从零预训练的算力与显存难题，但绝大多数研究者与企业面对的是另一个日常场景：手头已经有一座开源的 7B 或 70B 预训练/对齐大模型，只想在特定领域数据上做 SFT 或 DPO 后训练。如果做\textbf{全参数微调（Full Fine-Tuning）}，就必须承受全额 $16\Phi$ 字节的静态显存（7B 要 112\,GB！）且每个下游任务都得存一份完整的数十 GB 模型权重副本。

\subsection{内在维度假说、低秩分解与 $B = 0$ 初始化铁律}

Aghajanyan 等人的实证研究表明，预训练大模型在适应特定下游任务时，其有效参数更新位于一个维度极低的\textbf{内在子空间（Intrinsic Subspace）}之中。受此启发，Edward Hu 等人（2021）提出了 \textbf{LoRA（Low-Rank Adaptation，低秩适配）}：完全冻结庞大的预训练权重矩阵 $W_0 \in \R^{d \times k}$，仅在其旁路注入一对可训练的低秩分解矩阵 $B \in \R^{d \times r}$ 与 $A \in \R^{r \times k}$（其中秩 $r \ll \min(d, k)$），将前向变换写为：
\begin{equation}
  h \;=\; W_0 x \;+\; \Delta W\, x \;=\; W_0 x \;+\; \frac{\alpha}{r}\, B A x,
  \label{eq:lora-forward}
\end{equation}
其中常数 $\alpha$ 为固定缩放超参数（使 $\alpha / r$ 稳定不同秩 $r$ 下的梯度量级，调整 $r$ 时无需剧烈重调学习率）。

LoRA 之所以成为大模型后训练的绝对标配，归功于三条环环相扣的数学与工程优势：
\begin{enumerate}
  \item \textbf{零扰动初始化铁律（$A \sim \mathcal{N}(0, \sigma^2),\; B = \mathbf{0}$）}：训练第 0 步时，$\Delta W = \frac{\alpha}{r} B A = \mathbf{0}$，模型输出与原始预训练模型\textbf{逐比特完全一致}，丝毫不破坏底座已学得的知识！同时，记损失对输出 $h$ 的梯度为 $G = \nabla_h \mathcal{L}$，由链式法则：
  \begin{equation}
    \nabla_B \mathcal{L} \;=\; \frac{\alpha}{r}\, G (A x)^\top \;\neq\; \mathbf{0}, \qquad
    \nabla_A \mathcal{L} \;=\; \frac{\alpha}{r}\, B^\top G x^\top \Big|_{t=0} \;=\; \mathbf{0}.
    \label{eq:lora-grad}
  \end{equation}
  第 1 步更新立即使 $B$ 获得非零梯度而脱离零点，第 2 步起 $A$ 也获得非零梯度，完美打破对称性！（若错把 $A$ 与 $B$ 同时置零，则 $\nabla_A \mathcal{L} \equiv \mathbf{0}$ 且 $\nabla_B \mathcal{L} \equiv \mathbf{0}$，权重将永远锁死在零点，见习题~\textbf{6.6.5}！）
  \item \textbf{惊人的参数与显存压缩比}：单层全参数微调需训练 $dk$ 个参数，而 LoRA 仅需训练 $r(d + k)$ 个参数，精确参数压缩比为
  \begin{equation}
    \rho_{\mathrm{LoRA}} \;=\; \frac{r(d + k)}{d k}.
    \label{eq:lora-ratio}
  \end{equation}
  对于方阵 $d = k = 4096$、典型秩 $r = 16$，可训练参数占比仅为 $\rho = \frac{2 \times 16}{4096} = \mathbf{0.78\%}$！更重要的是，冻结的底座权重 $W_0$ 无需任何梯度、无需 FP32 主权重、也无需 AdamW 的 $m, v$ 状态，整网静态显存从 $16\Phi$ 断崖式暴跌至：
  \[
    M_{\mathrm{LoRA}} \;=\; 2\Phi + 14\,\Phi_{\mathrm{LoRA}} \;\approx\; 2.1\Phi\;\text{Bytes}
  \]
  （7B 模型从 $112\,\mathrm{GB}$ 直降至约 $15\,\mathrm{GB}$！若进一步将底座 $W_0$ 量化为 4-bit NormalFloat，即 \textbf{QLoRA}，7B 静态显存更是压至仅 $\sim 4\,\mathrm{GB}$，普通单张消费级显卡即可微调 65B 模型！）。
  \item \textbf{零推理延迟开销（Weight Merging）}：训练完毕部署上线时，只需离线做一次矩阵加和合并权重：
  \begin{equation}
    W_{\mathrm{merged}} \;=\; W_0 \;+\; \frac{\alpha}{r}\, B A.
    \label{eq:lora-merge}
  \end{equation}
  线上推理图与原始底座结构完全相同，\textbf{增加的推理延迟严格为零}！若需服务多租户任务，甚至只需在显存中常驻一份共享底座 $W_0$，毫秒级热切换几兆字节的 $(B_k, A_k)$ 适配器即可。
\end{enumerate}

\subsection{实测验证：LoRA 秩扫描与内在维度的 $100\times$ 误差悬崖}

在配套脚本 \texttt{code/vol6\_experiments.py} 的 \texttt{exp6\_quant\_and\_lora()} 实验中，我们在一个冻结的 $64\times 64$ 预训练线性层 $W_0$ 上构造了“真秩为 4 + 微弱全秩噪声（$\sigma = 0.03$）”的下游任务偏移 $W_{\mathrm{target}} = W_0 + U_t V_t + 0.03\varepsilon$，并依次扫描 LoRA 秩 $r \in \{1, 2, 4, 8, 16\}$ 与全参数微调（等价于 $r = 64$），实测数据保存于 \texttt{data/lora\_rank.csv}。

\begin{table}[htbp]
\centering
\caption{冻结 $64\times 64$ 线性层适配秩-4 任务偏移的 LoRA 实测结果（来源：\texttt{data/lora\_rank.csv}）。}
\label{tab:lora-rank}
\small
\begin{tabular}{@{}rrrrl@{}}
\toprule
\textbf{适配器秩 $r$} & \textbf{可训练参数量} & \textbf{参数占比} & \textbf{测试集 MSE} & \textbf{物理机制解读} \\
\midrule
$1$ & $128$ & $3.12\%$ & $4.96746$ & 欠拟合：表达秩严重低于任务内在秩 4 \\
$2$ & $256$ & $6.25\%$ & $2.89887$ & 欠拟合：仍缺失一半主奇异方向 \\
$\mathbf{4}$ & $\mathbf{512}$ & $\mathbf{12.50\%}$ & $\mathbf{0.04956}$ & \textbf{相变拐点：误差暴跌 $100\times$！精准捕获主子空间} \\
$8$ & $1024$ & $25.00\%$ & $0.04013$ & 收益递减：开始缓慢拟合微弱全秩背景噪声 \\
$16$ & $2048$ & $50.00\%$ & $0.02630$ & 收益递减：参数翻倍而误差仅微降 \\
$64$（全微调） & $4096$ & $100.00\%$ & $0.00000$ & 完整拟合全部 4096 个自由度（含背景噪声） \\
\bottomrule
\end{tabular}
\end{table}

图~\ref{fig:lora-rank-exp} 直观绘出了 \texttt{data/lora\_rank.csv} 中的 MSE 悬崖：仅需 $12.50\%$ 参数量的 $r = 4$ 就将测试误差从 $4.96746$ 瞬间劈落至 $0.04956$（下降整整两个数量级！），再往上加倍秩数至 $r = 8$（$25.00\%$ 参数，MSE $0.04013$）边际收益迅速饱和。这正是实践中对 7B/70B 模型选取 $r \in \{8, 16, 32\}$ 即可逼近全微调效果的数学本质。

\begin{figure}[htbp]
\centering
\begin{tikzpicture}
  \begin{axis}[
    width=0.78\linewidth,
    height=5.6cm,
    xmode=log,
    log basis x=2,
    xmin=0.8, xmax=85,
    ymin=-0.35, ymax=5.6,
    xlabel={LoRA 适配秩 $r$（$r=64$ 表示全参数微调）},
    ylabel={测试集均方误差（Test MSE）},
    title={\sffamily\bfseries\small LoRA 秩扫描实验：匹配任务内在维度（$r=4$）时的 $100\times$ 误差悬崖},
    grid=both,
    grid style={line width=0.25pt, draw=slate!20},
    major grid style={line width=0.4pt, draw=slate!35},
    tick label style={font=\scriptsize},
    label style={font=\footnotesize\sffamily},
    xtick={1,2,4,8,16,64},
    xticklabels={1 (3.1\%), 2 (6.3\%), 4 (12.5\%), 8 (25.0\%), 16 (50.0\%), 64 (100\%)},
    legend style={at={(0.97,0.95)}, anchor=north east, font=\scriptsize\sffamily, fill=white, fill opacity=0.9, draw=tealLine}
  ]
    \addplot[teal, very thick, mark=*, mark size=2.5pt, mark options={fill=tealLight, draw=tealDark, thick}]
      table[col sep=comma, x=rank, y=test_mse] {data/lora_rank.csv};
    \addlegendentry{实测 Test MSE (\texttt{lora\_rank.csv})}

    % Highlight intrinsic rank r=4
    \addplot[ochre, dashed, thick] coordinates {(4, -0.35) (4, 5.4)};
    \node[font=\scriptsize\sffamily\bfseries, color=ochre, anchor=west] at (axis cs:4.4, 3.6) {内在秩拐点 $r = 4$：};
    \node[font=\scriptsize\sffamily, color=ochre, anchor=west] at (axis cs:4.4, 2.9) {仅 $12.5\%$ 参数，MSE 从 $4.97 \to 0.050$（降 $100\times$）};
    \node[font=\tiny\sffamily, color=navy, anchor=south west] at (axis cs:1.05, 5.02) {$4.967$};
    \node[font=\tiny\sffamily, color=tealDark, anchor=south west] at (axis cs:4.1, 0.35) {$0.050$};
    \node[font=\tiny\sffamily, color=tealDark, anchor=south east] at (axis cs:62, 0.25) {$0.000$};
  \end{axis}
\end{tikzpicture}
\caption{基于 \texttt{data/lora\_rank.csv} 绘制的 LoRA 秩扫描测试误差曲线（受控合成设定：预先构造了主秩为 $4$ 的低秩任务偏移 $U_t V_t$ 叠加微弱全秩扰动，旨在直观展示当任务偏移近似低秩时 LoRA 为何奏效；真实大模型下游任务的内在有效秩则需视具体领域而定）。横轴括号内标出了各秩对应的可训练参数占比；一旦适配器秩触及内在秩 $r = 4$，泛化误差立刻坠落两个数量级。}
\label{fig:lora-rank-exp}
\end{figure}

代码~\ref{lst:lora-code} 给出了极简且可直接运行的 PyTorch LoRA 线性层与离线权重合并实现。

\begin{lstlisting}[caption={支持 $B=0$ 初始化与离线零延迟合并的 PyTorch LoRA 线性层核心实现。},label={lst:lora-code}]
import math, torch, torch.nn as nn, torch.nn.functional as F

class LoRALinear(nn.Module):
    def __init__(self, base_linear: nn.Linear, r: int = 8, alpha: float = 16.0):
        super().__init__()
        self.base = base_linear
        self.base.weight.requires_grad_(False)            # 冻结预训练底座
        if self.base.bias is not None:
            self.base.bias.requires_grad_(False)
        d_out, d_in = base_linear.weight.shape
        self.scale = alpha / r
        self.A = nn.Parameter(torch.empty(r, d_in))
        self.B = nn.Parameter(torch.zeros(d_out, r))      # B 严格零初始化!
        nn.init.kaiming_uniform_(self.A, a=math.sqrt(5))  # A 随机初始化

    def forward(self, x: torch.Tensor) -> torch.Tensor:
        return self.base(x) + (x @ self.A.T @ self.B.T) * self.scale

    @torch.no_grad()
    def merge_to_base(self) -> nn.Linear:
        """推理部署前离线合并，线上零额外延迟。"""
        self.base.weight.add_((self.B @ self.A) * self.scale)
        return self.base
\end{lstlisting}

\section{低比特后训练量化（PTQ）：INT8/INT4 数学分析与离群值难题}
\label{sec:quantization}

回到第~\ref{ch:hardware-roofline-flash} 章的 Roofline 定理~\ref{thm:prefill-vs-decode}：大模型在 Decode 阶段之所以只能跑出不到 $1\%$ 的 GPU 算力，罪魁祸首是每生成一个词元都必须把整个模型的 BF16 权重（$2\Phi$ 字节）从 HBM 搬进片上一次。如果我们能把权重从 16 比特压缩到 \textbf{8 比特（INT8）甚至 4 比特（INT4）}，便一箭双雕：
\begin{itemize}
  \item 显存占用直接腰斩 $2\times$ 或骤降 $4\times$（70B 模型从 $140\,\mathrm{GB}$ 压至 $35\,\mathrm{GB}$，单卡即可装下！）；
  \item Decode 阶段每步搬运的 HBM 字节数同步减少 $2\text{--}4\times$，吞吐率几乎直接线性翻倍！
\end{itemize}
这正是\textbf{后训练量化（Post-Training Quantization, PTQ）}的核心价值：无需昂贵的从头重训，直接将训练好的浮点权重映射为低比特整数网格。

\subsection{对称均匀量化与重尾离群通道灾难定理}

\begin{defBox}
\begin{definition}[$b$ 比特对称均匀量化（Symmetric Uniform Quantization）]
\label{def:sym-quant}
给定实数张量（或子块）$W$ 与目标有符号整数位宽 $b \in \{4, 8\}$，令最大正整数刻度为 $q_{\max} = 2^{b-1} - 1$（例如 INT8 对应 $q_{\max} = 127$，INT4 对应 $q_{\max} = 7$）。对称均匀量化的缩放因子 $s$、定点整数编码 $Q$ 与反量化重构值 $\hat{W}$ 定义为：
\begin{equation}
  s \;=\; \frac{\max |W|}{q_{\max}}, \qquad
  Q \;=\; \clip\!\left(\operatorname{round}\!\left(\frac{W}{s}\right),\, -q_{\max},\, q_{\max}\right), \qquad
  \hat{W} \;=\; s \cdot Q.
  \label{eq:sym-quant}
\end{equation}
\end{definition}
\end{defBox}

如果权重的所有元素都像标准高斯分布那样乖巧地集中在同一个尺度上，全矩阵共享一个全局缩放因子 $s$（称为\textbf{逐张量量化，Per-Tensor}）效果还算体面。但真实大语言模型的权重和激活却存在一个惊人的结构特征——\textbf{系统性通道离群值（Systematic Channel Outliers）}：在全部几千个隐藏通道中，极少数特定通道的数值幅度会比其余 $99\%$ 正常通道大出 $8\times$ 乃至 $100\times$！下面这条定理从数学上精准刻画了离群值对逐张量量化的毁灭性打击，以及逐通道/分组量化如何力挽狂澜。

\begin{thmBox}
\begin{theorem}[重尾离群通道对逐张量量化的破坏效应与逐通道误差隔离]
\label{thm:quant-outlier-error}
设权重矩阵 $W \in \R^{m \times n}$ 包含 $m - k$ 个正常输出通道与 $k \ll m$ 个离群输出通道：正常通道元素幅值上界为 $M_0$，离群通道元素幅值上界高达 $M_1 = \kappa M_0$（离群倍数 $\kappa \gg 1$）。假设在步长 $s$ 下，未截断元素的量化舍入误差 $e = w - \hat{w}$ 近似服从 $[-s/2,\, s/2]$ 上的均匀分布（期望均方误差 $\E[e^2] = s^2 / 12$），则：
\begin{enumerate}
  \item \textbf{逐张量量化（Per-Tensor）}被迫采用全局最大步长 $s_{\mathrm{pt}} = \frac{\kappa M_0}{q_{\max}}$，导致全部 $m \times n$ 个元素的单元素期望均方误差被离群通道放大 $\kappa^2$ 倍：
  \begin{equation}
    \mathrm{MSE}_{\mathrm{pt}} \;=\; \frac{\kappa^2 M_0^2}{12\, q_{\max}^2}.
    \label{eq:mse-pt}
  \end{equation}
  特别地，若 $\kappa > 2 q_{\max}$（对 INT4 即 $\kappa > 14$），正常通道中全部元素满足 $|w| < s_{\mathrm{pt}}/2$，经 $\operatorname{round}(w / s_{\mathrm{pt}})$ 后将\textbf{全体坍缩为零}（分辨率彻底归零）！
  \item \textbf{逐通道量化（Per-Channel）}为每一行 $i \in \{1, \dots, m\}$ 分配独立缩放因子 $s_i = \frac{\max_j |w_{ij}|}{q_{\max}}$，将离群步长严格隔离在 $k$ 个离群通道内部，使全矩阵平均均方量化误差骤降至
  \begin{equation}
    \mathrm{MSE}_{\mathrm{pc}} \;=\; \left(\frac{m - k}{m} + \frac{k}{m}\,\kappa^2\right) \frac{M_0^2}{12\, q_{\max}^2},
    \label{eq:mse-pc}
  \end{equation}
  当 $k / m \ll 1$ 且 $\kappa \gg 1$ 时，相对误差比值满足 $\sqrt{\mathrm{MSE}_{\mathrm{pc}} / \mathrm{MSE}_{\mathrm{pt}}} \approx \sqrt{\frac{k}{m} + \frac{1}{\kappa^2}} \ll 1$！同时，存储 $m$ 个 FP16 缩放因子带来的额外空间开销占比仅为 $\frac{16}{b \cdot n} = O(1/n)$。
\end{enumerate}
\end{theorem}
\end{thmBox}

\begin{proof}
（1）对任意位于 $[-s/2, s/2]$ 上均匀分布的随机变量 $E$，其二阶矩为 $\E[E^2] = \frac{1}{s}\int_{-s/2}^{s/2} e^2\,de = \frac{s^2}{12}$。在逐张量量化中，全矩阵共用同一缩放因子 $s_{\mathrm{pt}} = \frac{\max_{i,j}|w_{ij}|}{q_{\max}} = \frac{\kappa M_0}{q_{\max}}$，代入即得式~\eqref{eq:mse-pt}。当 $\kappa > 2 q_{\max}$ 时，对所有正常通道元素均有 $\frac{|w|}{s_{\mathrm{pt}}} \le \frac{M_0}{\kappa M_0 / q_{\max}} = \frac{q_{\max}}{\kappa} < \frac{1}{2}$，故四舍五入后 $Q_{ij} = \operatorname{round}(w_{ij}/s_{\mathrm{pt}}) \equiv 0$，正常通道信息全毁。

（2）在逐通道量化中，占比 $\frac{m-k}{m}$ 的正常通道使用精细步长 $s_0 = \frac{M_0}{q_{\max}}$（方差 $\frac{M_0^2}{12 q_{\max}^2}$），仅有占比 $\frac{k}{m}$ 的离群通道使用粗步长 $s_1 = \frac{\kappa M_0}{q_{\max}}$（方差 $\frac{\kappa^2 M_0^2}{12 q_{\max}^2}$），按通道占比加权平均立得式~\eqref{eq:mse-pc}。对于存储开销，原始 $b$-bit 权重每行占 $b n$ 比特，新增 1 个 FP16 缩放因子仅占 16 比特，额外比例为 $\frac{16}{bn}$（例如 $n = 4096, b = 4$ 时仅增加 $0.098\%$ 显存！）。若进一步按每组 $g = 128$ 个元素分配独立缩放因子（\textbf{分组量化，Group-Wise Quantization}），则能连同一通道内部的局部尖峰也进一步隔离，额外开销也仅约 $16/(4\times 128) = 3.1\%$。证毕。
\end{proof}

\begin{insightBox}{直通现代前沿量化算法：SmoothQuant、AWQ 与 GPTQ 的核心直觉}
定理~\ref{thm:quant-outlier-error} 完美解释了\textbf{仅量化权重（Weight-Only，如 W4A16）}为何只需分组量化即可大获全胜；但如果要把\textbf{激活值 $X$ 也量化为 INT8（W8A8）}以启用 INT8 Tensor Core 算力翻倍，问题立刻变得棘手——激活离群值出现在固定列维度上，按行（逐 Token）量化激活无法隔离列离群值！现代三大经典算法分别从不同维度破局：
\begin{itemize}
  \item \textbf{SmoothQuant（尺度平滑迁移）}：利用矩阵乘法对角等价恒等式
  \[
    Y \;=\; X W \;=\; \big(X \diag(s)^{-1}\big)\,\big(\diag(s) W\big),
  \]
  将激活 $X$ 在离群通道上的巨大尺度通过 $\diag(s)^{-1}$ 压平，并把对应的放大系数 $\diag(s)$ 优雅地“转嫁”给容易逐通道量化的权重矩阵 $W$！
  \item \textbf{AWQ（激活感知权重量化）}：发现仅仅保护对\textbf{大激活通道}起决定作用的顶层 $1\%$ 显著权重（通过等价尺度放大这 $1\%$ 权重再做 W4A16 量化），就能避免极低的 INT4 精度塌方，且无需反向传播微调。
  \item \textbf{GPTQ（二阶海森误差补偿）}：逐列量化权重 $W$ 时，利用校准集输入协方差近似海森矩阵 $H = 2 X^\top X$，每把一列权重舍入产生量化误差 $\delta w_q$，就立刻用 $H^{-1}$ 闭式调整尚未量化的剩余列权重来精确抵消输出误差 $\|X W - X \hat{W}\|_F^2$！
\end{itemize}
\end{insightBox}

\subsection{实测验证：离群矩阵上的 INT8 与 INT4 逐张量 vs 逐通道对比}

在 \texttt{code/vol6\_experiments.py} 的 \texttt{exp6\_quant\_and\_lora()} 实验中，我们构造了一个 $256 \times 256$ 的真实重尾权重矩阵 $W \sim \mathcal{N}(0, 0.25^2)$，并将其前 4 个输出通道（仅占 $1.56\%$ 的通道！）放大 $8\times$ 模拟离群通道，在随机输入 $X \in \R^{256\times 64}$ 上实测矩阵乘法输出 $Y = W X$ 的相对弗罗贝尼乌斯误差 $\|W_{\mathrm{q}} X - Y_{\mathrm{fp}}\|_F / \|Y_{\mathrm{fp}}\|_F$。代码~\ref{lst:quant-code} 与表~\ref{tab:quant-exp} 展示了触目惊心的实测对比。

\begin{lstlisting}[caption={对称均匀量化（逐张量 vs 逐通道）在重尾离群权重矩阵上的实测核验脚本。},label={lst:quant-code}]
import numpy as np

rng = np.random.default_rng(2026)
W = rng.normal(0, 0.25, size=(256, 256))
W[0:4, :] *= 8.0                             # 注入 4 个 8x 离群通道
X = rng.normal(0, 1.0, size=(256, 64))
Y_fp = W @ X

def quant_dequant(M: np.ndarray, bits: int, per_channel: bool = False):
    qmax = (1 << (bits - 1)) - 1             # INT8: 127, INT4: 7
    axis = 1 if per_channel else None        # 逐通道沿输出行独立求尺度
    scale = np.maximum(np.max(np.abs(M), axis=axis, keepdims=True) / qmax, 1e-8)
    Q = np.clip(np.round(M / scale), -qmax, qmax)
    return Q * scale

for bits in [8, 4]:
    rel_pt = np.linalg.norm(quant_dequant(W, bits, False) @ X - Y_fp) / np.linalg.norm(Y_fp)
    rel_pc = np.linalg.norm(quant_dequant(W, bits, True)  @ X - Y_fp) / np.linalg.norm(Y_fp)
    print(f"INT{bits}: per-tensor={rel_pt*100:.3f}%, per-channel={rel_pc*100:.3f}%")
\end{lstlisting}

\begin{table}[htbp]
\centering
\caption{含 4 个 $8\times$ 离群通道的 $256\times 256$ 权重矩阵实测输出相对误差（来源：\texttt{exp6\_quant\_and\_lora()}）。}
\label{tab:quant-exp}
\small
\begin{tabular}{@{}lcrrcl@{}}
\toprule
\textbf{量化位宽} & \textbf{整数刻度 $q_{\max}$} & \textbf{逐张量误差} & \textbf{逐通道误差} & \textbf{误差缩减倍数} & \textbf{结果诊断} \\
\midrule
\textbf{INT8} & $127$ & $4.603\%$ & $\mathbf{0.658\%}$ & $\mathbf{7.0\times}$ & 逐通道几乎无损复原浮点输出 \\
\textbf{INT4} & $7$ & $63.826\%$ & $\mathbf{12.138\%}$ & $\mathbf{5.3\times}$ & 逐张量彻底崩溃；逐通道成功挽救 \\
\bottomrule
\end{tabular}
\end{table}

实测结果与定理~\ref{thm:quant-outlier-error} 严丝合缝：仅仅 4 个离群通道（占总通道不足 $1.6\%$），就让 INT4 逐张量量化的输出误差飙升到灾难性的 $\mathbf{63.826\%}$（绝大多数正常权重被压扁归零，模型输出彻底胡言乱语）；仅仅改为逐通道分配缩放因子（每行多存 1 个标量），INT8 误差便暴降 $7.0\times$ 至 $\mathbf{0.658\%}$，INT4 误差也从彻底崩溃的 $63.826\%$ 骤降至 $\mathbf{12.138\%}$（若再进一步配合分组量化或 AWQ，INT4 误差还能继续压至 $1\text{--}2\%$ 以内，见习题~\textbf{6.6.7}！）。

\section{高吞吐大模型推理系统：连续批处理、PagedAttention 与投机解码}
\label{sec:llm-serving}

现在，我们已经把权重切好、微调好、量化好，终于来到了生产环境的最后一公里——\textbf{大模型在线推理服务引擎（如 vLLM、SGLang、TensorRT-LLM）}。在生产环境中，系统每天面对成千上万个提示长度各异、生成长度不可预测的并发请求，如何把 H100 的吞吐率压榨到极致？现代推理引擎建立在三大核心技术支柱之上。

\subsection{KV 缓存显存核算与连续批处理（Continuous Batching）}

自回归解码为了避免每吐一个新词都重算前面所有词元的键向量 $K$ 与值向量 $V$，必须把每个请求所有层的历史 $(K, V)$ 缓存在 GPU 显存中，称为 \textbf{KV Cache}。单条请求每增加 1 个上下文词元所消耗的精确 KV 缓存字节数为：
\begin{equation}
  M_{\mathrm{KV/tok}} \;=\; \underbrace{2}_{K\text{ 与 }V} \cdot L \cdot n_{\mathrm{kv}} \cdot d_{\mathrm{head}} \cdot b_{\mathrm{dtype}}\;\text{Bytes},
  \label{eq:kv-cache-formula}
\end{equation}
其中 $L$ 为层数，$n_{\mathrm{kv}}$ 为 KV 头数（标准 MHA 时 $n_{\mathrm{kv}} = h$；采用分组查询注意力 GQA 时大幅缩小，如 Llama-3 70B 取 $L = 80, n_{\mathrm{kv}} = 8, d_{\mathrm{head}} = 128$），$b_{\mathrm{dtype}}$ 为精度字节数（BF16 为 2）。代入一个 70B 标准 MHA 模型（$L=80, n_{\mathrm{kv}}=64, d_{\mathrm{head}}=128$, BF16），单词元 KV 缓存高达 $2.62\,\mathrm{MB}$——如果同时服务 $B = 32$ 个长度 $S = 4096$ 的对话，仅 KV 缓存就要吞掉惊人的 $\mathbf{343\,\mathrm{GB}}$ 显存（甚至超过模型权重本身两倍以上）！即使启用 8-组 GQA（$n_{\mathrm{kv}} = 8$），仍需 $43\,\mathrm{GB}$。

而在请求调度层面，传统深度学习框架采用\textbf{静态批处理（Static Batching）}：把 $B$ 条请求打包成一个批次送入 GPU，必须死等批次中\textbf{生成最长那条请求}吐出结束符 \texttt{<EOS>}，整批才能释放并接入下一批请求！短回复提前算完后，它占据的显存与计算槽位全部白白空转。

Yu 等人（Orca, 2022）提出的\textbf{连续批处理（Continuous Batching，亦称迭代级调度 Iteration-Level Scheduling）}彻底打破了请求级屏障：调度器不再以“整条对话结束”为周期，而是以\textbf{“每生成一个词元的单个 Decode 迭代步”}为粒度——一旦某条请求吐出 \texttt{<EOS>}，立刻在下一毫秒的迭代步中将其踢出批次、回收 KV 缓存，并无缝把等待队列中的新请求插入该空槽！仅仅这一项调度革新，就能把真实线上推理吞吐率提升 $5\text{--}10\times$！

\subsection{PagedAttention（vLLM）：408 操作系统虚拟内存分页的完美重现}

连续批处理把计算槽位填满了，但 KV 缓存的\textbf{显存管理灾难}却骤然爆发。因为系统根本无法预知用户请求最终会生成 20 个词还是 2048 个词，传统推理引擎（如早期 FasterTransformer）只能按最大上下文长度 $S_{\max}$ 为每条请求预先分配一块\textbf{物理连续的巨型显存槽}。实测表明，这种静态连续预分配导致了高达 $\mathbf{60\%\text{--}80\%}$ 的显存浪费——来源包括三类经典碎片：
\begin{enumerate}
  \item \textbf{预留浪费（Reserved）}：给将来的词元提前圈地，但请求往往早早结束根本用不到；
  \item \textbf{内部碎片（Internal Fragmentation）}：实际生成长度远小于 $S_{\max}$ 留下的尾部空洞；
  \item \textbf{外部碎片（External Fragmentation）}：不同长度请求的连续块释放后留下的零散空隙，拼不出新的大连续块。
\end{enumerate}

\begin{recallBox}{408 操作系统速忆：虚拟内存分页存储管理与写时复制（Copy-on-Write）}
四十年前，操作系统设计师面对多进程动态内存膨胀与外部碎片问题，发明了计算机科学史上最优雅的抽象之一——\textbf{分页虚拟内存（Virtual Memory Paging）}：进程的逻辑地址空间被切分为固定大小的页（Page），通过\textbf{页表（Page Table）}映射到物理内存中任意离散不连续的页框（Frame），外部碎片瞬间归零，内部碎片被严格封顶在每进程最后一页的不到 1 页之内！当父进程 \texttt{fork()} 子进程时，二者先共享同一批只读物理页框，仅在某一方尝试写入修改时才触发\textbf{写时复制（Copy-on-Write, CoW）}。
\end{recallBox}

2023 年，加州大学伯克利分校的 Kwon 等人把这套经典的 408 操作系统分页机制原封不动地搬进了 GPU HBM，诞生了震撼工业界的 \textbf{PagedAttention（vLLM）}（如图~\ref{fig:paged-attention} 所示）：
\begin{itemize}
  \item \textbf{逻辑 KV 块与离散物理块}：将每条请求的 KV 缓存切分为固定词元数的\textbf{逻辑 KV 块}（例如每块存 $B_{\mathrm{blk}} = 16$ 个词元的 K/V 向量）；在 GPU HBM 中预先划分出一个全局\textbf{物理 KV 块池}，按需动态分配，彻底允许同一请求的相邻逻辑块存放在\textbf{物理上完全不连续}的 HBM 地址中！
  \item \textbf{每请求块表（Block Table）}：维护一张极轻量的映射表，记录“逻辑块号 $\to$ 物理块号 + 当前已填词元数”。计算注意力时，定制的 PagedAttention CUDA 核函数通过块表间接寻址，分块聚合离散物理块中的 $K, V$ 向量完成点积与在线 Softmax。
  \item \textbf{碎片率从 $>60\%$ 骤降至 $<4\%$}：物理外部碎片彻底消失；唯一的内部碎片仅发生在每条请求的\textbf{最后一个未填满的块}里，平均每条请求仅浪费半个块（$B_{\mathrm{blk}}/2 = 8$ 个词元），在常规对话长度下总显存浪费被死死压在 $\mathbf{4\%}$ 以下——腾出来的显存允许并发批大小直接暴增 $2\text{--}4\times$，顺势借 Roofline 定理大幅推高 GPU 利用率！
  \item \textbf{写时复制（Copy-on-Write）实现零拷贝前缀共享与并行采样}：多条请求共享同一段长系统提示（Prefix Caching），或第 4 章中 Best-of-$N$、Beam Search、GRPO 针对同一提示并行采样 $G$ 条回复时，所有分支只需将各自的块表指向\textbf{同一组物理前缀块}（引用计数 $\text{ref\_cnt} = G$），显存零冗余！直到各分支分化生成不同新词元触及末尾块写入时，才对那唯一一块执行 Copy-on-Write 分配新物理块。
\end{itemize}

\begin{figure}[htbp]
\centering
\resizebox{0.97\linewidth}{!}{%
\begin{tikzpicture}[
  font=\sffamily\footnotesize,
  lblk/.style={rectangle, rounded corners=0.9mm, draw=navy, thick, fill=navyLight, minimum width=18mm, minimum height=7.2mm, align=center, font=\scriptsize, text=ink},
  sblk/.style={rectangle, rounded corners=0.9mm, draw=gold, very thick, fill=ochreLight, minimum width=18mm, minimum height=7.2mm, align=center, font=\scriptsize\bfseries, text=ink},
  pblk/.style={rectangle, rounded corners=0.9mm, draw=teal, thick, fill=tealLight, minimum width=22mm, minimum height=7.4mm, align=center, font=\scriptsize, text=ink},
  fblk/.style={rectangle, rounded corners=0.9mm, draw=slate!60, dashed, fill=codebg, minimum width=22mm, minimum height=7.4mm, align=center, font=\scriptsize, text=slate}
]
  % Column titles
  \node[font=\sffamily\small\bfseries, color=navy] at (-4.9, 3.2) {逻辑 KV 块视图（连续序列）};
  \node[font=\sffamily\small\bfseries, color=slate] at (0.05, 3.2) {块表（Block Table）};
  \node[font=\sffamily\small\bfseries, color=tealDark] at (4.7, 3.2) {GPU HBM 物理 KV 块池（离散非连续）};

  % Request A Logical Blocks
  \node[anchor=west, font=\scriptsize\bfseries, color=tealDark] at (-7.55, 2.48) {请求 A（共享系统提示 + 续写 A）};
  \node[sblk] (a0) at (-6.60, 1.80) {逻辑块 0\\{\tiny 系统提示}};
  \node[lblk] (a1) at (-4.60, 1.80) {逻辑块 1\\{\tiny 回答 A1}};
  \node[rectangle, rounded corners=0.9mm, draw=navy, thick, fill=white, minimum width=17mm, minimum height=7.2mm, font=\scriptsize, align=center] (a2real) at (-2.60, 1.80) {逻辑块 2\\{\tiny 尾部 5/16}};

  % Request B Logical Blocks
  \node[anchor=west, font=\scriptsize\bfseries, color=plum] at (-7.55, 0.35) {请求 B（同前缀并发采样 / CoW）};
  \node[sblk] (b0) at (-6.60, -0.35) {逻辑块 0\\{\tiny 系统提示}};
  \node[lblk, draw=plum, fill=plumLight] (b1) at (-4.60, -0.35) {逻辑块 1\\{\tiny 回答 B1}};

  % Block Tables
  \node[rectangle, rounded corners=1mm, draw=tealDark, thick, fill=white, inner sep=3.5pt, align=left, font=\tiny\ttfamily] (tblA) at (0.05, 1.80) {
    \textbf{Req A 块表}\\
    L0 $\to$ P1 (shared)\\
    L1 $\to$ P4 (16/16)\\
    L2 $\to$ P0 (\;5/16)
  };
  \node[rectangle, rounded corners=1mm, draw=plum, thick, fill=white, inner sep=3.5pt, align=left, font=\tiny\ttfamily] (tblB) at (0.05, -0.35) {
    \textbf{Req B 块表}\\
    L0 $\to$ P1 (shared)\\
    L1 $\to$ P3 (11/16)
  };

  % Physical Pool in HBM
  \draw[tealLine, thick, rounded corners=1.8mm, fill=codebg!50] (2.65, -1.68) rectangle (6.75, 2.75);
  \node[pblk] (p0) at (4.70, 2.15) {\textbf{物理块 P0}（A-L2）};
  \node[sblk, minimum width=26mm] (p1) at (4.70, 1.25) {\textbf{物理块 P1}（ref=2 共享）};
  \node[fblk] (p2) at (4.70, 0.35) {物理块 P2（空闲）};
  \node[pblk, draw=plum, fill=plumLight] (p3) at (4.70, -0.55) {\textbf{物理块 P3}（B-L1, CoW）};
  \node[pblk] (p4) at (4.70, -1.30) {\textbf{物理块 P4}（A-L1）};

  % Mapping Arrows
  \draw[-{Stealth[length=1.8mm]}, thick, gold] (tblA.east) ++(0, 0.08) to[out=0, in=180] (p1.west);
  \draw[-{Stealth[length=1.8mm]}, thick, tealDark] (tblA.east) ++(0, -0.14) to[out=-10, in=180] (p4.west);
  \draw[-{Stealth[length=1.8mm]}, thick, navy] (tblA.east) ++(0, -0.34) to[out=12, in=180] (p0.west);

  \draw[-{Stealth[length=1.8mm]}, thick, gold] (tblB.east) ++(0, 0.06) to[out=15, in=180] (p1.west);
  \draw[-{Stealth[length=1.8mm]}, thick, plum] (tblB.east) ++(0, -0.18) to[out=0, in=180] (p3.west);

  % Bottom note
  \node[font=\scriptsize, color=tealDark, align=center] at (-0.2, -2.18) {外部碎片严格归零，内部碎片仅限各请求末块（浪费 $<4\%$）；共享前缀块 \textbf{P1} 通过写时复制零冗余复用！};
\end{tikzpicture}%
}
\caption{PagedAttention（vLLM）将 408 操作系统分页虚拟内存引入 GPU KV 缓存管理：连续逻辑块通过每请求块表映射到离散物理块，不仅消除外部碎片，更通过引用计数与写时复制（CoW）实现前缀与并行采样分支的零拷贝共享。}
\label{fig:paged-attention}
\end{figure}

\subsection{投机解码（Speculative Decoding）：用并行验证打破自回归访存墙}

PagedAttention 与连续批处理极大地抬升了多用户并发吞吐率，但如果某位 VIP 用户发起的是\textbf{单请求低延迟（Low-Latency）长文本生成}呢？由定理~\ref{thm:prefill-vs-decode} 可知，单请求自回归逐词生成受限于串行因果依赖，每步 $T = 1$ 的算术强度仅 $1.00\,\mathrm{FLOP/B}$，H100 有 $99.66\%$ 的算力都在睡大觉！

Leviathan 等人与 Chen 等人（2023）提出了一个堪称“空手套白狼”的绝妙思路——\textbf{投机解码（Speculative Decoding）}：
\begin{quote}
\emph{既然把 $70\mathrm{B}$ 大模型权重搬进片上一次，算 $T = 1$ 个词元和并行算 $T = \gamma + 1$（例如 $\gamma = 5$）个词元的墙钟耗时几乎一模一样（都在深度访存受限区，见表~\ref{tab:roofline-data}！），那为什么不先让一个极小、极快的“草稿模型 $q$”（如 0.5B 或 1B 参数）飞快地自回归草拟出 $\gamma$ 个候选词元 $x_1, \dots, x_\gamma$，再让目标大模型 $p$ 在\textbf{单次并行前向传播}里一口气把这 $\gamma + 1$ 个位置的条件概率分布全部验算出来呢？}
\end{quote}

最震撼人心的问题来了：草稿模型 $q$ 毕竟比大模型 $p$ 弱得多，用 $q$ 提议、$p$ 审核，难道不会破坏大模型原本的输出质量吗？下面这条短小精悍的定理将证明：\textbf{通过极其巧妙的拒绝采样（Rejection Sampling）机制，投机解码输出序列的概率分布与纯目标大模型 $p$ 逐词自回归采样的分布数学上完全一致、零质量损失！}

\begin{thmBox}
\begin{theorem}[投机解码的精确无损分布保持定理]
\label{thm:speculative-lossless}
在任意解码步，设目标大模型的真实下一词概率分布为 $p(x)$，小草稿模型的提议分布为 $q(x)$（定义于同一词表 $\mathcal{V}$ 上）。执行如下单步投机验证算法：
\begin{enumerate}
  \item 从草稿模型采样候选词元 $x \sim q(\cdot)$，并独立采样 $u \sim \mathrm{Uniform}(0, 1)$；
  \item 若 $u \le \min\!\left(1,\, \frac{p(x)}{q(x)}\right)$，则\textbf{接受}并输出 $X = x$；
  \item 否则\textbf{拒绝} $x$，并从如下归一化残差分布中重新采样输出 $X \sim p'(\cdot)$：
  \begin{equation}
    p'(x) \;\triangleq\; \frac{\max\!\big(0,\; p(x) - q(x)\big)}{\sum_{y \in \mathcal{V}} \max\!\big(0,\; p(y) - q(y)\big)}.
    \label{eq:spec-residual}
  \end{equation}
\end{enumerate}
则最终输出词元 $X$ 的边缘概率分布\textbf{严格恒等于目标分布 $p$}：对任意词元 $x \in \mathcal{V}$，均有
\begin{equation}
  \Prob(X = x) \;=\; p(x).
  \label{eq:spec-exact}
\end{equation}
此外，若在连续 $\gamma$ 个草稿位置上的平均接受率均为 $\alpha \in [0, 1)$（注意全部接受时大模型还会免费多吐出第 $\gamma + 1$ 个词元），则目标大模型单次并行验证前向所产出的期望词元数为
\begin{equation}
  \E[L] \;=\; \frac{1 - \alpha^{\gamma + 1}}{1 - \alpha}.
  \label{eq:spec-expected-len}
\end{equation}
\end{theorem}
\end{thmBox}

\begin{proof}
首先计算一步采样中“接受并输出特定词元 $x$”的概率：
\[
  \Prob(\text{accept} \wedge X = x)
  \;=\; q(x) \cdot \min\!\left(1,\, \frac{p(x)}{q(x)}\right)
  \;=\; \min\!\big(q(x),\, p(x)\big).
\]
对全体词表 $y \in \mathcal{V}$ 求和，立得总接受概率与总拒绝概率：
\[
  \Prob(\text{accept}) = \sum_{y \in \mathcal{V}} \min\!\big(q(y),\, p(y)\big), \qquad
  \Prob(\text{reject}) = 1 - \sum_{y \in \mathcal{V}} \min\!\big(q(y),\, p(y)\big).
\]
现在利用实数初等恒等式 $a - \min(a, b) = \max(0, a - b)$，注意到 $\sum_{y} p(y) = 1$，因此总拒绝概率恰好等于残差分布的分母：
\[
  \Prob(\text{reject})
  \;=\; \sum_{y \in \mathcal{V}} p(y) - \sum_{y \in \mathcal{V}} \min\!\big(q(y),\, p(y)\big)
  \;=\; \sum_{y \in \mathcal{V}} \max\!\big(0,\; p(y) - q(y)\big)!
\]
将其代入全概率公式，分母与 $\Prob(\text{reject})$ \textbf{一刀抵消}：
\begin{align*}
  \Prob(X = x)
  &= \Prob(\text{accept} \wedge X = x) + \Prob(\text{reject}) \cdot p'(x) \\
  &= \min\!\big(q(x),\, p(x)\big) + \Prob(\text{reject}) \cdot \frac{\max\!\big(0,\; p(x) - q(x)\big)}{\Prob(\text{reject})} \\
  &= \min\!\big(q(x),\, p(x)\big) + \max\!\big(0,\; p(x) - q(x)\big)
  \;=\; p(x)!
\end{align*}
（最后一步只需分情形核对：若 $p(x) \le q(x)$，左边为 $p(x) + 0 = p(x)$；若 $p(x) > q(x)$，左边为 $q(x) + (p(x) - q(x)) = p(x)$！）

至于期望生成词元数，记产出长度为 $L \in \{1, 2, \dots, \gamma + 1\}$：只要前 $k$ 个草稿词元全部接受（事件概率为 $\alpha^k$），本次验证就至少产出 $k + 1$ 个词元（因为第 $k+1$ 个位置即便被拒也会从残差分布重采 1 个正确词元，若全接受则大模型免费附赠第 $\gamma + 1$ 个词元）。故由非负整数随机变量的尾概率求和公式：
\[
  \E[L] \;=\; \sum_{k=0}^{\gamma} \Prob(L \ge k + 1) \;=\; \sum_{k=0}^{\gamma} \alpha^k \;=\; \frac{1 - \alpha^{\gamma + 1}}{1 - \alpha}.
\]
证毕。
\end{proof}

直觉上，式~\eqref{eq:spec-expected-len} 的威力极其惊人：在代码补全、摘要改写或格式化输出等常见任务中，小草稿模型与大模型的平均一致率往往高达 $\alpha \approx 0.8$；取草稿步数 $\gamma = 5$ 时，大模型仅做 1 次并行验证就能一口气无损产出 $\E[L] = \frac{1 - 0.8^6}{1 - 0.8} \approx \mathbf{3.69}$ 个词元！扣除极轻量的草稿模型耗时，单请求生成延迟直接提速 $2\text{--}3\times$，而输出质量与原版大模型在概率论意义下\textbf{零差异}！

\section{分布式与推理系统排错指南及前沿架构拓展}
\label{sec:sys-pitfalls-ext}

\begin{failBox}{分布式训练、LoRA 与线上推理的三大致命系统陷阱}
\begin{enumerate}
  \item \textbf{把张量并行（TP）错误地跨节点部署到慢速以太网/InfiniBand 上}：初学者在双机 16 卡上运行 70B 模型时，常图省事直接设置 \texttt{tp\_size=16}。结果每层 Transformer 的 4 次层内阻塞式 AllReduce 全部被迫跨过慢一个数量级的机间网卡，GPU 算力利用率（MFU）瞬间从 $48\%$ 惨跌至不足 $8\%$！正确架构永远是：单机 8 卡内部跑 \texttt{tp\_size=8}（走 $900\,\mathrm{GB/s}$ NVLink），跨机方向只走 \texttt{pp\_size=2} 或 \texttt{ZeRO/FSDP}。
  \item \textbf{LoRA 忘记 $B = 0$ 初始化（或误将 $A, B$ 同时置零）}：如果按 PyTorch 默认随机初始化对 $A$ 和 $B$ 都用高斯/均匀分布初始化，第 0 步就会往冻结的预训练底座里注入巨大的随机矩阵噪声 $\frac{\alpha}{r}BA$，导致初始困惑度（Perplexity）暴涨、把预训练分布瞬间打碎；反之，若错把 $A$ 和 $B$ 同时初始化为全零，则由式~\eqref{eq:lora-grad} 可见 $\nabla_A \mathcal{L} \equiv \mathbf{0}$ 且 $\nabla_B \mathcal{L} \equiv \mathbf{0}$，整个适配器在优化中彻底“假死”纹丝不动！必须严格执行“$A$ 随机初始化、$B$ 全零初始化”的非对称规范。
  \item \textbf{长上下文服务忽视 KV 缓存爆发导致线上 OOM}：许多团队上线前只按权重大小（如 7B BF16 占 14\,GB）估算显存，觉得一张 24\,GB 或 40\,GB 显卡绰绰有余。结果一放开 $32\mathrm{K}$ 上下文与多并发请求，由式~\eqref{eq:kv-cache-formula} 决定的 KV 缓存瞬间吞噬几十 GB 显存直接触发 CUDA OOM 崩溃！工程上线前必须依据 $M_{\mathrm{KV/tok}} \times B \times S_{\max}$ 精确封顶最大在飞词元预算（\texttt{max\_num\_batched\_tokens}），并搭配 PagedAttention 与 FP8/INT8 KV-Cache 量化。
\end{enumerate}
\end{failBox}

\begin{extBox}{混合专家模型（MoE）专家并行与 Prefill/Decode 分离架构}
站在 2025--2026 年的超大规模 AI 系统最前沿，两项范式级演进正重新定义集群架构：
\begin{itemize}
  \item \textbf{稀疏混合专家（MoE）与专家并行（Expert Parallelism, EP）}：以 DeepSeek-V3（总参 671B、单词元仅激活 37B）与 Mixtral 为代表的 MoE 架构，把稠密 MLP 替换为数十至数百个小型专家网络，每个词元仅由门控路由器（Router）动态挑选 Top-$K$（如 $K = 2\text{--}8$）个专家计算，一举打破了“参数量与单步 FLOPs 严格成正比”的魔咒！在分布式训练上，这催生了第四维并行——\textbf{专家并行（EP）}：不同 GPU 存放不同专家子集，路由后通过 \textbf{All-to-All} 集合通信将词元分发给对应专家所在的 GPU、算完后再 All-to-All 传回。其核心系统挑战在于\textbf{负载均衡（Load Balancing）}：一旦热门专家被大量词元挤爆，就会引发严重的木桶短板效应（Straggler）甚至词元丢弃（Token Dropping），必须配合辅助负载均衡损失或无辅助损失动态偏置调节（Auxiliary-Loss-Free Balancing）。
  \item \textbf{Prefill / Decode 分离部署架构（Disaggregated Serving，如 Splitwise / DistServe）}：回到第~\ref{ch:hardware-roofline-flash} 章定理~\ref{thm:prefill-vs-decode} 的深刻洞察——Prefill 是彻底的\textbf{计算受限}任务（优化指标是首字延迟 TTFT），而 Decode 是彻底的\textbf{访存带宽受限}任务（优化指标是词间延迟 TPOT）。把它们混在同一张 GPU 上跑连续批处理时，一个突发长提示的 Prefill 请求会立刻卡死同一张卡上所有正在 Decode 的老用户，造成严重的尾延迟毛刺！分离部署架构将集群物理划分为专职 Prefill 节点池与专职 Decode 节点池：Prefill 节点吃满算力算完提示词后，通过高速 RDMA/NVLink 网络将生成好的 KV 缓存单向传输给 Decode 节点接力续写，彻底消除双阶段干扰，成为超大规模推理集群的标准蓝图。
\end{itemize}
\end{extBox}

\section*{本章分层习题}
\addcontentsline{toc}{section}{本章分层习题}

\exTier{A 层：概念与公式手算}{熟练核算 3D 并行/ZeRO 显存通信、LoRA 参数比与 KV 缓存账本}

\begin{enumerate}
  \item[\textbf{6.6.1}] \textbf{70B 模型在 64 卡集群上的 ZeRO 显存与 3D 并行综合核算}：现使用 8 台 8 卡 H100 服务器（共 $N_{\mathrm{GPU}} = 64$ 张 80\,GB 显卡）训练一个 $\Phi = 7\times 10^{10}$（70B）大模型，采用 BF16 + FP32 AdamW 混合精度。
  \begin{enumerate}
    \item 若纯用 64 路数据并行（$N_d = 64$），分别按表~\ref{tab:zero-ledger} 计算 DDP、ZeRO-1、ZeRO-2 与 ZeRO-3 下的单卡静态显存（GB）与每步单卡梯度/权重总通信量（GB）；判断哪几种方案连静态显存都装不进单张 80\,GB 显卡；
    \item 若改为经典 3D 并行：单机内张量并行 $N_t = 8$、跨机流水线并行 $N_p = 4$、顶层 ZeRO-1 数据并行 $N_d = 2$（满足 $8 \times 4 \times 2 = 64$），此时每张卡负责的局部模型参数量为 $\Phi_{\mathrm{local}} = \Phi / (N_t N_p)$。计算该配置下的单卡静态显存，并求出当微批数 $m = 32$ 时的流水线气泡率 $f_{\mathrm{bubble}}$。
  \end{enumerate}
  \exHint{(a) $\Psi = 2\Phi = 140\,\mathrm{GB}$；DDP $1120\,\mathrm{GB}$、ZeRO-1 $293.1\,\mathrm{GB}$、ZeRO-2 $155.3\,\mathrm{GB}$ 均超 80\,GB，唯有 ZeRO-3 的 $17.5\,\mathrm{GB}$ 能装下（通信 $3\Psi = 420\,\mathrm{GB}$）；(b) $\Phi_{\mathrm{local}} = 70 / 32 = 2.1875\mathrm{B}$，叠加 $N_d=2$ 的 ZeRO-1（每参数 $4 + 12/2 = 10\,\mathrm{B}$），单卡静态显存仅 $21.88\,\mathrm{GB}$！气泡率 $f_{\mathrm{bubble}} = \frac{4-1}{32+4-1} = \frac{3}{35} \approx 8.57\%$。}

  \item[\textbf{6.6.2}] \textbf{LoRA 显存预算与 GQA KV-Cache 精确核算}：考虑一座类似 Llama-3 8B 的模型：$L = 32$ 层、隐藏维度 $d = 4096$、查询头数 $h_q = 32$、KV 头数 $n_{\mathrm{kv}} = 8$（GQA-4）、头维度 $d_{\mathrm{head}} = 128$。
  \begin{enumerate}
    \item 若仅对每层的 $W_Q, W_V \in \R^{4096 \times 4096}$（实际 $W_V$ 为 $1024 \times 4096$，此处先按双 $4096 \times 4096$ 标准 MHA 近似估算）施加秩 $r = 16$ 的 LoRA 适配器，计算全网 32 层总可训练参数量、占 8B 模型的百分比，以及 BF16 底座 + LoRA 训练状态的总静态显存；
    \item 用式~\eqref{eq:kv-cache-formula} 计算该模型在 BF16 下的单词元 KV 缓存大小 $M_{\mathrm{KV/tok}}$；若在单张 80\,GB H100 上部署（BF16 权重占 16\,GB，预留 4\,GB 框架开销，余下 60\,GB 全部给 KV 缓存），在平均序列长度 $S = 4096$ 下，借助 PagedAttention 理论最多可并发服务多少条请求？
  \end{enumerate}
  \exHint{(a) 每层两个方阵共增 $2 \times 2 \times 4096 \times 16 = 262144$ 参数，32 层合计约 $8.39\mathrm{M}$ 参数（仅占 $0.105\%$！），静态显存仅比 $16\,\mathrm{GB}$ 底座多约 $0.13\,\mathrm{GB}$；(b) $M_{\mathrm{KV/tok}} = 2 \times 32 \times 8 \times 128 \times 2 = 131072\,\mathrm{B} = 128\,\mathrm{KiB/tok}$，单条 4096 序列占 $512\,\mathrm{MiB} \approx 536.9\,\mathrm{MB}$，故 $60\,\mathrm{GiB}$ 二进制口径可并发 $120$ 条（十进制 $60\,\mathrm{GB}$ 约 $111$ 条）！}
\end{enumerate}

\exTier{B 层：严谨数学证明}{严格推导环形归约延迟边界与投机解码期望加速比}

\begin{enumerate}
  \item[\textbf{6.6.3}] \textbf{Ring AllReduce 对比 Halving-Doubling 树形归约的延迟—带宽权衡}：由定理~\ref{thm:ring-allreduce} 知，Ring AllReduce 在 $P$ 张卡上传输 $M$ 字节张量的总耗时为 $T_{\mathrm{Ring}} = 2(P-1)\alpha + 2\frac{P-1}{P}\frac{M}{B}$。另一种经典的 Rabenseifner（递归减半—加倍，Recursive Halving-Doubling）算法在 $P = 2^k$ 时，用 $\log_2 P$ 轮完成 ReduceScatter、再用 $\log_2 P$ 轮完成 AllGather，第 $j$ 轮传输数据量为 $M / 2^j$。
  \begin{enumerate}
    \item 证明 Rabenseifner 算法每张 GPU 累计发送的总字节数同样精确等于 $2\frac{P-1}{P}M$，而固定延迟项从 $2(P-1)\alpha$ 锐减至 $2(\log_2 P)\alpha$；
    \item 当 $P = 1024$ 卡、延迟 $\alpha = 5\,\mu\mathrm{s}$、带宽 $B = 50\,\mathrm{GB/s}$ 时，计算不考虑拥塞差异时两种算法的固定延迟差 $T_{\mathrm{Ring}} - T_{\mathrm{Rab}}$，并求出 Ring AllReduce 自身延迟项与带宽项恰好相等的临界消息体积 $M_{\mathrm{crit}}$，解释为何超大集群同步小张量（如 LayerNorm 梯度）时 Ring 会因 $O(P)\alpha$ 延迟项急剧退化！
  \end{enumerate}
  \exHint{(a) 等比级数求和 $\sum_{j=1}^{\log_2 P} \frac{M}{2^j} = M\big(1 - \frac{1}{P}\big) = \frac{P-1}{P}M$，两阶段合计精确为 $2\frac{P-1}{P}M$；(b) Ring 延迟项高达 $2 \times 1023 \times 5\,\mu\mathrm{s} = 10.23\,\mathrm{ms}$（而 Rabenseifner 仅 $0.10\,\mathrm{ms}$，相差 $10.13\,\mathrm{ms}$），在 $50\,\mathrm{GB/s}$ 链路上等价于多传了约 $511.5\,\mathrm{MB}$ 通信量，对应临界张量大小 $M_{\mathrm{crit}} = P\alpha B = 256\,\mathrm{MB}$！}

  \item[\textbf{6.6.4}] \textbf{投机解码的多步联合分布精确性与真实墙钟加速比公式}：
  \begin{enumerate}
    \item 定理~\ref{thm:speculative-lossless} 证明了单个位置的边缘无损性。请利用概率链式法则 $\Prob(X_1, \dots, X_T) = \prod_{t=1}^T \Prob(X_t \mid X_{<t})$ 严格说明：在每轮迭代依次接受前 $k$ 个草稿词元、在第 $k+1$ 处拒绝并重采（随后立刻截断丢弃第 $k+2, \dots, \gamma$ 个草稿）的完整多步投机循环中，整条输出序列的联合分布依然与目标大模型 $p$ 严格等价；
    \item 设小草稿模型单步前向耗时与目标大模型单次并行验证耗时之比为 $c = t_{\mathrm{draft}} / t_{\mathrm{target}} \ll 1$（忽略验证 $\gamma+1$ 个词元相对 $1$ 个词元的微小增量）。结合式~\eqref{eq:spec-expected-len}，证明投机解码相对原始自回归解码的期望墙钟加速比为
    \[
      \mathrm{Speedup}(\gamma, \alpha, c) \;=\; \frac{1 - \alpha^{\gamma + 1}}{(1 - \alpha)(\gamma c + 1)},
    \]
    并在 $\alpha = 0.8, c = 0.05$ 下求出最大化加速比的最优草稿长度 $\gamma^*$。
  \end{enumerate}
  \exHint{(a) 因为一旦在第 $k+1$ 步拒绝，后续草稿被全部作废，新一轮草稿重新以已验证的精确前缀 $X_{\le k+1}$ 为条件启动，故每个已接受与重采位置都严格服从 $p(\cdot \mid X_{<t})$；(b) 每轮总耗时为 $\gamma t_{\mathrm{draft}} + t_{\mathrm{target}} = (\gamma c + 1) t_{\mathrm{target}}$，产出期望词元数 $\E[L] = \frac{1-\alpha^{\gamma+1}}{1-\alpha}$，作比即得！代入 $\alpha=0.8, c=0.05$ 逐个枚举 $\gamma = 6, 7, 8, 9$ 可见 $\gamma^* = 8$ 时加速比达到峰值约 $3.092\times$。}
\end{enumerate}

\exTier{C 层：机制剖析与反例设计}{诊断 LoRA 对称锁死陷阱与离群值 SmoothQuant 尺度迁移}

\begin{enumerate}
  \item[\textbf{6.6.5}] \textbf{反例剖析：为什么 LoRA 把 $A, B$ 同时零初始化会导致训练永久锁死？}：考虑极简标量线性模型 $y = (w_0 + b a) x$（秩 $r = 1$，$w_0$ 冻结），损失函数为平方误差 $\mathcal{L}(a, b) = \frac{1}{2}(y - y^*)^2$。
  \begin{enumerate}
    \item 写出 $\frac{\partial \mathcal{L}}{\partial a}$ 与 $\frac{\partial \mathcal{L}}{\partial b}$ 的解析式，证明若初始化 $a_0 = 0, b_0 = 0$，则原点 $(0, 0)$ 是损失函数的驻点（鞍点），梯度下降与 Adam 均会永远停在 $(0, 0)$ 无法动弹；
    \item 证明当取 $a_0 \neq 0, b_0 = 0$ 时，不仅第 0 步模型输出精确等于 $w_0 x$，而且第 1 步梯度更新后立即有 $b_1 \neq 0$，从而在第 2 步顺势激活 $\frac{\partial \mathcal{L}}{\partial a} \neq 0$；最后思考：如果反过来取 $a_0 = 0, b_0 \neq 0$ 在数学上是否同样可行？为何工程惯例偏好 $B = 0$？
  \end{enumerate}
  \exHint{(a) $\frac{\partial \mathcal{L}}{\partial a} = (y - y^*) b x$，$\frac{\partial \mathcal{L}}{\partial b} = (y - y^*) a x$，当 $a=b=0$ 时二者同时严格为零；(b) $A=0, B\neq 0$ 在数学上同样对称可行，只是输入维 $x$ 乘随机 $A$ 后再乘全零 $B$ 在多适配器拼接时前向输出恒为零、反向传播更直观。}

  \item[\textbf{6.6.6}] \textbf{重尾通道 INT4 量化手算与 SmoothQuant 迁移演练}：设某 $2\times 2$ 权重矩阵与单行激活向量为
  \[
    W = \begin{bmatrix} 0.4 & -0.3 \\ 7.0 & -5.0 \end{bmatrix}, \qquad
    x = \begin{bmatrix} 10.0 & 0.5 \end{bmatrix}.
  \]
  \begin{enumerate}
    \item 分别对 $W$ 执行 \textbf{INT4（$q_{\max} = 7$）逐张量量化}与\textbf{逐行（逐通道）量化}，手算量化后的重构矩阵 $\hat{W}_{\mathrm{pt}}$ 与 $\hat{W}_{\mathrm{pc}}$，观察正常第 1 行 $[0.4, -0.3]$ 在逐张量量化下变成了什么！
    \item 若还要同时对激活 $x$（第 1 维存在 $20\times$ 离群值！）做 INT4 逐张量量化，请设计一个对角平滑矩阵 $\diag(s_1, s_2)$（例如取 SmoothQuant 经典迁移指数 $\alpha = 0.5$，即 $s_j = \max|x_j|^{0.5} / \max|W_{j,:}|^{0.5}$），计算变换后的 $\tilde{x} = x \diag(s)^{-1}$ 与 $\tilde{W} = \diag(s) W$，验证乘积 $\tilde{x} \tilde{W} = x W$ 严格不变且两者动态范围同时大幅收窄！
  \end{enumerate}
  \exHint{(a) 逐张量 $s_{\mathrm{pt}} = 7.0 / 7 = 1.0$，于是第一行 $\operatorname{round}([0.4, -0.3]/1.0) = [0, 0]$ —— 整个通道彻底被清零！而逐行量化第一行取 $s_1 = 0.4 / 7 \approx 0.0571$，重构为 $[0.400, -0.286]$，误差极小！}
\end{enumerate}

\exTier{D 层：编程实验与系统延伸}{动手实现分组量化折中扫描与无损投机解码蒙特卡洛验算}

\begin{enumerate}
  \item[\textbf{6.6.7}] \textbf{编程实验：分组量化（Group-Wise Quantization）的 INT4 误差—显存帕累托前沿}：在代码~\ref{lst:quant-code}（\texttt{exp6\_quant\_and\_lora()}）基础上，将离群噪声进一步改为“随机分布在全矩阵不同通道内的 $2\%$ 孤立尖峰离群元素（$8\times$ 放大）”：
  \begin{enumerate}
    \item 实现通用函数 \texttt{quant\_groupwise(W, bits=4, group\_size=g)}，沿列维将每行切分为若干长度为 $g$ 的连续小组，每组独立计算对称缩放因子；
    \item 扫描组大小 $g \in \{256\text{（逐通道）}, 128, 64, 32, 16\}$，记录对应的输出相对误差与“有效平均每权重比特数” $b_{\mathrm{eff}} = 4 + 16 / g$，绘制误差与有效比特数的折中曲线，验证为何主流 4-bit 量化（GPTQ/AWQ）普遍选用 $g = 128$（有效 $4.125\,\mathrm{bit}$）作为甜点！
  \end{enumerate}
  \exHint{对行内孤立离群元素，逐通道（$g=256$）会被同一行里的单个尖峰拖累整行，而 $g=64$ 或 $128$ 把破坏范围严格锁死在一小组之内，且每权重仅多花 $0.125\text{--}0.25$ 比特！}

\par\pagebreak[4]
  \item[\textbf{6.6.8}] \textbf{编程实验：蒙特卡洛验证投机解码无损分布定理与加速比曲线}：编写纯 Python/NumPy 实验验证定理~\ref{thm:speculative-lossless}：
  \begin{enumerate}
    \item 在词表大小 $|\mathcal{V}| = 16$ 上随机生成一对有差异的概率分布 $(p, q)$（例如 $p = \softmax(z)$，$q = \softmax(0.6 z + \epsilon)$）；
    \item 实现定理~\ref{thm:speculative-lossless} 的拒绝采样与残差分布重采样过程，独立重复 $N = 2\times 10^5$ 次，统计经验频率分布 $\hat{p}$，验证总变差距离 $\frac{1}{2}\sum_x |\hat{p}(x) - p(x)| < 5\times 10^{-3}$（远小于 $q$ 与 $p$ 之间的原始距离！）；
    \item 计算理论一步接受率 $\alpha = \sum_x \min(p(x), q(x))$，并模拟连续 $\gamma = 5$ 步投机循环 $2\times 10^4$ 次，验证实测平均单轮产出词元数与理论公式~\eqref{eq:spec-expected-len} $\frac{1 - \alpha^6}{1 - \alpha}$ 吻合至小数点后两位！
  \end{enumerate}
  \exHint{注意构造残差分布时，若恰巧 $p \equiv q$（分母为 0），则每步必然 $100\%$ 接受，绝不会进入拒绝分支；实测经验总变差距离随 $O(1/\sqrt{N})$ 收敛到零，直观印证投机解码严格零损。}
\end{enumerate}


% ==================== Appendix: Complete Solutions ====================
\appendix
\ctexset{chapter={name={附录\ ,}, number=\Alph{chapter}}}
\renewcommand{\chapWatermarkNum}{\thechapter}
\renewcommand{\chaptermark}[1]{\markboth{附录 \thechapter\quad #1}{}}
\newcommand{\solItem}[2]{%
  \par\vspace{4.5mm}\noindent
  \colorbox{tealLight}{\strut\;\sffamily\bfseries\footnotesize\color{tealDark}习题\,#1\;}\;%
  {\normalsize\sffamily\bfseries\color{ink}#2}\par\nopagebreak\vspace{1.2mm}%
}

\chapter[全卷 48 道分层习题完整精解]{全卷 48 道分层习题完整精解}
\label{app:solutions}

\chapEpigraph{简洁是可靠性的先决条件。}{Edsger W. Dijkstra（艾兹格·迪杰斯特拉，1975）}

本附录给出第 1 章至第 6 章全部 48 道分层习题（A/B/C/D 四层，每章 8 题）的完整分步解答。所有手算题、反例算术、不等式常数与编程实验预期指标均经过独立的 Python 数值程序交叉验算，读者可直接对照自检。

% =====================================================================
\section*{第 1 章习题解答（共 8 题）}
\addcontentsline{toc}{section}{第 1 章习题解答（共 8 题）}

\solItem{6.1.1}{单条轨迹的 MC 与 TD(0) 逐步手算}

已知状态集合 $\{s_1, s_2, \mathsf{T}\}$，初值 $V(s_1) = 0$、$V(s_2) = 0$、$V(\mathsf{T}) = 0$，折扣因子 $\gamma = 0.9$，固定步长 $\alpha = 0.5$。

\textbf{(a) 第一条回合 $s_1 \xrightarrow{r=2} s_2 \xrightarrow{r=4} \mathsf{T}$ 的更新结果。}
\begin{enumerate}
  \item \textbf{首次访问蒙特卡洛（MC）}：先从回合末尾倒推计算各时刻完整折扣回报 $G_t$：
  \[
    G_1 = r_1 = 4, \qquad
    G_0 = r_0 + \gamma G_1 = 2 + 0.9 \times 4 = 5.6.
  \]
  回合结束后按式~\eqref{eq:mc-inc} 执行恒定步长增量更新：
  \begin{align*}
    V_{\mathrm{MC}}(s_1) &\leftarrow 0 + 0.5 \times (5.6 - 0) = \mathbf{2.8}, \\
    V_{\mathrm{MC}}(s_2) &\leftarrow 0 + 0.5 \times (4.0 - 0) = \mathbf{2.0}.
  \end{align*}
  \item \textbf{在线 TD(0)}：按时间步顺序逐转移即时更新（式~\eqref{eq:td0-update}）。
  \begin{itemize}
    \item 第 1 步（$s_1 \xrightarrow{r=2} s_2$）：此时 $V(s_2)$ 尚未更新，仍为初值 $0$，故 TD 误差 $\delta_0 = 2 + 0.9 \times 0 - 0 = 2.0$，更新得
    \[
      V^{(1)}(s_1) = 0 + 0.5 \times 2.0 = \mathbf{1.0}.
    \]
    \item 第 2 步（$s_2 \xrightarrow{r=4} \mathsf{T}$）：终止态价值 $V(\mathsf{T}) = 0$，TD 误差 $\delta_1 = 4 + 0.9 \times 0 - 0 = 4.0$，更新得
    \[
      V^{(1)}(s_2) = 0 + 0.5 \times 4.0 = \mathbf{2.0}.
    \]
  \end{itemize}
\end{enumerate}

\textbf{(b) 第二条相同回合的 TD(0) 更新结果。} 沿用第一回合后的价值表 $V^{(1)}(s_1) = 1.0, V^{(1)}(s_2) = 2.0$：
\begin{itemize}
  \item 第 1 步（$s_1 \xrightarrow{r=2} s_2$）：此刻 $V^{(1)}(s_2) = 2.0$ 首次通过自举目标传入 $s_1$，TD 目标为 $2 + 0.9 \times 2.0 = 3.8$，TD 误差 $\delta_0 = 3.8 - 1.0 = 2.8$，更新得
  \[
    V^{(2)}(s_1) = 1.0 + 0.5 \times 2.8 = \mathbf{2.4}.
  \]
  \item 第 2 步（$s_2 \xrightarrow{r=4} \mathsf{T}$）：TD 误差 $\delta_1 = 4 + 0 - 2.0 = 2.0$，更新得
  \[
    V^{(2)}(s_2) = 2.0 + 0.5 \times 2.0 = \mathbf{3.0}.
  \]
\end{itemize}
\textbf{物理洞察}：MC 在第 1 条回合结束时就一次性把远端奖励 $r_1 = 4$ 传回了 $s_1$；而单步在线 TD(0) 的奖励信息每经历 1 条回合才沿时间逆向回传 1 步（$n$ 步长的链至少需要 $n$ 条回合，远端奖励才会首次触及起点），这正是引入多步 TD 与 TD($\lambda$) 资格迹的直接动因。

\solItem{6.1.2}{SARSA、Q-Learning 与期望 SARSA 单步目标对比}

已知 $r_t = -1, \gamma = 0.9, \varepsilon = 0.15$，状态 $s_2$ 有 $m = 3$ 个动作，表值 $Q(s_2, a_1) = 10$、$Q(s_2, a_2) = 2$、$Q(s_2, a_3) = -20$。由式~\eqref{eq:eps-greedy}，在 $\varepsilon$-贪心策略下贪心最优动作 $a_1$ 的采样概率为
\[
  \pi(a_1 \mid s_2) = 1 - \varepsilon + \frac{\varepsilon}{3} = 1 - 0.15 + 0.05 = 0.90, \qquad
  \pi(a_2 \mid s_2) = \pi(a_3 \mid s_2) = \frac{\varepsilon}{3} = 0.05.
\]

\textbf{(a) Q-Learning 与期望 SARSA 的精确目标值。}
\begin{align*}
  y_{\mathrm{Q}} &= r_t + \gamma \max_{a \in \{a_1, a_2, a_3\}} Q(s_2, a) = -1 + 0.9 \times 10 = \mathbf{8.0}, \\
  \E_{a \sim \pi(\cdot \mid s_2)}[Q(s_2, a)] &= 0.90 \times 10 + 0.05 \times 2 + 0.05 \times (-20) = 9.0 + 0.1 - 1.0 = 8.1, \\
  y_{\mathrm{ExpSARSA}} &= -1 + 0.9 \times 8.1 = -1 + 7.29 = \mathbf{6.29}.
\end{align*}

\textbf{(b) SARSA 随机目标的分布、期望与方差。} 设随机变量 $Y = r_t + \gamma Q(s_2, A_{t+1})$，则其离散分布为：
\begin{center}
\footnotesize
\begin{tabular}{cccc}
\toprule
下一动作 $A_{t+1}$ & $a_1$（最优） & $a_2$（次优） & $a_3$（灾难动作） \\
\midrule
采样概率 $\Prob(A_{t+1})$ & $0.90$ & $0.05$ & $0.05$ \\
目标取值 $y = -1 + 0.9\,Q(s_2, A_{t+1})$ & $+8.0$ & $+0.8$ & $-19.0$ \\
\bottomrule
\end{tabular}
\end{center}
计算期望与方差：
\begin{align*}
  \E[Y] &= 0.90 \times 8.0 + 0.05 \times 0.8 + 0.05 \times (-19.0) = 7.20 + 0.04 - 0.95 = \mathbf{6.29} = y_{\mathrm{ExpSARSA}}, \\
  \Var(Y) &= 0.90(8.0 - 6.29)^2 + 0.05(0.8 - 6.29)^2 + 0.05(-19.0 - 6.29)^2 \\
          &= 0.90 \times 2.9241 + 0.05 \times 30.1401 + 0.05 \times 639.5841 \\
          &= 2.63169 + 1.507005 + 31.979205 = \mathbf{36.1179},
\end{align*}
标准差高达 $\sqrt{36.1179} \approx 6.01$。由此可见，期望 SARSA 与 SARSA 在给定下一状态 $s_2$ 下具有\textbf{完全相同的条件期望}，但期望 SARSA 直接对后继动作求解析加权平均，彻底消除了动作采样引入的方差 $36.1179$（尤其消除了以 $5\%$ 概率抽到灾难动作 $a_3$ 导致的剧烈更新震荡）。

\solItem{6.1.3}{Robbins--Monro 步长家族的边界检验}

\textbf{(a) 充要条件证明。} 令 $\alpha_t = (t+1)^{-p}$（$t \ge 0$），设 $n = t + 1 \ge 1$，则两个 Robbins--Monro 条件~\eqref{eq:rm-cond} 等价于两个 $p$-级数的敛散性：
\[
  \sum_{t=0}^{\infty} \alpha_t = \sum_{n=1}^{\infty} \frac{1}{n^p} = \infty \iff p \le 1, \qquad
  \sum_{t=0}^{\infty} \alpha_t^2 = \sum_{n=1}^{\infty} \frac{1}{n^{2p}} < \infty \iff 2p > 1 \iff p > \frac{1}{2}.
\]
取两者交集，当且仅当 $\frac{1}{2} < p \le 1$ 时同时满足条件~\eqref{eq:rm-cond}。特别地，$p = 1$（即 $\alpha_t = \frac{1}{t+1}$，对应算术平均）落在收敛区间右端点；若 $p > 1$（衰减过快），步长总和有限，算法可能停在远离真值之处；若 $p \le \frac{1}{2}$（衰减过慢），平方和发散，噪声方差无法被压制至零。

\textbf{(b) 固定常数步长 $\alpha \in (0, 1)$ 的稳态均方误差。} 记 $e_t = \theta_t - \mu$，更新式改写为 $e_{t+1} = (1 - \alpha) e_t + \alpha \xi_t$。因 $\xi_t$ 均值为 $0$、方差为 $\sigma^2$ 且与历史 $e_t$ 独立，两端取平方期望得均方误差递推式~\eqref{eq:rm-mse-rec}：
\[
  v_{t+1} \triangleq \E[e_{t+1}^2] = (1 - \alpha)^2 v_t + \alpha^2 \sigma^2.
\]
由于 $\alpha \in (0, 1)$，有 $0 < (1 - \alpha)^2 < 1$，上述线性收缩递推存在唯一不动点 $v_\infty$。设 $v_\infty = (1 - \alpha)^2 v_\infty + \alpha^2 \sigma^2$，展开得
\[
  \big[1 - (1 - 2\alpha + \alpha^2)\big] v_\infty = \alpha^2 \sigma^2 \implies \alpha(2 - \alpha) v_\infty = \alpha^2 \sigma^2 \implies
  v_\infty = \lim_{t \to \infty} \E[(\theta_t - \mu)^2] = \mathbf{\frac{\alpha}{2 - \alpha}\,\sigma^2}.
\]
展开显式通项 $v_t = (1 - \alpha)^{2t}(v_0 - v_\infty) + v_\infty$，可见初始误差以指数速率 $(1-\alpha)^{2t}$ 衰减，但最终不会归零，而是残留在宽度为 $O(\sqrt{\alpha}\sigma)$ 的稳态噪声球内。\textbf{在非平稳环境中受青睐的原因}：将 $\theta_t$ 展开为加权平均 $\theta_t = (1-\alpha)^t \theta_0 + \sum_{k=0}^{t-1} \alpha(1-\alpha)^{t-1-k}(\mu + \xi_k)$，常数步长赋予近期样本指数更高的权重、自动遗忘过时历史；在 RL 控制任务中，策略 $\pi$ 持续演进使得目标分布非平稳，递减步长 $\alpha_t \to 0$ 反而会令估值器“失去追踪新策略的能力”。

\solItem{6.1.4}{均匀噪声下的最大化偏差与 Double Q 对比}

\textbf{(a) 单表最大值的精确期望。} 设 $X_1, \dots, X_K \overset{\mathrm{i.i.d.}}{\sim} \mathrm{Uniform}(-1, 1)$，记 $M_K = \max_{1 \le a \le K} X_a$。对任意 $x \in [-1, 1]$，单个变量的 CDF 为 $F(x) = \frac{x + 1}{2}$，由独立性得最大值的 CDF 与 PDF 分别为
\[
  F_{\max}(x) = \Prob(M_K \le x) = \prod_{a=1}^K \Prob(X_a \le x) = \left(\frac{x + 1}{2}\right)^K, \qquad
  f_{\max}(x) = F'_{\max}(x) = \frac{K}{2}\left(\frac{x + 1}{2}\right)^{K-1}.
\]
作换元 $u = \frac{x + 1}{2} \in [0, 1]$（则 $x = 2u - 1, dx = 2\,du$），直接积分计算期望：
\[
  \E[M_K] = \int_{-1}^{1} x \cdot \frac{K}{2}\left(\frac{x + 1}{2}\right)^{K-1} dx
  = \int_{0}^{1} (2u - 1) K u^{K-1}\,du
  = 2K \cdot \frac{1}{K+1} - K \cdot \frac{1}{K}
  = \mathbf{\frac{K - 1}{K + 1}}.
\]
检验：$K = 1$ 时 $\E[M_1] = 0$（无偏差）；$K = 2$ 时为 $1/3 \approx 0.333$；$K = 10$ 时高达 $9/11 \approx 0.818$；当 $K \to \infty$ 时，$\E[M_K] \to \mathbf{1}$（逼近噪声分布的上确界），尽管所有动作的真实价值均为 $0$！这从理论上精确解释了为何动作空间越大，单表 Q-Learning 的正向高估偏差越严重。

\textbf{(b) Double Q-Learning 的解耦期望。} 令 $a^* = \argmax_{1 \le a \le K} Q_A(s, a)$。因为表 $Q_A$ 与表 $Q_B$ 相互独立，随机索引 $a^*$ 完全由 $Q_A$ 决定而与整张表 $Q_B$ 独立。由全期望公式对 $a^*$ 取条件期望：
\[
  \E\!\left[Q_B(s, a^*)\right]
  = \E_{Q_A}\!\Big[\,\E_{Q_B}\!\big[Q_B(s, a^*) \mid Q_A\big]\,\Big]
  = \E_{Q_A}\!\left[\,0\,\right] = \mathbf{0}.
\]
由此严格证明：将“动作选择（用 $Q_A$）”与“价值评估（用 $Q_B$）”解耦后，零均值估计噪声引发的系统性正向最大化偏差被彻底归零（代价是若不同动作真值不全相等，评估表可能引入轻微的负向低估，但通常远温和且不会导致发散）。

\solItem{6.1.5}{批量 MC 与批量 TD(0) 的非马尔可夫反例}

构造如下带\textbf{隐状态（部分可观测）}的数据生成过程：真实环境存在三个内部非终止状态 $\{A, B_{\mathrm{bad}}, B_{\mathrm{good}}\}$，但观测器将 $B_{\mathrm{bad}}$ 与 $B_{\mathrm{good}}$ 折叠为同一个观测符号 $B$（即状态观测不满足马尔可夫性）：
\begin{itemize}
  \item 回合起点：以概率 $1/8$ 从 $A$ 出发，以概率 $7/8$ 直接从 $B_{\mathrm{good}}$ 出发；
  \item 状态转移与奖励：从 $A$ 出发必以奖励 $0$ 确定性进入“损坏态” $B_{\mathrm{bad}}$；$B_{\mathrm{bad}}$ 以概率 $1$ 得奖励 $0$ 终止（$V^*(B_{\mathrm{bad}}) = 0$）；而 $B_{\mathrm{good}}$ 以概率 $6/7$ 得奖励 $1$ 终止、以概率 $1/7$ 得奖励 $0$ 终止（$V^*(B_{\mathrm{good}}) = 6/7$）。
\end{itemize}
在此真实机制下，取样 8 条回合恰好可以出现例~\ref{ex:batch-mc-td} 中的全部观测：1 条 $A \xrightarrow{0} B \xrightarrow{0} \mathsf{T}$（真实经过 $B_{\mathrm{bad}}$）、6 条 $B \xrightarrow{1} \mathsf{T}$ 和 1 条 $B \xrightarrow{0} \mathsf{T}$（真实来自 $B_{\mathrm{good}}$）。此时：
\begin{itemize}
  \item 状态 $A$ 的真实价值严格为 $V^*(A) = 0 + V^*(B_{\mathrm{bad}}) = \mathbf{0}$，这恰与批量 MC 给出的 $V_{\mathrm{MC}}(A) = 0$ 完全一致（因为 MC 直接统计从 $A$ 出发的完整轨迹回报，\textbf{不依赖任何马尔可夫假设}）；
  \item 相反，批量 TD(0) 强行将观测符号 $B$ 当作马尔可夫状态，把 $B_{\mathrm{good}}$ 的高奖励错误地自举传导给 $A$，得到严重失真的估值 $V_{\mathrm{TD}}(A) = 0.75$。
\end{itemize}
\textbf{结论}：批量 TD(0) 的高数据效率本质上是“用马尔可夫结构假设换取低方差”；当特征表示存在状态混叠（部分可观测 POMDP）时，MC 对非马尔可夫性天然免疫，而 TD 自举会引入不可消除的结构偏差。

\solItem{6.1.6}{GLIE 与步长耦合的探索陷阱}

回忆 GLIE（Greedy in the Limit with Infinite Exploration）包含两条缺一不可的条件：（i）每个状态--动作对 $(s, a)$ 都被访问无穷多次；（ii）策略在极限下收敛于贪心策略。

设动作空间大小 $m \ge 2$，在某个高频访问状态 $s$ 下，记事件 $E_k$ 为“在第 $k$ 回合首次访问状态 $s$ 时发生了 $\varepsilon_k$-随机探索”。由题意 $\varepsilon_k = 0.5^k$，有
\[
  \sum_{k=1}^{\infty} \Prob(E_k) = \sum_{k=1}^{\infty} 0.5^k = \frac{0.5}{1 - 0.5} = 1 < \infty.
\]
由概率论中的\textbf{第一 Borel--Cantelli 引理}，立得
\[
  \Prob\!\left(\limsup_{k \to \infty} E_k\right) = \Prob(\text{无穷多个 } E_k \text{ 同时发生}) = 0.
\]
也就是说，以概率 $1$，随机探索仅在有限多个回合中发生；存在某个随机时刻 $K_0 < \infty$，使得从第 $K_0 + 1$ 回合起，算法在状态 $s$ 永远只执行当前贪心动作。进一步地，考虑某个真实最优动作 $a^*$（假设初始估值 $Q_0(s, a^*) = 0$ 与其他动作相同，或略低于某次优动作的首次噪声估值），计算它在\emph{所有回合中一次都不被随机探索选中}的概率下界：每回合随机选中 $a^*$ 的概率为 $\frac{\varepsilon_k}{m} = \frac{0.5^k}{m}$，于是
\[
  \Prob(\text{永远不通过随机探索选中 } a^*) = \prod_{k=1}^{\infty}\!\left(1 - \frac{0.5^k}{m}\right) > 0
\]
（因为 $\sum_{k=1}^\infty \frac{0.5^k}{m} = \frac{1}{m} < \infty$，无穷乘积收敛于严格正数；例如 $m = 2$ 时该下界超过 $0.58$）。因此，指数衰减探索率 $\varepsilon_k = 0.5^k$ 违反了 GLIE 的无穷探索条件，会以严格正概率把智能体永久锁死在次优动作上。正确的 GLIE 衰减应满足 $\sum_{k=1}^\infty \varepsilon_k = \infty$ 且 $\varepsilon_k \to 0$，例如 $\varepsilon_k = 1/k$ 或 $\varepsilon_k = c / \sqrt{k}$。

\solItem{6.1.7}{编程实验：在悬崖行走与双状态带噪 MDP 上对比三种表格算法}

\textbf{(a) 悬崖行走中的 Expected SARSA 实测基准。} 在 \texttt{exp1\_cliff\_sarsa\_vs\_q()} 中，将更新目标替换为解析期望：对非终止态 $s'$，设 $a^* = \argmax_a Q[s', a]$，构造 $\varepsilon$-贪心分布 $\pi(a \mid s') = \frac{\varepsilon}{4} + (1 - \varepsilon)\mathbb{I}[a = a^*]$，则目标为 $y = r + \gamma \sum_{a=0}^3 \pi(a \mid s') Q[s', a]$。在相同 20 个随机种子、500 回合、$\varepsilon = 0.1$ 设定下，独立运行验证得到后 100 轮平均在线回报如下：
\begin{center}
\footnotesize
\begin{tabular}{lccc}
\toprule
算法（后 100 轮在线均回报） & $\alpha = 0.1$ & $\alpha = 0.5$ & $\alpha = 0.9$ \\
\midrule
标准 SARSA & $-22.68$ & $-26.22$ & $-63.30$（剧烈震荡） \\
Expected SARSA & $\mathbf{-20.66}$ & $\mathbf{-20.55}$ & $\mathbf{-20.85}$（高度稳健） \\
\bottomrule
\end{tabular}
\end{center}
\textbf{原因分析}：如习题~\textbf{6.1.2} 所证，标准 SARSA 的目标带有单次动作采样的随机方差；当 $\alpha = 0.9$ 时，偶尔一次 $\varepsilon$-探索跌落悬崖（$-100$）几乎以 $90\%$ 权重覆写当前 $Q$ 表，触发大面积策略抖动；而 Expected SARSA 彻底消除了后继动作采样方差，因此即使在极高步长 $\alpha = 0.9$ 下依然稳定在 $-20.85$。

\textbf{(b) Sutton--Barto 双状态 $A/B$ 带噪 MDP 实测基准。} 状态 $A$ 含动作“左”（进入 $B$，$r=0$）与“右”（直接终止，$r=0$），状态 $B$ 含 $K=10$ 个动作，每个动作均立即终止并给出独立噪声奖励 $r \sim \mathcal{N}(-0.1, 1)$。由于状态 $B$ 的真实期望价值为 $-0.1 < 0$，最优策略在 $A$ 处应选“右”，$\varepsilon = 0.1$ 下的最优探索选“左”下限为 $\varepsilon / 2 = 5\%$。在 $\alpha = 0.1, \varepsilon = 0.1$ 下重复 200 次独立试验（采用随机打破平局）：
\begin{itemize}
  \item {\rightskip=0pt plus 1fil\relax \textbf{标准 Q-Learning}：由于 10 个噪声动作的最大值期望显著大于 $0$（由极值理论 $\E[\max_{1 \le i \le 10} \mathcal{N}(-0.1, 1)] \approx -0.1 + 1.54 = +1.44$），前 50 回合内选“左”比例飙升至峰值约 $\mathbf{96\%}$（第 50 回合仍高达 $82\%$），第 300 回合才缓慢降至约 $13.5\%$；\par}
  \item \textbf{Double Q-Learning}：维护两张独立表 $Q_A, Q_B$，以 $0.5$ 概率轮流交叉更新（选动作按 $Q_A + Q_B$），前 50 回合选“左”比例很快回落至约 $\mathbf{14.5\%}$，第 300 回合稳定收敛至约 $10\%$ 并持续逼近理论下限 $5\%$，彻底验证了定理~\ref{thm:max-bias} 与解耦机制的有效性。
\end{itemize}

\solItem{6.1.8}{前瞻思考：从查表走向线性逼近的第一个裂缝}

记单步均方贝尔曼残差目标函数为 $L(w) = \frac{1}{2}\big[r_t + \gamma V_w(s_{t+1}) - V_w(s_t)\big]^2$，其中 $\delta_t = r_t + \gamma V_w(s_{t+1}) - V_w(s_t)$。

\textbf{1. 真实全梯度（残差梯度）与半梯度的差异。} 因为目标项中的 $V_w(s_{t+1})$ 同样依赖参数 $w$，按链式法则对 $L(w)$ 求关于 $w$ 的真实负梯度为：
\[
  -\nabla_w L(w)
  = -\delta_t \cdot \nabla_w\!\big[r_t + \gamma V_w(s_{t+1}) - V_w(s_t)\big]
  = \delta_t \Big(\nabla_w V_w(s_t) - \gamma \nabla_w V_w(s_{t+1})\Big).
\]
相比之下，半梯度 TD(0) 更新 $w \leftarrow w + \alpha \delta_t \nabla_w V_w(s_t)$ \textbf{刻意丢弃了第二项 $-\gamma \delta_t \nabla_w V_w(s_{t+1})$}（等价于把自举目标 $r_t + \gamma V_w(s_{t+1})$ 视作与 $w$ 无关的常量，在自动微分中施加 \texttt{.detach()}）。因此半梯度方向并不是任何标量损失函数的真实负梯度场（其旋度一般非零，不能写成某个势函数的梯度）。

\textbf{2. 为何会预示第~\ref{ch:dqn-pg-ac} 章的“致命三元组”与目标网络。} 在有限查表情形中，每个状态拥有独立参数，更新 $V(s_t)$ 不会牵动 $V(s_{t+1})$，贝尔曼算子在 $\|\cdot\|_\infty$ 下是严格压缩映射；但一旦引入函数逼近，不同状态共享同一组参数 $w$：当我们用半梯度把 $V_w(s_t)$ 拉向目标 $\gamma V_w(s_{t+1})$ 时，参数 $w$ 的改动可能同步把 $V_w(s_{t+1})$ 推得更远，在异策略采样下形成正反馈指数发散（即第 2 章例~\ref{ex:deadly-triad} 的致命三元组）。另一方面，为何不直接采用全残差梯度？因为在随机环境中，全残差梯度需要从同一状态 $s_t$ 采集\emph{两次独立后继样本}才能无偏估计期望的平方（双采样问题），且收敛极慢。DQN 的\textbf{目标网络} $Q_{\theta^-}$ 正是针对半梯度不稳定的工程折中：通过把目标中的参数冻结为滞后副本 $\theta^-$，在相邻同步周期 $C$ 步内把“追逐移动靶”转化为良态的静态监督回归问题！

% =====================================================================
\section*{第 2 章习题解答（共 8 题）}
\addcontentsline{toc}{section}{第 2 章习题解答（共 8 题）}

\solItem{6.2.1}{Softmax 策略与高斯策略的 Score Function 解析推导}

\textbf{(a) 离散 Softmax 策略。} 取对数得 $\ln \pi_\theta(a \mid s) = \theta^\top \phi(s, a) - \ln \sum_{a'} \exp(\theta^\top \phi(s, a'))$。对 $\theta$ 求梯度：
\begin{align*}
  \nabla_\theta \ln \pi_\theta(a \mid s)
  &= \phi(s, a) - \frac{\sum_{a'} \exp(\theta^\top \phi(s, a'))\,\phi(s, a')}{\sum_{a''} \exp(\theta^\top \phi(s, a''))} \\
  &= \phi(s, a) - \sum_{a'} \pi_\theta(a' \mid s)\,\phi(s, a')
  = \phi(s, a) - \E_{a' \sim \pi_\theta(\cdot \mid s)}\!\left[\phi(s, a')\right].
\end{align*}
直观含义：每当采到的动作 $a$ 具有正优势，参数 $\theta$ 就朝着“当前动作特征减去策略期望特征”的方向移动，从而精确提高 $a$ 的相对对数几率并压低其余动作的概率；且直接验证 $\E_{a \sim \pi_\theta}[\nabla_\theta \ln \pi_\theta(a \mid s)] = \E[\phi] - \E[\phi] = \mathbf{0}$。

\textbf{(b) 一维常数方差高斯策略。} 概率密度为 $\pi_\theta(a \mid s) = \frac{1}{\sqrt{2\pi}\sigma}\exp\!\Big(-\frac{(a - \mu_\theta(s))^2}{2\sigma^2}\Big)$，取对数得
\[
  \ln \pi_\theta(a \mid s) = -\frac{1}{2}\ln(2\pi\sigma^2) - \frac{\big(a - \mu_\theta(s)\big)^2}{2\sigma^2}.
\]
对 $\theta$ 求梯度（第一项为常数，导数为 $0$）：
\[
  \nabla_\theta \ln \pi_\theta(a \mid s) = \mathbf{\frac{a - \mu_\theta(s)}{\sigma^2}\,\nabla_\theta \mu_\theta(s)}.
\]
直观更新方向：在策略梯度上升 $\Delta\theta \propto \hat{A}_t \cdot \frac{a_t - \mu_\theta(s_t)}{\sigma^2}\nabla_\theta \mu_\theta(s_t)$ 中，若 $\hat{A}_t > 0$（动作优于平均）且 $a_t > \mu_\theta(s_t)$，梯度推动均值 $\mu_\theta(s_t)$ 向上朝 $a_t$ 靠拢；若 $\hat{A}_t < 0$，则把均值推离坏动作 $a_t$。同时注意分母 $\sigma^2$：当策略探索标准差 $\sigma$ 过小时，梯度幅值按 $1/\sigma^2$ 剧烈放大，极易引发数值爆炸，因此连续控制代码中通常对 $\ln \sigma$ 施加下界裁剪（如 \texttt{exp3\_ppo\_vs\_unclipped()} 中的 \texttt{clamp(log\_std, -2.0, 1.0)}）。

\solItem{6.2.2}{三步轨迹的 GAE($\lambda$) 手算}

已知 $(r_0, r_1, r_2) = (1, 0, 5)$，$(V(s_0), V(s_1), V(s_2), V(s_3)) = (2.0, 2.5, 4.0, 0)$，$\gamma = 0.9$。

\textbf{(a) 计算单步 TD 误差 $\delta_t = r_t + \gamma V(s_{t+1}) - V(s_t)$。}
\begin{align*}
  \delta_0 &= 1 + 0.9 \times 2.5 - 2.0 = 1 + 2.25 - 2.0 = \mathbf{1.25}, \\
  \delta_1 &= 0 + 0.9 \times 4.0 - 2.5 = 3.6 - 2.5 = \mathbf{1.10}, \\
  \delta_2 &= 5 + 0.9 \times 0 - 4.0 = \mathbf{1.00}.
\end{align*}

\textbf{(b) 按反向递推式~\eqref{eq:gae-rec} $\hat{A}_t = \delta_t + (\gamma\lambda)\hat{A}_{t+1}$（$\hat{A}_3 = 0$）逐项计算：}
\begin{itemize}
  \item \textbf{当 $\lambda = 0$（$\gamma\lambda = 0$，纯 1 步 TD 优势）}：
  \[
    \hat{A}_2 = 1.00, \qquad \hat{A}_1 = 1.10, \qquad \hat{A}_0 = 1.25.
  \]
  \item \textbf{当 $\lambda = 0.5$（$\gamma\lambda = 0.45$）}：
  \begin{align*}
    \hat{A}_2 &= \delta_2 = \mathbf{1.00}, \\
    \hat{A}_1 &= \delta_1 + 0.45 \times 1.00 = 1.10 + 0.45 = \mathbf{1.55}, \\
    \hat{A}_0 &= \delta_0 + 0.45 \times 1.55 = 1.25 + 0.6975 = \mathbf{1.9475}.
  \end{align*}
  \item \textbf{当 $\lambda = 1$（$\gamma\lambda = 0.90$）}：
  \begin{align*}
    \hat{A}_2 &= 1.00, \\
    \hat{A}_1 &= 1.10 + 0.90 \times 1.00 = \mathbf{2.00}, \\
    \hat{A}_0 &= 1.25 + 0.90 \times 2.00 = \mathbf{3.05}.
  \end{align*}
\end{itemize}
\textbf{验算 $\lambda = 1$ 与 $G_t - V(s_t)$ 的等价性}：直接计算蒙特卡洛回报：
\[
  G_2 = 5.0, \quad G_1 = 0 + 0.9 \times 5.0 = 4.5, \quad G_0 = 1 + 0.9 \times 4.5 = 5.05.
\]
于是 $G_2 - V(s_2) = 5.0 - 4.0 = 1.00$、$G_1 - V(s_1) = 4.5 - 2.5 = 2.00$、$G_0 - V(s_0) = 5.05 - 2.0 = 3.05$，三者与 $\lambda = 1$ 时的 $\hat{A}_t^{\mathrm{GAE}}$ 逐位完全一致（裂项相消恒等式的精确验证）。

\solItem{6.2.3}{单步标量策略的最优最小方差基线}

\textbf{(a) 严格凸二次性证明。} 首先，因 $\E[u(a)] = \sum_a \pi_\theta(a) \frac{\partial_\theta \pi_\theta(a)}{\pi_\theta(a)} = \partial_\theta \sum_a \pi_\theta(a) = 0$，估计量的期望为
\[
  \E[g(b)] = \E[u(a) R(a)] - b\,\E[u(a)] = \E[u(a) R(a)],
\]
完全不依赖于基线 $b$。因此方差展开为
\[
  \Var(g(b)) = \E[g(b)^2] - (\E[g(b)])^2 = \E[u(a)^2]\,b^2 - 2\,\E[u(a)^2 R(a)]\,b + \E[u(a)^2 R(a)^2] - \big(\E[u(a)R(a)]\big)^2.
\]
求一阶与二阶导数：
\[
  \frac{d}{db}\Var(g(b)) = 2\,\E[u(a)^2]\,b - 2\,\E[u(a)^2 R(a)], \qquad
  \frac{d^2}{db^2}\Var(g(b)) = 2\,\E[u(a)^2] > 0.
\]
故 $\Var(g(b))$ 是关于 $b$ 的严格凸二次函数，存在唯一全局极小值点。

\textbf{(b) 最优基线 $b^*$ 与状态价值 $\E[R(a)]$ 的比较。} 令一阶导数为零，立即解得
\[
  b^* = \mathbf{\frac{\E[u(a)^2 R(a)]}{\E[u(a)^2]}}.
\]
可见 $b^*$ 本质上是按得分函数平方 $u(a)^2$ 加权的奖励期望，而常规状态价值基线 $V = \E[R(a)]$ 是按策略概率等权平均。二者之差为 $b^* - \E[R(a)] = \frac{\operatorname{Cov}(u(a)^2, R(a))}{\E[u(a)^2]}$。只有当 $u(a)^2$ 与 $R(a)$ 不相关（例如所有动作的奖励均值相等、或对称二元策略下 $|u(a)|$ 恒定）时，$b^*$ 才与 $\E[R(a)]$ 严格重合。工程中普遍直接选用 $V(s) = \E[R(a)]$，是因为在多维神经网络参数 $\theta \in \R^d$ 下，每个参数分量 $i$ 的 $u_i(a)^2$ 不同导致最优基线因维度而异，无法用单一标量共享，而 $V(s)$ 与 $b^*$ 通常极为接近且易于统一拟合。

\solItem{6.2.4}{强化学习中的“因果性降方差”定量验证}

\textbf{1. 期望完全相同。} 由于奖励 $r_0$ 发生在动作 $a_1$ 之前，$r_0$ 与 $(u_1, r_1)$ 相互独立，且 $\E[u_1] = 0$。故
\[
  \E[g_{\mathrm{full}}] = \E[u_1 r_0 + u_1 r_1] = \E[u_1]\,\E[r_0] + \E[u_1 r_1] = 0 + \E[g_{\mathrm{rtg}}] = \E[g_{\mathrm{rtg}}].
\]

\textbf{2. 方差之差的精确推导。} 因为两者期望相同，方差之差即二阶矩之差：
\begin{align*}
  \Var(g_{\mathrm{full}}) - \Var(g_{\mathrm{rtg}})
  &= \E\!\left[u_1^2 (r_0 + r_1)^2\right] - \E\!\left[u_1^2 r_1^2\right] \\
  &= \E\!\left[u_1^2 r_0^2\right] + 2\,\E\!\left[u_1^2 r_0 r_1\right] \\
  &= \mathbf{\E[u_1^2]\,\E[r_0^2] + 2\mu_0\,\E[u_1^2 r_1]},
\end{align*}
其中最后一步再次使用了 $r_0$ 与 $(u_1, r_1)$ 的独立性及 $\E[r_0] = \mu_0$（且 $\E[r_0^2] = \sigma_0^2 + \mu_0^2$）。

\textbf{3. 交叉项与基线的精妙机制。}
\begin{itemize}
  \item 当 $\mu_0 = 0$ 时，交叉项为零，方差差严格化简为 $\E[u_1^2]\,\sigma_0^2 \ge 0$，即剔除历史噪声奖励 $r_0$ 严格降低方差；
  \item 当 $\mu_0 \neq 0$ 时，若不加基线，交叉项 $2\mu_0\,\E[u_1^2 r_1]$ 可正可负——极端情形下（例如 $r_1 > 0$ 且 $r_0$ 恰好是负常数 $-\E[u_1^2 r_1]/\E[u_1^2]$），$r_0$ 歪打正着充当了 $r_1$ 的负基线，反而可能偶然使 $\Var(g_{\mathrm{full}}) < \Var(g_{\mathrm{rtg}})$；
  \item 一旦对未来回报 $r_1$ 减去习题~\textbf{6.2.3} 的最优基线 $b^* = \frac{\E[u_1^2 r_1]}{\E[u_1^2]}$，令 $\tilde{r}_1 = r_1 - b^*$，便自动有 $\E[u_1^2 \tilde{r}_1] = 0$，交叉项彻底消失！此时
  \[
    \Var\!\big(u_1(r_0 + \tilde{r}_1)\big) - \Var\!\big(u_1 \tilde{r}_1\big) = \E[u_1^2]\,(\sigma_0^2 + \mu_0^2) \ge 0
  \]
  对\emph{任意} $\mu_0 \in \R$ 恒成立。这说明“未来回报（Reward-to-Go）”与“状态基线”是相辅相成、缺一不可的降方差黄金搭档。
\end{itemize}

\solItem{6.2.5}{致命三元组的拆解实验（基于例~\ref{ex:deadly-triad}）}

回顾例~\ref{ex:deadly-triad}：$\gamma = 0.9$，$\phi(s_1) = 1, \phi(s_2) = 2$，单步转移 $(s_1 \to s_2)$ 产生增量 $\Delta w_1 = \alpha(2\gamma - 1)w = 0.8\alpha w$，自环 $(s_2 \to s_2)$ 产生增量 $\Delta w_2 = 4\alpha(\gamma - 1)w = -0.4\alpha w$。

\textbf{(a) 去掉“异策略”：收敛采样概率阈值 $p_{\max}$。} 设以概率 $p \in [0, 1]$ 采样 $(s_1 \to s_2)$、以概率 $1 - p$ 采样 $(s_2 \to s_2)$，则条件期望更新为
\[
  \E[w_{t+1} \mid w_t] = \Big\{1 + \alpha\big[p(2\gamma - 1) + 4(1 - p)(\gamma - 1)\big]\Big\} w_t.
\]
代入 $\gamma = 0.9$，括号内的期望漂移系数为
\[
  c(p) = 0.8p - 0.4(1 - p) = 1.2p - 0.4.
\]
要使算法对充分小的步长 $\alpha > 0$ 期望收敛至真值 $w^* = 0$，必须满足 $c(p) < 0$，即
\[
  1.2p - 0.4 < 0 \iff p < \mathbf{\frac{1}{3} \approx 0.3333}.
\]
特别地，在真实马尔可夫链的同策略平稳分布下，智能体离开 $s_1$ 后永久停留在 $s_2$，故 $d(s_1) = 0, d(s_2) = 1$（即 $p = 0 < 1/3$），此时 $c(0) = -0.4 < 0$，只要 $0 < \alpha < 2 / 0.4 = 5$ 即保证指数收敛（事实上由 Tsitsiklis--Van Roy 1997 定理，线性同策略半梯度 TD 在任意有限 MDP 下均保证收敛）。

\textbf{(b) 去掉“自举”：改用蒙特卡洛目标。} 由于该马尔可夫过程所有转移奖励均为 $0$，从任意状态出发的真实完整回报恒为 $G_t \equiv 0$。MC 线性回归更新为 $w_{t+1} = w_t + \alpha\big(0 - w_t \phi(s_t)\big)\phi(s_t) = \big(1 - \alpha \phi(s_t)^2\big) w_t$。取期望：
\[
  \E[w_{t+1} \mid w_t] = \Big(1 - \alpha\big[p \cdot 1^2 + (1 - p)\cdot 2^2\big]\Big) w_t = \big(1 - \alpha(4 - 3p)\big) w_t.
\]
对任意 $p \in [0, 1]$，均有 $4 - 3p \in [1, 4]$ 严格为正！只要步长满足 $0 < \alpha < \frac{2}{4 - 3p}$（对任意 $p$ 取 $\alpha < 0.5$ 即可），收缩因子 $|1 - \alpha(4 - 3p)| < 1$，因此无论采样分布如何异策略，$w_t$ 均以指数速率收敛至 $0$。这证明“致命三元组”缺一不可。

\solItem{6.2.6}{DQN 目标网络更新周期 $C$ 的两难权衡}

\textbf{1. 极端情况 $C = 1$（每步立即同步 $\theta^- \leftarrow \theta$）。} 目标网络彻底失效，式~\eqref{eq:dqn-loss} 退化为无目标网络的原始非线性异策略半梯度 Q-Learning。参数更新立即改变同一批状态下的自举目标 $y_i = r_i + \gamma \max_{a'} Q_\theta(s'_i, a')$，导致目标剧烈漂移、$Q$ 值过估计正反馈与致命三元组震荡发散。

\textbf{2. 极端情况 $C$ 极大（如 $C = 10^6$ 步同步一次）。} 从\textbf{训练稳定性}角度看，在任意两个同步时刻之间的 $C$ 步内，目标 $y_i^{\mathrm{DQN}}$ 完全固定，优化退化为良态的静态有监督回归，绝不会在周期内发散。但从\textbf{偏差传播速度（样本效率）}角度看，奖励信息沿时间逆向回传遵循动态规划的 Bellman 算子迭代节奏：假设任务从起点到奖励终点需要 $H$ 步，每完成 1 次目标网络硬同步，远端奖励信息最多只能沿时间反向传播 1 步！因此至少需要 $H \cdot C$ 个梯度步（当 $H = 50, C = 10^6$ 时高达 $5 \times 10^7$ 步）才能让起点的 $Q$ 值感知到终点奖励，训练几乎停滞不前，且在漫长周期内过度拟合陈旧的过期目标。

\textbf{3. 工程权衡。} 因此 $C$ 必须落在“足以让在线网络大致拟合当前静止目标、又不至于拖慢贝尔曼回传”的甜点区（Atari 经典 DQN 取 $C = 10^4$ 帧），或改用连续平滑的 Polyak 软更新 $\theta^- \leftarrow \tau \theta + (1 - \tau)\theta^-$（$\tau \approx 0.005$），相当于有效滞后窗口约为 $1/\tau \approx 200$ 步。

\solItem{6.2.7}{编程实验：验证 \texttt{.detach()} 缺失与奖励平移对 REINFORCE 的冲击}

\textbf{(a) 奖励常量平移对比（偏移量 $0.0$ vs $+4.0$ vs $+20.0$）。} 设奖励偏移为 $c$，12 步折扣回报随之平移 $\Delta G_t = c \frac{1 - 0.95^{12-t}}{1 - 0.95} \approx 9.2c$（在 $t=0$ 处）。
\begin{itemize}
  \item \textbf{无基线 REINFORCE}：由式~\eqref{eq:pg-raw}，策略梯度方差随 $c^2$ 平方级暴涨。当 $c = 0.0$ 时回报尚有正有负，无基线版本仍能勉强学出较好策略；当 $c = +4.0$ 时方差显著恶化（配套实测终段均回报落后带基线版本一大截）；当 $c = +20.0$ 时，好坏动作的 $G_t$ 均在 $+180$ 左右、相对差异不足 $3\%$，所有动作都被盲目大幅拉高概率，梯度几乎被纯噪声淹没，终段有效位置得分接近随机策略；
  \item \textbf{带基线 REINFORCE}：Critic 网络在最初十几回合内即可把均方误差中的常量偏移 $c \frac{1 - 0.95^{12-t}}{1 - 0.95}$ 精确吸收进偏置项，此后优势 $A_t = G_t - V_\phi(s_t)$ 对 $c$ 完全不变，因此在 $c \in \{0.0, 4.0, 20.0\}$ 三种设定下除去常量基准后的真实跟踪误差几乎完全重合！
\end{itemize}

\textbf{(b) 合并优化器并删去 \texttt{vals.detach()} 的崩坏机制。} 若删去 \texttt{.detach()} 且联合最小化总损失
\[
  \mathcal{L}_{\mathrm{tot}} = -\frac{1}{T}\sum_{t} \ln \pi_\theta(a_t \mid s_t)\big(G_t - V_\phi(s_t)\big) + \mathcal{L}_{\mathrm{MSE}}(V_\phi, G),
\]
则 Critic 参数 $\phi$ 会额外收到来自策略损失的虚假梯度：
\[
  \frac{\partial \mathcal{L}_{\mathrm{actor}}}{\partial V_\phi(s_t)} = +\frac{1}{T}\ln \pi_\theta(a_t \mid s_t) < 0.
\]
由于对数概率 $\ln \pi_\theta(a_t \mid s_t)$ 恒为负值，最小化 $\mathcal{L}_{\mathrm{tot}}$ 会驱动 Critic 不断\textbf{虚高抬升 $V_\phi(s_t)$}（甚至冲向正无穷）以人为拉低 $\mathcal{L}_{\mathrm{actor}}$！结果 Critic 的 MSE 损失剧烈飙升、$V_\phi(s_t)$ 彻底偏离真实回报，导致所有动作的优势 $G_t - V_\phi(s_t)$ 全变为巨大负数，Actor 策略随之崩坏。

\solItem{6.2.8}{从普通 Actor-Critic 走向多轮样本复用的困境预览}

若在 \texttt{ac} 分支中把 \texttt{loss\_p = -(logps\_t * adv).mean()} 对同一批旧轨迹循环执行 10 次 \texttt{opt\_p.step()}，实测会迅速观察到：策略熵急剧坍缩至接近 $0$（输出概率趋近独热分布），随后回合回报剧烈跳水并再也无法恢复。原因在于两重根本性数学破坏：
\begin{enumerate}
  \item \textbf{分布失配（失去无偏性）}：策略梯度定理~\ref{thm:pg-theorem} 要求期望下标 $a_t \sim \pi_\theta(\cdot \mid s_t)$ 与求导参数 $\theta$ 严格一致。第 1 步更新后，$\theta$ 已变为 $\theta_1 \neq \theta_{\mathrm{old}}$，继续对旧轨迹求 $\nabla_\theta \ln \pi_\theta(a_t \mid s_t) \hat{A}_t$ 却不乘以重要性采样权重 $r_t(\theta) = \frac{\pi_\theta(a_t \mid s_t)}{\pi_{\theta_{\mathrm{old}}}(a_t \mid s_t)}$，梯度已不再指向目标 $J(\theta)$ 的上升方向；
  \item \textbf{无界过优化与策略坍缩}：旧轨迹中某个偶然抽中且 $\hat{A}_t > 0$ 的动作会被连续拉高 10 步直至 $\pi_\theta(a_t \mid s_t) \to 1$；而某个 $\hat{A}_t < 0$ 的动作其 $\ln \pi_\theta$ 会被狂压向 $-\infty$，由于 $\ln p$ 在 $p \to 0^+$ 处斜率发散，哪怕步长很小也会产生灾难性梯度，彻底摧毁所有状态的 logits。
\end{enumerate}
要安全地实现多轮样本复用，必须在数学上同时引入第~\ref{ch:trpo-ppo-grpo} 章的两大支柱：\textbf{重要性采样概率比 $r_t(\theta)$}（修正分布偏移）与\textbf{信任域 KL 约束或 PPO 裁剪边界 $[1-\epsilon, 1+\epsilon]$}（封顶单批轨迹允许引发的策略位移）。

% =====================================================================
\section*{第 3 章习题解答（共 8 题）}
\addcontentsline{toc}{section}{第 3 章习题解答（共 8 题）}

\solItem{6.3.1}{PPO-Clip 六种边界情形逐一验算}

已知 $\epsilon = 0.2$，信任区间为 $[0.8, 1.2]$，目标函数 $f(r_t; \hat{A}_t) = \min\big(r_t \hat{A}_t,\; \clip(r_t, 0.8, 1.2)\hat{A}_t\big)$。对照命题~\ref{prop:ppo-cases} 逐一计算 6 个样本的精确结果如下表：

\begin{center}
\footnotesize
\setlength{\tabcolsep}{4.5pt}
\begin{tabular}{@{}ccccccc@{}}
\toprule
样本 & $(r_t, \hat{A}_t)$ & 未裁剪 $r_t \hat{A}_t$ & 裁剪 $\clip(r_t)\hat{A}_t$ & $f(r_t; \hat{A}_t)$ & $\partial f / \partial r_t$ & 是否截断 \\
\midrule
(1) & $(0.6, +2.0)$ & $+1.20$ & $0.8 \times 2.0 = +1.60$ & $\mathbf{+1.20}$ & $+2.0$ & 否（保留纠偏） \\
(2) & $(1.1, +2.0)$ & $+2.20$ & $1.1 \times 2.0 = +2.20$ & $\mathbf{+2.20}$ & $+2.0$ & 否（信任域内） \\
(3) & $(1.5, +2.0)$ & $+3.00$ & $1.2 \times 2.0 = +2.40$ & $\mathbf{+2.40}$ & $\mathbf{0}$ & \textbf{是}（防贪婪过推） \\
(4) & $(0.6, -2.0)$ & $-1.20$ & $0.8 \times (-2.0) = -1.60$ & $\mathbf{-1.60}$ & $\mathbf{0}$ & \textbf{是}（防过度压杀） \\
(5) & $(1.1, -2.0)$ & $-2.20$ & $1.1 \times (-2.0) = -2.20$ & $\mathbf{-2.20}$ & $-2.0$ & 否（信任域内） \\
(6) & $(1.5, -2.0)$ & $-3.00$ & $1.2 \times (-2.0) = -2.40$ & $\mathbf{-3.00}$ & $-2.0$ & 否（保留纠偏） \\
\bottomrule
\end{tabular}
\end{center}

\textbf{核心观察}：样本 (1) 与样本 (6) 虽已落在 $[0.8, 1.2]$ 之外，但外层 $\min$ 选中的恰好是\textbf{未裁剪分支} $r_t \hat{A}_t$，从而保留了非零梯度把被误压的好动作拉回、把被误抬的坏动作压回；唯有样本 (3) 与 (4) 属于“沿优势方向已走过头”，梯度被安全归零。

\solItem{6.3.2}{GRPO 群体优势与 $k_3$ KL 估计量手算}

\textbf{(a) GRPO 组内相对优势。} 给定组内奖励 $(R_1, R_2, R_3, R_4) = (1, 0, 1, 0)$：
\[
  \bar{R} = \frac{1 + 0 + 1 + 0}{4} = \mathbf{0.5}, \qquad
  \sigma_R = \sqrt{\frac{(1-0.5)^2 + (0-0.5)^2 + (1-0.5)^2 + (0-0.5)^2}{4}} = \sqrt{0.25} = \mathbf{0.5}.
\]
代入式~\eqref{eq:grpo-adv} $\hat{A}_i = (R_i - \bar{R}) / \sigma_R$，得
\[
  (\hat{A}_1, \hat{A}_2, \hat{A}_3, \hat{A}_4) = \mathbf{(+1.0,\; -1.0,\; +1.0,\; -1.0)}.
\]

\textbf{(b) 朴素 $k_1$ 与 Schulman 非负 $k_3$ KL 估计量对比。} 记 $\rho = \pi_{\mathrm{ref}} / \pi_\theta$，则 $k_1 = \ln(\pi_\theta / \pi_{\mathrm{ref}}) = -\ln \rho$，$k_3 = \rho - \ln \rho - 1$：
\begin{itemize}
  \item \textbf{情形一（$\pi_\theta = 0.5, \pi_{\mathrm{ref}} = 0.25$）}：$\rho = 0.5$，于是
  \[
    k_1 = \ln 2 \approx \mathbf{0.6931}, \qquad
    k_3 = 0.5 - \ln(0.5) - 1 = \ln 2 - 0.5 \approx \mathbf{0.1931} > 0.
  \]
  \item \textbf{情形二（$\pi_\theta = 0.25, \pi_{\mathrm{ref}} = 0.5$）}：$\rho = 2.0$，于是
  \[
    k_1 = \ln(0.5) = -\ln 2 \approx \mathbf{-0.6931} < 0, \qquad
    k_3 = 2 - \ln 2 - 1 = 1 - \ln 2 \approx \mathbf{0.3069} > 0.
  \]
\end{itemize}
由基本不等式 $e^x \ge 1 + x$（令 $x = \ln \rho$）即知 $\rho - \ln \rho - 1 \ge 0$ 对任意 $\rho > 0$ 逐点恒成立，且 $\E_{a \sim \pi_\theta}[k_3] = \sum_a \pi_{\mathrm{ref}}(a) - 1 + \E_{a \sim \pi_\theta}[\ln(\pi_\theta / \pi_{\mathrm{ref}})] = \DKL(\pi_\theta \| \pi_{\mathrm{ref}})$，因此 $k_3$ 兼具\textbf{无偏性、逐点非负性与低方差}三大优势。

\solItem{6.3.3}{KL 散度的 Hessian 严格等于 Fisher 信息矩阵}

为简洁记 $\pi_\theta(a) = \pi_\theta(a \mid s)$，在归一化恒等式 $\sum_a \pi_\theta(a) \equiv 1$ 两端连续对 $\theta$ 求一次和两次梯度（允许微分与求和交换）：
\begin{equation*}
  \sum_a \nabla_\theta \pi_\theta(a) = \mathbf{0}, \qquad
  \sum_a \nabla_\theta^2 \pi_\theta(a) = \mathbf{0}.
\end{equation*}

\textbf{(a) 一阶梯度在 $\theta = \theta_0$ 处为零。} 展开 $D(\theta) = \sum_a \pi_{\theta_0}(a) \ln \pi_{\theta_0}(a) - \sum_a \pi_{\theta_0}(a) \ln \pi_\theta(a)$，第一项与 $\theta$ 无关，求导得
\[
  \nabla_\theta D(\theta) = -\sum_a \pi_{\theta_0}(a)\,\nabla_\theta \ln \pi_\theta(a) = -\sum_a \pi_{\theta_0}(a)\,\frac{\nabla_\theta \pi_\theta(a)}{\pi_\theta(a)}.
\]
代入 $\theta = \theta_0$，分母与权重精确相消：
\[
  \nabla_\theta D(\theta)\big|_{\theta=\theta_0} = -\sum_a \nabla_\theta \pi_{\theta_0}(a) = \mathbf{0}.
\]

\textbf{(b) 二阶 Hessian 等于 Fisher 信息矩阵。} 对上式再次求导得
\[
  \nabla_\theta^2 D(\theta) = -\sum_a \pi_{\theta_0}(a)\,\nabla_\theta^2 \ln \pi_\theta(a).
\]
利用二阶对数导数恒等式
\[
  \nabla_\theta^2 \ln \pi_\theta(a)
  = \nabla_\theta\!\left(\frac{\nabla_\theta \pi_\theta(a)}{\pi_\theta(a)}\right)
  = \frac{\nabla_\theta^2 \pi_\theta(a)}{\pi_\theta(a)} - \frac{\nabla_\theta \pi_\theta(a)\,\nabla_\theta \pi_\theta(a)^\top}{\pi_\theta(a)^2}
  = \frac{\nabla_\theta^2 \pi_\theta(a)}{\pi_\theta(a)} - \nabla_\theta \ln \pi_\theta(a)\,\nabla_\theta \ln \pi_\theta(a)^\top.
\]
代入 $\theta = \theta_0$ 并对 $a \sim \pi_{\theta_0}$ 求期望，第一项化为 $-\sum_a \nabla_\theta^2 \pi_{\theta_0}(a) = \mathbf{0}$，故立得
\[
  \nabla_\theta^2 D(\theta)\big|_{\theta=\theta_0}
  = -\E_{a \sim \pi_{\theta_0}}\!\left[\nabla_\theta^2 \ln \pi_\theta(a)\right]\Big|_{\theta=\theta_0}
  = \E_{a \sim \pi_{\theta_0}}\!\left[\nabla_\theta \ln \pi_{\theta_0}(a)\,\nabla_\theta \ln \pi_{\theta_0}(a)^\top\right] = F(\theta_0).
\]

\textbf{(c) 反向 KL 散度 $\tilde{D}(\theta) = \sum_a \pi_\theta(a)\ln\frac{\pi_\theta(a)}{\pi_{\theta_0}(a)}$ 具有相同的二阶 Hessian。} 对 $\tilde{D}(\theta)$ 求一阶梯度：
\[
  \nabla_\theta \tilde{D}(\theta)
  = \sum_a \nabla_\theta \pi_\theta(a)\left(\ln\frac{\pi_\theta(a)}{\pi_{\theta_0}(a)} + 1\right)
  = \sum_a \nabla_\theta \pi_\theta(a)\ln\frac{\pi_\theta(a)}{\pi_{\theta_0}(a)},
\]
其中用到了 $\sum_a \nabla_\theta \pi_\theta(a) = \mathbf{0}$。在 $\theta = \theta_0$ 处 $\ln(\pi_{\theta_0}/\pi_{\theta_0}) = 0$，故 $\nabla_\theta \tilde{D}(\theta_0) = \mathbf{0}$。再次对 $\theta$ 求导：
\[
  \nabla_\theta^2 \tilde{D}(\theta)
  = \sum_a \nabla_\theta^2 \pi_\theta(a)\ln\frac{\pi_\theta(a)}{\pi_{\theta_0}(a)}
  + \sum_a \nabla_\theta \pi_\theta(a)\,\frac{\nabla_\theta \pi_\theta(a)^\top}{\pi_\theta(a)}.
\]
代入 $\theta = \theta_0$，第一项因对数为 $0$ 而严格消失，第二项化为
\[
  \nabla_\theta^2 \tilde{D}(\theta)\big|_{\theta=\theta_0}
  = \sum_a \pi_{\theta_0}(a)\,\nabla_\theta \ln \pi_{\theta_0}(a)\,\nabla_\theta \ln \pi_{\theta_0}(a)^\top = F(\theta_0).
\]
由此证明：在局部二阶近似下，正、反向 KL 散度完全等价，均等于 $\frac{1}{2}\Delta\theta^\top F(\theta_0)\Delta\theta$（差异仅出现在三阶及更高阶项）。

\solItem{6.3.4}{TRPO 自然梯度步长的参数重参数化不变性}

\textbf{(a) 拉格朗日乘子法推导闭式解。} 在约束 $\frac{1}{2}\Delta\theta^\top F \Delta\theta \le \delta$（$F \succ 0, \delta > 0, g \neq \mathbf{0}$）下最大化线性目标 $g^\top \Delta\theta$。因目标是线性的、可行域是紧凸椭球，最大值必在边界取到。构造拉格朗日函数（乘子 $\lambda > 0$）：
\[
  \mathcal{L}(\Delta\theta, \lambda) = g^\top \Delta\theta - \lambda\left(\frac{1}{2}\Delta\theta^\top F \Delta\theta - \delta\right).
\]
令 $\nabla_{\Delta\theta}\mathcal{L} = g - \lambda F \Delta\theta = \mathbf{0}$，得 $\Delta\theta = \frac{1}{\lambda} F^{-1} g$。代入紧约束边界 $\frac{1}{2}\Delta\theta^\top F \Delta\theta = \delta$：
\[
  \frac{1}{2\lambda^2} g^\top F^{-1} F F^{-1} g = \frac{1}{2\lambda^2} g^\top F^{-1} g = \delta
  \implies \lambda = \sqrt{\frac{g^\top F^{-1} g}{2\delta}} > 0.
\]
回代即得式~\eqref{eq:npg-step}：
\[
  \Delta\theta^* = \sqrt{\frac{2\delta}{g^\top F^{-1} g}}\,F^{-1} g.
\]

\textbf{(b) 可逆线性重参数化 $\psi = A\theta$ 下的严格不变性。} 记 $\tilde{\pi}_\psi = \pi_{A^{-1}\psi}$。由多变量链式法则，对任意标量函数 $h(\theta) = \tilde{h}(A\theta)$ 有 $\nabla_\theta h = A^\top \nabla_\psi \tilde{h}$，故 $\nabla_\psi \tilde{h} = A^{-\top} \nabla_\theta h$。据此立即得到新坐标系下的梯度与 Fisher 矩阵：
\begin{align*}
  g_\psi &= \nabla_\psi J = A^{-\top} g, \\
  F_\psi &= \E\!\left[\big(\nabla_\psi \ln \tilde{\pi}_\psi\big)\big(\nabla_\psi \ln \tilde{\pi}_\psi\big)^\top\right]
          = \E\!\left[A^{-\top}(\nabla_\theta \ln \pi_\theta)(\nabla_\theta \ln \pi_\theta)^\top A^{-1}\right]
          = A^{-\top} F A^{-1}.
\end{align*}
对 $F_\psi$ 求逆得 $F_\psi^{-1} = A F^{-1} A^\top$。于是自然梯度方向与二次型标量分别为：
\begin{align*}
  F_\psi^{-1} g_\psi &= (A F^{-1} A^\top)(A^{-\top} g) = A F^{-1} g, \\
  g_\psi^\top F_\psi^{-1} g_\psi &= (g^\top A^{-1})(A F^{-1} g) = g^\top F^{-1} g.
\end{align*}
代入闭式步长公式，严格得出
\[
  \Delta\psi^* = \sqrt{\frac{2\delta}{g_\psi^\top F_\psi^{-1} g_\psi}}\,F_\psi^{-1} g_\psi
  = A \left(\sqrt{\frac{2\delta}{g^\top F^{-1} g}}\,F^{-1} g\right) = A\,\Delta\theta^*.
\]
因此新旧坐标系下的更新点严格满足 $\psi + \Delta\psi^* = A(\theta + \Delta\theta^*)$，对应完全相同的概率分布！相比之下，普通欧氏梯度 $\Delta\psi_{\mathrm{GD}} = \eta g_\psi = \eta A^{-\top} g \neq A(\eta g)$（除非 $A$ 是正交矩阵），说明欧氏步长严重依赖参数尺度，而自然梯度直接在概率分布流形的黎曼几何上移动。

\solItem{6.3.5}{反例构造：删去 PPO-Clip 外层 $\min$ 的灾难后果}

\textbf{(a) 删去 $\min$ 后的分段函数及其与真实 PPO-Clip 的分歧区间。} 记 $\tilde{f}(r; A) = \clip(r, 1-\epsilon, 1+\epsilon)\,A$：
\begin{itemize}
  \item 当 $A > 0$ 时：$\tilde{f}$ 在 $r < 1-\epsilon$ 为水平平台 $(1-\epsilon)A$、在 $[1-\epsilon, 1+\epsilon]$ 为斜率 $+A$ 的直线、在 $r > 1+\epsilon$ 为水平平台 $(1+\epsilon)A$。\textbf{分歧区间在 $r < 1-\epsilon$}：真实 PPO-Clip 保持斜率 $+A$ 继续向下延伸（惩罚把好动作降得过低，从而产生拉回梯度），而 $\tilde{f}$ 将其拉平、梯度归零；
  \item 当 $A < 0$ 时：$\tilde{f}$ 在 $r < 1-\epsilon$ 为水平平台 $(1-\epsilon)A$、在 $[1-\epsilon, 1+\epsilon]$ 为斜率 $A < 0$ 的下降直线、在 $r > 1+\epsilon$ 为水平平台 $(1+\epsilon)A$。\textbf{分歧区间在 $r > 1+\epsilon$}：真实 PPO-Clip 保持斜率 $A < 0$ 继续暴跌（严厉惩罚把坏动作抬高），而 $\tilde{f}$ 把损失封底在常数 $(1+\epsilon)A$、梯度归零！
\end{itemize}

\textbf{(b) 两样本一维参数反例。} 构造两个独立状态 $s_1, s_2$，各自对数概率由同一标量参数 $\theta$ 线性控制（设当前旧参数 $\theta_{\mathrm{old}} = 0$，样本 1 有 $\hat{A}_1 = +1$、$\ln\pi_\theta(a_1\mid s_1) = \ln p_0 + \theta$，样本 2 有 $\hat{A}_2 = -1$、$\ln\pi_\theta(a_2\mid s_2) = \ln p_0 + 2\theta$，取 $\epsilon = 0.2$）。于是概率比为 $r_1(\theta) = e^\theta, r_2(\theta) = e^{2\theta}$。

假设某步较大梯度上升把参数推到 $\theta_1 = \frac{1}{2}\ln 1.5 \approx 0.2027$，此时：
\[
  r_2(\theta_1) = 1.5 > 1.2, \qquad r_1(\theta_1) = \sqrt{1.5} \approx 1.2247 > 1.2.
\]
或者考虑更危险的耦合：$r_1(\theta_1) = 1.1 \in [0.8, 1.2]$（好动作尚在信任域内），而坏动作已被推到 $r_2(\theta_1) = 1.5 > 1.2$（例如 $\ln\pi_\theta(a_1\mid s_1) = \ln p_0 + 0.5\theta$、$\ln\pi_\theta(a_2\mid s_2) = \ln p_0 + 2.13\theta$，在 $\theta_1 \approx 0.19$ 处有 $r_1 \approx 1.10, r_2 \approx 1.50$）。此时考察删去 $\min$ 的目标 $\tilde{L}(\theta) = \frac{1}{2}\big[\clip(r_1)\cdot(+1) + \clip(r_2)\cdot(-1)\big]$：
\begin{itemize}
  \item 对坏动作样本 2，因为 $r_2 = 1.5 > 1.2$，$\clip(r_2, 0.8, 1.2)(-1) \equiv -1.2$ 为常数，其梯度 $\frac{\partial}{\partial \theta}\tilde{f}_2 = \mathbf{0}$——目标函数对“坏动作概率飙升 $1.5\times$”完全装瞎；
  \item 对好动作样本 1，仍有正梯度 $+0.5 r_1 > 0$，于是下一步梯度上升会继续增大 $\theta$，把坏动作 $r_2$ 进一步推向 $2.0, 5.0, \dots$ 直至彻底失控！
\end{itemize}
而保留外层 $\min$ 后，真实 PPO-Clip 在 $r_2 = 1.5$ 处命中未裁剪分支 $f_2 = -r_2 = -1.5$，立即产生强烈的负向纠偏梯度 $-2.13 r_2 = -3.195$，迅速把参数拉回安全信任域。

\solItem{6.3.6}{GRPO 在二元奖励 $R_i \in \{0, 1\}$ 下的难度加权效应分析}

\textbf{(a) 归一化优势的精确闭合式。} 当题目通过率为 $p \in (0, 1)$ 且 $G$ 充分大时，$\bar{R} \approx p$，$\sigma_R \approx \sqrt{p(1 - p)}$。代入式~\eqref{eq:grpo-adv}：
\[
  \hat{A}_+ = \frac{1 - p}{\sqrt{p(1 - p)}} = \mathbf{\sqrt{\frac{1 - p}{p}}}, \qquad
  \hat{A}_- = \frac{0 - p}{\sqrt{p(1 - p)}} = \mathbf{-\sqrt{\frac{p}{1 - p}}}.
\]
显然可检验零均值与单位方差：$p \hat{A}_+ + (1 - p)\hat{A}_- = \sqrt{p(1-p)} - \sqrt{p(1-p)} = 0$，$p\hat{A}_+^2 + (1-p)\hat{A}_-^2 = (1-p) + p = 1$。

\textbf{(b) 单道题的总正权重与极端难度下的梯度冲击。} 单道题上，所有答对样本（占比 $p$）贡献的总正权重为
\[
  W_+(p) = p \cdot \hat{A}_+ = \mathbf{\sqrt{p(1 - p)}}
\]
（所有答错样本贡献的总负权重恰好为 $-W_+(p) = -\sqrt{p(1-p)}$）。
\begin{itemize}
  \item \textbf{整题权重视角}：$W_+(p) = \sqrt{p(1-p)}$ 在中等难度 $p = 0.5$ 处达到最大 $0.5$，在极难 $p \to 0^+$ 或极易 $p \to 1^-$ 处趋近于 $0$（全错或全对的题目整组优势为零）；
  \item \textbf{单条样本梯度幅值视角}：当 $p = 0.01$（极难题）时，虽然答对样本极少，但单条答对轨迹的优势高达 $\hat{A}_+ = \sqrt{99} \approx \mathbf{+9.95}$！若这条唯一“答对”的轨迹恰好是推理胡扯但末尾蒙对数字的假阳性样本，它会以近 $10\times$ 的剧烈梯度放大污染模型权重；对称地，当 $p = 0.99$（极易题）时，偶尔一条笔误答错样本的优势高达 $\hat{A}_- \approx \mathbf{-9.95}$，引发强烈过度惩罚；
  \item \textbf{对比 Dr. GRPO（去掉除以 $\sigma_R$）}：若改用无标准差归一化的 $\tilde{A}_i = R_i - \bar{R}$，则 $\tilde{A}_+ = 1 - p \in (0, 1)$、$\tilde{A}_- = -p \in (-1, 0)$，单条轨迹的优势幅值严格被锁在 $[-1, 1]$ 内，不再在 $p \to 0^+$ 或 $p \to 1^-$ 处发散，且整题总权重变为 $p(1 - p)$，更平滑地抑制极端题目的噪声梯度。
\end{itemize}

\solItem{6.3.7}{编程实验：PPO 裁剪阈值 $\epsilon$ 与复用轮数 $K$ 的相图扫描}

\textbf{(a) 扫描 $\epsilon \in \{0.05, 0.1, 0.2, 0.4, 0.8, \infty\}$（固定 $K = 10$）的相图规律。} 在 \texttt{exp3\_ppo\_vs\_unclipped()} 的 1D 连续高斯 LQR 任务（初始回报约 $-12.6$，理想平稳最优约 $-3.0$）上运行扫描：
\begin{itemize}
  \item $\epsilon = 0.05$：信任域极窄，中后期几乎零回撤，但每个 epoch 极早触发梯度截断，前 10 轮爬升明显迟缓，样本效率偏低；
  \item $\epsilon = 0.1\text{--}0.2$：\textbf{黄金甜点区}——前 10 轮能够充分借助 $K = 10$ 次复用快速爬升，同时严格封住大幅概率比偏移，末 10 轮均值稳定在约 $-3.0\text{--}-3.3$ 且全程无剧烈回撤；
  \item $\epsilon = 0.4\text{--}0.8$：前几轮爬升略快，但中后期单批高方差优势偶尔推动 $r_t$ 突破 $1.5$ 以上，开始出现显著的局部回撤尖峰；
  \item $\epsilon = \infty$（等价于未裁剪 IS-PG）：在 $K = 10$ 下彻底失控，高斯标准差与均值被反复过推，出现灾难性回撤（配套实测最差跌至 $-10$ 以下，末 10 轮均值远差于 $\epsilon = 0.2$）。
\end{itemize}

\textbf{(b) 加入 KL 早停（\texttt{max\_kl = 0.05}）的实测现象。} 在内层 10 轮循环中，每次前向计算近似 KL $\widehat{\DKL} = \operatorname{mean}\big((r_t - 1) - \ln r_t\big)$（即非负 $k_3$ 估计量），若 $\widehat{\DKL} > 0.05$ 则 \texttt{break}。实测会发现：训练早期策略梯度大，往往仅执行 $2\text{--}4$ 个 epoch 即触发早停；训练中后期优势幅度变小、策略趋于平稳，可执行 $6\text{--}10$ 个 epoch。KL 早停相当于在 PPO-Clip 的逐样本软保护之外，再套上一道全局分布级硬保险丝。

\solItem{6.3.8}{编程实验：在上下文赌博机上亲手实现无 Critic 的 GRPO}

\textbf{实现骨架与核心逻辑。} 构造 32 个独热或随机向量上下文状态 $s \in \R^{16}$，每个状态对应 8 个离散动作，仅固定 1 个动作奖励 $R = 1$、其余 $7$ 个为 $0$。参考策略 $\pi_{\mathrm{ref}}$ 取均匀分布（$\pi_{\mathrm{ref}}(a \mid s) = 1/8$）。每步对每个状态采样 $G = 8$ 个动作，计算组内均值与标准差得到 $\hat{A}_i = (R_i - \bar{R}) / (\sigma_R + 10^{-8})$（若全对或全错则该组优势自动为 $0$），调用 \texttt{grpo\_loss} 执行 PPO-Clip 更新并叠加 $\beta \cdot k_3(\pi_\theta, \pi_{\mathrm{ref}})$ 正则。

\textbf{不同 $\beta$ 下的理论精确解与实测验证。} 该上下文赌博机的 KL 正则目标在每个状态下具有闭式玻尔兹曼最优策略（正是第 4 章定理~\ref{thm:rlhf-optimal} 的直接特例）：
\[
  \pi_\beta^*(a^* \mid s) = \frac{\frac{1}{8} e^{1/\beta}}{\frac{1}{8} e^{1/\beta} + \frac{7}{8} e^{0}} = \frac{e^{1/\beta}}{e^{1/\beta} + 7}.
\]
代入三种 $\beta$ 取值对照预期指标：
\begin{itemize}
  \item $\beta = 0.0$：无 KL 约束，策略收敛至接近 $100\%$ 确定性选中最优动作，与均匀分布的 KL 散度逼近 $\ln 8 \approx \mathbf{2.079\;\text{nats}}$；
  \item $\beta = 0.05$：$e^{20} \gg 7$，最优动作概率仍接近 $100\%$，KL 接近 $2.08\;\text{nats}$（正则几乎不压低准确率，仅防止 logits 无穷发散）；
  \item $\beta = 0.5$：理论最优准确率停在 $\frac{e^2}{e^2 + 7} \approx \mathbf{51.35\%}$，对应精确 KL 散度约 $\mathbf{0.474\;\text{nats}}$（注意若 GRPO 中优势经过了组内除以 $\sigma_R$ 归一化，实际有效温度会因 $\sigma_R$ 缩放而按 $\beta \sigma_R$ 折算，收敛准确率略高于纯非归一化值）。
\end{itemize}
\textbf{桥接第~\ref{ch:posttrain-rlhf-dpo} 章}：若不再由环境给出绝对标量 $R \in \{0, 1\}$，而是仅给出一对回答的相对人类偏好 $y_w \succ y_l$，把上述玻尔兹曼闭式解反解出隐式奖励 $r(x, y) = \beta \ln\frac{\pi^*(y\mid x)}{\pi_{\mathrm{ref}}(y\mid x)} + \beta \ln Z(x)$ 代入 Bradley--Terry 似然，配分函数 $Z(x)$ 神奇相消，直接诞生了无需采样也无需显式奖励模型的 \textbf{DPO}！

% =====================================================================
\section*{第 4 章习题解答（共 8 题）}
\addcontentsline{toc}{section}{第 4 章习题解答（共 8 题）}

\solItem{6.4.1}{DPO 损失、隐式奖励与错判权重手算}

已知 $\beta = 0.5$，$\ln \pi_{\mathrm{ref}}(y_w \mid x) = -12.0$、$\ln \pi_{\mathrm{ref}}(y_l \mid x) = -10.0$、$\ln \pi_\theta(y_w \mid x) = -11.0$、$\ln \pi_\theta(y_l \mid x) = -8.0$。

\textbf{(a) 隐式奖励计算（式~\eqref{eq:implicit-reward}）。}
\begin{align*}
  \hat{r}_\theta(x, y_w) &= \beta \big(\ln \pi_\theta(y_w \mid x) - \ln \pi_{\mathrm{ref}}(y_w \mid x)\big) = 0.5 \times (-11.0 - (-12.0)) = \mathbf{+0.5}, \\
  \hat{r}_\theta(x, y_l) &= \beta \big(\ln \pi_\theta(y_l \mid x) - \ln \pi_{\mathrm{ref}}(y_l \mid x)\big) = 0.5 \times (-8.0 - (-10.0)) = \mathbf{+1.0}.
\end{align*}

\textbf{(b) 偏好预测概率、DPO 损失与梯度错判权重。} 隐式奖励差为 $\Delta \hat{r} = \hat{r}_w - \hat{r}_l = 0.5 - 1.0 = -0.5$：
\begin{align*}
  P_\theta(y_w \succ y_l \mid x) &= \sigma(-0.5) = \frac{1}{1 + e^{0.5}} \approx \mathbf{0.3775}, \\
  \mathcal{L}_{\mathrm{DPO}} &= -\ln \sigma(-0.5) = \ln(1 + e^{0.5}) \approx \mathbf{0.9741}, \\
  w_\theta &= \sigma(\hat{r}_l - \hat{r}_w) = \sigma(+0.5) = \frac{1}{1 + e^{-0.5}} \approx \mathbf{0.6225}.
\end{align*}
可检验恒等式 $P_\theta(y_w \succ y_l) + w_\theta = \sigma(-0.5) + \sigma(0.5) = 1$。

\textbf{(c) 判别状态与权重直观含义。} 由于 $\hat{r}_l = 1.0 > \hat{r}_w = 0.5$（即模型认为落败回答 $y_l$ 相对参考策略的增幅比胜出回答 $y_w$ 更大），模型当前\textbf{判反了}（预测正确概率仅 $37.75\% < 50\%$）。因此式~\eqref{eq:dpo-grad} 的自适应错判权重 $w_\theta = 1 - P_\theta(y_w \succ y_l) = 0.6225 > 0.5$，自动加大对该错判样本的梯度回传力度，强力拉高 $y_w$ 并压低 $y_l$。

\solItem{6.4.2}{RLHF 闭式最优策略的两动作数值验证}

已知 $\pi_{\mathrm{ref}}(y_1) = 0.2, \pi_{\mathrm{ref}}(y_2) = 0.8$，$r(y_1) = 3.0, r(y_2) = 1.0$。由定理~\ref{thm:rlhf-optimal}，$Z(x) = \sum_i \pi_{\mathrm{ref}}(y_i) e^{r(y_i)/\beta}$，$\pi^*(y_i) = \frac{1}{Z}\pi_{\mathrm{ref}}(y_i) e^{r(y_i)/\beta}$。

\textbf{(a) 两种温度 $\beta$ 下的精确数值计算。}
\begin{enumerate}
  \item \textbf{当 $\beta = 1.0$ 时}：
  \begin{align*}
    Z(x) &= 0.2 e^3 + 0.8 e^1 \approx 0.2 \times 20.085537 + 0.8 \times 2.718282 \approx \mathbf{6.191733}, \\
    \pi^*(y_1) &= \frac{0.2 e^3}{0.2 e^3 + 0.8 e^1} = \frac{e^2}{e^2 + 4} \approx \mathbf{0.648786}, \qquad
    \pi^*(y_2) = 1 - \pi^*(y_1) \approx \mathbf{0.351214}, \\
    \DKL(\pi^* \| \pi_{\mathrm{ref}}) &= 0.648786 \ln\frac{0.648786}{0.2} + 0.351214 \ln\frac{0.351214}{0.8} \\
    &\approx 0.763401 - 0.289045 = \mathbf{0.474356\;\text{nats}}, \\
    \beta \ln Z(x) &= 1.0 \times \ln(6.191733) \approx \mathbf{1.823215}.
  \end{align*}
  \item \textbf{当 $\beta = 0.2$ 时}（奖励放大 $1/\beta = 5$ 倍，故 $r/\beta$ 分别为 $15$ 与 $5$）：
  \begin{align*}
    Z(x) &= 0.2 e^{15} + 0.8 e^{5} \approx 653803.47 + 118.73 \approx \mathbf{6.539222 \times 10^5}, \\
    \pi^*(y_1) &= \frac{e^{10}}{e^{10} + 4} \approx \mathbf{0.999818}, \qquad
    \pi^*(y_2) = \frac{4}{e^{10} + 4} \approx \mathbf{0.000182}, \\
    \DKL(\pi^* \| \pi_{\mathrm{ref}}) &= 0.999818 \ln\frac{0.999818}{0.2} + 0.000182 \ln\frac{0.000182}{0.8} \approx \mathbf{1.607441\;\text{nats}}, \\
    \beta \ln Z(x) &= 0.2 \times \ln(653922.205) \approx 0.2 \times 13.390744 \approx \mathbf{2.678149}.
  \end{align*}
  注意当 $\beta \to 0^+$ 时，$\pi^*(y_1) \to 1$，$\DKL \to \ln(1/0.2) = \ln 5 \approx 1.6094\;\text{nats}$，而 $\beta = 0.2$ 时策略已经几乎完全贪婪化！
\end{enumerate}

\textbf{(b) 目标恒等式精确验算。}
\begin{itemize}
  \item $\beta = 1.0$：$\E_{\pi^*}[r] = 0.648786 \times 3 + 0.351214 \times 1 = 2.297571$，于是
  \[
    \E_{\pi^*}[r] - \beta \DKL = 2.297571 - 1.0 \times 0.474356 = \mathbf{1.823215} = \beta \ln Z(x).
  \]
  \item $\beta = 0.2$：$\E_{\pi^*}[r] = 0.999818 \times 3 + 0.000182 \times 1 = 2.999637$，于是
  \[
    \E_{\pi^*}[r] - \beta \DKL = 2.999637 - 0.2 \times 1.607441 = 2.999637 - 0.321488 = \mathbf{2.678149} = \beta \ln Z(x).
  \]
\end{itemize}
两者精确相等至小数点后六位，完美印证了变分自由能恒等式。

\solItem{6.4.3}{Bradley--Terry 奖励等价类与最优策略的唯一对应}

\textbf{(a) 等价奖励函数诱导相同的偏好分布与最优策略。} 设 $r_2(x, y) = r_1(x, y) + c(x)$。
\begin{enumerate}
  \item \textbf{偏好分布不变性}：代入 Bradley--Terry 式~\eqref{eq:bt-model}，同一提示 $x$ 下两回答的奖励差为
  \[
    r_2(x, y_w) - r_2(x, y_l) = \big(r_1(x, y_w) + c(x)\big) - \big(r_1(x, y_l) + c(x)\big) = r_1(x, y_w) - r_1(x, y_l),
  \]
  故 $P_{r_2}(y_w \succ y_l \mid x) \equiv P_{r_1}(y_w \succ y_l \mid x)$。
  \item \textbf{闭式最优策略不变性}：将 $r_2$ 代入式~\eqref{eq:rlhf-opt-policy}，分子分母同时提取仅依赖 $x$ 的正因子 $e^{c(x)/\beta}$ 并消去：
  \[
    \pi_{r_2}^*(y \mid x)
    = \frac{\pi_{\mathrm{ref}}(y \mid x)\exp\!\left(\frac{r_1(x, y) + c(x)}{\beta}\right)}{\sum_{y'} \pi_{\mathrm{ref}}(y' \mid x)\exp\!\left(\frac{r_1(x, y') + c(x)}{\beta}\right)}
    = \frac{e^{c(x)/\beta}\,\pi_{\mathrm{ref}}(y \mid x)\exp\!\left(\frac{r_1(x, y)}{\beta}\right)}{e^{c(x)/\beta}\sum_{y'} \pi_{\mathrm{ref}}(y' \mid x)\exp\!\left(\frac{r_1(x, y')}{\beta}\right)}
    = \pi_{r_1}^*(y \mid x).
  \]
\end{enumerate}

\textbf{(b) 从策略到奖励等价类的双射（DPO 不损失表达能力）。} 设 $\pi_{\mathrm{ref}}(y \mid x) > 0$ 对一切 $(x, y)$ 成立，$\beta > 0$。
\begin{itemize}
  \item \textbf{存在性}：对任意处处为正的合法策略 $\pi(\cdot \mid x)$，构造代表元奖励函数 $r_\pi(x, y) \triangleq \beta \ln\frac{\pi(y \mid x)}{\pi_{\mathrm{ref}}(y \mid x)}$。检验其配分函数：
  \[
    Z_{r_\pi}(x) = \sum_{y} \pi_{\mathrm{ref}}(y \mid x) \exp\!\left(\frac{r_\pi(x, y)}{\beta}\right) = \sum_y \pi(y \mid x) = \mathbf{1}.
  \]
  因此 $\pi_{r_\pi}^*(y \mid x) = \frac{1}{1}\pi_{\mathrm{ref}}(y \mid x)\frac{\pi(y \mid x)}{\pi_{\mathrm{ref}}(y \mid x)} = \pi(y \mid x)$，说明每个合法策略都至少对应一个奖励函数（且恰是配分函数恒为 $1$ 的规范代表元）。
  \item \textbf{等价类唯一性}：假设两个奖励函数 $r_1, r_2$ 诱导了同一个最优策略 $\pi_{r_1}^* = \pi_{r_2}^* = \pi$，则对所有 $(x, y)$ 有
  \[
    \frac{1}{Z_1(x)}\pi_{\mathrm{ref}}(y \mid x) e^{r_1(x,y)/\beta} = \frac{1}{Z_2(x)}\pi_{\mathrm{ref}}(y \mid x) e^{r_2(x,y)/\beta}.
  \]
  消去 $\pi_{\mathrm{ref}}(y \mid x) > 0$ 并两边取对数再乘以 $\beta$，得
  \[
    r_2(x, y) - r_1(x, y) = \beta \ln\frac{Z_2(x)}{Z_1(x)} \triangleq c(x),
  \]
  仅依赖于 $x$ 而与 $y$ 无关，即 $r_1 \sim r_2$。综上，在给定 $(\pi_{\mathrm{ref}}, \beta)$ 下，\textbf{Bradley--Terry 奖励等价类与合法策略空间之间存在严格一一对应}，因此 DPO 用 $\beta\ln\frac{\pi_\theta}{\pi_{\mathrm{ref}}}$ 重参数化奖励等价类绝不会损失任何表达能力。
\end{itemize}

\solItem{6.4.4}{推理时缩放的优美定理：Best-of-$N$ 的精确 KL 散度公式~\eqref{eq:bon-kl}}

\textbf{(a) 最大分位数 $U^*$ 的概率密度。} 因 $\pi_{\mathrm{ref}}$ 连续无原子且奖励函数几乎处处无结，随机变量 $R = r(Y)$ 的累积分布函数 $F(r)$ 连续严格递增（在支撑集上），由概率积分变换知 $U_i = F(r(Y_i)) \overset{\mathrm{i.i.d.}}{\sim} \mathrm{Uniform}(0, 1)$。由于 $F$ 单调递增，挑选奖励最高者等价于挑选分位数最高者 $U^* = \max(U_1, \dots, U_N)$。对任意 $u \in [0, 1]$，其累积分布函数为
\[
  F_{U^*}(u) = \Prob(U^* \le u) = \prod_{i=1}^N \Prob(U_i \le u) = u^N.
\]
对 $u$ 求导即得概率密度函数：
\[
  p_N(u) = \frac{d}{du}(u^N) = \mathbf{N u^{N-1}}, \qquad u \in [0, 1].
\]

\textbf{(b) 精确 KL 散度积分证明。} 注意 Best-of-$N$ 仅依赖奖励值 $r(Y)$ 的顺序挑选样本，在给定奖励值（或等价地给定分位数 $U = u$）的条件下，Best-of-$N$ 输出 $Y^*$ 的条件分布与原参考策略 $\pi_{\mathrm{ref}}(Y \mid U = u)$ \textbf{完全一致}！由 KL 散度的链式法则：
\begin{align*}
  \DKL(\pi_{\mathrm{BoN}} \| \pi_{\mathrm{ref}})
  &= \DKL\!\big(p_{U^*} \,\|\, p_U\big) + \E_{u \sim p_{U^*}}\!\Big[\underbrace{\DKL\!\big(\pi_{\mathrm{BoN}}(\cdot \mid u) \,\|\, \pi_{\mathrm{ref}}(\cdot \mid u)\big)}_{=\,0}\Big] \\
  &= \int_0^1 p_N(u) \ln\frac{p_N(u)}{1}\,du
  = \int_0^1 N u^{N-1} \ln\!\big(N u^{N-1}\big)\,du.
\end{align*}
拆分被积函数并分部积分计算 $\int_0^1 u^{N-1}\ln u\,du$：
\[
  \int_0^1 u^{N-1} \ln u\,du
  = \left[\frac{u^N}{N}\ln u\right]_0^1 - \int_0^1 \frac{u^N}{N} \cdot \frac{1}{u}\,du
  = 0 - \frac{1}{N}\left[\frac{u^N}{N}\right]_0^1 = -\frac{1}{N^2}.
\]
代回原积分，立即得出与奖励分布完全无关的普适优美公式~\eqref{eq:bon-kl}：
\[
  \DKL(\pi_{\mathrm{BoN}} \| \pi_{\mathrm{ref}})
  = \ln N \underbrace{\int_0^1 N u^{N-1}\,du}_{=\,1} + N(N - 1)\left(-\frac{1}{N^2}\right)
  = \mathbf{\ln N - \frac{N - 1}{N}}.
\]

\solItem{6.4.5}{奖励过优化的数学定量诊断}

已知 $R_{\mathrm{gold}}(d) = d(\alpha_{\mathrm{RM}} - \beta_{\mathrm{RM}} d) = 2.0\,d - 0.25\,d^2$，其中 $d = \sqrt{\DKL(\pi \| \pi_{\mathrm{ref}})} \ge 0$。

\textbf{(a) 真实人类质量峰值与最优 KL 散度。} $R_{\mathrm{gold}}(d)$ 是开口向下的二次函数，令一阶导数为零：
\[
  \frac{d R_{\mathrm{gold}}}{d d} = 2.0 - 0.5\,d = 0 \implies d^* = \mathbf{4.0}.
\]
由此得最优 KL 散度与最高真实质量得分：
\[
  \DKL^* = (d^*)^2 = 4.0^2 = \mathbf{16.0\;\text{nats}}, \qquad
  R_{\mathrm{gold}}^{\max} = 4.0 \times (2.0 - 0.25 \times 4.0) = \mathbf{4.0}.
\]

\textbf{(b) 无约束过优化后果与 Best-of-$N$ 采样数估算。}
\begin{itemize}
  \item 当 $\DKL = 64\;\text{nats}$ 时，$d = \sqrt{64} = 8.0$，代入标度律得
  \[
    R_{\mathrm{gold}}(8.0) = 8.0 \times (2.0 - 0.25 \times 8.0) = \mathbf{0.0}.
  \]
  此时代理奖励模型分数虽一路狂飙，真实人类质量却已把全部收益吐尽归零（若 $\DKL > 64$，真实质量甚至跌为负数，彻底不如微调前的 SFT 底座！）；
  \item 对 Best-of-$N$ 采样，欲使 $\DKL(\pi_{\mathrm{BoN}} \| \pi_{\mathrm{ref}}) = \ln N - \frac{N-1}{N} = 16$，因 $N$ 极大时 $\frac{N-1}{N} \approx 1$，故
  \[
    \ln N \approx 17 \implies N \approx e^{17} \approx \mathbf{2.42 \times 10^7}.
  \]
  需要采样两千四百万条回答才能把 KL 推到过优化拐点 $16\;\text{nats}$！由于 Best-of-$N$ 的 KL 随采样数仅呈对数 $\ln N$ 增长，它在工程中具有天然的抗过优化鲁棒性；而在线 RLHF 仅需几百步梯度更新就可能把 KL 推过 $64\;\text{nats}$，因此必须严格监控 KL 散度与施加 $\beta$ 正则。
\end{itemize}

\solItem{6.4.6}{为什么确定性偏好会让 DPO 绕开 $\beta$ 发生过拟合？}

对唯一无噪声偏好对 $(y_w, y_l)$，记隐式奖励边际为
\[
  z(\theta) = \beta\!\left(\ln\frac{\pi_\theta(y_w \mid x)}{\pi_{\mathrm{ref}}(y_w \mid x)} - \ln\frac{\pi_\theta(y_l \mid x)}{\pi_{\mathrm{ref}}(y_l \mid x)}\right).
\]
DPO 损失函数 $\ell(z) = -\ln \sigma(z) = \ln(1 + e^{-z})$ 在整个实数轴 $z \in \R$ 上严格单调递减且严格大于 $0$，其下确界 $\inf_z \ell(z) = 0$ \textbf{只能在 $z \to +\infty$ 时渐近达到}，在任何有限参数处梯度 $-\sigma(-z) \neq 0$ 永不归零！

因此，只要不加早停或额外正则，无约束优化无论 $\beta > 0$ 取多大（哪怕 $\beta = 10$），都会无休止地把 $z(\theta)$ 推向 $+\infty$。因为 $\ln \pi_\theta(y_w \mid x) \le 0$ 有上界 $0$，让 $z(\theta) \to +\infty$ 的唯一途径就是把落败回答的概率推向零：
\[
  \lim_{t \to \infty} \pi_\theta(y_l \mid x) = \mathbf{0} \qquad (\ln \pi_\theta(y_l \mid x) \to -\infty),
\]
同时 $\DKL(\pi_\theta \| \pi_{\mathrm{ref}}) \to \infty$（若词表支撑集外还有其他回答，甚至可能同时把 $\pi_\theta(y_w \mid x)$ 也拖低，只要 $\pi_\theta(y_l \mid x)$ 跌得更快即可）。可见在确定性经验偏好 $\hat{P}(y_w \succ y_l) = 1$ 下，Bradley--Terry 的对数似然会把真实隐式奖励差推向 $+\infty$，令 $\beta$ 的正则保护形同虚设！

\textbf{IPO（Identity Preference Optimization）的根治原理}：Azar 等人（2024）将交叉熵 $-\ln \sigma(z)$ 替换为有限目标边际的\textbf{二次回归损失}：
\[
  \mathcal{L}_{\mathrm{IPO}}(\theta) = \left[\left(\ln\frac{\pi_\theta(y_w \mid x)}{\pi_{\mathrm{ref}}(y_w \mid x)} - \ln\frac{\pi_\theta(y_l \mid x)}{\pi_{\mathrm{ref}}(y_l \mid x)}\right) - \frac{1}{2\tau}\right]^2.
\]
该损失在有限的对数比差 $\frac{1}{2\tau}$ 处即达到严格全局极小值 $0$ 且梯度恰好归零；继续把 $\pi_\theta(y_l \mid x)$ 推向 $0$ 反而会使损失平方级上升，从而从根本上锁死了 KL 漂移与确定性偏好过拟合。

\solItem{6.4.7}{编程实验：绘制完整的 DPO 奖励—KL 帕累托前沿并监测似然位移}

\textbf{(a) 扫描 $\beta \in \{0.02, 0.05, 0.1, 0.3, 0.6, 1.5, 5.0\}$ 的帕累托前沿特征。} 在 \texttt{exp4\_dpo\_alignment()} 中固定 120 轮 Adam 更新并扫描 $\beta$，以最终 \texttt{KL} 为横轴、真实期望奖励 \texttt{exp\_rew} 为纵轴，将得到一条典型的\textbf{递增凹形帕累托曲线}：
\begin{itemize}
  \item 大 $\beta$ 端（$\beta = 5.0, 1.5$）：隐式奖励差被强力放大，损失迅速饱和，模型几乎不敢离开 $\pi_{\mathrm{ref}}$，最终 KL 极小（$< 0.15\;\text{nats}$），但真实奖励提升受限；
  \item 中间拐角（$\beta \in [0.1, 0.6]$）：以温和的 KL 代价（约 $0.5\text{--}1.4\;\text{nats}$）换取绝大部分真实奖励收益，性价比最高；
  \item 极低 $\beta$ 端（$\beta = 0.02, 0.05$）：由于每对样本仅靠有限噪声偏好拟合，极低 $\beta$ 迫使对数概率差剧烈拉伸才能达到相同边际，KL 飙升至 $2\;\text{nats}$ 以上，且在噪声反转的偏好对上开始出现过拟合，真实奖励边际收益趋零甚至回落。
\end{itemize}

\textbf{(b) 似然位移（Likelihood Displacement）与轻量 SFT 正则效果。} 在训练中额外追踪 $\bar{\ell}_w = \E[\ln \pi_\theta(y_w \mid x)]$ 与 $\bar{\ell}_l = \E[\ln \pi_\theta(y_l \mid x)]$：当 $\beta$ 较小且训练轮数较长时，DPO 仅约束差值 $\bar{\ell}_w - \bar{\ell}_l$ 增大，往往通过“猛压 $\bar{\ell}_l$、同时把概率质量挤向未出现在数据对中的其他回复”来降低损失，导致 $\bar{\ell}_w$ 在初期上升后停滞甚至掉头向下（即著名的\textbf{似然位移}现象）。在 DPO 损失中加入轻量负对数似然正则 $+0.1 \cdot (-\bar{\ell}_w)$（即 DPO+NLL / RPO 技巧）后，直接对胜出回答施加向上的绝对概率锚点，$\bar{\ell}_w$ 保持稳步抬升，既防止了似然崩塌，又进一步稳定了高置信度偏好对的生成质量。

\solItem{6.4.8}{编程实验：在模拟偏好数据上实现 Best-of-$N$ 并验证公式~\eqref{eq:bon-kl}}

取候选回复数 $K_{\mathrm{resp}} = 500$（保证 $K_{\mathrm{resp}} \gg N$ 使独立采样碰撞概率可忽略），参考策略取均匀分布 $\pi_{\mathrm{ref}}(y) = 1 / 500$，给 500 个回复赋予互异的真实奖励值（例如从 $\mathcal{N}(0, 1)$ 采样并固定）。对每个 $N \in \{2, 4, 8, 16, 32\}$，执行 $10^4$ 次蒙特卡洛抽样（每次有放回均匀抽 $N$ 个回复并输出奖励最高者），用频率估计 $\hat{\pi}_{\mathrm{BoN}}$ 并计算经验散度 $\DKL(\hat{\pi}_{\mathrm{BoN}} \| \pi_{\mathrm{ref}})$（或直接按离散无碰撞精确分布计算）。与习题~\textbf{6.4.4} 的解析理论值 $\ln N - \frac{N-1}{N}$ 逐项对照如下表：

\begin{center}
\footnotesize
\begin{tabular}{cccccc}
\toprule
Best-of-$N$ 采样数 $N$ & $2$ & $4$ & $8$ & $16$ & $32$ \\
\midrule
理论精确 KL：$\ln N - \frac{N-1}{N}$（nats） & $0.1931$ & $0.6363$ & $1.2044$ & $1.8351$ & $2.4970$ \\
$K_{\mathrm{resp}} = 500$ 时离散精确 KL & $0.1931$ & $0.6363$ & $1.2044$ & $1.8350$ & $2.4968$ \\
\bottomrule
\end{tabular}
\end{center}
注意细节：用有限样本频率 $\hat{\pi}_{\mathrm{BoN}}$ 估计高维（500 类）分布的 KL 时存在经典的 Miller--Madow 正偏差 $\approx \frac{K_{\mathrm{eff}}-1}{2 M}$（$M = 10^4$ 时约 $+0.025$）；若希望蒙特卡洛验算精确吻合至小数点后两位，应增大抽样次数至 $M = 10^5$ 或在连续均匀分位数 $U^* = \max(U_1, \dots, U_N)$ 上做分箱密度估计，实测与表内理论值完美吻合。

% =====================================================================
\section*{第 5 章习题解答（共 8 题）}
\addcontentsline{toc}{section}{第 5 章习题解答（共 8 题）}

\solItem{6.5.1}{Roofline 算术强度与跨芯片拐点核算}

已知 $P_{\mathrm{peak}} = 2000\,\mathrm{TFLOP/s} = 2.0 \times 10^{15}\,\mathrm{FLOP/s}$，$B_{\mathrm{HBM}} = 8.0\,\mathrm{TB/s} = 8.0 \times 10^{12}\,\mathrm{B/s}$；BF16 线性层 $Y = XW$，其中 $X \in \R^{T \times d_{\mathrm{in}}}, W \in \R^{d_{\mathrm{in}} \times d_{\mathrm{out}}}, Y \in \R^{T \times d_{\mathrm{out}}}$，$d_{\mathrm{in}} = 4096, d_{\mathrm{out}} = 16384$。

\textbf{(a) 硬件拐点 $I^*$。} 由式~\eqref{eq:ridge-point}：
\[
  I^* = \frac{P_{\mathrm{peak}}}{B_{\mathrm{HBM}}} = \frac{2000\,\mathrm{TFLOP/s}}{8.0\,\mathrm{TB/s}} = \mathbf{250\,\mathrm{FLOP/B}}.
\]

\textbf{(b) 非方阵线性层的精确算术强度 $\mathrm{AI}(T)$。}
\begin{itemize}
  \item 总浮点运算数（每个输出元素做 $d_{\mathrm{in}}$ 次乘加，即 $2 d_{\mathrm{in}}$ 次 FLOP）：$F = 2 T d_{\mathrm{in}} d_{\mathrm{out}}$；
  \item 最低 HBM 访存字节数（BF16 每元素 2 字节，读 $X$、读 $W$、写 $Y$）：$Q = 2\big(T d_{\mathrm{in}} + d_{\mathrm{in}} d_{\mathrm{out}} + T d_{\mathrm{out}}\big)$。
\end{itemize}
作比消去因子 $2$，得精确算术强度公式：
\[
  \mathrm{AI}(T) = \frac{F}{Q} = \mathbf{\frac{T\,d_{\mathrm{in}} d_{\mathrm{out}}}{d_{\mathrm{in}} d_{\mathrm{out}} + T(d_{\mathrm{in}} + d_{\mathrm{out}})}}.
\]
（当 $d_{\mathrm{in}} = d_{\mathrm{out}} = d$ 时化简为 $\frac{Td}{d + 2T}$，与式~\eqref{eq:gemm-ai} 完全吻合。）

\textbf{(c) $T = 1$ 与 $T = 512$ 的可达算力、利用率及临界词元数 $T^*$。} 代入 $d_{\mathrm{in}} d_{\mathrm{out}} = 4096 \times 16384 = 67{,}108{,}864$、$d_{\mathrm{in}} + d_{\mathrm{out}} = 20{,}480$：
\begin{itemize}
  \item \textbf{当 $T = 1$（单请求 Decode）}：
  \[
    \mathrm{AI}(1) = \frac{67{,}108{,}864}{67{,}108{,}864 + 20{,}480} \approx \mathbf{0.9997\,\mathrm{FLOP/B}} \ll 250\,\mathrm{FLOP/B}
  \]
  （深度访存受限）。可达算力 $P(1) = B_{\mathrm{HBM}} \cdot \mathrm{AI}(1) = 8.0 \times 0.9997 \approx \mathbf{7.998\,\mathrm{TFLOP/s}}$，硬件利用率仅为 $7.998 / 2000 \approx \mathbf{0.40\%}$！
  \item \textbf{当 $T = 512$（Prefill）}：
  \[
    \mathrm{AI}(512) = \frac{512 \times 67{,}108{,}864}{67{,}108{,}864 + 512 \times 20{,}480}
    = \frac{512}{1 + 512 \times \frac{5}{16384}} = \frac{512}{1.15625} \approx \mathbf{442.81\,\mathrm{FLOP/B}} > 250\,\mathrm{FLOP/B}
  \]
  （进入计算受限区）。可达算力撞顶 $P(512) = \mathbf{2000\,\mathrm{TFLOP/s}}$，硬件利用率 $\mathbf{100\%}$。
  \item \textbf{最小词元数 $T^*$}：令 $\mathrm{AI}(T^*) = I^* = 250$，解线性方程：
  \[
    T^*\big(d_{\mathrm{in}} d_{\mathrm{out}} - I^*(d_{\mathrm{in}} + d_{\mathrm{out}})\big) = I^* d_{\mathrm{in}} d_{\mathrm{out}}
    \implies
    T^* = \frac{I^*}{1 - I^*\left(\frac{1}{d_{\mathrm{in}}} + \frac{1}{d_{\mathrm{out}}}\right)}
    = \frac{250}{1 - 250 \times \frac{5}{16384}}
    \approx \mathbf{270.65}.
  \]
  故向上取整至少需要 $T^* = \mathbf{271}$ 个词元才能让该非方阵线性层跨越拐点进入计算受限区。
\end{itemize}

\solItem{6.5.2}{13B 大模型训练静态与激活显存全账本}

已知 $\Phi = 1.3 \times 10^{10}$，$L = 40, d = 5120, h = 40, B = 2, S = 4096$（按表~\ref{tab:16phi-ledger} 约定以 $1\,\mathrm{GB} = 10^9\,\mathrm{Bytes}$ 计，括号内同时给出二进制 $\mathrm{GiB} = 2^{30}\,\mathrm{Bytes}$ 对照）。

\textbf{(a) $16\Phi$ 静态显存逐项明细。}
\begin{itemize}
  \item BF16 模型权重（$2\Phi$）：$2 \times 13 = \mathbf{26\,\mathrm{GB}}$；
  \item BF16 梯度张量（$2\Phi$）：$2 \times 13 = \mathbf{26\,\mathrm{GB}}$；
  \item FP32 主权重副本（$4\Phi$）：$4 \times 13 = \mathbf{52\,\mathrm{GB}}$；
  \item FP32 AdamW 一阶矩与二阶矩（$8\Phi$）：$8 \times 13 = \mathbf{104\,\mathrm{GB}}$；
  \item \textbf{静态显存总计（$16\Phi$）}：$26 + 26 + 52 + 104 = \mathbf{208\,\mathrm{GB}}$（已超过 2 张 80\,GB GPU 总容量）。
\end{itemize}

\textbf{(b) 全网 40 层总激活显存（式~\eqref{eq:act-memory-layer}）。} 先算公因子与二次注意力项系数：
\[
  B S d = 2 \times 4096 \times 5120 = 41{,}943{,}040\,\mathrm{B} \approx 41.943\,\mathrm{MB}, \qquad
  \frac{5 h S}{d} = \frac{5 \times 40 \times 4096}{5120} = \mathbf{160}.
\]
注意 $O(S^2)$ 项的系数 $160$ 竟是线性项 $34$ 的 $4.7\times$！
\begin{itemize}
  \item \textbf{未启用 FlashAttention}：单层 $M_{\mathrm{act,1}} = 41.943\,\mathrm{MB} \times (34 + 160) = 41.943 \times 194 \approx 8.137\,\mathrm{GB}$，40 层总计
  \[
    M_{\mathrm{act,no\text{-}flash}} = 40 \times 8.13695\,\mathrm{GB} \approx \mathbf{325.48\,\mathrm{GB}} \quad (303.13\,\mathrm{GiB}).
  \]
  \item \textbf{启用 FlashAttention（消去 $160$ 项）}：单层 $M_{\mathrm{act,1}}^{\mathrm{FA}} = 41.943\,\mathrm{MB} \times 34 \approx 1.426\,\mathrm{GB}$，40 层总计
  \[
    M_{\mathrm{act,flash}} = 40 \times 1.42606\,\mathrm{GB} \approx \mathbf{57.04\,\mathrm{GB}} \quad (53.13\,\mathrm{GiB}),
  \]
  激活显存立降 $5.7\times$！
\end{itemize}

\textbf{(c) 叠加均匀激活检查点后的峰值激活显存。} 由定理~\ref{thm:activation-checkpointing} 式~\eqref{eq:ckpt-mem}，峰值保留 $L/k + k$ 层激活。对 $L = 40$，最优连续值 $k^* = \sqrt{40} \approx 6.32$，取整除 $40$ 的两个邻近块长 $k = 5$ 或 $k = 8$ 恰好均给出相同的最小层数：
\[
  \frac{40}{5} + 5 = 8 + 5 = 13, \qquad \frac{40}{8} + 8 = 5 + 8 = 13.
\]
因此启用 FlashAttention 并配合均匀激活检查点后，峰值激活显存进一步压缩至
\[
  M_{\mathrm{peak}} = 13 \times 1.42606\,\mathrm{GB} \approx \mathbf{18.54\,\mathrm{GB}} \quad (17.27\,\mathrm{GiB}),
\]
相比原始 $325.48\,\mathrm{GB}$ 合计压低了整整 $17.6\times$！

\solItem{6.5.3}{激活检查点定理的推广与极限边界}

\textbf{(a) 两级嵌套检查点的显存极小值与算力代价。} 将 $L$ 层分为 $L / k_1$ 个大段（每大段含 $k_1$ 层），每大段内部又按 $k_2$ 层划分为 $k_1 / k_2$ 个子段（$1 \le k_2 \le k_1 \le L$）：
\begin{enumerate}
  \item \textbf{峰值激活显存}：反向传播退回到最深处的某个子段内部时，HBM 中同时驻留三部分激活：（i）全网 $L / k_1$ 个大段起点的顶级检查点；（ii）当前大段重算前向时保存的 $k_1 / k_2$ 个子段起点的二级检查点；（iii）当前正在反向求导的单个子段内部全部 $k_2$ 层的完整中间激活。故峰值显存上界为
  \[
    M_{\mathrm{peak}}(k_1, k_2) = \left(\frac{L}{k_1} + \frac{k_1}{k_2} + k_2\right) M_1.
  \]
  应用三元算术--几何均值不等式（AM-GM）：
  \[
    \frac{L}{k_1} + \frac{k_1}{k_2} + k_2 \;\ge\; 3\left(\frac{L}{k_1} \cdot \frac{k_1}{k_2} \cdot k_2\right)^{1/3} = \mathbf{3 L^{1/3}}.
  \]
  等号当且仅当 $\frac{L}{k_1} = \frac{k_1}{k_2} = k_2$ 时成立，解得最优分段为 $k_2^* = L^{1/3},\; k_1^* = L^{2/3}$，最小峰值显存降至 $\mathbf{3 L^{1/3} M_1}$。
  \item \textbf{额外前向传播次数与总算力增幅}：每个最底层除了原始 1 次前向传播外，还需要：（i）在当前大段重建二级检查点时额外重算 1 次前向；（ii）在当前子段反向求导时再额外重算 1 次前向——共计最多 \textbf{2 次额外前向传播}！设单层前向耗时 $t_F$、反向耗时 $t_B \approx 2 t_F$（基线总耗时 $T_0 = L(t_F + 2 t_F) = 3 L t_F$），则总算力相对增量上界达到
  \[
    \frac{\Delta T}{T_0} \;\le\; \frac{2 L t_F}{3 L t_F} = \frac{2}{3} \approx \mathbf{+66.7\%}
  \]
  （比单级检查点的 $+33.3\%$ 整整翻倍）。
\end{enumerate}

\textbf{(b) 工程界普遍采用单级（或每层边界）检查点的原因。}
\begin{itemize}
  \item \textbf{边际收益递减与瓶颈转移}：以常见 $L = 80$ 层大模型为例，单级检查点已把激活层数从 $80$ 压至 $2\sqrt{80} \approx 18$ 层（在 3D 并行与 FlashAttention 下往往已小于 $10\,\mathrm{GB}$），甚至每层边界全检查点（每层只存输入张量 $2 B S d$，仅约每层完整激活 $34 B S d$ 的 $1/17$）即可把激活降到几乎可忽略的量级，且仅花 1 次额外前向（$+33.3\%$）；
  \item \textbf{算力极其昂贵}：三级以上嵌套重算虽然在理论上把显存压至 $O(L^{1/d})$ 乃至 $O(\log L)$，但每深一层递归就多烧掉一整遍前向算力（$+66.7\%, +100\%, \dots$），直接令千万美元 GPU 集群的有效吞吐率（MFU）腰斩，在工程上彻底得不偿失。
\end{itemize}

\solItem{6.5.4}{FlashAttention-2 延迟归一化形式与反向传播 $L_i$ 重构证明}

\textbf{(a) 延迟归一化与逐步归一化的严格代数等价性。} 我们对块索引 $j = 0, 1, \dots, T_c$ 用数学归纳法证明恒等式 $\tilde{O}_i^{(j)} = \ell_i^{(j)} O_i^{(j)}$：
\begin{itemize}
  \item \textbf{奠基（$j = 0$）}：$\ell_i^{(0)} = 0$，$\tilde{O}_i^{(0)} = \mathbf{0} = 0 \cdot O_i^{(0)}$，成立。
  \item \textbf{归纳步}：假设对 $j - 1$ 有 $\tilde{O}_i^{(j-1)} = \ell_i^{(j-1)} O_i^{(j-1)}$。按题意延迟归一化更新：
  \[
    \tilde{O}_i^{(j)} = \alpha_i^{(j)} \tilde{O}_i^{(j-1)} + \tilde{P}_{ij} V_j
    = \Big(\alpha_i^{(j)} \ell_i^{(j-1)}\Big) O_i^{(j-1)} + \tilde{P}_{ij} V_j.
  \]
  对照式~\eqref{eq:flash-o} $O_i^{(j)} = \frac{\big(\alpha_i^{(j)} \ell_i^{(j-1)}\big) O_i^{(j-1)} + \tilde{P}_{ij} V_j}{\ell_i^{(j)}}$，两端同乘 $\ell_i^{(j)}$ 即得 $\tilde{O}_i^{(j)} = \ell_i^{(j)} O_i^{(j)}$！
\end{itemize}
因此循环结束 $j = T_c$ 后统一除以 $\ell_i^{(T_c)}$ 所得 $\tilde{O}_i^{(T_c)} / \ell_i^{(T_c)} \equiv O_i^{(T_c)}$，与逐步除法在代数上完全恒等，同时省去了内循环中 $T_c - 1$ 次昂贵的逐元素浮点除法。

\textbf{(b) 仅存 $L_i$ 的反向传播精确重构与行同步消除证明。}
\begin{enumerate}
  \item \textbf{概率块精确重构}：前向结束时，每行最终最大值记为 $m_i = m_i^{(T_c)}$，分母记为 $\ell_i = \ell_i^{(T_c)} = \sum_{k=1}^N \exp(S_{ik} - m_i)$，仅存单标量 log-sum-exp $L_i = m_i + \ln \ell_i$。反向重算任意块 $S_{ij} = \tau Q_i K_j^\top$ 后，直接计算
  \[
    \exp(S_{ij} - L_i) = \exp(S_{ij} - m_i - \ln \ell_i) = \frac{\exp(S_{ij} - m_i)}{\ell_i} = \mathbf{P_{ij}},
  \]
  精确还原了全序列 Softmax 注意力概率块，无需在 HBM 中存储 $N \times N$ 矩阵！
  \item \textbf{Softmax 反向与 $D_i$ 局域化推导}：由前向 $O = P V$，得 $dP_{ij} = dO_i V_j^\top$（与 $dV_j = \sum_i P_{ij}^\top dO_i$）。对第 $r$ 行（为简洁略去块标，以单行向量 $p_r = \softmax(s_r)$ 记），Softmax 的雅可比矩阵为 $\frac{\partial p_{r,c}}{\partial s_{r,k}} = p_{r,c}(\delta_{ck} - p_{r,k})$，故链式法则给出
  \[
    dS_{r,k} = \sum_{c=1}^N dP_{r,c}\,\frac{\partial p_{r,c}}{\partial s_{r,k}}
    = p_{r,k}\!\left(dP_{r,k} - \sum_{c=1}^N p_{r,c}\,dP_{r,c}\right).
  \]
  最惊艳的一步在于化简括号内的整行求和 $\sum_{c=1}^N p_{r,c}\,dP_{r,c}$：代入 $dP_{r,c} = \langle dO_r,\, V_c \rangle$，利用内积线性性交换求和顺序：
  \[
    \sum_{c=1}^N p_{r,c}\,dP_{r,c}
    = \sum_{c=1}^N p_{r,c}\,\langle dO_r,\, V_c \rangle
    = \left\langle dO_r,\; \sum_{c=1}^N p_{r,c} V_c \right\rangle
    = \langle dO_r,\, O_r \rangle = \operatorname{rowsum}(dO_r \odot O_r) \triangleq D_r.
  \]
  回到块矩阵形式即得 $D_i = \operatorname{rowsum}(dO_i \odot O_i) \in \R^{B_r \times 1}$ 与 $dS_{ij} = P_{ij} \odot (dP_{ij} - D_i)$。这一恒等式彻底消除了反向计算 $dS_{ij}$ 时跨列块 $j$ 的全局归约依赖，使得 $D_i$ 可在反向主循环前独立预处理一次，所有块计算完全保留在高速片上 SRAM 中！
\end{enumerate}

\solItem{6.5.5}{定量诊断算子融合对访存受限算子的加速比}

已知张量形状 $(B, S, d) = (16, 2048, 4096)$，精度 BF16（每元素 2 字节）。单张量体积为：
\[
  V = 16 \times 2048 \times 4096 \times 2\,\mathrm{B} = 268{,}435{,}456\,\mathrm{B} = 256\,\mathrm{MiB} \approx \mathbf{268.435\,\mathrm{MB}}.
\]

\textbf{(a) 融合前后 HBM 总字节读写量。}
\begin{itemize}
  \item \textbf{未融合（3 个独立核函数）}：
  \begin{enumerate}
    \item 残差加 $H_1 = X + R$：读 $X$、读 $R$、写 $H_1$ $\implies 3V$；
    \item 归一化 $H_2 = \mathrm{RMSNorm}(H_1)$：读 $H_1$、写 $H_2$（忽略 $4096 \times 2\,\mathrm{B} = 8\,\mathrm{KB}$ 的 $\gamma$）$\implies 2V$；
    \item 激活函数 $H_3 = \mathrm{SiLU}(H_2)$：读 $H_2$、写 $H_3$ $\implies 2V$。
  \end{enumerate}
  总计 $Q_{\mathrm{unfused}} = (3 + 2 + 2)V = \mathbf{7V \approx 1.8790\,\mathrm{GB}}$（其中中间张量 $H_1, H_2$ 白白往返 HBM 两次）。
  \item \textbf{完全融合（单核片上流水）}：从 HBM 读入 $X, R$，在寄存器/SRAM 中完成加法、行归约 RMSNorm 与 SiLU，最后仅将最终结果 $H_3$ 写回 HBM。总计
  \[
    Q_{\mathrm{fused}} = \mathbf{3V \approx 0.8053\,\mathrm{GB}}.
  \]
\end{itemize}

\textbf{(b) H100（$B_{\mathrm{HBM}} = 3.35\,\mathrm{TB/s} = 3.35 \times 10^{12}\,\mathrm{B/s}$）上的理论墙钟耗时与加速比。}
\begin{align*}
  t_{\mathrm{unfused}} &= \frac{1.879048 \times 10^9\,\mathrm{B}}{3.35 \times 10^{12}\,\mathrm{B/s}} \approx \mathbf{0.5609\,\mathrm{ms}}, \\
  t_{\mathrm{fused}} &= \frac{0.805306 \times 10^9\,\mathrm{B}}{3.35 \times 10^{12}\,\mathrm{B/s}} \approx \mathbf{0.2404\,\mathrm{ms}}, \\
  \mathrm{Speedup} &= \frac{t_{\mathrm{unfused}}}{t_{\mathrm{fused}}} = \frac{7V}{3V} = \frac{7}{3} \approx \mathbf{2.33\times}.
\end{align*}
不仅直接提速 $2.33\times$，还省去了两次内核启动（Kernel Launch）开销，这正是 Triton 与 \texttt{torch.compile} 带来巨大加速的核心物理机制。

\solItem{6.5.6}{数值反例：为什么纯 BF16 权重更新会导致优化彻底“冻结”？}

\textbf{(a) 区间 $[1, 2)$ 内的 BF16 步长（ULP）。} 在区间 $[1, 2) = [2^0, 2^1)$ 内，浮点数的无偏指数恒为 $e = 0$，BF16 有 7 位显式尾数，故二进制最低有效位对应的数值间距为
\[
  \mathrm{ULP}_{\mathrm{BF16}}([1, 2)) = 2^{0 - 7} = 2^{-7} = \frac{1}{128} = \mathbf{0.0078125}.
\]
因此 $1.0$ 之后下一个可精确表示的 BF16 数是 $1 + 2^{-7} = 1.0078125$。

\textbf{(b) $10000$ 步更新后的结果与破局机制。} 在“舍入到最近偶数（Round-to-Nearest-Even）”规则下，$1.0$ 与 $1.0078125$ 的中点为
\[
  1.0 + \frac{1}{2}\mathrm{ULP} = 1.0 + 2^{-8} = 1.00390625.
\]
第 1 步精确更新值为 $w - \eta g = 1.0 - 10^{-3}\times(-1.0) = 1.001$。由于单步增量 $0.001 < 0.00390625$，$1.001$ 严格更靠近 $1.0$，舍入后严格退回：
\[
  w_1 = \mathrm{BF16}(1.001) = 1.0.
\]
因为每一步都从完全相同的 $1.0$ 出发，数学归纳法立知跑完 $10000$ 步后 $w_{10000} = \mathbf{1.0}$（与理想值 $1.0 + 10000 \times 10^{-3} = 11.0$ 相差整整 $10.0$，权重彻底冻结！）。

\textbf{破局方案}：（1）\textbf{FP32 主权重（标准混合精度）}：FP32 拥有 23 位尾数，在 $[1, 2)$ 内 $\mathrm{ULP}_{\mathrm{FP32}} = 2^{-23} \approx 1.19 \times 10^{-7} \ll 10^{-3}$，微小梯度在 FP32 主权重中无损累加，每约 8 步跨越半个 BF16 ULP 时 BF16 副本自然跳升至下一个网格点；（2）\textbf{随机舍入（Stochastic Rounding）}：以概率 $p = \frac{0.001}{0.0078125} = 0.128$ 向上舍入到 $1.0078125$、以概率 $1 - p = 0.872$ 向下舍入到 $1.0$，单步期望更新精确等于 $0.128 \times 0.0078125 = 0.001$，无偏打破冻结陷阱（代价是引入微小舍入方差）。

\solItem{6.5.7}{编程实验：为分块在线 Softmax 添加因果掩码（Causal Mask）与上三角跳过加速}

\textbf{(a)--(b) 实现要点。} 在代码~\ref{lst:flash-online}（\texttt{exp5\_flash\_\allowbreak online\_and\_roofline()}）中，对标准注意力矩阵上三角（\texttt{k=1}）施加 $-\infty$ 因果掩码（调用 \texttt{np.triu(..., k=1)}）得到基准输出 $O_{\mathrm{causal}}$。在分块在线 Softmax 的内层 \texttt{for j in range(0, N, Bc):} 循环内加入两行工程优化：
\begin{enumerate}
  \item \textbf{纯上三角块零计算跳过}：若 \texttt{j > i + Br - 1}，则该块内所有列索引均严格大于行索引，权重全为 $0$，直接 \texttt{continue}；
  \item \textbf{主对角线块局部掩码}：当 $B_r = B_c$ 且 \texttt{j == i} 时，对局部瓦片分数 $S_{ij}$ 施加局部上三角 $-\infty$ 掩码（即 \texttt{Sij[np.triu(\allowbreak np.ones((Br, Bc), bool), 1)] = -np.inf}；若 $B_r \neq B_c$，则对所有跨越主对角线的块按全局坐标 \texttt{col\_idx > row\_idx} 掩码）；下三角块（\texttt{j < i}）无需任何掩码操作。
\end{enumerate}

\textbf{(c) 独立数值验证结果。} 在 $N = 256, d = 64, B_r = B_c = 32$ 下运行验证：
\begin{itemize}
  \item 与标准因果注意力的最大绝对误差实测为 $\max|O_{\mathrm{causal}} - O_{\mathrm{flash}}| = \mathbf{5.55 \times 10^{-16}} < 10^{-15}$（纯机器双精度浮点舍入级）；
  \item 总瓦片数为 $8 \times 8 = 64$，实际仅计算下三角与主对角线共 $\frac{8 \times 9}{2} = \mathbf{36}$ 个瓦片，跳过 $28$ 个瓦片，精确省去 $\frac{28}{64} = \mathbf{43.75\%}$ 的块乘加计算；一般地，当块数为 $T = N / B$ 时，跳过比例为 $\frac{T(T-1)/2}{T^2} = \frac{1}{2}\big(1 - \frac{1}{T}\big) \xrightarrow{T \to \infty} \mathbf{50\%}$，实现近 $2\times$ 因果加速！
\end{itemize}

\solItem{6.5.8}{编程实验：Roofline 分析标准 MHA 与分组查询注意力（GQA）在长上下文 Decode 中的带宽瓶颈}

已知单层 Decode 阶段批量 $B = 16$，$d = 4096$，查询头数 $h_q = 32$，头维度 $d_h = 128$（满足 $h_q d_h = d$），历史上下文长度 $S \in [512, 32768]$，精度 BF16（2 字节）。

\textbf{(a) 总 FLOPs 与 HBM 访存字节数推导。}
\begin{itemize}
  \item \textbf{计算量}：每个请求、每个查询头需计算 $q \in \R^{1 \times d_h}$ 与 $K_{\mathrm{cache}}^\top \in \R^{d_h \times S}$ 的内积（$2 S d_h$ FLOPs），以及权重 $p \in \R^{1 \times S}$ 与 $V_{\mathrm{cache}} \in \R^{S \times d_h}$ 的加权求和（$2 S d_h$ FLOPs）。汇总 $B$ 个请求与 $h_q$ 个查询头，忽略低阶的 $O(B h_q S)$ Softmax 标量运算，总计
  \[
    F_{\mathrm{attn}} = B \cdot h_q \cdot (2 S d_h + 2 S d_h) = 4 B S h_q d_h = \mathbf{4 B S d}\;\text{FLOPs}.
  \]
  \item \textbf{HBM 访存量}：每个请求拥有\textbf{彼此独立}的历史 KV 缓存，无法跨批量共享！读入 $K_{\mathrm{cache}}$ 与 $V_{\mathrm{cache}}$（各含 $B \cdot h_{\mathrm{kv}} \cdot S \cdot d_h$ 个 BF16 元素）占 $2 \times (B h_{\mathrm{kv}} S d_h) \times 2 = 4 B h_{\mathrm{kv}} S d_h$ 字节；若再计入读入当前步 $Q$（$2Bd$ 字节）与写出输出 $O$（$2Bd$ 字节），总访存为
  \[
    Q_{\mathrm{attn}} = 4 B h_{\mathrm{kv}} S d_h + 4 B d\;\text{Bytes}.
  \]
\end{itemize}

\textbf{(b) 算术强度 $\mathrm{AI}_{\mathrm{attn}}$ 与 MHA/GQA-8 对比。} 作比消去 $4 B d_h$：
\[
  \mathrm{AI}_{\mathrm{attn}} = \frac{h_q S}{h_{\mathrm{kv}} S + h_q} = \frac{h_q / h_{\mathrm{kv}}}{1 + \frac{h_q}{h_{\mathrm{kv}} S}}
  \;\xrightarrow{\;S \gg 1\;}\; \mathbf{\frac{h_q}{h_{\mathrm{kv}}}\,\mathrm{FLOP/B}}.
\]
代入 $S \in \{512, 4096, 32768\}$ 的精确数值如下表：
\begin{center}
\footnotesize
\begin{tabular}{lcccc}
\toprule
注意力结构 & $S = 512$ & $S = 4096$ & $S = 32768$ & 渐近极限（与 $B, S$ 无关） \\
\midrule
标准 MHA（$h_{\mathrm{kv}} = 32$） & $0.998\,\mathrm{FLOP/B}$ & $1.000\,\mathrm{FLOP/B}$ & $1.000\,\mathrm{FLOP/B}$ & $\mathbf{1.0\,\mathrm{FLOP/B}}$ \\
GQA-8（$h_{\mathrm{kv}} = 4$） & $7.877\,\mathrm{FLOP/B}$ & $7.984\,\mathrm{FLOP/B}$ & $7.998\,\mathrm{FLOP/B}$ & $\mathbf{8.0\,\mathrm{FLOP/B}}$（$8\times$ 提升） \\
\bottomrule
\end{tabular}
\end{center}
\textbf{深刻系统结论}：公式中批量大小 $B$ 分子分母完全约去！与线性投影层可以通过增大批量 $B$ 提高权重复用不同，Decode 阶段的注意力层因每条请求 KV 缓存独立，标准 MHA 无论 $B$ 和 $S$ 多大，算术强度都被死死钉在 $1\,\mathrm{FLOP/B}$（在 H100 上利用率不足 $0.34\%$）；而 GQA-8（以及 DeepSeek 的 MLA）通过让 $8$ 个查询头共享同一份 KV 头，不仅把 KV 显存砍掉 $8\times$，更将 Decode 注意力的 Roofline 算术强度与实际吞吐率直接抬升 $8\times$！

% =====================================================================
\section*{第 6 章习题解答（共 8 题）}
\addcontentsline{toc}{section}{第 6 章习题解答（共 8 题）}

\solItem{6.6.1}{70B 模型在 64 卡集群上的 ZeRO 显存与 3D 并行综合核算}

已知 $\Phi = 7 \times 10^{10}$（70B），BF16 权重体积 $\Psi = 2\Phi = 140\,\mathrm{GB}$，$N_{\mathrm{GPU}} = 64$，单卡显存 $80\,\mathrm{GB}$。

\textbf{(a) 64 路纯数据并行（$N_d = 64$）下的 DDP 与 ZeRO-1/2/3 账本。} 对照表~\ref{tab:zero-ledger} 逐项代入 $\Phi = 70\,\mathrm{B}, \Psi = 140\,\mathrm{GB}, N_d = 64$：
\begin{center}
\footnotesize
\setlength{\tabcolsep}{4.5pt}
\begin{tabular}{@{}lcccc@{}}
\toprule
并行方案 & 单卡静态显存公式 & 单卡静态显存 & 每步单卡通信量 & 能否装入 80\,GB \\
\midrule
DDP & $16\Phi$ & $1120.00\,\mathrm{GB}$ & $2\Psi = 280\,\mathrm{GB}$ & 否（超 $14\times$） \\
ZeRO-1 & $4\Phi + 12\Phi / 64$ & $280 + 13.125 = \mathbf{293.13\,\mathrm{GB}}$ & $2\Psi = 280\,\mathrm{GB}$ & 否 \\
ZeRO-2 & $2\Phi + 14\Phi / 64$ & $140 + 15.3125 = \mathbf{155.31\,\mathrm{GB}}$ & $2\Psi = 280\,\mathrm{GB}$ & 否 \\
ZeRO-3 / FSDP & $16\Phi / 64$ & $\mathbf{17.50\,\mathrm{GB}}$ & $3\Psi = 420\,\mathrm{GB}$ & \textbf{能}（余 $62.5\,\mathrm{GB}$） \\
\bottomrule
\end{tabular}
\end{center}
可见 DDP、ZeRO-1 与 ZeRO-2 甚至连静态模型状态都装不进单张 80\,GB GPU，纯数据并行下唯有 ZeRO-3 / FSDP 可行。

\textbf{(b) 3D 并行（$N_t = 8, N_p = 4, N_d = 2$，顶层 ZeRO-1，$m = 32$）核算。}
\begin{itemize}
  \item \textbf{单卡局部参数量}：张量并行与流水线并行将模型无冗余切分为 $N_t N_p = 8 \times 4 = 32$ 份，每卡分管
  \[
    \Phi_{\mathrm{local}} = \frac{70\,\mathrm{B}}{32} = \mathbf{2.1875\,\mathrm{B}}.
  \]
  \item \textbf{单卡静态显存}：在 $N_d = 2$ 组数据并行副本上对 $\Phi_{\mathrm{local}}$ 启用 ZeRO-1（每参数占 $4 + 12 / 2 = 10\,\mathrm{B}$）：
  \[
    M_{\mathrm{static}} = 10 \times 2.1875\,\mathrm{GB} = \mathbf{21.875\,\mathrm{GB}} \ll 80\,\mathrm{GB},
  \]
  余下近 $58\,\mathrm{GB}$ 留给激活与缓冲区，且所有高频 TP AllReduce 均局限在单机 8 卡 NVLink 内部，跨机仅传递轻量的 PP 边界激活与 $N_d = 2$ 梯度！
  \item \textbf{流水线气泡率}：由定理~\ref{thm:pipeline-bubble} 式~\eqref{eq:pp-bubble}，代入 $N_p = 4, m = 32$：
  \[
    f_{\mathrm{bubble}} = \frac{N_p - 1}{m + N_p - 1} = \frac{4 - 1}{32 + 4 - 1} = \frac{3}{35} \approx \mathbf{8.57\%}.
  \]
\end{itemize}

\solItem{6.6.2}{LoRA 显存预算与 GQA KV-Cache 精确核算}

已知 $L = 32, d = 4096, h_q = 32, n_{\mathrm{kv}} = 8, d_{\mathrm{head}} = 128$，底座规模 $\Phi = 8 \times 10^9$（8B）。

\textbf{(a) LoRA 参数量、占比与总静态显存。} 按题目近似，$W_Q, W_V \in \R^{4096 \times 4096}$，施加秩 $r = 16$ 的低秩旁路（每个矩阵增设 $A \in \R^{r \times d}$ 与 $B \in \R^{d \times r}$，共 $2dr$ 参数）：
\begin{itemize}
  \item 单层可训练参数：$2 \times (2 \times 4096 \times 16) = 262{,}144$；
  \item 全网 32 层总可训练参数：
  \[
    \Phi_{\mathrm{LoRA}} = 32 \times 262{,}144 = \mathbf{8{,}388{,}608 \approx 8.39\,\mathrm{M}},
  \]
  仅占 8B 底座参数量的 $\frac{8.3886 \times 10^6}{8 \times 10^9} \approx \mathbf{0.105\%}$！
  \item 总静态显存：冻结底座仅需保存 BF16 权重（无梯度、无主权重、无 AdamW 状态），占 $2 \times 8 = 16.0\,\mathrm{GB}$；LoRA 可训练参数需完整 $16\Phi_{\mathrm{LoRA}}$ 状态，仅占 $16 \times 8.3886\,\mathrm{MB} \approx \mathbf{134.2\,\mathrm{MB} \approx 0.134\,\mathrm{GB}}$。故静态总显存仅约 $\mathbf{16.13\,\mathrm{GB}}$（相比全量微调 $16 \times 8 = 128\,\mathrm{GB}$ 暴降近 $8\times$！）。
\end{itemize}

\textbf{(b) GQA KV 缓存核算与 PagedAttention 最大并发数。} 代入式~\eqref{eq:kv-cache-formula}（BF16 取 $b_{\mathrm{dtype}} = 2$）：
\[
  M_{\mathrm{KV/tok}} = 2 \cdot L \cdot n_{\mathrm{kv}} \cdot d_{\mathrm{head}} \cdot b_{\mathrm{dtype}}
  = 2 \times 32 \times 8 \times 128 \times 2 = \mathbf{131{,}072\,\mathrm{B}} = \mathbf{128\,\mathrm{KiB/tok}} \approx 131.07\,\mathrm{KB/tok}.
\]
单条长度 $S = 4096$ 的请求占用 KV 缓存为：
\[
  M_{\mathrm{seq}} = 4096 \times 131{,}072\,\mathrm{B} = 536{,}870{,}912\,\mathrm{B} = \mathbf{512\,\mathrm{MiB}} \approx \mathbf{536.87\,\mathrm{MB}}.
\]
可用 KV 显存池为 $60\,\mathrm{GB}$（PagedAttention 下内部碎片仅半个块 $< 0.2\%$，可忽略）：
\begin{itemize}
  \item 按硬件 HBM 惯用的二进制口径（$60\,\mathrm{GiB} = 60 \times 1024\,\mathrm{MiB}$）：最多可并发服务 $\frac{60 \times 1024}{512} = \mathbf{120}$ 条请求；
  \item 按严格十进制口径（$60\,\mathrm{GB} = 60 \times 10^9\,\mathrm{B}$）：$\lfloor 60 \times 10^9 / 536{,}870{,}912 \rfloor = \mathbf{111}$ 条请求（工程速算常直接以 $60{,}000 / 512 \approx 117$ 估算，均在 110--120 条并发量级）。
\end{itemize}

\solItem{6.6.3}{Ring AllReduce 对比 Halving-Doubling 树形归约的延迟—带宽权衡}

\textbf{(a) Rabenseifner 算法通信量与延迟证明。} 设 $P = 2^k$（即 $k = \log_2 P$）。
\begin{itemize}
  \item \textbf{ReduceScatter 阶段}：共执行 $k = \log_2 P$ 轮两两交换，第 $j$ 轮（$j = 1, \dots, k$）每张卡发送的数据量为 $M / 2^j$。由等比级数求和公式，每卡累计发送字节数为
  \[
    V_{\mathrm{RS}} = \sum_{j=1}^{k} \frac{M}{2^j} = M \cdot \frac{\frac{1}{2}\left(1 - \frac{1}{2^k}\right)}{1 - \frac{1}{2}} = M\left(1 - \frac{1}{P}\right) = \frac{P - 1}{P}\,M.
  \]
  \item \textbf{AllGather 阶段}：完全逆序执行 $k = \log_2 P$ 轮，每轮发送数据量依次为 $M/P, 2M/P, \dots, M/2$，累计发送字节数同样精确等于 $V_{\mathrm{AG}} = \frac{P-1}{P} M$。
\end{itemize}
合并两阶段，每张 GPU 的累计发送总字节数为 $\mathbf{2\frac{P-1}{P}M}$（与 Ring AllReduce 完全一致，均达到理论带宽下界），而串行通信步数从 $2(P-1)$ 锐减至 $2\log_2 P$，故总耗时模型为
\[
  T_{\mathrm{Rab}} = 2(\log_2 P)\alpha + 2\frac{P - 1}{P}\frac{M}{B}.
\]

\textbf{(b) 千卡集群（$P = 1024, \alpha = 5\,\mu\mathrm{s}, B = 50\,\mathrm{GB/s}$）的延迟差与临界消息体积。} 由于两种算法在无拥塞模型下拥有完全相同的带宽项 $2\frac{P-1}{P}\frac{M}{B}$，两者耗时的差值恒为固定延迟差：
\begin{align*}
  T_{\mathrm{Ring,lat}} &= 2(1024 - 1) \times 5\,\mu\mathrm{s} = 2046 \times 5\,\mu\mathrm{s} = \mathbf{10.23\,\mathrm{ms}}, \\
  T_{\mathrm{Rab,lat}} &= 2(\log_2 1024) \times 5\,\mu\mathrm{s} = 20 \times 5\,\mu\mathrm{s} = \mathbf{0.10\,\mathrm{ms}}, \\
  T_{\mathrm{Ring}} - T_{\mathrm{Rab}} &= 10.23 - 0.10 = \mathbf{10.13\,\mathrm{ms}} \quad (\text{相差 } 102\times).
\end{align*}
令 Ring AllReduce 自身的延迟项与带宽传输项恰好相等，解得临界消息体积 $M_{\mathrm{crit}}$：
\[
  2(P - 1)\alpha = 2\frac{P - 1}{P}\frac{M_{\mathrm{crit}}}{B}
  \implies M_{\mathrm{crit}} = P \alpha B = 1024 \times (5 \times 10^{-6}\,\mathrm{s}) \times (50 \times 10^9\,\mathrm{B/s}) = \mathbf{256\,\mathrm{MB}}
\]
（对应链路上的双向总数据流量 $2\frac{P-1}{P} M_{\mathrm{crit}} \approx 511.5\,\mathrm{MB}$）。\textbf{工程含义}：对于几 KB 的 LayerNorm 梯度张量（$M \ll 256\,\mathrm{MB}$），在 1024 卡上跑 Ring AllReduce 时 $99.99\%$ 以上的时间全部白白消耗在 $2046$ 跳的 $10.23\,\mathrm{ms}$ 串行握手延迟上，网络带宽利用率近乎归零！因此现代集合通信库（NCCL）在千卡集群上同时采用两项对策：（1）用\textbf{梯度分桶（Gradient Bucketing）}把碎小张量拼装成数十上百 MB 的大桶；（2）对中小消息自动切换至\textbf{双二叉树（Double Binary Tree）}或 Halving-Doubling 树形拓扑以将延迟压至 $O(\log P)$。

\solItem{6.6.4}{投机解码的多步联合分布精确性与真实墙钟加速比公式}

\textbf{(a) 多步投机循环的联合分布严格无损证明。} 对长度为 $T$ 的输出序列运用概率链式法则
\[
  \Prob(X_1, \dots, X_T) = \prod_{t=1}^T \Prob(X_t \mid X_{<t}).
\]
只需证明：无论某位置 $t$ 是在某轮投机循环中被草稿“直接接受”，还是因草稿被拒而触发“残差重采样”，其条件分布 $\Prob(X_t = x \mid X_{<t})$ 均严格等于目标大模型分布 $p(x \mid X_{<t})$。
\begin{itemize}
  \item 在每轮循环开始时，前缀 $X_{<t_0}$ 已完全确立。小模型依次生成 $\gamma$ 个草稿 $(\tilde{x}_1, \dots, \tilde{x}_\gamma)$，大模型单次并行前向即可算出全部条件分布 $p(\cdot \mid X_{<t_0}, \tilde{x}_{<j})$（$j = 1, \dots, \gamma + 1$）；
  \item 按顺序从 $j = 1$ 开始逐位检验：只要前 $j - 1$ 个草稿全部被接受，当前条件前缀就是已确立的精确历史 $X_{<t_0 + j - 1} = (X_{<t_0}, \tilde{x}_1, \dots, \tilde{x}_{j-1})$。对第 $j$ 位应用定理~\ref{thm:speculative-lossless}，无论接受还是拒绝重采，输出词元 $X_{t_0 + j - 1}$ 严格服从 $p(\cdot \mid X_{<t_0 + j - 1})$；
  \item \textbf{关键截断规则}：一旦在第 $k + 1$ 位发生拒绝并重采出新词元 $X_{t_0 + k} \neq \tilde{x}_{k+1}$，后续基于错误前缀 $\tilde{x}_{k+1}$ 生成的草稿 $\tilde{x}_{k+2}, \dots, \tilde{x}_\gamma$ \textbf{立刻全部作废丢弃}，下一轮重新以已验证前缀 $X_{\le t_0 + k}$ 为条件生成新草稿；而若全部 $\gamma$ 个草稿均获接受，第 $\gamma + 1$ 位直接由大模型自身的 $p(\cdot \mid X_{<t_0}, \tilde{x}_{1:\gamma})$ 精确采样。
\end{itemize}
因此每个位置 $t \in \{1, \dots, T\}$ 均严格满足 $\Prob(X_t = x \mid X_{<t}) = p(x \mid X_{<t})$，按链式法则相乘即证整条序列的联合分布与纯目标大模型严格一致！

\textbf{(b) 期望墙钟加速比推导与最优草稿长度 $\gamma^*$。} 每轮投机循环依次执行 $\gamma$ 次小草稿模型串行前向与 $1$ 次目标大模型并行验证前向，单轮总墙钟耗时为
\[
  T_{\mathrm{round}} = \gamma t_{\mathrm{draft}} + t_{\mathrm{target}} = (\gamma c + 1)\,t_{\mathrm{target}}.
\]
由式~\eqref{eq:spec-expected-len}，每轮产出的期望已验证词元数为 $\E[L] = \frac{1 - \alpha^{\gamma + 1}}{1 - \alpha}$。而原始自回归解码生成同样 $\E[L]$ 个词元耗时为 $\E[L]\,t_{\mathrm{target}}$。两者相除即得精确墙钟加速比公式：
\[
  \mathrm{Speedup}(\gamma, \alpha, c)
  = \frac{\E[L]\,t_{\mathrm{target}}}{(\gamma c + 1)\,t_{\mathrm{target}}}
  = \mathbf{\frac{1 - \alpha^{\gamma + 1}}{(1 - \alpha)(\gamma c + 1)}}.
\]
代入 $\alpha = 0.8, c = 0.05$ 逐个枚举整数 $\gamma$：
\begin{center}
\footnotesize
\begin{tabular}{ccccccccc}
\toprule
草稿长度 $\gamma$ & $3$ & $4$ & $5$ & $6$ & $7$ & $\mathbf{8}$ & $9$ & $10$ \\
\midrule
期望产出 $\E[L]$ & $2.952$ & $3.362$ & $3.689$ & $3.951$ & $4.161$ & $\mathbf{4.329}$ & $4.463$ & $4.571$ \\
轮耗倍数 $\gamma c + 1$ & $1.15$ & $1.20$ & $1.25$ & $1.30$ & $1.35$ & $\mathbf{1.40}$ & $1.45$ & $1.50$ \\
墙钟加速比（$\times$） & $2.567$ & $2.801$ & $2.951$ & $3.040$ & $3.082$ & $\mathbf{3.092}$ & $3.078$ & $3.047$ \\
\bottomrule
\end{tabular}
\end{center}
由表可见加速比先升后降，在 $\gamma^* = \mathbf{8}$ 处取得唯一全局最大值 $\mathbf{3.092\times}$！直觉上，当 $\gamma$ 过大时，后续草稿被前面拒绝所截断作废的概率急剧增大（边际产出 $\alpha^{\gamma+1}$ 指数衰减），而草稿串联耗时 $\gamma c$ 却线性增加，导致收益反被吞噬。

\solItem{6.6.5}{反例剖析：为什么 LoRA 把 $A, B$ 同时零初始化会导致训练永久锁死？}

已知标量模型 $y = (w_0 + b a)x$，损失 $\mathcal{L}(a, b) = \frac{1}{2}(y - y^*)^2$（记残差 $e = y - y^*$）。

\textbf{(a) 全零初始化 $(a_0, b_0) = (0, 0)$ 的永久锁死证明。} 由链式法则直接求偏导：
\[
  \frac{\partial \mathcal{L}}{\partial a} = e \cdot \frac{\partial y}{\partial a} = (y - y^*)\,b x, \qquad
  \frac{\partial \mathcal{L}}{\partial b} = e \cdot \frac{\partial y}{\partial b} = (y - y^*)\,a x.
\]
当 $(a_0, b_0) = (0, 0)$ 时，即便残差 $e_0 = w_0 x - y^* \neq 0$，两项偏导数均严格等于 $0$！
\begin{itemize}
  \item 在二阶曲率上，原点处 Hessian 矩阵为 $H(0, 0) = \begin{bmatrix} 0 & e_0 x \\ e_0 x & 0 \end{bmatrix}$，其特征值为 $\pm |e_0 x|$ 一正一负，故 $(0, 0)$ 是严格的\textbf{鞍点}（驻点）；
  \item 对标准梯度下降与 Adam（初始动量 $m_0 = 0$，每步收到精确零梯度 $g_t = 0$），所有迭代恒有 $a_t \equiv 0, b_t \equiv 0$，参数永远锁死在原点无法动弹。
\end{itemize}

\textbf{(b) 非对称初始化 $a_0 \neq 0, b_0 = 0$ 的破局机制与工程惯例。}
\begin{itemize}
  \item \textbf{第 0 步无扰动}：$y_0 = (w_0 + 0 \cdot a_0)x = w_0 x$，与预训练底座输出严格一致；
  \item \textbf{第 1 步激活 $b$}：此时 $\frac{\partial \mathcal{L}}{\partial a}\big|_0 = 0$（故 $a_1 = a_0 \neq 0$），但 $\frac{\partial \mathcal{L}}{\partial b}\big|_0 = e_0 a_0 x \neq 0$，于是梯度下降立刻更新出 $b_1 = -\eta e_0 a_0 x \neq 0$；
  \item \textbf{第 2 步全面唤醒 $a$}：由于 $b_1 \neq 0$，第 2 步立刻产生非零梯度 $\frac{\partial \mathcal{L}}{\partial a}\big|_1 = e_1 b_1 x \neq 0$，矩阵 $A$ 与 $B$ 协同进入正常优化轨道！
\end{itemize}
\textbf{为何反取 $a_0 = 0, b_0 \neq 0$ 也数学可行、但工程偏好 $B = 0$？} 从偏导公式可见 $(a, b)$ 在数学上完全对称，$A = 0, B \neq 0$ 同样可在第 1 步唤醒 $A$、第 2 步唤醒 $B$。工程界（Hu et al., 2022）统一采用 $A \sim \mathcal{N}(0, \sigma^2)$、$B = \mathbf{0}$ 的深层原因有二：（1）前向数据流 $x \mapsto A x \mapsto B(Ax)$ 中，随机正交/高斯投影 $A$ 先将高维输入特征映射到 $r$ 维低秩子空间，$B = 0$ 保证输出端精确零初始化且在多 LoRA 适配器动态热插拔时接口统一；（2）按 $\mu$P（Maximal Update Parametrization，LoRA+ 分析）在超大宽度 $d \to \infty$ 下的特征学习标度律，$A$ 用 $O(1/\sqrt{d})$ 随机初始化、$B = 0$（并给 $B$ 略高学习率）能获得更稳定的梯度范数与特征演化。

\solItem{6.6.6}{重尾通道 INT4 量化手算与 SmoothQuant 迁移演练}

已知 $W = \begin{bmatrix} 0.4 & -0.3 \\ 7.0 & -5.0 \end{bmatrix}$，$x = \begin{bmatrix} 10.0 & 0.5 \end{bmatrix}$，精确乘积为 $y_{\mathrm{true}} = x W = [10\times 0.4 + 0.5\times 7.0,\; 10\times(-0.3) + 0.5\times(-5.0)] = \mathbf{[7.5,\; -5.5]}$。对称 INT4 取 $q_{\max} = 7$。

\textbf{(a) 逐张量 vs 逐行（逐通道）INT4 量化手算。}
\begin{enumerate}
  \item \textbf{逐张量量化}：全矩阵最大绝对值为 $7.0$，缩放因子 $s_{\mathrm{pt}} = \frac{7.0}{7} = \mathbf{1.0}$。逐元素除以 $1.0$ 并四舍五入：
  \[
    Q_{\mathrm{pt}} = \operatorname{round}\!\begin{bmatrix} 0.4 & -0.3 \\ 7.0 & -5.0 \end{bmatrix}
    = \begin{bmatrix} 0 & 0 \\ 7 & -5 \end{bmatrix}, \qquad
    \hat{W}_{\mathrm{pt}} = s_{\mathrm{pt}} Q_{\mathrm{pt}} = \mathbf{\begin{bmatrix} 0 & 0 \\ 7.0 & -5.0 \end{bmatrix}}.
  \]
  \textbf{灾难后果}：正常第 1 行 $[0.4, -0.3]$ 因为被第 2 行的离群值 $7.0$ 拉爆缩放因子，整行被彻底压扁清零！
  \item \textbf{逐行（逐通道）量化}：第 1 行取 $s_1 = \frac{0.4}{7} \approx 0.05714$，第 2 行取 $s_2 = \frac{7.0}{7} = 1.0$。逐行量化：
  \[
    Q_{\mathrm{pc},1} = \operatorname{round}\!\left(\frac{[0.4, -0.3]}{0.4/7}\right) = \operatorname{round}([7, -5.25]) = [7, -5], \qquad
    Q_{\mathrm{pc},2} = [7, -5].
  \]
  反量化重构矩阵为
  \[
    \hat{W}_{\mathrm{pc}} = \begin{bmatrix} 7 \times \frac{0.4}{7} & -5 \times \frac{0.4}{7} \\ 7.0 & -5.0 \end{bmatrix}
    = \mathbf{\begin{bmatrix} 0.4000 & -0.2857 \\ 7.0000 & -5.0000 \end{bmatrix}},
  \]
  第 1 行得到高精度完整保留（误差仅 $0.0143$）！
\end{enumerate}

\textbf{(b) SmoothQuant（$\alpha = 0.5$）对角迁移与 W4A4 精度暴涨验算。} 激活向量第 1 维为 $10.0$、第 2 维为 $0.5$（动态范围高达 $20\times$），权重行最大绝对值分别为 $0.4$ 与 $7.0$（动态范围高达 $17.5\times$）。按 SmoothQuant 公式计算平滑缩放因子 $s_j = |x_j|^{0.5} / (\max|W_{j,:}|)^{0.5}$：
\[
  s_1 = \sqrt{\frac{10.0}{0.4}} = \sqrt{25} = \mathbf{5.0}, \qquad
  s_2 = \sqrt{\frac{0.5}{7.0}} = \frac{1}{\sqrt{14}} \approx \mathbf{0.26726}.
\]
执行等价缩放变换 $\tilde{x} = x \diag(s)^{-1}$ 与 $\tilde{W} = \diag(s) W$：
\[
  \tilde{x} = \left[\frac{10.0}{5.0},\; \frac{0.5}{1/\sqrt{14}}\right] = \mathbf{[2.0000,\; 1.8708]}, \qquad
  \tilde{W} = \begin{bmatrix} 5.0 \times 0.4 & 5.0 \times (-0.3) \\ \frac{7.0}{\sqrt{14}} & \frac{-5.0}{\sqrt{14}} \end{bmatrix}
  \approx \mathbf{\begin{bmatrix} 2.0000 & -1.5000 \\ 1.8708 & -1.3363 \end{bmatrix}}.
\]
\begin{itemize}
  \item \textbf{代数等价性}：$\tilde{x}\tilde{W} = x \diag(s)^{-1}\diag(s) W = x W = [7.5, -5.5]$ 严格恒等；
  \item \textbf{动态范围剧烈收窄}：$\tilde{x}$ 的通道幅值比从 $20\times$ 暴降至 $2.0 / 1.8708 \approx \mathbf{1.07\times}$，$\tilde{W}$ 的行幅值比也从 $17.5\times$ 暴降至 $\mathbf{1.07\times}$！
  \item \textbf{W4A4 逐张量量化对比}：迁移前若同时对 $x$ 和 $W$ 做 INT4 逐张量量化，$\hat{x}_{\mathrm{pt}} = [10, 0]$、$\hat{W}_{\mathrm{pt}}$ 第一行全零，输出竟为 $\hat{x}_{\mathrm{pt}}\hat{W}_{\mathrm{pt}} = \mathbf{[0, 0]}$（相对误差高达 $\mathbf{100\%}$！）；而 SmoothQuant 迁移后再做 W4A4 逐张量量化，输出为 $[8.00, -5.71]$，相对误差骤降至仅 $\mathbf{5.85\%}$！
\end{itemize}

\solItem{6.6.7}{编程实验：分组量化（Group-Wise Quantization）的 INT4 误差—显存帕累托前沿}

\textbf{(a) 通用分组量化函数实现要点。} 给定权重矩阵 $W \in \R^{M \times N}$ 且组大小 $g \mid N$，只需三行 NumPy 向量化代码即可实现 \texttt{quant\_groupwise(W, bits=4, group\_size=g)}：先将 $W$ 重塑为三维张量 \texttt{W\_g = W.reshape(M, N // g, g)}，沿最后一维计算每组独立缩放因子 $s = \max(\max|W_g| / q_{\max},\; 10^{-8})$（用 \texttt{np.max(np.abs(W\_g), axis=-1)} 实现），完成对称量化反量化后再 \texttt{reshape(M, N)}。每组额外存储 1 个 FP16/BF16 缩放因子（16 比特），均摊到组内 $g$ 个权重上，有效平均比特数为 $b_{\mathrm{eff}} = 4 + 16 / g$。

\textbf{(b) $2\%$ 孤立 $8\times$ 尖峰离群元素下的实测帕累托前沿基准。} 在 $256 \times 256$ 权重矩阵（$\mathcal{N}(0, 0.25^2)$ 并以 $2\%$ 概率随机乘以 $8\times$ 尖峰）上运行扫描，独立实测结果如下表：
\begin{center}
\footnotesize
\begin{tabular}{cccccc}
\toprule
分组大小 $g$ & $256$（逐通道） & $\mathbf{128}$（主流默认） & $64$ & $32$ & $16$ \\
\midrule
有效比特数 $b_{\mathrm{eff}} = 4 + 16/g$ & $4.0625$ & $\mathbf{4.1250}$ & $4.2500$ & $4.5000$ & $5.0000$ \\
INT4 输出相对误差（$\%$） & $35.10\%$ & $\mathbf{28.81\%}$ & $22.50\%$ & $16.89\%$ & $12.72\%$ \\
\bottomrule
\end{tabular}
\end{center}
\textbf{规律与甜点分析}：当离群尖峰随机散落于各行内部时，逐通道量化（$g=256$）每行平均落入约 $5$ 个孤立尖峰，拖垮整行全部 $256$ 个权重；缩小 $g$ 把单颗尖峰的“污染半径”严格锁死在 $g$ 个相邻权重内（对真实大模型更低的离群密度 $\sim 0.1\%$，$g=128$ 即可令绝大多数小组不含尖峰，将困惑度损失压至几乎为零），且仅付出 $+0.125\,\mathrm{bit}$ 极低显存代价（相对纯 4-bit 仅增重 $3.1\%$）。相比之下，继续压到 $g = 16$ 虽精度更高，但有效比特已涨至 $5.0\,\mathrm{bit}$（显存膨胀 $25\%$）且破坏了向量化内存对齐，因此 GPTQ 与 AWQ 工业界普遍把 $g = 128$ 作为最佳甜点。

\solItem{6.6.8}{编程实验：蒙特卡洛验证投机解码无损分布定理与加速比曲线}

\textbf{(a)--(b) 单步无损分布验证基准。} 在词表大小 $|\mathcal{V}| = 16$ 上按题意构造一对有显著差异的分布 $(p, q)$（设随机种子 2026，$z \sim \mathcal{N}(0, 1.5^2)$，$p = \softmax(z)$，$q = \softmax(0.6z + \mathcal{N}(0, 0.3^2))$），两者原始总变差距离高达 $\mathrm{TV}(p, q) = \frac{1}{2}\sum_x |p(x) - q(x)| = \mathbf{0.2529}$，理论一步接受率 $\alpha = \sum_x \min(p(x), q(x)) = 1 - \mathrm{TV}(p, q) = \mathbf{0.7471}$。

实现拒绝采样与残差重采样（先计算未归一化残差 $r = \max(0, p - q)$，若其和 $> 0$ 则归一化为 $p' = r / \sum r$，否则永不触发拒绝分支）。独立重复 $N = 2 \times 10^5$ 次后统计经验分布 $\hat{p}$：
\[
  \mathrm{TV}(\hat{p}, p) = \frac{1}{2}\sum_{x=1}^{16} |\hat{p}(x) - p(x)| = \mathbf{0.00253} < 5 \times 10^{-3},
\]
比原始草稿偏差 $0.2529$ 缩小了整整 \textbf{100 倍}，且随抽样数按标准蒙特卡洛误差律 $O(1/\sqrt{N})$ 趋于零，从实验上确凿印证了定理~\ref{thm:speculative-lossless} 的严格零损性。

\textbf{(c) $\gamma = 5$ 步投机循环的期望产出词元数验证。} 理论一步接受率 $\alpha = 0.7471$ 时，代入式~\eqref{eq:spec-expected-len} 得单轮期望产出词元数的精确理论值为
\[
  \E[L]_{\mathrm{theo}} = \frac{1 - \alpha^6}{1 - \alpha} = \frac{1 - 0.7471^6}{1 - 0.7471} = \mathbf{3.2667}\;\text{tokens/step}.
\]
对上述 $\gamma = 5$ 步截断投机循环独立模拟 $2 \times 10^4$ 轮，统计实测平均单轮产出词元数为 $\E[L]_{\mathrm{emp}} = \mathbf{3.2631}$ 词元/步，与理论值吻合至小数点后两位（相对误差仅 $0.11\%$）。至此，全卷从无模型 RL 到现代大模型推理系统的 48 道理论与工程闭环习题全部验算完毕！


\end{document}
